\documentclass[11pt,twoside,openany]{book}
\usepackage[a4paper,inner=31mm,outer=27mm,top=28mm,bottom=31mm,marginparwidth=17mm]{geometry}
\usepackage{fontspec}
\defaultfontfeatures{Ligatures=TeX,Scale=MatchLowercase}
\setmainfont{Malgun Gothic}[
  ItalicFont={Malgun Gothic},
  ItalicFeatures={FakeSlant=0.2},
  BoldItalicFont={Malgun Gothic Bold},
  BoldItalicFeatures={FakeSlant=0.2}
]
\setsansfont{Malgun Gothic}
\setmonofont{TeX Gyre Cursor}
\XeTeXlinebreaklocale "ko"
\XeTeXlinebreakskip=0pt plus 1pt
\usepackage{microtype}
\usepackage{graphicx}
\usepackage{mathtools}
\usepackage{amsthm}
\usepackage{unicode-math}
\setmathfont{Latin Modern Math}
\usepackage{bookmark}
\usepackage{hyperref}
\hypersetup{hidelinks,unicode=true,pdfencoding=auto}

\setcounter{secnumdepth}{3}
\setcounter{tocdepth}{2}
\numberwithin{equation}{section}
\renewcommand{\contentsname}{목차}
\makeatletter
\renewcommand{\@pnumwidth}{2.2em}
\renewcommand{\@tocrmarg}{3.2em}
\makeatother
\pdfstringdefDisableCommands{\def\textsuperscript#1{#1}}

\newcommand{\sourceunit}{front}
\newcommand{\src}[1]{%
  \phantomsection\hypertarget{src:\sourceunit:#1}{}%
  \marginpar{\raggedleft\footnotesize\sffamily #1}%
}
\newcommand{\sid}[1]{\phantomsection\hypertarget{fga:#1}{}}
\newcommand{\Osh}{\underline{O}}
\newcommand{\Zsh}{\underline{\mathbf Z}}
\newcommand{\Ccat}{\mathcal C}
\newcommand{\Homsh}{\underline{\operatorname{Hom}}}
\newcommand{\Extsh}{\underline{\operatorname{Ext}}}
\newcommand{\hsh}{\underline{h}}
\newcommand{\fquote}[1]{\textquotedbl#1\textquotedbl}
\newcommand{\seminar}[1]{%
  \begin{flushleft}\sffamily 부르바키 세미나\\#1\end{flushleft}%
}
\newcommand{\seminaren}[1]{%
  \begin{flushleft}\sffamily 부르바키 세미나\\#1\end{flushleft}%
}
\newcommand{\paperhead}[2]{%
  \begin{center}
    {\Large\bfseries #1\par}
    \vspace{0.6\baselineskip}
    #2
  \end{center}
  \vspace{0.8\baselineskip}
}

\theoremstyle{plain}
\newtheorem{proposition}{명제}[section]
\newtheorem{theoreme}[proposition]{정리}
\newtheorem{theorem}[proposition]{정리}
\newtheorem{corollaire}[proposition]{따름정리}
\newtheorem{corollary}[proposition]{따름정리}
\newtheorem{lemme}[proposition]{보조정리}
\newtheorem{lemma}[proposition]{보조정리}

\theoremstyle{definition}
\newtheorem{definition}[proposition]{정의}
\newtheorem{definitionen}[proposition]{정의}

\theoremstyle{remark}
\newtheorem*{remarque}{주}
\newtheorem*{remark}{주}

\renewcommand{\proofname}{증명}
\newcommand{\unit}[1]{%
  \chapter{#1}%
  \markboth{#1}{#1}%
}
\newcommand{\fsection}[2]{%
  \section*{#1. #2}%
  \addcontentsline{toc}{section}{#1. #2}%
}
\newcommand{\result}[2]{%
  \par\medskip\noindent\textbf{#1%
  \if\relax\detokenize{#2}\relax\else\space#2\fi. ---}\enspace
}

\setlength{\emergencystretch}{1.5em}
\hypersetup{
  pdftitle={대수기하학의 기초 -- 여덟 강연과 1962년 논평 -- 한국어 누적판},
  pdfauthor={알렉상드르 그로텐디크},
  pdfsubject={프랑스어 원전과 정오표를 바탕으로 한 한국어 AI 번역 교정판}
}

\begin{document}
\frontmatter
\begin{titlepage}
\centering
\vspace*{0.16\textheight}
{\Huge\bfseries 대수기하학의 기초\par}
\vspace{1.2\baselineskip}
{\Large 알렉상드르 그로텐디크\par}
\vfill
{\large 한국어 누적 교정판\par}
\vspace{0.5\baselineskip}
{\normalsize 여덟 강연, 각 강연의 정오표, 1962년 논평\par}
\end{titlepage}
\chapter*{이 판에 관하여}
\addcontentsline{toc}{chapter}{이 판에 관하여}
이 책은 알렉상드르 그로텐디크의 부르바키 세미나 강연 149, 182, 190,
195, 212, 221, 232, 236과 각 강연의 정오표 및 보충, 1962년 논평을
한 권에 수록한다. 정오표와 후대의 수학적 정정을 반영한 교정판이며,
필요한 원문 독법과 수학적 설명은 주석에 제시한다.

한국어 본문은 AI가 번역하였고, 교정과 편집 및 조판도 AI가 수행하였다.
출판자는 본문을 번역하지 않았다. 프랑스어 원전이 일차 자료이며,
영어판은 의미를 대조하는 이차 자료이다. 원문의 쪽 표지는 바깥 여백에
표시되어 있다.
\clearpage

\chapter*{출처와 권리}
\addcontentsline{toc}{chapter}{출처와 권리}
원저자는 알렉상드르 그로텐디크이며, 원전은
\emph{Fondements de la géométrie algébrique}의 여덟 부르바키 강연과
정오표 및 1962년 논평이다. 프랑스어 원전은
\href{https://www.numdam.org/}{NUMDAM}에서 열람할 수 있다.

대조에 사용된 \href{https://github.com/KokunoYumeto/fga-en}{영어판}에는
\href{https://github.com/thosgood/fga}{Tim Hosgood의 인간 번역}과 공유되는
영어 표현이 있다. 확인된 대조 버전은
\href{https://github.com/thosgood/fga/tree/9976dd5f8e2ddcb388f0a07f3aa334f120d1019f}{\nolinkurl{9976dd5f8e2ddcb388f0a07f3aa334f120d1019f}}이다.
이는 대조 버전을 식별하며 모든 문장이 그 버전에서 유래했다는 주장은 아니다.
원래 번역자의 주석은 Steve Hnizdur의 교정과 Jon Sterling의 기술 지원도
명시한다. 이 공헌은 영어 참고 번역에 관한 것이며, 한국어 AI 번역이나
별도의 AI 교정을 이들에게 귀속시키지 않는다.

원전과 이 번역 및 전사에 포함된 보호되는 기여의 권리는 각각의 권리자에게
있다. 이 판에서 발행자가 권리를 가진 모든 것은 CC0 1.0(https://creativecommons.org/publicdomain/zero/1.0/deed.ko)에 따라 퍼블릭 도메인에 헌정된다. 기존의 개별 권리 고지와 이미 부여된
허가는 계속 적용된다. 출처 표시는 이용 허가나 보증을 뜻하지 않는다.

이 판은 비상업적이며 순수한 교육 목적을 위해 무료로 제공된다. 이는
프랑스어 원전과 함께 번역 및 전사를 무료로 열람할 수 있게 하는 데
권리자들이 이의를 제기하지 않을 것이라는 추정에 따른 것이다.
이 추정은 허가를 받았다는 뜻이 아니며, 원전의 무료 열람 자체가 번역,
개작 또는 재배포를 허용하는 것은 아니다.

수록 자료의 권리자로서 공개에 이의가 있다면
\href{https://github.com/KokunoYumeto/fga-ko/issues}{GitHub 이슈}를 통해
알려 주시기 바란다. 출판자는 해당 권리자의 요청에 따라 해당 자료를 내린다.
이 판은 역사적 저자, 원 출판기관 또는 인용된 제삼자의 승인이나 보증을
뜻하지 않는다.
\clearpage

\tableofcontents
\mainmatter
\renewcommand{\sourceunit}{149}
\unit{강연 149 — 연접인 대수적 층에 대한 쌍대성 정리들}
\sid{149}
\src{169}
\seminar{(1957년 5월)}
\paperhead{연접인 대수적 층에 대한 쌍대성 정리들}%
  {알렉상드르 그로텐디크 지음}

이하의 결과들은 세르의 \fquote{대수적 쌍대성 정리}에서 영감을 받아
1955년 겨울과 1956년 겨울에 얻은 것이다. 이 결과들은 사영공간의
코호몰로지에 관한 상당히 초등적인 결과들 [3]과,
[2]에 제시된 형태의 카르탕--에일렌베르크 호몰로지 대수를 집중적으로
사용하여 매우 간단히 증명된다.

\fsection{1}{가군층의 Ext ({[2]}, 제3장 및 제4장)}

$X$를 단위원을 갖는 환(반드시 가환일 필요는 없음)의 층 $\Osh$가 주어진
위상공간이라 하자. $\Osh$-가군층들로 이루어진 아벨 범주
$\Ccat^{\Osh}$를 생각하자. 그 대상들을 $\Osh$-가군이라고도 한다.
이 범주의 모든 대상은 단사 분해를 갖는다는 것이 알려져 있으므로, 이 범주에서
잘 알려진 형식적 성질들을 지닌 $\operatorname{Ext}$ 함자들을 정의할 수
있다. 혼동을 피하기 위해 좀 더 정확히 말하면,
$\operatorname{Hom}_{\Osh}(X;A,B)$ 또는 간단히
$\operatorname{Hom}(X;A,B)$로 $A$에서 $B$로 가는 $\Osh$-준동형들의 아벨군을
나타내고, $\Homsh_{\Osh}(A,B)$로 $A$에서 $B$로 가는 준동형들의
싹의 층을 나타내겠다($A,B\in\Ccat^{\Osh}$). 고정된
$A\in\Ccat^{\Osh}$에 대하여, 각각 아벨군들의 범주
$\Ccat$와 $X$ 위의 아벨군층들의 범주
$\Ccat^{\Zsh}$에 값을 갖는 함자 $h_A$, $\hsh_A$를 다음 식으로 정의한다:
\begin{equation}\tag{1.1}
  h_A(B)=\operatorname{Hom}_{\Osh}(X;A,B),
  \qquad
  \hsh_A(B)=\Homsh_{\Osh}(A,B).
\end{equation}

$h_A$와 $\hsh_A$는 공변 왼쪽 완전 함자이다.
그 오른쪽 유도함자를 생각하고, 각각
$\operatorname{Ext}_{\Osh}^{p}(X;A,B)$와 $\Extsh_{\Osh}^{p}(A,B)$로 나타낸다. 그러므로
정의에 따라
\begin{equation}\tag{1.2}
\left\{
\begin{aligned}
\operatorname{Ext}_{\Osh}^{p}(X;A,B)
  &= (R^{p}h_A)(B)
   = H^{p}\bigl(\operatorname{Hom}_{\Osh}(X;A,C(B))\bigr),\\
\Extsh_{\Osh}^{p}(A,B)
  &= (R^{p}\hsh_A)(B)
   = H^{p}\bigl(\Homsh_{\Osh}(A,C(B))\bigr),
\end{aligned}
\right.
\end{equation}

여기서 기호 $R^{p}$는 오른쪽 유도함자를 취하는 것을 나타내고,
$C(B)$는 $\Ccat^{\Osh}$에서 $B$의 임의의 단사 분해를 나타낸다.
다음으로 나타낸
\[
  \Gamma:\Ccat^{\Zsh}\longrightarrow\Ccat
\]
\fquote{단면} 함자를 생각하자. 그 오른쪽 유도함자는 다음으로
나타낸다는 것을 상기하자:
\[
  B\longmapsto H^{p}(X,B).
\]

\src{170}
\begin{equation}\tag{1.3}
  H^{p}(X,B)=(R^{p}\Gamma)(B)=H^{p}\bigl(\Gamma(C(B))\bigr).
\end{equation}
물론 $h_A=\Gamma\hsh_A$이다. 한편 $\hsh_A$가
단사 대상을 $\Gamma$-비순환 대상으로 보냄을 확인할 수 있다. 따라서
잘 알려진 방식으로 다음을 얻는다:

\result{명제}{1} 모든 $\Osh$-가군 $A$에 대하여 $\Ccat^{\Osh}$ 위의
코호몰로지형 스펙트럼 함자가 존재하며, 이는 등급 함자
$\bigl(\operatorname{Ext}_{\Osh}^{*}(X;A,B)\bigr)$로 수렴하고 그 초항은
\begin{equation}\tag{1.4}
  E_{2}^{p,q}(A,B)=H^{p}\bigl(X,\Extsh_{\Osh}^{q}(A,B)\bigr).
\end{equation}
이로부터 \fquote{가장자리 준동형}들과 우리가 적지 않을 5항
완전열을 얻는다.

\result{따름정리}{1} $A$가 국소적으로 $\Osh^{n}$과 동형이면, 다음 표준
동형들이 성립한다:
\begin{equation}\tag{1.5}
  \operatorname{Ext}_{\Osh}^{p}(X;A,B)
  \overset{\sim}{\longleftarrow}
  H^{p}\bigl(X,\Homsh_{\Osh}(A,B)\bigr)
\end{equation}
(스펙트럼 열의 가장자리 준동형들이 주는 동형들이다). 특히 다음 표준
동형이 성립한다:
\begin{equation}\tag{1.6}
  \operatorname{Ext}_{\Osh}^{p}(X;\Osh,B)=H^{p}(X,B).
\end{equation}

이 결과들을 사용하려면 $\Extsh_{\Osh}^{p}(A,B)$를 구체적으로
계산할 수 있어야 한다. 이 함자들은 국소적으로 계산된다.
즉 $U$가 $X$의 열린집합이면 다음이 성립한다:
\[
  \Extsh_{\Osh}^{p}(A,B)|U
  =\Extsh_{\Osh|U}^{p}(A|U,B|U),
\]
이는 단사 $\Osh$-가군의 $U$로의 제한이
단사 $(\Osh|U)$-가군이기 때문이다. 또한 고정된 $x\in X$에 대하여
다음 함자적인 준동형들이 있다:
\begin{equation}\tag{1.7}
  \Homsh_{\Osh}(A,B)_x
  \longrightarrow
  \operatorname{Hom}_{\Osh_x}(A_x,B_x),
\end{equation}
이들은 $B$에 관한 코호몰로지 함자 사이의 하나의 준동형으로 유일하게
연장된다:
\begin{equation}\tag{1.8}
  \Extsh_{\Osh}^{p}(A,B)_x
  \longrightarrow
  \operatorname{Ext}_{\Osh_x}^{p}(A_x,B_x).
\end{equation}

\result{명제}{2} $A$가 $x$의 어떤 근방에서
준동형 $\Osh^{m}\longrightarrow\Osh^{n}$의 여핵과 동형이라고 하자. 그러면 $A$가 국소적으로 유한 생성 자유 가군들로 이루어진 차수 $p+1$까지의 부분 분해를 가지면 (1.8)은 차수 $p$에서 동형이고, $A$가 의사연접이면 모든 $p$에 대하여 동형이다. 하나의 유한 표시만으로는 차수 0에서의 비교만을 얻는다.\footnote{인쇄본에는 “(1.7)”로 되어 있으나,
$p$가 나타나고 앞 문장이 있으므로 “(1.8)”이어야 한다.} 특히
$A$가 연접 $\Osh$-가군이면 그러하다 [3].

\src{171}
\result{명제}{3} $L_{*}=(L_i)$를
$\Osh$-가군 $A$의 왼쪽 분해라 하자. 여기서 각 항은 국소적으로 어떤
$\Osh^{n}$과 동형인 $\Osh$-가군이다. 그러면 $\Extsh_{\Osh}^{*}(A,B)$는
$H^{*}\bigl(\Homsh_{\Osh}(L_{*},B)\bigr)$와 동일시되고,
$\operatorname{Ext}_{\Osh}^{*}(X;A,B)$는 복합체
$\Homsh_{\Osh}(L_{*},B)$에 관한 $X$의 초코호몰로지와 동일시된다.

증명은 표준적이다. $C(B)$가 $B$의 단사 분해일 때 이중복합체
$\Homsh_{\Osh}(L_{*},C(B))$와, $\Homsh_{\Osh}(L_{*},B)$ 및
$\Homsh_{\Osh}(A,C(B))$에서 이 이중복합체로 가는 자연
준동형들을 고려한다.

마지막으로, (1.2)에서 도입한 두 $\operatorname{Ext}$ 함자는 $B$에 관한
코호몰로지 함자일 뿐만 아니라, $B$에 관해서는 공변이고 $A$에
관해서는 반변인 코호몰로지 쌍함자이기도 하다.

\fsection{2}{Ext의 합성 법칙}

이 절의 결과는 CARTIER와 YONEDA가 서로 독립적으로 얻었다. 자세한
내용은 CARTIER의 강연 [1]을 보라. $\Ccat$를 아벨 범주라 하자.
$K$와 $L$을 $\Ccat$의 두 등급 대상이라 하자.
$\operatorname{Hom}(K,L)$로 다음 등급 아벨군을 나타내기로 한다. 그 차수
$n$의 성분은 $K$에서 $L$로 가는 차수 $n$의 동차 준동형들, 즉
준동형 $K^{i}\longrightarrow L^{i+n}$의 족 $(u_i)$들로 이루어진다.

$K$와 $L$이 복합체라고 하자(구체적으로 미분 연산자의 차수는 $+1$이라
하자). $\operatorname{Hom}(K,L)$에 다음과 같이 주어지는 미분
연산자를 정의한다.
\begin{equation}\tag{2.1}
  \delta(u)=du+(-1)^{n+1}ud,\qquad\text{여기서 }n=\deg(u),
\end{equation}
이로써 $\operatorname{Hom}(K,L)$은 차수 $+1$의 미분 연산자를 갖는
복합체가 된다. 차수 $n$의 순환은 (동차 연산자로서) $d$와 반가환하는
차수 $n$의 연산자이다.\footnote{인쇄본에는 “$u$와”로 되어 있으나, 식
(2.1)에 따르면 “$d$와”이어야 한다.} 따라서
$H^{*}\bigl(\operatorname{Hom}(K,L)\bigr)$을 생각할 수 있으며, 이는
$K$와 $L$의 호모토피형들의 불변량이다. 이를 $H^{*}(K,L)$로 나타내도 좋다.

세 번째 복합체 $M$이 주어지면, 준동형의 합성은 짝지음
\[
  \operatorname{Hom}(K,L)\times\operatorname{Hom}(L,M)
  \longrightarrow\operatorname{Hom}(K,M)
\]
을 정의하며 이는 미분 연산자들과 양립한다. 따라서 코호몰로지로
넘어가면 짝지음
\begin{equation}\tag{2.2}
  H^{*}(K,L)\times H^{*}(L,M)\longrightarrow H^{*}(K,M),
\end{equation}
을 얻으며, 이를 $(u,v)\longmapsto vu$로 나타낸다. 이 짝지음들은 자명한
결합법칙을 만족한다. 특히 $H^{*}(K,K)$는 단위원을 갖는 결합적
등급환이고, $H^{*}(K,L)$ (각각 $H^{*}(L,K)$)는 이 환 위의 오른쪽
(각각 왼쪽) 등급가군이다. 차수 $0$에서 (2.2)는 복합체 준동형들의
합성으로 귀착된다. 끝으로, 복합체의 완전열
\[
  0\longrightarrow K'\longrightarrow K\longrightarrow K''\longrightarrow 0,
\]
\src{172}
에서 모든 $i$에 대하여 $K'^{i}$가 $K^{i}$의 직접 인자와 동일시되면,
$\operatorname{Hom}(K'',L)$ 등을 비롯한 군의 복합체들의 완전열이
생기며, 따라서 연결 연산자
\[
  H^{i}(K',L)\longrightarrow H^{i+1}(K'',L)
\]
를 얻는다. $L$의 완전열에 관한 연결 연산자도 같은 방식으로 정의한다.
짝지음 (2.2)는 통상적인 의미에서 이 연결 연산자들과 양립한다.

이제 $\Ccat$의 모든 대상 $A$가 단사 분해 $C(A)$를 갖는다고 하자.
이중복합체 정리의 여러 변형 중 하나를 사용하면 다음이 성립한다.
\[
  H^{*}\bigl(C(A),C(B)\bigr)
  =H^{*}\bigl(\operatorname{Hom}(C(A),C(B))\bigr)
\]
은 다음과 표준적으로 동형이다.
\[
  H^{*}\bigl(\operatorname{Hom}(A,C(B))\bigr)
  =\operatorname{Ext}^{*}(A,B).
\]
위에서 살펴본 연결 연산자들은 $\operatorname{Ext}$의 연결 준동형들을
준다. 또한 여기서 짝지음 (2.2)는 연결 준동형들과 양립하는 결합적
짝지음
\begin{equation}\tag{2.3}
  \operatorname{Ext}^{*}(A,B)\times\operatorname{Ext}^{*}(B,C)
  \longrightarrow\operatorname{Ext}^{*}(A,C).
\end{equation}
을 준다. 특히 $\operatorname{Ext}^{*}(A,A)$는 단위원을 갖는 결합적
등급환이다. (같은 방식으로 $\operatorname{Ext}$ 함자들이 임의의 함자의
유도함자들에 작용함을 보일 수 있으나, 여기서는 이 사실을 사용하지 않는다.)

고려하는 범주가 $X$ 위의
$\Osh$-가군들의 범주 $\Ccat^{\Osh}$인 경우에는 다음 짝지음을 얻는다.
\begin{equation}\tag{2.4}
  \operatorname{Ext}_{\Osh}^{p}(X;A,B)
  \times\operatorname{Ext}_{\Osh}^{q}(X;B,C)
  \longrightarrow\operatorname{Ext}_{\Osh}^{p+q}(X;A,C),
\end{equation}
이 짝지음들은 앞에서 말한 대로 구체적으로 나타낼 수 있다. 더구나 같은
전개에서 아벨군의 범주를 $X$ 위의 아벨층의
범주로, $\operatorname{Hom}$ 함자들을
$\Homsh$ 함자들로 바꾸면, 같은 형식적 성질을 가지면서
이번에는 \fquote{국소적 성격}을 띠는 다음 짝지음도 정의된다:
\begin{equation}\tag{2.5}
  \Extsh_{\Osh}^{p}(A,B)\times\Extsh_{\Osh}^{q}(B,C)
  \longrightarrow\Extsh_{\Osh}^{p+q}(A,C).
\end{equation}
이 후자의 짝지음은 (1.8)의 준동형들이
$\operatorname{Ext}$ 사이의 짝지음과 양립한다는 사실을 관찰하면 더 구체화된다.

\src{173}
끝으로 $H^{p}(X,A)$ 함자들 사이에도 곱셈
구조(컵곱)가 있음을 상기하자. 그러면 명제 1의 스펙트럼 열들이
곱셈 구조와 양립함을 알 수 있다. 더
정확히 말하면, 스펙트럼 열 $E(A,B)$와 스펙트럼 열
$E(B,C)$에서 스펙트럼 열 $E(A,C)$로 가는 짝지음이 있으며, 수렴항에서는
전역 $\operatorname{Ext}$ 사이의 짝지음으로 귀착되고,
초기항에서는 (1.4)의 둘째 항에 나타나는
컵곱과 국소 $\operatorname{Ext}$의 짝지음으로부터 유도되는 짝지음으로 귀착된다. 이에 따라
특히 다음 \fquote{에지 준동형}들은
\begin{equation}\tag{2.6}
  \operatorname{Ext}_{\Osh}^{n}(X;A,B)
  \longrightarrow H^{0}\bigl(X,\Extsh_{\Osh}^{n}(A,B)\bigr)
\end{equation}
그리고
\begin{equation}\tag{2.7}
  H^{n}\bigl(X,\Homsh_{\Osh}(A,B)\bigr)
  \longrightarrow\operatorname{Ext}_{\Osh}^{n}(X;A,B)
\end{equation}
곱셈 구조와 양립한다. 따라서
국소적으로 어떤 $\Osh^{m}$과 동형인 층들로 한정하면, 이는
컵곱을 사용하여 전역 $\operatorname{Ext}$ 사이의 합성을 완전히 구체적으로 나타내며,
여기서는 동형 (1.5)를 감안하였다.

\fsection{3}{국소 코호몰로지의 결과}

$J$를 아이디얼로 갖는 단위원이 있는 가환환 $A$를 취하자. 우리는
가변 $A$-가군 $M$에 대하여 다음 함자적 준동형들을
\begin{equation}\tag{3.1}
\left\{
\begin{aligned}
\operatorname{Ext}_{A}^{p}(A/J,M)
  &\longrightarrow
  \operatorname{Hom}_{A}\!\left(\bigwedge^{p}J/J^{2},M\otimes A/J\right),\\
\operatorname{Tor}_{p}^{A}(A/J,M)
  &\longleftarrow
  \left(\bigwedge^{p}J/J^{2}\right)
  \otimes\operatorname{Hom}_{A}(A/J,M),
\end{aligned}
\right.
\end{equation}
여기서 텐서곱과 외승은 환 $A$ 위에서 취한다. 또한
$J/J^{2}$는 사실 $A/J$-가군이고 그
외승들은 이를 $A$-가군으로 보든 $A/J$-가군으로 보든 서로 같음에 유의하자.
준동형 (3.1)의 정의는 다음 데이터의 정의로 귀착된다. 즉,
$J$의 원소들로 이루어진 임의의 계 $\underline{x}=(x_1,\ldots,x_p)$마다
다음 준동형 $\varphi_{\underline{x}}$를 정의한다:
\begin{equation}\tag{3.2}
\left\{
\begin{aligned}
\varphi_{\underline{x}}:
  \operatorname{Ext}_{A}^{p}(A/J,M)&\longrightarrow M\otimes A/J,\\
\varphi_{\underline{x}}:
  \operatorname{Hom}_{A}(A/J,M)&\longrightarrow
  \operatorname{Tor}_{p}^{A}(A/J,M),
\end{aligned}
\right.
\end{equation}
이들은 다음 조건을 만족해야 한다:

\src{174}
\begin{enumerate}
\item[\textup{i.}] $\varphi_{x_1,\ldots,x_p}$는
계 $(x_i)$에 대해 교대 $A$-다중선형이다. 여기서 $x_i\in J$이다.
\item[\textup{ii.}] $\varphi_{x_1,\ldots,x_p}$는 $x_i$ 가운데 하나라도
$J^{2}$에 속하면 $0$이다.
\end{enumerate}
실제로 ii는 i에서 따른다. $a\in J$이면 $a\varphi_{\underline{x}}=0$이기 때문이다.
이는 (3.2)에 나오는 모든 가군이
$J$에 의해 소멸된다는 사실에서 알 수 있다.

$\varphi_{\underline{x}}$를 정의하기 위해, 복합체
$K_{\underline{x}}$를 생각하자. 그 바탕 $A$-가군은 $\bigwedge A^{p}$이고,
그 미분 연산자는 $A^{p}$ 위의 선형형식으로 본
성분 $x_1,\ldots,x_p$의 $\underline{x}$에 의한 내부곱 $i_{\underline{x}}$이다.
미분 연산자의 차수는 $-1$이고, 차수들은
음이 아니며, 이 복합체의 차수 $0$ 호몰로지는
$A/(x_1A+\cdots+x_pA)$이다. $x_i$들이 $J$에 속하므로,
증대사상 $K_{\underline{x},0}\longrightarrow A/J$를 얻는다. 따라서
$K_{\underline{x}}$는 증대된 자유 복합체로 나타나며, 그 증대 가군은
$A/J$이다. 잘 알려진 방법으로 다음 준동형들이 유도된다.
\[
\operatorname{Ext}_{A}^{*}\bigl(H_0(K_{\underline{x}}),M\bigr)
\longrightarrow
H^{*}\bigl(\operatorname{Hom}_{A}(K_{\underline{x}},M)\bigr)
\]
그리고
\[
\operatorname{Tor}_{*}^{A}\bigl(H_0(K_{\underline{x}}),M\bigr)
\longleftarrow
H_{*}\bigl(K_{\underline{x}}\otimes_A M\bigr),
\]
따라서 이를 증대 준동형으로부터 유도되는 $\operatorname{Ext}$와
$\operatorname{Tor}$의 준동형들과 합성한다. 그 증대 준동형은
$H_0(K_{\underline{x}})\longrightarrow A/J$이며, 다음 준동형들을 얻는다:
\begin{equation}\tag{3.3}
\left\{
\begin{aligned}
\psi_{\underline{x}}:
\operatorname{Ext}_{A}^{*}(A/J,M)
&\longrightarrow H^{*}\bigl(\operatorname{Hom}_{A}(K_{\underline{x}},M)\bigr),\\
\psi_{\underline{x}}:
\operatorname{Tor}_{*}^{A}(A/J,M)
&\longleftarrow H_{*}\bigl(K_{\underline{x}}\otimes_A M\bigr).
\end{aligned}
\right.
\end{equation}
그런데 최고 차수 $p$에서 첫째 우변의 코호몰로지는
$M/(x_1M+\cdots+x_pM)$이고, 둘째 우변의 호몰로지는 $M$에서
모든 $x_i$에 의해 소멸되는 원소들의 집합임을 즉시 알 수 있다. $x_i$들이 $J$에 속하므로
다음 준동형들을 얻는다.
\begin{equation}\tag{3.4}
\left\{
\begin{aligned}
H^{p}\bigl(\operatorname{Hom}_{A}(K_{\underline{x}},M)\bigr)
&\longrightarrow M\otimes A/J,\\
H_{p}\bigl(K_{\underline{x}}\otimes_A M\bigr)
&\longleftarrow\operatorname{Hom}_{A}(A/J,M).
\end{aligned}
\right.
\end{equation}
(3.3)과 (3.4)의 준동형들을 합성하면, 정의하고자 했던 (3.2)의
준동형들을 얻는다. i.를 확인하는 일은 번거롭지만
어렵지는 않다.

\src{175}
\result{명제}{4} $A$를 단위원을 갖는 가환환이라 하고,
$(x_1,\ldots,x_p)$를 $A$의 원소들의 열이라 하자.
모든 $1\leq i\leq p$에 대하여, $A$를 $(x_1,\ldots,x_{i-1})$이 생성하는 아이디얼로 나눈 몫에서
$x_i$의 상이 영인자가 아니라고 하자. $J$를 $x_i$들이 생성하는 아이디얼이라 하자.
그러면 $J/J^{2}$는 $x_i$들의 표준상들을 기저로 갖는 자유 $(A/J)$-가군이고,
복합체 $K_{\underline{x}}$는 $A/J$의 자유 분해이며, 모든
$A$-가군 $M$에 대하여 차수 $p$에서 (3.1)의 준동형들은
전단사이다. 또한 $J.M=0$이면 임의의 차수 $i$에서 정의된 유사한
준동형들도 전단사이다.

(다른 모든 결론이 따르는 핵심은 $K_{\underline{x}}$의 양의 차수
호몰로지가 모두 $0$이라는 사실이며, 이는 주어진 조건 아래 잘 알려져 있다.)

\result{따름정리}{1} $A$와 $J$가 위와 같다고 하고, 더 나아가
$A$가 완전체 $k$ 위의 차원 $n$인 정칙 아핀 대수이며
$A/J$도 정칙 아핀 대수라고 하자. $\Omega^{1}(A)$와
$\Omega^{1}(A/J)$로 각각의 켈러 미분 가군을 나타내자. 그러면
국소화와 양립하는 다음 표준 동형사상이 있다.
\begin{equation}\tag{3.5}
  \operatorname{Ext}_{A}^{p}\bigl(\Omega^{n-p}(A/J),\Omega^{n}(A)\bigr)=A/J.
\end{equation}
실제로 $\Omega^{n-p}(A/J)$는 계수(階數) $1$인 사영 $(A/J)$-가군이고, 마찬가지로
$\Omega^{n}(A)$는 계수(階數) $1$인 사영 $A$-가군이다.\footnote{인쇄본에는
「계수(階數) $n$」이라고 되어 있다. 최고차 형식의 두 가군은 계수(階數)
$1$이며, 이는 아래의 $\bigwedge^{p}(J/J^{2})$와의 동일시에도 필요하다.}
따라서 좌변은 다음과 같다.
\[
  \operatorname{Ext}_{A}^{p}(A/J,A)
  \otimes\Omega^{n-p}(A/J)'\otimes\Omega^{n}(A)
\]
(여기서 기호 $'$는 $(A/J)$-가군의 쌍대를 뜻한다.) 마지막 두 인자의
텐서곱은 $\bigwedge^{p}(J/J^{2})$와 동일시되므로, 전체는
다음과 동일시된다.
\[
  \operatorname{Ext}_{A}^{p}\!\left(A/J,\bigwedge^{p}(J/J^{2})\right),
\]
따라서 명제에 의해 다음과 동일시된다.
\[
  \operatorname{Hom}_{A}\!\left(\bigwedge^{p}J/J^{2},
  \bigwedge^{p}J/J^{2}\right),
\]
즉 $A/J$와 동일시된다. 특히
$\operatorname{Ext}_{A}^{p}\bigl(\Omega^{n-p}(A/J),\Omega^{n}(A)\bigr)$에는
$A/J$의 단위원에 대응하는 특별히 정해진 원소가 있으며, 이를
\emph{$A$ 안의 아이디얼 $J$의 기본류}라 한다. (실제로 이는
훨씬 더 넓은 조건 아래에서도 정의할 수 있다.) 따름정리 1은
더 기하적이고 더 전역적인 다음 형태로 나타낼 수 있다:

\result{따름정리}{2} $X$를 대수적으로 닫힌 체 $k$ 위의 비특이
다양체라 하고, $Y$를 $X$의 닫힌 비특이 부분다양체라 하자.
$\Osh_X$를 $X$의 국소환들의 층, $\Osh_Y$를 $Y$의 국소환들의 층이라
하고, 후자를 $\Osh_X$의 몫층으로 간주하자. $X$의 차원을 $n$,
$Y$의 차원을 $n-p$라 하자.

\src{176}
$\Omega_X$(각각 $\Omega_Y$)를 $X$(각각 $Y$) 위의 정칙
미분형식의 싹의 층이라 하자. 그러면 다음 표준 동형사상들이
있다:
\begin{equation}\tag{3.6}
  \Extsh_{\Osh_X}^{p}\bigl(\Omega_Y^{n-p},\Omega_X^{n}\bigr)=\Osh_Y,
\end{equation}
또는 같은 말로,
\begin{equation}\tag{3.6 bis}
  \Extsh_{\Osh_X}^{p}\bigl(\Osh_Y,\Omega_X^{n}\bigr)=\Omega_Y^{n-p}.
\end{equation}
식 (3.6 bis)는 $Y$가 특이 다양체일 때 $\Omega_Y^{n-p}$의 정의로
삼을 수 있다. 더 정확히 말하면 다음과 같다:

\result{명제}{5} $Y$를 차원이 $q=n-p$인, 차원 $n$의 비특이
대수다양체 $X$ 안의 대수적 부분집합이라 하자. $F$를 지지집합이
$Y$에 포함되는 $X$ 위의 연접인 대수적 층으로, $L$을 $X$ 위의
국소 자유인 대수적 층으로 하자. 그러면 층
$\Extsh_{\Osh_X}^{i}(F,L)$은 $i<p$일 때 영이고, $i=p$일 때에는 다음
표준 동형사상이 있다.
\begin{equation}\tag{3.7}
  \Extsh_{\Osh_X}^{p}(F,L)
  =\Homsh_{\Osh_X}\!\left(F,
  \Extsh_{\Osh_X}^{p}(\Osh_X/J,L)\right),
\end{equation}
여기서 $J$는 $F$를 소멸시키고 영점집합이 $Y$인 $X$ 위의 임의의
아이디얼층이다. 특히 $F$가 $Y$ 위의 연접인 대수적 층이면 다음이 성립한다.
\begin{equation}\tag{3.7 bis}
  \Extsh_{\Osh_X}^{p}(F,L)
  =\Homsh_{\Osh_Y}\!\left(F,\Extsh_{\Osh_X}^{p}(\Osh_Y,L)\right).
\end{equation}
끝으로 여전히 $F$가 $Y$ 위의 연접인 대수적 층이라고 할 때, 층
\[
  E^{i}(F)=\Extsh_{\Osh_X}^{p+i}(F,\Omega_X^{n})
\]
들은 대수적 공간 $Y$를 비특이 다양체 $X$ 안에 몰입시키는 방식에
의존하지 않는다.

문제는 국소적이므로 $X$가 아핀이고 $L=\Osh_X$라고 가정할 수 있다.
이로써 문제는 가환대수의 문제, 더 나아가 국소대수의 문제로 귀착된다. 즉
$A$가 정칙 국소환이고 $M$이 지지집합의 차원이
$\leq q=n-p$인 $A$-가군일 때, $\operatorname{Ext}_{A}^{i}(M,A)=0$이 $i<p$에 대해 성립하고 다음이 성립함을 보이면 된다.
\[
  \operatorname{Ext}_{A}^{p}(M,A)
  =\operatorname{Hom}_{A}\!\left(M,\operatorname{Ext}_{A}^{p}(A/J,A)\right),
\]
여기서 $J$는 $M$을 소멸시키고 \fquote{차원}이 $\leq q$인 임의의 아이디얼이다.
첫 번째 주장에는 $q$에 대한 귀납법을 쓴다. 간단한 단계적 분해 논법으로
$M$이 $A/J$ 꼴인 경우로 귀착된다. 이어서 $J$를 더 작은 아이디얼로
바꾸고 귀납 가정과 $\operatorname{Ext}$의 완전열을 이용하면, $J$가
명제 4에서와 같은 \fquote{매개변수계}로 생성되는 경우로 귀착되며, 이때 결과는 즉시 성립한다. 앞의 결과에 따르면 $J$가

\src{177}
\fquote{차원}이 $\leq q$인 고정된 아이디얼이면, $A/J$-가군의 범주에서
정의되는 반변 함자 $E(M)=\operatorname{Ext}_{A}^{p}(M,A)$는 왼쪽
완전하다. 또한 이 함자는 직합을 직접곱으로 보내므로 다음이 쉽게 따른다.
\[
  E(M)=\operatorname{Hom}_{A}\bigl(M,E(A/J)\bigr).
\]
\footnote{인쇄본에는 $E(A)$로 되어 있다. 여기서 다루는 범주에서 이를
표현하는 대상은 $A/J$이고, (3.7)에 필요한 등식은 $E(A/J)$이다.} 끝으로
명제 5의 마지막 주장은 더 미묘하며, 여기서는 제시할 수 없는 국소
쌍대성 정리를 이용하여 $E^{i}(F)$를 내재적으로 특징짓는 데서 따른다.

\result{따름정리}{} $\omega_Y^{q}$로 층
$\Extsh_{\Osh_X}^{p}(\Osh_Y,\Omega_X^{n})$을 나타내자. 그러면
$F$가 $Y$ 위의 연접 대수적 층이면 다음 함자적 동형사상이 성립한다:
\begin{equation}\tag{3.8}
  \Extsh_{\Osh_X}^{p}(F,\Omega_X^{n})
  =\Homsh_{\Osh_X}(F,\omega_Y^{q}).
\end{equation}

\fsection{4}{부분다양체에 대응하는 코호몰로지류}

이하 전체에서 $X$는 차원 $n$이고 체 $k$ 위에 정의된 대수적 집합을
나타내며, 간단히 하기 위해 이 체는 대수적으로 닫혔다고 가정한다. 6절을
제외하면 $X$는 비특이라고 가정한다. $\Osh_X$로 $X$ 위의 국소환들의 층을,
\[
  \Omega_X^{*}=\bigoplus_p\Omega_X^{p}
\]
로 $X$ 위의 미분형식들의 싹의 층을 나타낸다.\footnote{인쇄본은
합집합 기호를 사용한다. 그러나 등급 구조는 층 $\Omega_X^p$들의 직합
이다.} $Y$가 $X$의 닫힌 부분집합이면, 그 위의 연접 대수적 층을
$Y$ 밖에서 영인 $X$ 위의 연접층과 동일시한다. 특히 $\Osh_Y$와
$\Omega_Y$도 그러하다.

\result{보조정리}{1} $F$를 $X$ 위의 연접 대수적 층으로서 그 지지집합의
차원이 $\leq n-p$인 것으로 하고, $L$을 $X$ 위의 국소 자유 연접 대수적 층으로
하자. 그러면 $\operatorname{Ext}_{\Osh_X}^{i}(X;F,L)$은 $i<p$일 때 영이고,
다음 표준 동형사상이 성립한다.
\begin{equation}\tag{4.1}
  \operatorname{Ext}_{\Osh_X}^{p}(X;F,L)
  =H^{0}\bigl(X,\Extsh_{\Osh_X}^{p}(F,L)\bigr).
\end{equation}
또한 $F$가 닫힌 부분집합 $Y$ 위의 연접 대수적 층이고, 이 닫힌집합은 $X$ 안에서
차원이 $\leq n-p$이면 다음 표준 동형사상이 성립한다.
\begin{equation}\tag{4.1 bis}
  \Extsh_{\Osh_X}^{p}(F,L)
  =\Homsh_{\Osh_X}\!\left(
    F\otimes L'\otimes\Omega_X^{n},\omega_Y^{n-p}
  \right),
\end{equation}
여기서 $\omega_Y^{n-p}$는 명제 5의 따름정리에서 정의한 $Y$ 위의 층이다
($\Omega_Y^{n-p}$와 동일시되는데, 이는 $Y$가 비특이일 때이다).

공식 (4.1)은 명제 1의 스펙트럼 열과 명제 5에서 곧바로 따른다. 공식
(3.8)에 따라 다음과 같이 쓸 수 있다.

\src{178}
\begin{align*}
\Extsh_{\Osh_X}^{p}(F,L)
&=L\otimes(\Omega_X^{n})'\otimes
  \Extsh_{\Osh_X}^{p}(F,\Omega_X^{n})\\
&=L\otimes(\Omega_X^{n})'\otimes
  \Homsh_{\Osh_X}(F,\omega_Y^{q})\\
&=\Homsh_{\Osh_X}\!\left(
  F\otimes L'\otimes\Omega_X^{n},\omega_Y^{q}\right),
\end{align*}
여기서 $q=n-p$이고, 이로부터 공식 (4.1 bis)를 얻는다.

특히 $F=\Osh_Y$, $L=\Omega_X^{p}$로 두자.
$\Omega_X^{n}\otimes(\Omega_X^{p})'=\Omega_X^{n-p}$임을 고려하면 다음 표준
동형사상을 얻는다.
\begin{equation}\tag{4.2}
  \operatorname{Ext}_{\Osh_X}^{p}(X;\Osh_Y,\Omega_X^{p})
  =\operatorname{Hom}_{\Osh_X}
  (X;\Omega_X^{n-p},\omega_Y^{n-p}).
\end{equation}
이제 간단히 하기 위해 $Y$가 비특이라고 가정하자. 그러면
$\omega_Y^{n-p}=\Omega_Y^{n-p}$이다. $\Omega_X^{n-p}$에서
$\Omega_Y^{n-p}$로 가는 자연스러운 준동형이 있으므로, 층
$\Extsh_{\Osh_X}^{p}(\Osh_Y,\Omega_X^{p})$의 표준 단면 $s_Y$를 얻는다.
$Y$의 모든 성분의 차원이 $n-p$라면 이 단면을 층
$\Extsh_{\Osh_X}^{p}(\Osh_Y,\Omega_X^{p})$의 \emph{기본 단면}이라 부르겠다.
이 단면은 (4.1)에 따라
$\operatorname{Ext}_{\Osh_X}^{p}(X;\Osh_Y,\Omega_X^{p})$의 한 원소를 정한다. 한편
자연스러운 준동형 $\Osh_X\longrightarrow\Osh_Y$는 다음 준동형을 정한다:
\[
  \operatorname{Ext}_{\Osh_X}^{p}(X;\Osh_Y,\Omega_X^{p})
  \longrightarrow
  \operatorname{Ext}_{\Osh_X}^{p}(X;\Osh_X,\Omega_X^{p})
  =H^{p}(X,\Omega_X^{p}).
\]
그리하여 $H^{p}(X,\Omega_X^{p})$의 한 원소를 얻으며, 이를 $P_X(Y)$로 나타내고
\emph{$Y$가 $X$ 안에서 갖는 코호몰로지류}라 부른다. 따라서 이 원소는 단면
$s_Y$가 $\Extsh_{\Osh_X}^{p}(\Osh_Y,\Omega_X^{p})$에서 정하는 원소로부터 다음
준동형 도식을 통해 얻어진다:
\begin{equation}\tag{4.3}
\begin{aligned}
H^{p}(X,\Omega_X^{p})
&=\operatorname{Ext}_{\Osh_X}^{p}(X;\Osh_X,\Omega_X^{p})\\
&\longleftarrow\operatorname{Ext}_{\Osh_X}^{p}(X;\Osh_Y,\Omega_X^{p})
\xrightarrow{\ \sim\ }
H^{0}\!\left(X,\Extsh_{\Osh_X}^{p}(\Osh_Y,\Omega_X^{p})\right).
\end{aligned}
\end{equation}

차원 $n-p$인 비특이 사이클이라 함은 차원 $n-p$인 $X$ 안의 비특이 기약
부분다양체들로 생성되는 자유가환군의 한 원소를 뜻한다. 그러면 함수
$Y\longmapsto P(Y)$는 차원 $n-p$인 $X$ 위의 비특이 사이클들의 군에서
군 $H^{p}(X,\Omega_X^{p})$으로 가는 준동형으로 연장된다.

$Z^{n-p}$와 $Z'^{n-p'}$를 각각 차원이 $n-p$와 $n-p'$인 비특이 사이클이라
하자. $Z$의 모든 성분이 $Z'$의 모든 성분과 횡단적으로 교차하면, 두 사이클이
\emph{횡단적으로 교차한다}고 한다. 그러면 사이클 $Z.Z'$가 정의되며, 이는
차원이 $n-p-p'$인 비특이 사이클이다. 이때 다음 정리가 성립한다.

\src{179}
\result{정리}{1} $Z^{n-p}$와 $Z'^{n-p'}$가 횡단적으로 교차하는 비특이
사이클이면 다음이 성립한다.
\begin{equation}\tag{4.4}
  P_X(Z.Z')=P_X(Z).P_X(Z'),
\end{equation}
여기서 오른쪽 항의 곱은 컵곱
\[
  H^{p}(X,\Omega_X^{p})\times H^{p'}(X,\Omega_X^{p'})
  \longrightarrow H^{p+p'}(X,\Omega_X^{p+p'})
\]
이다. (여기서는 $X$가 어떤 사영 공간의 국소 닫힌 부분집합과 동형이라고
가정한다.)

마지막 가정은 $X$ 위의 모든 연접 대수적 층이 국소 자유인 연접 대수적
층의 몫임을 결론내리는 데에만 사용된다(세르). 따라서 그러한 층은
국소 자유층에 의한 왼쪽 분해를 갖는다. 정리 1을 증명하기 위해
$Z$와 $Z'$가 횡단적으로 교차하는 기약 비특이 부분다양체 $Y$와 $Y'$라고
가정해도 된다. $L_{*}$를 $\Osh_Y$의 국소 자유층에 의한 왼쪽
분해라 하자. 그러면 명제 3에 따라 준동형 그림 (4.3)은 다음 그림과
동일시된다.
\[
R^{p}\Gamma(\Omega_X^{p})
\xleftarrow{\alpha}
\mathbb R^{p}\Gamma\!\left(\Homsh_{\Osh_X}(L_{*},\Omega_X^{p})\right)
\xrightarrow[\sim]{\beta}
\Gamma\!\left(H^{p}\bigl(\Homsh_{\Osh_X}(L_{*},\Omega_X^{p})\bigr)\right).
\]
여기서 $\Gamma$는 $X$ 위의 아벨군의 층들의 범주에서 ``단면군''을 취하는
함자이고, $\mathbb R^{p}\Gamma$는 이 함자의 $p$차 초코호몰로지를,
$R^{p}\Gamma$는 그 $p$번째 유도함자를 나타낸다. 단순히 하기 위해
$L_0=\Osh_X$이고 증대 준동형 $L_0\longrightarrow\Osh_Y$가 자연스러운
준동형이라고 가정해도 된다. 그러면 $\alpha$는 복합체의 준동형
$\Osh_X\longrightarrow L_{*}$로부터 얻어진다. 여기서 $\Osh_X$는 차수
$0$에만 놓인 복합체로 간주하며, 층들의 복합체 $K$가 차수 $0$에만
놓이면 $\mathbb R^{p}\Gamma(K)=R^{p}\Gamma(K_0)$임을 고려한다. 준동형
$\beta$는 잘 알려진 ``가장자리 사상(edge-homomorphism)''이다.
$\Osh_{Y'}$의 국소 자유층에 의한 분해 $L'_{*}$에 관한 유사한 그림을 생각하고,
다음과 같이 짝짓기로 이루어진 가환 그림을 생각하자.
\begin{equation}\tag{4.5}
\resizebox{\textwidth}{!}{$
\begin{array}{ccc}
R^{p}\Gamma(\Omega_X^{p})
&\longleftarrow
\mathbb R^{p}\Gamma\!\left(\Homsh_{\Osh_X}(L_{*},\Omega_X^{p})\right)
&\xrightarrow{\sim}
\Gamma H^{p}\!\left(\Homsh_{\Osh_X}(L_{*},\Omega_X^{p})\right)
\\[3pt]
\times&\times&\times\\[3pt]
R^{p'}\Gamma(\Omega_X^{p'})
&\longleftarrow
\mathbb R^{p'}\Gamma\!\left(\Homsh_{\Osh_X}(L'_{*},\Omega_X^{p'})\right)
&\xrightarrow{\sim}
\Gamma H^{p'}\!\left(\Homsh_{\Osh_X}(L'_{*},\Omega_X^{p'})\right)
\\[3pt]
\downarrow&\downarrow&\downarrow\\[3pt]
R^{p+p'}\Gamma(\Omega_X^{p+p'})
&\longleftarrow
\mathbb R^{p+p'}\Gamma\!\left(
\Homsh_{\Osh_X}(L_{*}\otimes L'_{*},\Omega_X^{p+p'})\right)
&\xrightarrow{\sim}
\Gamma H^{p+p'}\!\left(
\Homsh_{\Osh_X}(L_{*}\otimes L'_{*},\Omega_X^{p+p'})\right).
\end{array}
$}
\end{equation}

\src{180}
오른쪽 두 열의 짝짓기는 층 복합체들의 짝짓기
\[
\Homsh_{\Osh_X}(L_{*},\Omega_X^{p})
\times\Homsh_{\Osh_X}(L'_{*},\Omega_X^{p'})
\longrightarrow
\Homsh_{\Osh_X}(L_{*}\otimes L'_{*},\Omega_X^{p+p'}),
\]
으로부터 얻어진다. 이 짝짓기는 외적
$\Omega_X^{p}\times\Omega_X^{p'}\longrightarrow\Omega_X^{p+p'}$을 사용하여
정의한다. 첫째 열의 짝짓기는 (외적에 관한) 컵곱이다. 나는 (4.5)의
마지막 행이, (4.3)과 유사하되 $Y$를 $Y\cap Y'$으로, $p$를 $p+p'$으로
바꾼 동형들의 그림과 동일시된다고 주장한다. 이를 위해서는
$L\otimes L'$이 $\Osh_{Y\cap Y'}$의 (명백히 국소 자유인) 분해임을
보이면 충분하다. 실제로
\[
  H_0(L\otimes L')=\Osh_Y\otimes\Osh_{Y'}=\Osh_{Y\cap Y'}
\]
이고
\[
  H_i(L\otimes L')=\operatorname{Tor}_{i}^{\Osh_X}(\Osh_Y,\Osh_{Y'})=0
  \qquad(i>0),
\]
인데, 이는 $Y$와 $Y'$가 횡단적으로 교차하기 때문이다. 이제 다음 식에서
정리 1이 따른다.
\begin{equation}\tag{4.6}
  s_Y.s_{Y'}=s_{Y.Y'}
\end{equation}
(여기서 왼쪽 항의 곱은 (4.5)의 마지막 열에 나타난 곱이다.) 순전히
국소적인 식 (4.6)은 $L_{*}$와 $L'_{*}$로 명제 4에서 고려한 분해들을
택하면 어렵지 않게 확인된다. 마찬가지로 (더 쉬운 증명으로)
$Z\longmapsto P_X(Z)$가 데카르트 곱과 양립함을 보인다.
\begin{equation}\tag{4.7}
  P_{X\times X'}(Z\times Z')=P_X(Z)\otimes P_{X'}(Z')
\end{equation}
(이 식은 $Z$와 $Z'$가 각각 비특이 다양체 $X$와 $X'$ 위의 비특이
사이클이고, $Z\times Z'$를 $X\times X'$ 위의 비특이 사이클로 간주할 때
성립한다.) (4.4)와 (4.7)에서 $P_X(Z)$가 비특이 다양체의 사상
$f:X\longrightarrow X'$에 의한 ``역상'' 연산과도 양립함이 따른다.
\begin{equation}\tag{4.8}
  P_X\bigl(f^{-1}(Z')\bigr)=f^{*}\bigl(P_{X'}(Z')\bigr),
\end{equation}
이 식은 $Z'$가 $X'$ 위의 비특이 사이클이고 $f$가 $Z'$에 ``횡단적''일 때,
즉 $f$의 그래프가 $X\times X'$ 안의 사이클 $X\times Z'$와 횡단적으로
교차할 때 성립한다.

\src{181}
\result{따름정리}{1} $X$, $X'$을 사영공간 안에서 국소 닫힌 두 비특이
다양체라 하고, $X'$은 완비라고 하자. $U$를 $X\times X'$ 위의 비특이
사이클이라 하고, $a$, $b$를 $X'$의 두 점이라 하되 $U$가 사이클
$X\times(a)$ 및 $X\times(b)$와 횡단적으로 교차한다고 하자. $Z$, $Z'$을
다음을 만족하는 $X$ 위의 비특이 사이클이라 하자.
\[
  Z\times(a)=(X\times(a)).U,
  \qquad
  Z'\times(b)=(X\times(b)).U.
\]
그러면
\[
  P_X(Z)=P_X(Z').
\]

실제로 $f_a:X\longrightarrow X\times X'$를 $f_a(x)=(x,a)$로 정의하자.
그러면 (4.8)에 따라 $P(Z)=f_a^{*}(P(U))$이고, 마찬가지로
$P(Z')=f_b^{*}(P(U))$이다. 한편 K\"unneth 공식
\[
  H^{*}(X\times X',\Omega_{X\times X'}^{*})
  =H^{*}(X,\Omega_X^{*})\otimes H^{*}(X',\Omega_{X'}^{*})
\]
과 $H^{0}(X',\Osh_{X'})$가 스칼라들로만 이루어져 있다는 사실을
이용하면 $f_a^{*}=f_b^{*}$를 쉽게 얻으며, 이로써 결론이 따른다.

각 $x\in X$에 대하여 점 $(x)$는 $X$의 여차원 $n$인 비특이
부분다양체이므로 $H^{n}(X,\Omega_X^{n})$의 원소 $\varepsilon_x$를 정한다.
$X$가 비특이 사영 다양체이면 따름정리 1에 따라 $\varepsilon_x$는 선택한
점 $x$에 의존하지 않는다. 이를 $\varepsilon_X$로 나타내고
$H^{n}(X,\Omega_X^{n})$의 \emph{기본류}라 부른다.

\result{비고}{} 만족스러운 이론을 얻으려면 임의의 사이클 $Z$에 대하여
$P_X(Z)$를 정의하고, 사이클들의 고유 교차에 대하여 정리 1을 증명해야 할
것이다. (이 강연록을 쓸 당시에는 아직 일반적인 경우에 이 일이 이루어지지
않았다.) 이것이 이루어졌다고 가정하면 따름정리 1은 다음과 같이 된다.
$Z$와 $Z'$이 대수적으로 동치인 두 사이클이면 $P_X(Z)=P_X(Z')$이다
(이 명제는 $Z$와 $Z'$이 비특이라 하더라도 앞의 결과에서 나오는 것 같지
않다).

\fsection{5}{쌍대성 정리}

이 절에서 $X$는 차원 $n$인 비특이 사영 다양체를 나타낸다.

\result{정리}{2} $H^{n}(X,\Omega_X^{n})$의 기본류 $\varepsilon_X$는 이
벡터공간의 기저를 이룬다.

(증명은 아래에서 제시한다.) 이 정리에 따라
$H^{n}(X,\Omega_X^{n})$를 체 $k$와 동일시할 수 있다. 이제 제2절에서
고찰한 짝짓기들을 생각하자. 이들은 특히 다음 짝짓기를 준다.

\src{182}
\[
\operatorname{Ext}_{\Osh_X}^{p}(X;\Osh_X,F)
\times\operatorname{Ext}_{\Osh_X}^{n-p}(X;F,\Omega_X^{n})
\longrightarrow
\operatorname{Ext}_{\Osh_X}^{n}(X;\Osh_X,\Omega_X^{n}),
\]
즉,
\begin{equation}\tag{5.1}
  H^{p}(X,F)
  \times\operatorname{Ext}_{\Osh_X}^{n-p}(X;F,\Omega_X^{n})
  \longrightarrow H^{n}(X,\Omega_X^{n}).
\end{equation}
정리 2를 고려하면 이 짝짓기는 다음 준동형을 정한다.
\begin{equation}\tag{5.2}
  \operatorname{Ext}_{\Osh_X}^{n-p}(X;F,\Omega_X^{n})
  \longrightarrow\bigl(H^{p}(X,F)\bigr)'.
\end{equation}
이 준동형은 $F$에 관해 함자적이며, 완전열
$0\longrightarrow F'\longrightarrow F\longrightarrow F''
\longrightarrow0$에 딸린 연결 준동형들과 가환한다.

\result{정리}{3} 준동형 (5.2)는 동형이다.

특히 세르의 다음 결과를 다시 얻는다.

\result{따름정리}{} $E$를 $X$ 위의 대수적 벡터 다발이라 하고,
$\Osh_X(E)$를 $E$의 정칙 단면의 싹들로 이루어진 층이라 하자. 그러면
다음 표준 동형들이 성립한다.
\begin{equation}\tag{5.3}
  \bigl(H^{p}(X,\Osh_X(E))\bigr)'
  =H^{n-p}\bigl(X,\Omega_X^{n}\otimes\Osh_X(E')\bigr).
\end{equation}
명제 1에 정리 3과 따름정리 1을 적용하면 충분하다.

정리 2와 정리 3은 다음 진술로부터 따른다.

\noindent\textbf{(D) 다음 준동형}
\begin{equation}\tag{5.2 bis}
  \operatorname{Ext}_{\Osh_X}^{n-p}(X;F,\Omega_X^{n})
  \longrightarrow\bigl(H^{p}(X,F)\bigr)'\otimes L,
  \qquad L=H^{n}(X,\Omega_X^{n}),
\end{equation}
은 짝짓기 (5.1)에서 유도되며 동형이다.

실제로 (D)가 정리 2를 함의함을 보이자. $k_x=\Osh_{(x)}$를 한 점
$x\in X$로 이루어진 축소 다양체의 국소환의 층이라 하자. 표준 준동형
$\Osh_X\longrightarrow k_x$와 이에 딸린 준동형
\begin{equation}\tag{5.3 bis}
  H^{0}(X,\Osh_X)\longrightarrow H^{0}(X,k_x),
\end{equation}
을 생각하자. 그 쌍대 준동형은 다음 준동형과 동일시된다.
\begin{equation}\tag{5.4}
  \operatorname{Ext}_{\Osh_X}^{n}(X;k_x,\Omega_X^{n})\otimes L'
  \longrightarrow
  \operatorname{Ext}_{\Osh_X}^{n}(X;\Osh_X,\Omega_X^{n})\otimes L'
\end{equation}
이는 $\Osh_X\longrightarrow k_x$에 의해 $\operatorname{Ext}^{n}$에
유도되는 준동형
\begin{equation}\tag{5.5}
  \operatorname{Ext}_{\Osh_X}^{n}(X;k_x,\Omega_X^{n})
  \longrightarrow
  \operatorname{Ext}_{\Osh_X}^{n}(X;\Osh_X,\Omega_X^{n})
  =H^{n}(X,\Omega_X^{n}).
\end{equation}
에서 유도된다.

\src{183}
(5.3 bis)가 동형이므로\footnote{인쇄본은 182쪽에서 두 번째 공식
“(5.3)”을 제시하고 여기서 그것을 참조한다. 식별자 “(5.3 bis)”는 이
공식을 첫 번째 (5.3)과 구별한다.} (5.4)도 동형이고, 따라서 (5.5)도
동형이다. 한편 (4.2)에 따라 $s_{(x)}$는
$\operatorname{Ext}_{\Osh_X}^{n}(X;k_x,\Omega_X^{n})$의 기저이므로, 그
상인 $\varepsilon_X$는 $H^{n}(X,\Omega_X^{n})$의 기저이다.

이제 진술 (D)를 증명하는 일만 남았다. 이는 다음 보조정리들에 요약된
초등적 사실들로부터 순전히 형식적으로 따를 것이다. 이 보조정리들에서는
$X$를 차원 $r$인 사영공간 $P$의 닫힌 부분집합으로 가정한다. $X$는
특이일 수도 있고 비특이일 수도 있다. [3]에서 $\Osh(m)$으로 나타낸 $P$
위의 층을 $\Osh_P(m)$으로 나타내고, $X$ 위의 그와 같은 층을
$\Osh_X(m)$으로 나타낸다.

\result{보조정리}{2} $X=P$이고 $F=\Osh_P(m)$이면 진술 (D)가 성립한다.

이 보조정리는 직접 계산으로 확인된다. $H^{i}(P,\Osh_P(m))$의 명시적
계산은 [3]에 있지만 더 초등적으로도 할 수 있다. 짝짓기 (5.1)를
명시하는 데 필요한 컵곱
\[
  H^{i}(P,\Osh_P(m))\times H^{j}(P,\Osh_P(m'))
  \longrightarrow H^{i+j}(P,\Osh_P(m+m'))
\]
의 계산에는 어려움이 없다.\footnote{인쇄본에는 $\Osh_P(r+r')$로 되어
있으나, 두 인자 $\Osh_P(m)$과 $\Osh_P(m')$에 따라
$\Osh_P(m+m')$이어야 한다.}

\result{보조정리}{3} $X$ 위의 모든 연접 대수적 층 $F$는
$\Osh_X(-m)^{k}$의 어떤 몫층과 동형이다. 여기서 $m$은 원하는 만큼
크게 잡을 수 있다.

$m$이 충분히 클 때 $F\otimes\Osh_X(m)$이 ``그 단면들로 생성된다''는
사실에서 따른다([3] 참조).

\result{보조정리}{4} $i>0$이라 하자. 그러면
$m$이 충분히 클 때 $H^{r-i}(P,\Osh_P(-m))=0$이고, $X$ 위의 임의의
연접 대수적 층 $B$에 대하여 다음이
\[
  \operatorname{Ext}_{\Osh_X}^{i}(X;\Osh_X(-m),B)=0
\]
$m$이 충분히 클 때 성립한다.

첫 번째 명제는 위에서 언급한 명시적 계산으로부터 따른다. 두 번째 명제에
대해서는 다음 동형이 있음을 주목한다.
\[
  \operatorname{Ext}_{\Osh_X}^{i}(X;\Osh_X(-m),B)
  =H^{i}(X,B\otimes\Osh_X(m))
\]
(명제 1, 따름정리 1). 따라서 [3]의 잘 알려진 결과에 의해 결론이
따른다. 앞의 두 보조정리를 결합하면 다음 따름정리를 얻는다.

\result{따름정리}{} $i>0$이라 하자. 그러면 $P$ 위의 연접 대수적 층의
범주에서 정의된 함자 $F\longmapsto H^{r-i}(P,F)$는 에피사상에 의해
소거 가능하고, $X$ 위의 연접 대수적 층의 범주에서 정의된 함자
$\operatorname{Ext}_{\Osh_X}^{i}(X;F,B)$도 마찬가지이다.

\src{184}
\result{보조정리}{5} $A$와 $B$를 $X$ 위의 두 연접 대수적 층이라 하고,
$A(m)=A\otimes\Osh_X(m)$이라 놓자. 그러면 $m$이 충분히 클 때 표준
준동형
\[
  \operatorname{Ext}_{\Osh_X}^{1}(X;A(-m),B)
  \longrightarrow H^{0}\bigl(X,\Extsh_{\Osh_X}^{1}(A(-m),B)\bigr)
  =H^{0}\bigl(X,\Extsh_{\Osh_X}^{1}(A,B)(m)\bigr)
\]
은 동형이다.

이는 명제 1의 스펙트럴 수열을 $A(-m)$과 $B$에 적용하면 곧바로 따른다.
실제로 다음을 얻기 때문이다.
\[
  E_{2}^{p,q}(A(-m),B)
  =H^{p}\bigl(X,\Extsh_{\Osh_X}^{q}(A(-m),B)\bigr)
  =H^{p}\bigl(X,\Extsh_{\Osh_X}^{q}(A,B)(m)\bigr),
\]
$p>0$이고 $m$이 충분히 크면 이는 0이다.

이제 $X=P$인 경우에 (D)를 증명하자. 먼저 $p=n$일 때 (5.2 bis)가
동형임을 보인다. 이때 양변은 모두 왼쪽 완전 함자이다
($H^{r+1}(P,F)=0$이기 때문이다). 따라서 보조정리 3에 의해
$F=\Osh_P(-m)$인 경우에 명제를 증명하면 충분한데, 이 경우는
보조정리 2에 포함되어 있다. 준동형들 (5.2 bis)는 함자적이고
연결 준동형들과 양립하며, $p<n$일 때 (5.2 bis)의 양변은 모두
$F$에 관하여 에피사상에 의해 소거 가능한 함자이므로(보조정리 4의 따름정리),
표준적인 논법에 따라 (5.2 bis)는 모든 $p$에 대하여 동형이다.
이로써 사영공간에 대한 쌍대성 정리가 증명되었다. 이제 $X$는 임의이되
비특이라고 하자. $P$에 대한 쌍대성 정리에 의해 다음 동형이 있다.
\[
  H^{n}(X,F)'=H^{n}(P,F)'
  =\operatorname{Ext}_{\Osh_P}^{r-n}(P;F,\Omega_P^{r}).
\]
4절의 보조정리 1에 의해 마지막 항은 다음과 동일시된다.
\[
  \operatorname{Hom}_{\Osh_P}(P;F,\omega_X^{n})
  =\operatorname{Hom}_{\Osh_X}(X;F,\Omega_X^{n})
  =\operatorname{Ext}_{\Osh_X}^{0}(X;F,\Omega_X^{n}),
\]
따라서 다음 동형을 얻는다.
\begin{equation}\tag{5.6}
  H^{n}(X,F)'=\operatorname{Hom}_{\Osh_X}(X;F,\Omega_X^{n})
  =\operatorname{Ext}_{\Osh_X}^{0}(X;F,\Omega_X^{n}).
\end{equation}
$F=\Omega_X^{n}$으로 놓으면 다음 동형을 얻는다.
\begin{equation}\tag{5.7}
  \gamma:H^{n}(X,\Omega_X^{n})\xrightarrow{\sim}k.
\end{equation}

\src{185}
동형 (5.6)은 (5.7)을 사용하여 $L=k$로 놓았을 때 $p=n$에 대한
(5.2 bis)와 다름없음을 확인할 수 있다. 따라서 (5.2 bis)는 $p=n$일 때
동형이다. 모든 $p$에 대하여 동형임을 증명하려면, $p<n$일 때
(5.2 bis)의 양변이 $F$에 관하여 에피사상에 의해 소거 가능한 함자임을
보이면 충분하고, 보조정리 3을 고려하면 더 강하게 양변이
$F=\Osh_X(-m)$이고 $m$이 충분히 클 때 0임을 보이면 충분하다. 첫 번째
변에 대해서는 보조정리 4로부터 따르고, 두 번째 변에 대해서는
$P$ 위의 쌍대성 정리를 사용하여 다음과 같이 쓴다.
\[
  H^{p}(X,\Osh_X(-m))'
  =\operatorname{Ext}_{\Osh_P}^{r-p}
    (P;\Osh_X(-m),\Omega_P^{r}).
\]
두 번째 변은 $p<n$이고 $m$이 충분히 클 때 0이다. 이는 보조정리 5에서
$X=P$로 놓고, $P$ 위의 연접 대수적 층으로서 $\Osh_X$가 코호몰로지
차원 $\leq r-n$을 갖는다는 사실($X$가 비특이이므로)을 사용하면 따른다.
따라서
\[
  \Extsh_{\Osh_P}^{r-p}(\Osh_X,\Omega_P^{r})=0
  \qquad\text{$p<n$일 때}.
\]
\fsection{6}{특이 다양체에 대한 쌍대성 정리}

$X$를 차원이 $r$인 사영공간 $P$의 차원이 $n$인 닫힌 부분집합이라
하자. 여기서는 공식 (5.6)을 다음과 같이 써야 한다.
\begin{equation}\tag{6.1}
  H^{n}(X,F)'
  \simeq\operatorname{Hom}_{\Osh_X}(X;F,\omega_X^{n})
  =\operatorname{Ext}_{\Osh_X}^{0}(X;F,\omega_X^{n}),
\end{equation}
여기서 다음과 같이 놓았다.
\begin{equation}\tag{6.2}
  \omega_X^{n}=E^{0}(\Osh_X)
  =\Extsh_{\Osh_P}^{r-n}(\Osh_X,\Omega_P^{r}).
\end{equation}
명제 5에서 말했듯이, 이렇게 도입한 층은 실제로 $X$를 비특이 다양체
$P$에 넣는 데 선택한 몰입 사상에 의존하지 않는다.
(6.1)에서 $F=\omega_X^{n}$로 놓으면 다음을 얻는다.
\begin{equation}\tag{6.2 bis}
  H^{n}(X,\omega_X^{n})'
  \simeq\operatorname{Hom}_{\Osh_X}(X;\omega_X^{n},\omega_X^{n}),
\end{equation}
\footnote{인쇄본도 이 공식에 ``(6.2)''라는 번호를 붙였다. 접미사
``bis''는 두 출현을 구별한다.} 따라서 다음에 대응하는
$H^{n}(X,\omega_X^{n})'$의 특별한 원소가 존재한다.
$\omega_X^{n}$에서 자기 자신으로 가는 항등 준동형이다:
\begin{equation}\tag{6.3}
  \gamma:H^{n}(X,\omega_X^{n})\longrightarrow k.
\end{equation}

\src{186}
이제 $\operatorname{Ext}$의 합성으로 정의되는 짝지음들을 생각하자:
\begin{equation}\tag{6.4}
  H^{p}(X,F)\times
  \operatorname{Ext}_{\Osh_X}^{n-p}(X;F,\omega_X^{n})
  \longrightarrow H^{n}(X,\omega_X^{n}),
\end{equation}
이를 (6.3)의 준동형 $\gamma$와 합성하면 경계 준동형과
양립하는 다음 함자적 준동형들을 얻는다:
\begin{equation}\tag{6.5}
  \operatorname{Ext}_{\Osh_X}^{n-p}(X;F,\omega_X^{n})
  \longrightarrow H^{p}(X,F)'
\end{equation}
(이는 (5.2)를 일반화한다). $p=n$일 때에는 이렇게 하여 (6.1)의
동형사상을 얻음을 확인할 수 있다. 이제 다음 정리가 성립한다.

\result{정리}{3 bis} 주어진 정수 $k\geq0$에 대하여 $X$에 관한
다음 네 조건은 서로 동치이다:
\begin{enumerate}
  \item[\textit{i.}] 함자적 준동형 (6.5)는 $n-k\leq p\leq n$일 때
    동형사상이다.
  \item[\textit{ii.}] 충분히 큰 $m$과 $n-k\leq p<n$에 대하여
    $H^{p}(X,\Osh_X(-m))=0$이다.
  \item[\textit{iii.}] $X$ 위의 연접 대수적 층의 범주에서 함자
    $H^{p}(X,F)$는 $n-k\leq p<n$일 때 여소거 가능하다.
  \item[\textit{iv.}] $0<i\leq k$에 대하여
    $E^{i}(\Osh_X)=\Extsh_{\Osh_P}^{r-n+i}(\Osh_X,\Omega_P^{r})=0$
    이다.
\end{enumerate}

\noindent\textbf{증명.} --- \textit{i.}$\Rightarrow$\textit{ii.}는
보조정리 4에서, \textit{ii.}$\Rightarrow$\textit{iii.}은 보조정리 3에서,
\textit{iii.}$\Rightarrow$\textit{i.}은 잘 알려진 표준 논증에서 따른다.
실제로 그때 (6.5)의 양변은 $n-k\leq p<n$에서 여소거 가능한 함자이고
(첫째는 보조정리 4에 따라 그러하기 때문이다). 끝으로
\textit{ii.}$\Leftrightarrow$\textit{iv.}를 증명하자. 이는 다음 명제의
따름정리에서 나온다:

\result{명제}{6} $F$를 $X$ 위의 연접 대수적 층이라 하고,
$i$를 정수라 하자. 그러면 충분히 큰 $m$에 대하여 다음 동형사상이 성립한다:
\begin{equation}\tag{6.6}
  H^{i}(X,F(-m))'
  \simeq H^{0}\bigl(X,E^{n-i}(F)(m)\bigr),
\end{equation}
여기서 다음과 같이 놓는다(n\textsuperscript{o} 3, 명제 5와 비교):
\begin{equation}\tag{6.7}
  E^{j}(F)=\Extsh_{\Osh_P}^{r-n+j}(F,\Omega_P^{r}).
\end{equation}
실제로 $P$에 대한 쌍대성 정리에 의해 (6.6)의 좌변은
$\operatorname{Ext}_{\Osh_P}^{r-i}(P;F(-m),\Omega_P^{r})$와 동형이므로,
(6.6)은 보조정리 5에서
따른다.

\src{187}
\result{따름정리}{} 충분히 큰 $m$에 대하여 $H^{i}(X,F(-m))=0$이
성립할 필요충분조건은 $E^{n-i}(F)=0$인 것이다.

$E^{j}(F)$들은 실제로 고려 중인 사영 몰입 사상에 의존하지 않음을 상기하자.
따름정리의 조건은 순전히 국소적이다. 따라서
이 조건이 $F$에 대해 성립하면 국소적으로
어떤 층 $F^{n}$과 동형인 모든 층에 대해서도 성립한다. 특히 이 조건이
$\Osh_X$에 대해 성립하면 모든 국소 자유 연접 대수적 층에 대해
성립한다. 예를 들어 $X$가 비특이면 모든 $i<n$에 대해 그러하고,
$X$의 어느 성분도 한 점만으로 이루어지지 않으면 $i=0$에 대해,
$X$가 정규이고 모든 성분의 차원이 $>1$이면 $i=0,1$에 대해 그러하다(참조
[3]).\footnote{인쇄본은 여기서 ``$S$가 정규이다''라고 되어 있지만, 이 단락은
다양체 $X$만을 다룬다.} 모든 $i<k$에 대해 그러할 필요충분조건은
정의에 따라 국소환
$\Osh_x$ ($x\in X$)들이 ``호몰로지적 여차원'' $\geq k$를 갖는 것이다
(이 개념은 [4] 참조). $k=n$이면 정리 3 bis에 의해 이는
쌍대성 정리가 $X$에 대해 성립한다는 것, 즉 (6.5)가
모든 $p$와 모든 $F$에 대해 동형사상이라는 뜻이다. 이에 대한 국소환
$\Osh_x$의 동치 조건을 많이 제시할 수 있는데(Nagata), 예를 들어
코언--매콜리의 등차원성 정리를 만족하는 조건이 그러하다. 이는
예를 들어 $X$가 국소적으로 비특이 주위 다양체 안의 ``완전교차''이면
성립한다.

\fsection{7}{푸앵카레 쌍대성}

$X$를 차원 $n$인 비특이 사영 다양체라 하자. 그러면
\[
  H^{**}(X)=H^{*}(X,\Omega_X^{*})
\]
는 유한 차원의 반가환 이중등급 대수이며, 우리는 여기에 전체 차수에
따른 등급을 준다. 따라서 $H^{p,q}(X)=H^{p}(X,\Omega_X^{q})$의 차수는
$p+q$이다. $H^{*}(X)$에서 나타나는 차수들은 $0$에서 $2n$까지이다.
정리 2와 정리 3의 따름정리에 따라 $H^{*}(X)$는 차원 $2n$인
“푸앵카레 대수”이다. 즉 $H^{2n}(X)$에는 기초체 $k$와의 동형사상이
주어져 있고, 곱
\[
  H^{m}(X)\times H^{2n-m}(X)\longrightarrow H^{2n}(X)=k
\]
은 $H^{m}(X)$와 $H^{2n-m}(X)$ 사이의 쌍대성을 이룬다. 또 $Y$가
다른 비특이 사영 다양체이면, 연접 대수적 층에 대한 퀸네트 공식은
다음을 준다.
\begin{equation}\tag{7.1}
  H^{*}(X\times Y)=H^{*}(X)\otimes H^{*}(Y),
\end{equation}
을 얻으며, 이 동형사상은 푸앵카레 대수 구조들과 양립한다. 또한
$H^{*}(X)$는 가환 대수로서 $X$에 관한 반변 함자이고, 사상
$f:Y\longrightarrow X$는 명백한 방식으로 등급 대수의 준동형

\src{188}
\begin{equation}\tag{7.2}
  f^{*}:H^{*}(X)\longrightarrow H^{*}(Y).
\end{equation}
을 정의한다. 이들이 푸앵카레 대수이므로, 전치를 취하여 다음 벡터
공간의 준동형을 얻는다:
\begin{equation}\tag{7.3}
  f_{*}:H^{*}(Y)\longrightarrow H^{*}(X).
\end{equation}
4절에서, 비특이 순환에 대응하는 코호몰로지류에 대한 $f^{*}$의
작용은 그 순환들의 역상에 대응하는 코호몰로지류를 취하는 것으로
기하학적으로 해석됨을 보았다. 여기서는 (7.3)이 이와 마찬가지로
순환의 “직상” 연산에 대응함을 보이는 것이 중요하다. 적어도 적절한
비특이성 조건 아래에서는, 이는 다음 특수한 경우로부터 따른다:

\result{정리}{4} $f$가 $X^{n}$ 안의 비특이 부분다양체 $Y^{m}$에서
$X^{n}$으로 가는 포함 사상이고 $1_Y$가 $H(Y)$의 단위원을 나타내면,
다음이 성립한다.
\begin{equation}\tag{7.4}
  f_{*}(1_Y)=P_X(Y),
\end{equation}
여기서 우변은 $Y$에 대응하는 $X$의 코호몰로지류이다.

이 식은 다음과 동치이다.
\begin{equation}\tag{7.4 bis}
  \bigl\langle\xi^{m,m},P_X(Y)\bigr\rangle\varepsilon_Y
  =f^{*}(\xi^{m,m})
  \qquad
  \bigl(\xi^{m,m}\in H^{m}(X,\Omega_X^{m})\bigr),
\end{equation}
\footnote{인쇄본에는 $f_{*}(\xi^{m,m})$라고 되어 있다. 그런데
$\xi^{m,m}\in H^{m}(X,\Omega_X^{m})$이고 양변이
$H^{m}(Y,\Omega_Y^{m})$에 속하므로 필요한 사상은 $f^{*}$이다.} 여기서
$\varepsilon_Y$는 $H^{m}(Y,\Omega_Y^{m})$의 기본 원소이다.
비특이 사영 다양체의 경우 이 식은 $Y$에 대응하는 코호몰로지류의 새로운
정의를 준다. 정리 4를 증명하기 위해, 정리 3을 이용하여 준동형
\[
  H^{m}(X,\Omega_X^{m})
  \longrightarrow H^{m}(Y,\Omega_Y^{m})
  =H^{m}(X,\Omega_Y^{m})
\]
의 전치를 다음 준동형으로 해석한다.
\begin{equation}\tag{7.5}
\begin{split}
  H^{n-m}(X,\Omega_X^{n-m})
  &\simeq
  \operatorname{Ext}_{\Osh_X}^{n-m}(X;\Omega_X^{m},\Omega_X^{n})\\
  &\longleftarrow
  \operatorname{Ext}_{\Osh_X}^{n-m}(X;\Omega_Y^{m},\Omega_X^{n})
  \simeq
  \operatorname{Hom}_{\Osh_X}(X;\Omega_Y^{m},\Omega_Y^{m}).
\end{split}
\end{equation}
$H^{m}(Y,\Omega_Y^{m})$의 쌍대에 속하는 원소 $1_Y$는 우변에서
$\Omega_Y^{m}$의 항등 자기준동형에 대응하는 원소와 동일시됨을
확인할 수 있다. 또한 이 원소의 $H^{n-m}(X,\Omega_X^{n-m})$에서의
상은 바로 $P_X(Y)$임도 확인할 수 있다. 이 양립성들은 이미 4절에서
제시할 수도 있었으며, 실제로 임의의 비특이 다양체에 대하여도
서술할 수 있고 성립한다. 그중 둘째 양립성은

\src{189}
예컨대 다음 표준 자기준동형들의 도표가 가환한다는 사실에서 따른다:
\begin{equation}\tag{7.6}
\resizebox{\textwidth}{!}{$
\left\{
\begin{array}{ccccc}
\operatorname{Ext}_{\Osh_X}^{n-m}(X;\Omega_X^{m},\Omega_X^{n})
&\longleftarrow&
\operatorname{Ext}_{\Osh_X}^{n-m}(X;\Omega_Y^{m},\Omega_X^{n})
&\simeq&
\operatorname{Hom}_{\Osh_X}(X;\Omega_Y^{m},\Omega_Y^{m})\\[3pt]
\Big\downarrow&&&&\Big\downarrow\\[3pt]
\operatorname{Ext}_{\Osh_X}^{n-m}(X;\Osh_X,\Omega_X^{n-m})
&\longleftarrow&
\operatorname{Ext}_{\Osh_X}^{n-m}(X;\Osh_Y,\Omega_X^{n-m})
&\simeq&
\operatorname{Hom}_{\Osh_X}(X;\Omega_X^{m},\Omega_Y^{m})
\end{array}
\right.
$}
\end{equation}

이로써 방향이 주어진 콤팩트 다양체에서의 푸앵카레 쌍대성 형식론에
대한 정확한 대응물을 얻는다. 특히 정리 4로부터 $X\times X$의
대각선에 대응하는 코호몰로지류를 결정할 수 있다. 잘 알려진
논법으로부터, 예컨대 다음 렙셰츠 공식을 얻는다:

\result{정리}{5} $f$를, 그 모든 고정점의 중복도가 1인 비특이 사영
다양체 $X$의 자기준동형이라 하자. 그러면 이러한 고정점들의 개수는
$k$의 표수를 법으로,
$f$가 $H^{i}(X)$에 유도하는 자기준동형들의 대각합의 교대합과
합동이다.

$f$에 둔 이 제한은 4절의 비고에서 언급한 어려움 때문이다.
그러나 다음 사실에 유의하자. $f$가 모든 고정점의 중복도가 1인
자기준동형과 “호모토픽”이면 렙셰츠 공식은
여전히 성립한다.

\fsection{8}{쌍대성 정리의 일반화}

$X$를 다음 성질을 갖는 비특이 대수다양체라 하자. $X$ 위의 모든
연접인 대수적 층 $F$는 국소 자유인 어떤 연접인 대수적 층의 몫층과
동형이다(이는 $X$가 어떤 사영공간 안에서 국소 닫혀 있을 때 성립한다).
그러면 $X$ 위의 모든 연접인 대수적 층 $F$는 국소 자유층들로 이루어진
유한 분해 $L$을 갖는다. 또한 그러한 두 분해가 주어지면 언제나 세 번째
분해와, 증대 준동형들과 양립하는 그 세 번째 분해에서 앞의 두 분해로
가는 준동형들을 찾을 수 있다. 마찬가지로 $L$이 $F$의 유한 국소 자유
분해이고 준동형 $F'\longrightarrow F$가 주어졌다면, $F'$의 유한 국소
자유 분해 $L'$과 $F'\longrightarrow F$에 양립하는 준동형
$L'\longrightarrow L$이 존재한다. $F'\longrightarrow F$가 전사이면 이
준동형도 전사라고 가정할 수 있다. 따라서 두 정수 $r,s\geq0$가 주어지면,
$A_1,\ldots,A_r$; $B_1,\ldots,B_s$를 인수로 하는 두 코호몰로지적
다중함자를 정의할 수 있다.

\src{190}
이 인수들은 $X$ 위의 연접인 대수적 층들의 범주에서 변하며, 두 함자의
값은 각각 $X$ 위의 연접인 대수적 층들의 범주와
$H^{0}(X,\Osh_X)$ 위의 가군들의 범주에 속한다. 그 정의식은 다음과 같다.
\begin{equation}\tag{8.1}
\resizebox{\textwidth}{!}{$
\left\{
\begin{aligned}
\underline{T}_{r}^{s*}(A_1,\ldots,A_r;B_1,\ldots,B_s)
  &={}H^{*}\!\left(
    \Homsh_{\Osh_X}\!\left(
      L(A_1)\otimes\cdots\otimes L(A_r),
      L(B_1)\otimes\cdots\otimes L(B_s)
    \right)\right),\\[3pt]
T_{r}^{s*}(A_1,\ldots,A_r;B_1,\ldots,B_s)
  &={}\mathbf R^{*}\Gamma\!\left(
    \Homsh_{\Osh_X}\!\left(
      L(A_1)\otimes\cdots\otimes L(A_r),
      L(B_1)\otimes\cdots\otimes L(B_s)
    \right)\right).
\end{aligned}
\right.
$}
\end{equation}
이 식에서 $L(F)$는 연접인 대수적 층 $F$의 유한 국소 자유 분해를
나타내고, $\mathbf R^{*}\Gamma(K)$는 층 복합체 $K$에 관한 공간 $X$의
초코호몰로지를 나타낸다. $r$ 또는 $s$가 0이면 $L(A_i)$들의 텐서곱,
각각 $L(B_j)$들의 텐서곱을 $\Osh_X$로 바꾼다. 특히
$\underline{T}_{0}^{0}$와 $T_{0}^{0}$는 인수가 없는 등급함자이다.
$\underline{T}_{0}^{0}$는 0차에 집중되어 있고 그 차수에서 층
$\Osh_X$이며, $T_{0}^{0}$는 $H^{*}(X,\Osh_X)$와 같다. (8.1)의 우변들이
선택한 분해에 의존하지 않는다는 사실은 첫째 줄에 대해서는 어쨌든
명백하다(이 경우 문제는 국소적이다). 둘째 줄에 대해서는 앞의 일반적
고찰과 다음 층 복합체의 초코호몰로지 스펙트럼 열을 고려하면 이 사실을
얻는다.
\[
  K=\Homsh_{\Osh_X}\bigl(L(A_1)\otimes\cdots,
      L(B_1)\otimes\cdots\bigr),
\]
이 스펙트럼 열은 $K$에 관한 $X$의 초코호몰로지로 수렴하며, 그 초기항은
$H^{p}(X,H^{q}(K))$, 즉 다음과 같다.
\begin{equation}\tag{8.2}
  E_{2}^{p,q}
  =H^{p}\!\left(X,
    (\underline{T}_{r}^{s})^{(q)}
      (A_1,\ldots,A_r;B_1,\ldots,B_s)\right).
\end{equation}
그러므로 이 스펙트럼 열 자체도 선택한 분해에 의존하지 않으며, 그
수렴대상은 (8.1)의 좌변이다. 여러 인수 $A_i$, $B_j$에 관한 연결
연산자들도 쉽게 정의된다. 실제로 모든 완전열
$0\longrightarrow F'\longrightarrow F\longrightarrow F''\longrightarrow0$은
유한 국소 자유 복합체들로 이루어진 완전열로 분해할 수 있다.

함자계 $\underline{T}_{r}^{s*}$, 각각 $T_{r}^{s*}$에는 텐서 계산의
연산들과 비슷한 연산들을 정의한다. 그 정의는 정의식 (8.1)에서 곧바로
나온다. 이를테면 제2절에서 다룬 것을 일반화하는 다음 합성이 있다.
\begin{equation}\tag{8.3}
\resizebox{\textwidth}{!}{$
\begin{aligned}
&T_{r}^{s*}(A_1,\ldots,A_r;B_1,\ldots,B_s)
\times T_{r'}^{s'*}(A'_1,\ldots,A'_{r'};B'_1,\ldots,B'_{s'})\\
&\qquad\longrightarrow
T_{r+r'}^{s+s'*}
(A_1,\ldots,A_r,A'_1,\ldots,A'_{r'};
 B_1,\ldots,B_s,B'_1,\ldots,B'_{s'}),
\end{aligned}
$}
\end{equation}
이 합성은 명백한 결합법칙을 만족하며 함자적 준동형, 연결 준동형,
스펙트럼 열과 양립한다. 대칭 연산도 마찬가지로 있으며, 이를 명시하는
일은 독자에게 맡긴다. 또한 한 인수 $A_i$가 한 인수 $B_j$와 같을 때마다
축약 연산이 있다.

\src{191}
\begin{equation}\tag{8.4}
\resizebox{\textwidth}{!}{$
\begin{aligned}
&T_{r}^{s*}(A_1,\ldots,A_{i-1},C,A_{i+1},\ldots,A_r;
 B_1,\ldots,B_{j-1},C,B_{j+1},\ldots,B_s)\\
&\qquad\longrightarrow
T_{r-1}^{s-1*}(A_1,\ldots,\widehat{A_i},\ldots,A_r;
 B_1,\ldots,\widehat{B_j},\ldots,B_s).
\end{aligned}
$}
\end{equation}
더 나아가 한 인수 $A_i$가 국소 자유층이면, 한 $A_j$ ($j\neq i$)를
$A_j\otimes A_i$로 바꾸거나 한 $B_k$를 $B_k\otimes A_i'$로 바꾸는 대신
$A_i$를 없앨 수 있다. 여기서
$A_i'=\Homsh_{\Osh_X}(A_i,\Osh_X)$이다. 한 인수 $B_i$가 국소 자유층인
경우에도 이와 비슷한 계산 규칙이 있다. 특히 $\Osh_X$와 같은 인수는
언제나 없앨 수 있다. 모든 인수가 국소 자유이고 $B_i$들 가운데 많아야
하나만 예외라면, 방금 말한 규칙은 다음 함자적 동형을 준다.
\begin{equation}\tag{8.5}
  T_{r}^{s*}(A_1,\ldots,A_r;B_1,\ldots,B_s)
  =H^{*}\bigl(X,A_1'\otimes\cdots\otimes A_r'
    \otimes B_1\otimes\cdots\otimes B_s\bigr)
\end{equation}
(실제로 이는 $r=0$, $s=1$인 경우로 귀착되며 그 경우에는 명백하다.
초기항이 (8.2)인 스펙트럼 열을 직접 사용해도 된다.) 앞의 연산들에
대응하는 연산들을 $\underline{T}_{r}^{s}$에 대해서도 정의할 수 있다.
이렇게 도입한 여러 연산 사이의 관계는 텐서 계산의 대응하는 연산 사이의
관계와 같다.

$X$의 차원을 $n$이라 하자. 반복되는 인자들에 텐서 합성 (8.3)과
축약 (8.4)을 차례로 적용하면 다음 짝짓기를 얻는다.
\begin{equation}\tag{8.6}
\resizebox{\textwidth}{!}{$
  (T_{r}^{s})^{p}(A_1,\ldots;B_1,\ldots)
  \times
  (T_{s}^{r})^{n-p}(B_1,\ldots;A_1,\ldots,A_r\otimes\Omega_X^{n})
  \longrightarrow H^{n}(X,\Omega_X^{n}).
$}
\end{equation}

\result{정리}{6} $X$가 비특이 사영 대수다양체이면, 짝짓기 (8.6)은
쌍대성을 이룬다.

이는 제3정리의 따름정리로부터 순전히 형식적으로 나온다. 실제로 그
따름정리에 따르면, $K$가 국소 자유인 연접 대수적 층들의 복합체일 때
복합체 $K$에 대한 $X$의 초코호몰로지는 자연스러운 짝짓기
\begin{equation}\tag{8.7}
  \mathbf R^{p}\Gamma(K)
  \times\mathbf R^{n-p}\Gamma(K'\otimes\Omega_X^{n})
  \longrightarrow\mathbf R^{n}\Gamma(\Omega_X^{n})
  =H^{n}(X,\Omega_X^{n})
\end{equation}
를 통하여 복합체 $K'\otimes\Omega_X^{n}$에 대한 $X$의 초코호몰로지와
쌍대를 이룸을 쉽게 알 수 있다. 이는 첫 항이
$H^{p}(X,H^{q}(K))$인 스펙트럼 열\footnote{인쇄본에는
$H^{p}(H^{q}(X,K))$라고 되어 있으나, 첫 항은 스펙트럼 열 (8.2)의
것, 곧 $H^{p}(X,H^{q}(K))$이다.}과
$K'\otimes\Omega_X^{n}$에 대한 유사한 스펙트럼 열을 사용하면 알 수
있다. 앞의 결과에서 정의 (8.1)을 사용하면 정리 6이 따른다.

\src{192}
\subsection*{비고}

\noindent 1\textsuperscript{o} 정리 6 앞의 정의들에서는 $X$가
비특이일 필요가 없었다. 유한 분해만을 다루는 것이 필수적이지 않았기
때문이다. 그러나 $X$가 특이이면
$(\underline{T}_{r}^{s})^{p}(A_1,\ldots;B_1,\ldots)$이 연접층이라고
미리 단언할 수는 없다. 실제로 층들의 복합체
\[
  \Homsh_{\Osh_X}\bigl(L(A_1)\otimes\cdots,L(B_1)\otimes\cdots\bigr)
\]
에는 주어진 전체 차수마다 무한히 많은 성분이 생기기 때문이다.

\medskip
\noindent 2\textsuperscript{o} 식 (8.1)에서 $L(B_i)$ 가운데 하나를
$B_i$로 바꾸어도 됨은 쉽게 확인된다. 따라서 명제 3을 함께 고려하면
다음을 얻는다.
\begin{equation}\tag{8.8}
\left\{
\begin{aligned}
  \underline{T}_{1}^{1*}(A;B)&=\Extsh_{\Osh_X}^{*}(A,B),\\
  T_{1}^{1*}(A;B)&=\operatorname{Ext}_{\Osh_X}^{*}(X;A,B).
\end{aligned}
\right.
\end{equation}
특히 (8.6)에서 $r=s=1$, $A_1=\Osh_X$로 놓으면 정리 3을 다시 얻는다.
식 (8.8)은 또한 다음을 함의한다.
\[
  T_{0}^{1*}(B)=H^{*}(X,B),
  \qquad
  T_{1}^{0*}(A)=\operatorname{Ext}_{\Osh_X}^{*}(X;A,\Osh_X).
\]

\medskip
\noindent 3\textsuperscript{o} 식 (8.1)을 보면 일반적으로 함자
$\underline{T}_{r}^{s*}$와 $T_{r}^{s*}$는 양의 차수와 음의 차수의
성분을 모두 가진다. 비고 2를 사용하면 $X$의 차원이 $n$일 때
$\underline{T}_{r}^{s*}$의 영이 아닌 성분들은 $s>0$이면
$-(s-1)n$과 $rn$ 사이에, $s=0$이면 $0$과 $rn$ 사이에 있음을 알 수
있다. 한편 $T_{r}^{s*}$의 영이 아닌 성분들은 $s>0$이면
$-(s-1)n$과 $(r+1)n$ 사이에, $s=0$이면 $0$과 $(r+1)n$ 사이에 있다
(그리고 내가 잘못 보지 않았다면, $r>0$일 때에는 각각
$-(s-1)n$과 $rn$ 사이, $0$과 $rn$ 사이에 있다).

\section*{참고문헌}

\noindent[1] Cartier (Pierre). --- Les groupes
$\operatorname{Ext}^{s}(A,B)$, Séminaire A. Grothendieck : Algèbre homologique,
제1권, 1957, 제3호.

\noindent[2] Grothendieck (Alexandre). --- Sur quelques points d'algèbre
homologique, \emph{Tohoku Math. J.}, 제9권, 1957, 119--183쪽.

\noindent[3] Serre (Jean-Pierre). --- Faisceaux algébriques cohérents,
\emph{Annals of Math.}, 제61권, 1955, 197--278쪽.

\noindent[4] Serre (Jean-Pierre). --- Sur la dimension homologique des anneaux
et des modules noethériens, \emph{Proc. Intern. Symp. on Alg. Number Theory}
[1955, Tokyo et Nikko]. --- Tokyo, Science Council of Japan, 1956;
175--189쪽.

\clearpage
\src{193}
\section*{추기}

13쪽의 비고에서 지적한 어려움들은 쌍대성 정리를 임의의 대수다양체
(임의의 특이점을 가진 대수다양체)로 확장함으로써 이제 완전히
해결되었다. ``스킴 이론''의 일반적 틀에서 정식화할 수 있는 이 이론에
관해서는 현재 준비 중인 J.~Dieudonné와 A.~Grothendieck의
``Éléments de Géométrie algébrique''를 참조한다. 몇 가지 실마리는 다음
글에서 찾을 수 있다.

\begin{quote}
Grothendieck (A.). --- Cohomology theory of algebraic varieties,
\emph{Congrès int. des Mathématiciens}, 1958, Edinburgh (출간 예정).
\end{quote}

\begin{flushright}
[1959년 4월]
\end{flushright}

\begin{center}\rule{36mm}{0.5pt}\end{center}

\par\smallskip
\begingroup\footnotesize




\noindent\textit{명제 2의 인쇄본은 하나의 유한 표시만으로 모든 차수의 비교가 동형이라고 주장한다. 차수 $p$의 비교에는 차수 $p+1$까지 유한 자유 가군들로 이루어진 부분 분해가 필요하며, 모든 차수에 대해서는 의사연접성으로 충분하다.}\par
\endgroup

\par
\renewcommand{\sourceunit}{e149}
\clearpage
\chapter*{강연 149의 정오표}
\addcontentsline{toc}{chapter}{강연 149의 정오표}
\setcounter{footnote}{0}
\src{298}

\noindent\textbf{1. 연접인 대수적 층에 대한 쌍대성 정리.}

\begin{center}
[부르바키 세미나, 제9권, 1956/57, 제149호, 25쪽.]
\end{center}

\noindent\textbf{14쪽에 관한 주.} --- 내가 1958년 국제수학자대회에서 한
강연에서 지적했듯이,\footnote{
Grothendieck (Alexander). --- The cohomology theory of abstract algebraic
varieties, \emph{Proceedings of the International Congress of Mathematicians}
[1958, Edinburgh], p. 103--118. --- Cambridge, at the University Press, 1960.}
여기서 제기된 문제들은 이제 완전히 해결되었다.

독자는 연접층의 쌍대성에 관한 더 자세한 내용을
\emph{loco citato}\footnotemark[1], 112--115쪽, EGA III 제2부(준비 중),
그리고 그때까지는 SGA 1962에서 찾을 수 있다. 더 체계적인 논의는 EGA의
뒤 이은 한 장(계획상 제IX장)에 실릴 것이다.

\par
\renewcommand{\sourceunit}{182}
\unit{강연 182 — 형식기하학과 대수기하학}
\sid{182}
\src{193}
\seminar{(1959년 5월)}
\paperhead{형식기하학과 대수기하학}%
  {알렉상드르 그로텐디크 지음}

\fsection{1}{스킴}

체 $k$ 위에 정의된 아핀 대수공간은 그 아핀 대수 $A$($k$ 위에 정의된
정칙함수들의 환)에 의해 본질적으로 결정된다는 것은 잘 알려져 있다.
대수공간의 사상 $X\longrightarrow Y$는 $k$-대수 준동형
$A(Y)\longrightarrow A(X)$와 일대일로 대응한다. 한 대수공간에 대응하는
아핀 대수는 유한 생성 $k$-대수이며, ``고전적'' 관점에서는 멱영원을 갖지
않는다. 역으로 이러한 유형의 모든 대수는 $k$ 위에 정의된 한 대수공간의
아핀 대수로 얻어진다. 따라서 아핀 대수공간에 관한 상황을 가환대수의
말로 해석하게 하는 잘 알려진 사전이 있다. 그렇게 하면 더 일반적인
명제들을 얻게 된다는 사실은 오래전부터 알려져 있었다. 실제로 문제의
환들이 방금 살펴본 유형이라고 가정할 필요가 대개 없어졌으며, 흔히 뇌터
가정이면 충분했기 때문이다. 특히 주어진 바탕체가 있든 없든, 이 환들이
멱영원을 포함하는 경우를 배제할 이유는 없었다. 지금까지 기하학자들은
이러한 암시를 받아들이기를 거부하고 멱영원을 갖지 않는 아핀 대수, 즉
구조층에 멱영원이 없는 대수공간(더구나 대개는 ``절대 기약''인
대수공간)만을 고집해 왔다. 강연자는 이러한 사고방식이 대수기하학에서
참으로 자연스러운 방법들의 발전에 심각한 장애가 되었다고 생각한다.

$A$를 가환환이라 하자. $A$의 소아이디얼들의 집합
$X=\operatorname{Spec}(A)$에는 자연스러운 위상, 곧 ``자리스키 위상''
또는 스펙트럼 위상이 주어진다는 것은 잘 알려져 있다. 한편 $X$ 위에는
가환환의 층 $\Osh_X$가 있으며, $\mathfrak p\in X$에서 그 줄기는
국소화환 $A_{\mathfrak p}$이고 그 단면환은 $A$와 동일시된다. 이로써
$X$는 \emph{환 달린 공간}이 되며, 이를 $A$의 \emph{소 스펙트럼}이라
한다. 환 준동형 $f:A\longrightarrow B$는 환 달린 공간의 사상
$f':\operatorname{Spec}(B)\longrightarrow\operatorname{Spec}(A)$를
정의하며, 그 바탕 집합 사상은 다름 아닌
$\mathfrak p\longmapsto f^{-1}(\mathfrak p)$이다. 이런 방식으로 얻어지는

\src{194}
$\operatorname{Spec}(B)$에서 $\operatorname{Spec}(A)$로 가는 환 달린
공간의 사상들은 준동형 $\Osh_x\longrightarrow\Osh_y$ ($x=f'(y)$)가
국소 준동형인 것들, 그리고 오직 그런 것들뿐이다(즉 극대 아이디얼의
역상은 극대 아이디얼이다).

$\operatorname{Spec}(A)$와 동형인 환 달린 공간을 \emph{아핀 스킴}이라 하고,
국소적으로 아핀인 환 달린 공간, 즉 모든 점이 유도된 구조에 관하여 아핀
스킴인 열린 근방을 갖는 환 달린 공간을 \emph{준스킴}이라 한다. 준스킴의
\emph{사상}은 자명한 방식으로 정의하며, 국소적으로 환 준동형에 대응한다.

준스킴 $S$를 고정하고 준스킴의 사상 $X\longrightarrow S$를 생각하면,
$S$는 기초체나 기초환(더 정확히는 섬유화의 밑공간)의 역할을 한다. 이때
$X$를 $S$-준스킴이라 한다. $S=\operatorname{Spec}(A)$이면 이는
$\Osh_X$가 $A$-대수의 층이라는 뜻이기도 하다. 따라서 모든 준스킴은
유일한 방식으로 $\mathbf Z$-준스킴으로 볼 수 있다. 물론 $S$-준스킴들은
범주를 이루며, 더 나아가 이 범주에서 두 대상 $X,Y$의 곱은 항상 존재하고
$X\times_S Y$로 나타낸다. 이 곱 개념으로 사상
$S'\longrightarrow S$에 대응하는 $S$-준스킴의 밑준스킴 변경을 정의할
수 있다. 실제로 $X\times_S S'$은 $S'$-준스킴으로 볼 수 있다.

$X\times_S X$의 대각선이 닫혀 있으면 $X$가 $S$ \emph{위에서 분리}되어
있다고 한다. $\mathbf Z$ 위에서 분리된 준스킴을 \emph{스킴}이라 한다.
그러면 이는 어떤 밑 대상 위에서도 분리되어 있다. 단순히 하기 위해 이제부터
스킴만을 다루고, 더 나아가 이들이 \emph{뇌터}라고 가정한다. 즉 이들은
뇌터 환의 스펙트럼인 아핀 열린집합들의 유한 합집합이다. $S$의 모든 아핀
열린집합 $U$에 대하여 $X$에서의 그 역상이 유한 개의 아핀 열린집합의
합집합이고, 이 열린집합들의 환이 $U$의 환 위의 유한 생성 대수이면 $X$가
$S$ \emph{위에서 유한형}이라고 한다. 이러한 $S$-스킴들이 본래 의미의
기하학적 연구에 알맞다. 특히 모든 $s\in S$에 대하여 $s$ 위에서 $X$의
섬유 $f^{-1}(s)$는 $S$에서 $s$의 국소환 $\Osh_s$의 잉여체
$\kappa(s)$ 위의 대수적 스킴이다. 그러므로 $X$는 어느 정도는 매개변수
$s$가 $S$를 움직이는 ``대수공간'' $f^{-1}(s)$들의 족으로 볼 수 있다.
(국소적 관점에서 $s$는 주어진 환의 소아이디얼들을 움직인다.) 물론
$\kappa(s)$들은 서로 다른 표수를 가질 수 있다. $S=\operatorname{Spec}(k)$이고
$k$가 체이면, 이제는 구조층이 멱영원을 가질 수 있다는 차이만 빼고 통상적인
``대수공간'' 개념을 본질적으로 되찾는다.

\src{195}
잘 알려진 개념들에서 착안하여 \emph{사영 사상}, 더 일반적으로는
\emph{고유 사상}의 개념을 정의한다. 이러한 사상은 유한형이다. 또한 닫힌
부분집합을 닫힌 부분집합으로 보내며, 임의의 밑준스킴 변경 뒤에도 이 성질을
유지한다.

$X$가 스킴(언제나 그렇듯 뇌터 스킴)이면, 층 $\Osh_X$는 [2]의 의미에서
\emph{연접인 환의 층}이다. 따라서 $X$ 위의 연접 가군층은 국소적으로 어떤
준동형 $\Osh_X^{m}\longrightarrow\Osh_X^{n}$의 여핵과 동형인 층들이기도
하다.

\fsection{2}{형식 스킴}

$X$를 스킴, $X'$을 $X$의 닫힌 부분집합이라 하자. 그러면
$X'=\operatorname{Supp}(\Osh_X/J)$인 $\Osh_X$의 연접 부분층 $J$가
존재한다(더구나 그러한 것 가운데 가장 큰 것도 존재한다). $X'$에
$\Osh_X/J$를 갖추면 $X_0$로 나타내는 스킴이 되며, 이러한 스킴을
``$X$의 닫힌 부분스킴''이라 한다. 또한 모든 $n$에 대하여 $X'$에
$\Osh_X/J^{n+1}$를 갖춘 것을 생각하고 $X_n$으로 나타낼 수 있다. 이는
바탕 집합은 여전히 $X'$이지만 구조층
$\Osh_{X_n}=\Osh_X/J^{n+1}$이 다른 $X$의 닫힌 부분준스킴이다. 분명히
$\Osh_{X_n}$들은 $X'$ 위의 환의 층들의 사영계를 이루며, 그 사영극한
$\overline{\Osh}_X$를 \emph{$X'$을 따른 $\Osh_X$의 형식적 완비화}라
한다. 이 환의 층을 갖춘 $X'$을 \emph{$X'$을 따른 $X$의 형식적
완비화}라 한다. 따라서 이는 환 달린 공간이지만 일반적으로 스킴은 아니다.
$X$ 위의 모든 연접층 $F$에 대하여 마찬가지로 $X'$을 따른 $F$의 완비화
\[
\overline F=\varprojlim_n F_n,
\qquad F_n=F\otimes_{\Osh_X}\Osh_X/J^{n+1}
\]
를 생각할 수 있다. 이는 $\overline X$ 위의 가군층이다. 그 단면들을
``$X'$을 따른 $F$의 형식적 단면''이라 하며,
$\varprojlim_n\Gamma(X',F_n)$의 원소들과 동일시한다. $F=\Osh_X$이면
자리스키의 의미에서 ``$X'$을 따른 $X$의 정칙함수''를 얻는다. 다만 이
용어는 고전적 용어와 충돌하므로 여기서는 따르지 않겠다.

\emph{형식 스킴}(뇌터라고 이해한다)이란 다음 조건을 만족하는 위상환의 층
$\Osh_{\mathfrak X}$가 주어진 위상공간 $\mathfrak X$를 말한다. 위상환의
층들의 동형
\[
\Osh_{\mathfrak X}=\varprojlim_n\Osh_n
\]
이 있고, 여기서 $\Osh_n$들은 $\mathfrak X$ 위의 환층들의 사영계를
이루어 각각 $\mathfrak X$를 스킴 $\mathfrak X_n$으로 만든다. 또한
$m\geq n$일 때 준동형 $\Osh_m\longrightarrow\Osh_n$은 전사이고 그 핵은
$J_m^{n+1}$이며, 여기서 $J_m$은

\src{196}
$\Osh_m\longrightarrow\Osh_0$의 핵이다. 그러면 $\Osh_{\mathfrak X}$가
연접이고 줄기들이 뇌터 국소환인 환의 층임을 보일 수 있다.

정의에 따라 앞에서와 같은 형식적 완비화 $\overline X$는 형식 스킴이고,
역으로 모든 형식 스킴은 \emph{국소적으로} 이 유형이다. 실제로 \emph{아핀}
형식 스킴(즉 $\mathfrak X_0$가 아핀이며, 따라서 모든 $\mathfrak X_n$이
아핀인 형식 스킴)을 주는 것은 $J$-아딕이고 분리·완비인 뇌터 위상환을
주는 것과 동치이다.

통상적인 스킴에 대한 사상, 유한형 사상, 고유 사상 등의 정의는 어려움
없이 형식 스킴으로 확장된다.

\fsection{3}{세 기본정리}

$f:X\longrightarrow Y$를 스킴들의 고유 사상(언제나 그렇듯 뇌터
스킴들의 사상), $Y'$을 $Y$의 닫힌 부분집합, $X'$을 $X$에서의 그
역상이라 하고, 이에 대응하는 형식적 완비화 $\overline Y$와
$\overline X$를 생각하자. 그러면 $f$는 형식 스킴들의 사상
$\overline f:\overline X\longrightarrow\overline Y$를 유도하며, 이
사상도 고유이다. $F$가 $X$ 위의 연접층이면 $\overline F$는
$\overline X$ 위의 연접층이다. 정리 1에서는 $X,Y,F$를 잊고 형식
스킴들의 고유 사상 $\overline f$와 $\overline X$ 위의 연접층
$\overline F$만을 본다. (다만 강연자가 완전한 증명을 써 둔 것은
$X,Y,f,F$에서 출발하는 경우뿐이다.)

\result{정리}{1 (유한성 정리)}
\begin{enumerate}
  \item[\textit{i.}] $R^{q}\overline f_{*}(\overline F)$는
    $\overline Y$ 위의 연접층이다.
  \item[\textit{ii.}] 자연스러운 준동형
    \[
      R^{q}\overline f_{*}(\overline F)
      \longrightarrow \varprojlim_n R^{q}f_{n*}(F_n)
    \]
    은 동형이다.
\end{enumerate}

이 명제에서는 $Y'$을 정의하는 $\Osh_Y$의 연접 부분층 $J$를 하나
선택했다고 가정한다. 그 역상으로 $X'$을 정의하는 $\Osh_X$의 연접
부분층을 얻고, 따라서 제2절에서처럼 $F_n,X_n,Y_n$과
$f_n:X_n\longrightarrow Y_n$을 정의한다. 임의의 두 형식 스킴 사이의
고유 사상에서 출발했다면 이 표기들을 명시하는 데 필요한 사소한 변경은
자명하다.

정리 1은 ``형식 코호몰로지''만을 다루었다. 다음 정리는

\src{197}
이를 ``대수적 코호몰로지''와 관련짓고, 대수적 코호몰로지와 해석적
코호몰로지의 비교에 관한 세르의 잘 알려진 정리 [4]와 유사하다.

\result{정리}{2 (제1비교정리)} $R^{q}f_{*}(F)$는 $Y$ 위의
연접층이며(이는 정리 1의 특수한 경우이다), 자연스러운 준동형
\[
  \overline{R^{q}f_{*}(F)}
  \longrightarrow \varprojlim_n R^{q}f_{n*}(F_n)
\]
은 동형이다.

\result{따름정리}{1} 다음 표준 동형들이 있다.
\[
  \overline{R^{q}f_{*}(F)}=R^{q}\overline f_{*}(\overline F).
\]

$q=0$일 때 이 따름정리는 자리스키의 ``정칙함수의 기본정리''를 일반화한
것이며, 우리는 여기서 자리스키의 ``연결성 정리''의 일반화를 이끌어 낼
것이다. 한편 $q=0$일 때 정리 1의 \textit{(ii)}는 자명하지만 정리 2나
그와 동치인 정식화인 따름정리 1은 전혀 그렇지 않음에 유의하자. 실제로
증명은 $q$에 대한 내림 귀납법으로 진행된다. $q$가 크면 양변이 0이므로
자명하고, 따라서 $q=0$인 경우는 귀납법의 마지막 단계, 말하자면 ``가장
어려운'' 경우로 나타난다.

\result{따름정리}{2} $Y=\operatorname{Spec}(A)$이고 $Y'$은 $A$의 아이디얼
$J$로 정의된다고 하자. 그러면 $X$ 위의 모든 연접층 $F$에 대하여
$H^{q}(X,F)$는 유한 생성 $A$-가군이며, 그 $J$-아딕 완비화는
$H^{q}(\overline X,\overline F)$이다.

끝으로 이 따름정리를
$H=\Homsh_{\Osh_X}(F,G)$에 적용하면 다음을 얻는다.

\result{따름정리}{3} $Y=\operatorname{Spec}(A)$이고 $Y'$은 $A$의 아이디얼
$J$로 정의된다고 하자. $F,G$가 $X$ 위의 두 연접층이면
$\operatorname{Hom}(F,G)$는 유한 생성 $A$-가군이고, 그 $J$-아딕
완비화는 $\operatorname{Hom}(\overline F,\overline G)$와 동일시된다.

물론 자연스러운 사상
$\operatorname{Hom}(F,G)\longrightarrow
\operatorname{Hom}(\overline F,\overline G)$은 준동형
$u:F\longrightarrow G$에 그 ``연속성에 의한'' 연장
$\overline u:\overline F\longrightarrow\overline G$를 대응시키는
사상이다(이 규약 아래에서 $\overline F$는 $F$에 관한 함자가 된다).

이제 $A$가 그 $J$-아딕 위상에 관하여 분리·완비라고 하자. 그러면 앞의
따름정리 2와 3에서 다음을 얻는다.
\[
  H^{q}(X,F)=H^{q}(\overline X,\overline F),
  \qquad
  \operatorname{Hom}(F,G)
  =\operatorname{Hom}(\overline F,\overline G).
\]

\src{198}
이 마지막 항등식은 $X$ 위의 연접층의 범주가 $\overline X$ 위의 연접층의
범주의 한 \emph{부분범주}(그 사상들은 유도된 사상들이다)와 동일시됨을
보여 준다. 사실 다음의 더 강한 명제가 성립한다.

\result{정리}{3} $\overline X$ 위의 가군층이 연접일 필요충분조건은
$X$ 위의 어떤 연접층 $F$에 대하여 $\overline F$ 꼴의 층과 동형인
것이다(정리 2의 따름정리 3에 따라 $F$는 정준 동형까지 정해진다).
[이제 $Y=\operatorname{Spec}(A)$이며, $A$는 $J$-아딕이고 분리·완비인
뇌터 위상환임을 상기하자.]

\result{따름정리}{1} $X$의 닫힌 부분스킴들은 $\overline X$의 닫힌 형식적
부분스킴들과 일대일로 대응한다.

실제로 이들은 각각 $\Osh_X$와 $\overline{\Osh}_X$의 연접 부분층들과
대응한다. 사상의 그래프를 닫힌 부분스킴으로 보면, 따름정리 1에서 다음을
얻는다.

\result{따름정리}{2} $X,Z$가 $A$ 위에서 고유인 두 스킴이라고 하자
($A$는 $J$-아딕이고 분리·완비인 뇌터 환이다). 그러면
$g\longmapsto\overline g$는 $X$에서 $Z$로 가는 $Y$-사상들과
$\overline X$에서 $\overline Z$로 가는 $\overline Y$-사상들 사이의
일대일 대응을 정의한다.

다시 말해 $A$ 위에서 고유인 대수적 스킴들은 $\overline Y$ 위에서 고유인
형식적 스킴들의 범주의 한 \emph{부분범주}(그 사상들은 유도된 사상들이다)로
나타난다. 그러나 $\overline Y$ 위에서 고유이면서 ``대수화 가능''하지
않은 형식적 스킴, 즉 $A$ 위에서 고유인 어떤 $X$에 대해서도 $\overline X$와
동형이 아닌 형식적 스킴이 존재한다는 점에 주의해야 한다(복소수체 위에
정의된 대수다양체에서 오지 않는 콤팩트 복소 해석다양체가 존재하는 것과
마찬가지이다). 이러한 형식적 스킴은 ``모듈라이 이론''에서 자연스럽게
나타난다. 다만 형식적 스킴이 대수화 가능한 다음의 흥미로운 특수한 경우를
기록해 두자.

\result{정리}{4} $A$를 잉여체 $k$를 갖는 완비 뇌터 국소환이라 하고,
$\mathfrak X$를 $A$ 위에서 고유인 형식적 스킴이라 하자(그
$\mathfrak m(A)$-아딕 위상을 갖춘다). 다음을 가정한다.
\begin{enumerate}
  \item[\textit{i.}] $\Osh_{\mathfrak X}$의 국소환들은 평탄
    $A$-가군이다. 동치로, $\Osh_{\mathfrak X}$와 $A$에 $A$의 극대
    아이디얼의 거듭제곱들이 정의하는 여과를 주면, 연관 등급 대상들 사이에
    다음 관계가 성립한다.
    \[
      \operatorname{gr}(\Osh_{\mathfrak X})
      \simeq
      \operatorname{gr}^{0}(\Osh_{\mathfrak X})
      \otimes_k\operatorname{gr}(A).
    \]
  \item[\textit{ii.}]
    $\mathfrak X_0=\mathfrak X\otimes_A k$를 $k$ 위의 대수적 스킴으로 보면
    \[
      H^2(\mathfrak X_0,\Osh_{\mathfrak X_0})=0
    \]
    이다.
  \item[\textit{iii.}] $\mathfrak X_0$는 사영적이다.
\end{enumerate}
이 조건들 아래에서 $\mathfrak X$는 대수화 가능하며, 더 정확히는 사영
$A$-스킴 $X$의 형식적 완비화 $\overline X$와 동형이다.

조건 \textit{(ii)}와 \textit{(iii)}는 특히 $\mathfrak X_0$가 $k$ 위의
비특이 곡선이면 성립한다. 따라서 정리 4는 특히 주어진 종수의 곡선에 관한
``모듈라이 이론''에 적용할 수 있다\ldots{} 정리 4를 어떻게 증명하는지
설명하자. 조건 \textit{(i)}와 \textit{(ii)}가 다음 사실을 함의함을
보인다. $\mathfrak X_0$ 위의 모든 연접층 가운데 국소적으로 구조층과
동형인 것은 $\mathfrak X$ 위의 같은 종류의 층을 환원하여 얻어진다
(아래의 명제 3 참조).
이제 $\mathfrak X_0$ 위의 ``풍부한'' 층에서 출발하여(\textit{(iii)}에
따라 그러한 층이 존재한다) 이를 $\mathfrak X$ 위의 가역층으로 올린다.
그런 다음 정리 1을 사용하여, 이 가역층의 어떤 양의 텐서 거듭제곱이
$\mathfrak X$에서 스킴 $\mathbf P_A^r$의 형식적 완비화로 가는 몰입
사상을 정의함을 증명한다($\mathbf P_A^r$은 $A$ 위의 차원 $r$인
``표준 사영공간''이다).

정리 1부터 3까지의 증명은 [1]을 참조하라.

\src{199}
\fsection{4}{연결성 정리와 자리스키의 ``주정리''에의 응용}

$f:X\longrightarrow Y$를 스킴들의 고유 사상이라 하자. 유한성 정리에 따라
$f_*(\Osh_X)=\mathcal A$는 $Y$ 위의 연접층이다. 더구나 이는 가환 대수의
층이므로 $Y$ 위에서 유한인 $Y$-스킴 $g:Y'\longrightarrow Y$에 대응한다.
이 스킴은 $Y$ 위에서 아핀이라는 조건, 즉 아핀 열린집합의 역상이 아핀이라는
조건과 $g_*(\Osh_{Y'})=\mathcal A$라는 조건으로 정의된다. 이제 $f$가
표준적으로 $f=gf'$으로 인수분해됨은 바로 알 수 있다. 여기서
$f':X\longrightarrow Y'$은 $X$에서 $Y'$으로 가는 사상이며
$f'_*(\Osh_X)=\Osh_{Y'}$을 만족한다. 이 $f$의 인수분해를 $f$의
\emph{스타인 인수분해}라 한다. 제1비교정리의 따름정리 1을 $f'$과
$Y'$에서 점 $y'$ 하나로 이루어진 부분에 적용하면
$f'^{-1}(y')=X'$이 연결임을 얻는다. 실제로 그렇지 않다면 $X'$을 따른
$X$의 형식적 단면들이 국소환을 이루지 않을 터인데,
$f'_*(\Osh_X)_{y'}=\Osh_{Y',y'}$의 완비화는 국소환이기 때문이다! 따라서
다음을 증명하였다.

\result{정리}{5 (자리스키의 ``연결성 정리'')} $f:X\longrightarrow Y$를
고유 사상이라 하자. 그러면 $f$는 동형을 제외하고 유일하게 $f=gf'$으로
인수분해된다. 여기서 $g:Y'\longrightarrow Y$는 유한 사상이고
$f':X\longrightarrow Y'$은 $f'_*(\Osh_X)=\Osh_{Y'}$을 만족한다(따라서
$g_*(\Osh_{Y'})=f_*(\Osh_X)$이다). $f'$의 섬유들은 연결이다. 즉 $f$의
섬유 $f^{-1}(y)$의 연결성분들의 집합은 $y$ 위에 놓인 $Y'$의 점들의
집합, 다시 말해 $f_*(\Osh_X)_y$의 극대 아이디얼들의 집합과 일대일로
대응한다.

여기서 연결성 정리의 통상적인 변형들을 곧바로 얻는다. 다음의 결과만
기록하자.

\src{200}
\result{따름정리}{1} $X$의 점 $x$가 그 섬유 $f^{-1}(y)$에서 고립되어
있을 필요충분조건은 섬유 $f'^{-1}(y')$($y'=f'(x)$)이 $x$ 하나로
이루어져 있는 것이다. 동치로, $f'$이 $x$의 한 근방에서 $y'$의 한
근방으로 가는 동형을 유도하는 것이다. 이러한 점들의 집합은 열린집합
$U$이고, $f'$은 $U$에서 $Y'$의 한 열린집합으로 가는 동형을 유도한다.

$f'$이 $x$에서 국소 동형사상임을 보이기 위하여, $f'$이 동형
$\Osh_{y'}\longrightarrow\Osh_x$를 유도함을 주목하자. 이는
$f'_*(\Osh_X)=\Osh_{Y'}$와 다음 사실에서 알 수 있다. 즉 $f'$은 닫힌
사상이고 $y'$ 위의 그 섬유는 $x$ 하나로 이루어져 있으므로, $V$가
$y'$의 근방들의 기본계를 움직일 때 $f'^{-1}(V)$는 $x$의 근방들의
기본계를 움직인다. 이로부터 ``기하학적'' 경우에는 슈발레에게서 비롯된
다음 결과를 곧바로 얻는다.

\result{따름정리}{2} $f$가 유한 사상일 필요충분조건은 $f$가 고유이고
그 섬유들이 유한한 것이다.

실제로 이 경우에는 앞의 결과에 따라 $f'$이 동형이다.

이제 $f:X\longrightarrow Y$를 반드시 고유일 필요는 없는 사상이라 하자.
다만 $X$가 고유 $Y$-스킴
$\widetilde f:\widetilde X\longrightarrow Y$의 열린 부분집합으로 들어
있다고 가정하자(특히 $f$가 준사영이면 그러하다). 따름정리 1을 적용하면,
$\widetilde f'$이 자기 섬유에서 고립된 $X$의 점들의 집합 $U$에서
$Y'$의 한 열린집합으로 가는 동형을 유도함을 알 수 있다(따라서 $U$는
실제로 열린집합이다). 이로부터 자리스키의 ``주정리''의 다음 전역 형태를
얻는다.

\result{정리}{6 (자리스키의 ``주정리'')} $f:X\longrightarrow Y$를
유한형 사상이라 하자. 그러면 자신이 속한 섬유에서 고립된 $X$의 점들의 집합
$U$는 열린집합이다. 또한 $f$가 분리 사상이면,\footnote{1962년의 정오표는
여기서 인쇄본의 ``$f$는 준사영이다''라는 가정을 더 약한 ``$f$는 분리
사상이다''라는 가정으로 바꾼다. 이는 SGA VIII, 6.2, 즉 모든 준유한 분리
사상은 준사영이라는 사실에 따른 것이다.} $U$는 $Y$ 위에서 유한인 어떤
스킴 $Y'$의 열린 부분과 $Y$-동형이다.

유한형 사상은 국소적으로 아핀이며, 따라서 \textit{a fortiori}
국소적으로 준사영이므로, 정리 6에서 주정리의 통상적인 국소 형태들이
곧바로 따른다.
\fsection{5}{고유 평탄 사상의 코호몰로지적 연구에의 응용}

$f:X\longrightarrow Y$를 고유 사상이라 하고, $F$를 $X$ 위의 연접층이라
하자. $F$는 $Y$ 위에서 평탄하다고 가정한다. 즉 $F_x$들은 환
$\Osh_y$ 위의 평탄 가군이다($y=f(x)$). 이는 또한 다음을 뜻한다. 모든
$y\in Y$에 대하여 섬유 $f^{-1}(y)$을 따라 $F$에
$\mathfrak m_y^nF$들로 이루어진 여과를 주면($\mathfrak m_y$는
$\Osh_y$의 극대 아이디얼이다), 그 연관 등급 대상은
$(F/\mathfrak m_yF)\otimes_{\kappa(y)}\operatorname{gr}(\Osh_y)$와
동형이다. 다시 말해 모든 정수 $n$에 대하여
\[
  \mathfrak m_y^nF/\mathfrak m_y^{n+1}F
  =F_y\otimes_{\kappa(y)}
  (\mathfrak m_y^n/\mathfrak m_y^{n+1})
\]
이다. 여기서 $X_y$는 섬유 $f^{-1}(y)$을 $y$의 잉여체 $\kappa(y)$
위에서 고유인 스킴으로 본 것이고, $F_y$는 $F$가 $X_y$ 위에 유도하는
층 $F/\mathfrak m_yF$이다. 이 동형과 정리 2를 고려하면, $F_y$를
계수로 하는 $X_y$의 코호몰로지를 앎으로써 $y$의 근방에서
$R^qf_*(F)$의 상계를 얻고 때로는 이를 계산할 수 있다. 여기서 정리 2는
다음 꼴이 된다.
\[
  \widehat{R^qf_*(F)_y}
  =\varprojlim_n H^q(X_y,F/\mathfrak m_y^nF).
\]
여기서는 그 따름정리 하나만 기록하겠다.

\src{201}
\result{명제}{1} $f:X\longrightarrow Y$를 고유 사상, $F$를 $X$ 위의
연접층이라 하고, $F$가 $Y$ 위에서 평탄하다고 하자. $y\in Y$이고 $q$가
정수이며 $H^q(X_y,F_y)=0$이라고 가정하자. 그러면 $R^qf_*(F)$는 $y$의
한 근방에서 영이고, 모든 $n$에 대하여 자연스러운 준동형
\[
  R^{q-1}f_*(F)_y\longrightarrow
  H^{q-1}(X_y,F/\mathfrak m_y^nF)
\]
은 전사이다.

특히 $f$가 평탄 사상이면(즉 $\Osh_X$가 $Y$ 위에서 평탄하면), $X$ 위의
모든 국소 자유 연접층 $F$는 $Y$ 위에서 평탄하다. 이러한 두 층 $F,G$에
대하여 명제 1을 $\Homsh_{\Osh_X}(F,G)$와 $q=1$에 적용하면 다음을 얻는다.

\result{정리}{7} $f$를 고유 평탄 사상, $F,G$를 $X$ 위의 국소 자유인 두
연접층, $y\in Y$라 하고
\[
  H^1\bigl(X_y,\Homsh_{\Osh_{X_y}}(F_y,G_y)\bigr)=0
\]
이라고 가정하자. 그러면 모든 준동형 $u_0:F_y\longrightarrow G_y$는
준동형
\[
  u:F|V\longrightarrow G|V,
  \qquad V=f^{-1}(U),
\]
에서 유도된다. 여기서 $U$는 $y$의 한 근방이다.

\result{따름정리}{} $u_0$가 동형사상(각각 단사사상, 에피사상)이면,
$U$를 충분히 작게 잡았을 때 $u$도 마찬가지이다.

특히 다음이 성립한다.

\src{202}
\result{따름정리}{2} $E_0$를 $X_y$ 위의 국소 자유 연접층이라 하고
\[
  H^1\bigl(X_y;\Homsh_{\Osh_{X_y}}(E_0,E_0)\bigr)=0
\]
이라고 하자. 그러면 $X_y$에 제한했을 때 $E_0$와 동형인 두 국소
자유층은 $X_y$의 한 근방에서 서로 동형이다.

따라서 다음이 성립한다.

\result{따름정리}{3} $H^1(X_y,\Osh_{X_y})=0$이라고 가정하자. 그러면
$X$ 위의 두 가역층(즉 $\Osh_X$와 국소적으로 동형인 층)은 $X_y$에 대한
그 제한들이 서로 동형이면 $X_y$의 한 근방에서 서로 동형이다.\footnote{원문에는 근방 조건이 명시되어 있지 않다. 여기서는 따름정리2의 국소적 결론을 적용하며, 명제2의 증명에서도 이 결론을 사용한다.}

이로부터 다음을 얻는다.

\result{명제}{2} $Y$를 연결 스킴, $E$를 $Y$ 위의 계수가 적어도 $2$인
국소 자유 연접층이라 하자. $E$로 정의되는 사영 다발
$X=\mathbf P(E)$를 그 잘 알려진 가역층 $\Osh_X(1)$과 함께 생각하자.
그러면 $X$ 위의 모든 가역층 $L$은
\[
  f^*(L')\otimes\Osh_X(n)
\]
꼴의 층과 동형이다. 여기서 $L'$은 $Y$ 위의 가역층이고 $n$은 정수이다.
이 $n$은 유일하게 정해지고, $L'$은 동형을 제외하고 정해진다.

앞의 따름정리 3에 따라 $L$은 각 섬유의 한 근방에서 어떤
$\Osh_X(n)$과 동형이다. 나머지는 거의 형식적으로 따른다.

명제 2를 사용하면 $X=\mathbf P(E)$에서 다른 사영 다발로 가는
$Y$-사상들을 결정할 수 있다. 특히 다음을 얻는다.

\result{따름정리}{} $u$를 $X=\mathbf P(E)$의 자동사상이라 하자. 그러면
$Y$ 위의 가역층 $L'$과 $E$에서 $E\otimes L'$으로 가는 동형사상 $v$가
존재하여, $u$는 이에 대응하는 동형사상
\[
  \mathbf P(E)\xrightarrow{\sim}\mathbf P(E\otimes L')=\mathbf P(E)
\]
이 된다. 쌍 $(v,L')$은 동형을 제외하고 정해진다.

$\Gamma$를 $E\otimes L'$이 $E$와 동형이 되는 $Y$ 위의 가역층 $L'$의
동형류들로 이루어진 집합이라 하자. 그 원소들은 모두 유한 위수를 갖는다.
실제로 $n$이 $E$의 계수이면, $n$차 외승을 취하여
$L'^{\otimes n}\simeq\Osh_Y$이어야 함을 알 수 있다. 따라서 이
따름정리는 다음 군의 완전열이 존재한다는 말로 표현할 수 있다.
\[
  e\longrightarrow
  \operatorname{Aut}(E)/\Gamma(Y,\Osh_Y^*)
  \longrightarrow \operatorname{Aut}_Y(X)
  \longrightarrow \Gamma\longrightarrow e
\]
(사실 이는 다음 군의 층들의 완전열에서 유도되는 코호몰로지 완전열로부터도
얻어진다.
\[
  e\longrightarrow \Osh_Y^*
  \longrightarrow \underline{\operatorname{Aut}}(E)
  \longrightarrow \underline{\operatorname{Aut}}_Y(X)
  \longrightarrow e,
\]
여기서 $\Osh_Y^*$는 $\Osh_Y$의 ``가역원''들의 층이며
$\underline{\operatorname{Aut}}(E)$의 중심과 동일시된다).

\src{203}
\fsection{6}{\texorpdfstring{$J$}{J}-아딕 완비환 위의 층과 스킴에 대한 존재 및 유일성 정리의 응용}

정리 7에서는 정리 1과 2를 이용하여 국소 자유 연접층에 관한 유일성 결과를
주었다. 이제 정리 3을 이용하여 층, 스킴의 사상 또는 스킴 자체에 관한 존재
정리들을 도출한다. 이하에서 $A$는 분리·완비인 뇌터 국소환을 나타낸다.
일반적인 방법은 언제나 형식적 구성을 수행하는 데 있다. 이는 본질적으로
아르틴 환 위에서 대수기하학을 수행하고, 세 기본정리를 이용하여 그로부터
``대수적'' 성질의 결론을 끌어내는 것이다.

\result{명제}{3} $\mathfrak X$를 $A$ 위에서 고유이고 평탄인 형식적
스킴이라 하고, $F_0$를 $X_0$ 위의 국소 자유층이라 하자. 다음을 가정한다.
\[
  H^2\bigl(X_0,\Homsh_{\Osh_{X_0}}(F_0,F_0)\bigr)=0.
\]
그러면 $X_0$ 위에 $F_0$와 동형인 층을 유도하는 $\mathfrak X$ 위의 국소
자유층 $F$가 존재한다. (더구나
$H^1(X_0,\Homsh_{\Osh_{X_0}}(F_0,F_0))=0$이면 이 $F$는 동형까지
유일하다.)

$X_n$ 위의 국소 자유층 $F_n$들을 각각 앞의 것을 유도하도록 차례로
구성한다. $F_n$을 구성하는 데 따르는 장애는
\[
  H^2\bigl(X_0,\Homsh_{\Osh_{X_0}}(F_0,F_0)\bigr)
  \otimes_{A/J}(J^n/J^{n+1})
\]
안에 놓이므로, 가설에 따라 사라진다. 이제 정리 3을 이용하면 다음을 얻는다.

\result{따름정리}{1} $X$를 $A$ 위에서 고유이고 평탄인 스킴이라 하고,
$F_0$를 위와 같이 잡자. 그러면 $X_0$ 위에 $F_0$와 동형인 층을 유도하는
$X$ 위의 국소 자유층 $F$가 존재한다. 더구나
$H^1(X_0,\Homsh_{\Osh_{X_0}}(F_0,F_0))=0$이면 이 $F$는 동형까지
유일하다.

$X_0$를 체 $k$ 위의 유한형 스킴이라 하자. $X_0$는 $k$ 위에서
비특이(여기서는 절대 비특이라는 뜻이다)라고 가정하지만, 반드시 $k$ 위에서
고유일 필요는 없다. $A$를 잉여체가 $k$인 아르틴 국소환이라 하자. 우리는
$A$ 위에서 평탄이고
$X\otimes_A k=X_0$인 스킴 $X$를 찾는 데 관심이 있다(이것이 $X_0$의
``모듈라이 이론'', 또는 ``구조의 변형''의 출발점이다). 그러한 $X$를 주는
것과, 위상공간 $X_0$ 위의 층 $\Osh_X$에 다음 구조들을 주는 것은 같다.
\begin{enumerate}
  \item[\textit{i.}] 이는 $A$-대수의 층이다.
  \item[\textit{ii.}] $A$-대수 구조들과 양립하는 증대 준동형
    $\Osh_X\longrightarrow\Osh_{X_0}$을 갖는다. 이 데이터에는 다음 조건들이
    부과된다. 증대 준동형은 동형사상
    $\Osh_X\otimes_A k\xrightarrow{\sim}\Osh_{X_0}$을 유도한다. 또한
    $\Osh_X$는 $A$ 위에서 평탄하다. 즉 $A$의 극대 아이디얼
    $\mathfrak m$의 거듭제곱들로 여과한 $\Osh_X$의 연관 등급 대상은
    $\operatorname{gr}^0(\Osh_X)\otimes_k\operatorname{gr}(A)$와 동형이다.
    다시 말해 다음 동형사상들이 존재한다.
    \[
      \mathfrak m^n\Osh_X/\mathfrak m^{n+1}\Osh_X
      =\Osh_{X_0}\otimes_k(\mathfrak m^n/\mathfrak m^{n+1}).
    \]
\end{enumerate}
기본 사실은 다음과 같다.

\src{204}
\result{정리}{8} $X_0$를 체 $k$ 위의 유한형 비특이 스킴이라 하고,
$X_0$가 아핀이라고 가정하자. $A$를 잉여체가 $k$인 아르틴 국소환이라
하자. 그러면 $A$ 위에서 평탄이고 $X\otimes_A k=X_0$인 $A$-스킴 $X$가
존재하며, 그러한 두 스킴은 반드시 서로 동형이다.

여기서 말한 동형사상은 표준적이지 않음에 유의하자. 일반적으로 $X$에는
$X_0$ 위에서 항등을 유도하는 자명하지 않은 $A$-자기동형사상들이 있기
때문이다. 한편 주어진 조건들을 만족하는 $X$의 ``표준적'' 선택은 일반적으로
존재하지 않는다. 다만 $A$가 $k$-대수인 경우(동표수인 경우)에는
$X=X_0\otimes_k A$, 즉 $\Osh_X=\Osh_{X_0}\otimes_k A$로 잡을 수 있다
($X_0$가 아핀인지 여부와 무관하다). 이표수인 경우에는 $X_0$가 아핀이 아닐
때 $X_0$를 $A$ 위에서 정의된 $X$로 ``올릴'' 수 있는지 나는 일반적으로
알지 못한다. 그러나 $n>0$인 정수 $n$을 잡고
$A_{n-1}=A/\mathfrak m^n$으로 놓자. $X_0$를 평탄 $A_{n-1}$-스킴
$X_{n-1}$로 올렸다고 가정하고, $X_{n-1}$을 평탄 $A_n$-스킴 $X_n$으로
올리고자 한다. 정리 8에 따라 이것이 국소적으로 가능함은 이미 알고 있다.
한편 $U_n$이 $X_{n-1}$의 열린집합 $U_{n-1}$을 올린다면, $U_{n-1}$ 위에서
항등을 유도하는 $U_n$의 자기동형군의 층은
\[
  \Theta_{X_0/k}\otimes_k(\mathfrak m^n/\mathfrak m^{n+1})
\]
을 $U_{n-1}$에 제한한 층과 표준적으로 동형이며, 특히 가환임을 쉽게
확인할 수 있다. (여기서 $\Theta_{X_0/k}$는 $X_0$ 위의 $k$-미분들의
싹의 층을 나타낸다.) 따라서 $X_{n-1}$을 올리는 $X_n$을 구성하는 데
따르는 장애가 존재하며, 이는
\[
  H^2(X_0,\Theta_{X_0/k})\otimes_k
  \mathfrak m^n/\mathfrak m^{n+1}
\]
안에 놓인다. 따라서 다음을 얻는다.

\result{따름정리}{1} $X_0$를 $k$ 위의 유한형 비특이 스킴이라 하고,
$H^2(X_0,\Theta_{X_0/k})=0$이라고 가정하자. 그러면 잉여체가 $k$인 모든
아르틴 국소환 $A$에 대하여 $A$ 위에서 평탄이고
$X\otimes_A k=X_0$인 $A$-스킴 $X$가 존재한다.

더구나 $X_0$를 올리는, $A$ 위에서 평탄인 $X$를 하나 찾았다면, 정리 8에
따라 $X_0$를 올리는 평탄 $A$-스킴들의 동형류의 집합은
\[
  H^1\bigl(X_0,\underline{\operatorname{Aut}}(X)\bigr)
\]
과 동일시된다. 여기서 물론 $\underline{\operatorname{Aut}}(X)$는 증대와
양립하는 $A$-대수의 층 $\Osh_X$의 자기동형사상들의 싹의 층을 나타낸다.
$\Osh_X$의 여과는 $\underline{\operatorname{Aut}}(X)$의 여과를 정의하며,
이 층을 그 여과의 $n$번째 부분군으로 나눈 몫은
$\underline{\operatorname{Aut}}(X_n)$이다. 이 여과의 연관 등급 대상은
가환이고 $\Theta_{X_0/k}\otimes_k\operatorname{gr}(A)$와 동일시된다.
특히 $\mathfrak m^{n+1}$이 처음으로 0이 되는 $\mathfrak m$의 거듭제곱이면,
$F^n(\underline{\operatorname{Aut}}(X))$(여과의 마지막 단계)는
$\underline{\operatorname{Aut}}(X)$의 중심에 들어가며
$\Theta_{X_0/k}\otimes_k\mathfrak m^n$과 동형이다. 이것은 또한
$X_{n-1}=X\otimes_A A/\mathfrak m^n$ 위에서 항등을 유도하는 $X$의
자기동형사상들의 싹의 층이다. 이 결과들을 이용하면 다음 명제들을 곧바로
얻는다.

\src{205}
\result{따름정리}{2} $X_0$를 $k$ 위의 유한형 비특이 스킴이라 하고,
$A$를 잉여체 $k$와 극대 아이디얼
$\mathfrak m$을 갖는 아르틴 국소환이라 하자. $\mathfrak m^{n+1}=0$이라 가정하고,
$A_{n-1}=A/\mathfrak m^n$이라 하며, 끝으로 $X_{n-1}$을 평탄
$A_{n-1}$-스킴으로서 $X_{n-1}\otimes_A k=X_0$인 것이라 하자. 그러면
$X_{n-1}$ 위에서 항등사상을 유도하는 동형을 제외한, 평탄
$A$-스킴 $X_n$ 중 $X_n\otimes_A A_{n-1}=X_{n-1}$인 것들의
동형류 집합은 공집합이거나
다음 군에 대한 주동차공간이다.
\[
  H^1(X_0,\Theta_{X_0/k})\otimes_k\mathfrak m^n.
\]

(일반적으로 이 마지막 공간에는 특별히 정해진 원점이 없음을 유의하라.
$X_{n-1}$을 $X_n$으로 올리는 특별히 정해진 방법이 없기
때문이다.)

\result{따름정리}{3} $X_0$를 $k$ 위의 유한형 비특이 스킴이라 하고,
$H^1(X_0,\Theta_{X_0/k})=0$이라 가정하자. 그러면 잉여체 $k$를 갖는 모든
아르틴 국소환 $A$에 대하여, $X\otimes_A k=X_0$인 평탄
$A$-스킴 $X$는 동형을 제외하면 많아야 하나 존재한다.

따름정리 1과 3에서, $A$를 잉여체 $k$를 갖는 완비 뇌터
국소환으로만 가정하고 $X$를
$A$ 위의 형식적 스킴으로 취했을 때의, 겉보기에는 더 일반적인 명제들이 즉시 따른다.

\result{정리}{9} $k$를 체, $X_0$를 $k$ 위의 유한형 비특이 스킴이라
하자. 잉여체가 $k$인 모든 완비 뇌터 국소환 $A$에 대하여,
$F(A)$를 $X_0$ 위에서 항등사상을 유도하는 동형을 제외한, $A$ 위의
유한형 평탄 형식적 스킴 $\mathfrak X$ 중
$\mathfrak X\otimes_A k=X_0$인 것들의 동형류 집합이라 하자. 이 표기 아래에서 다음이 성립한다.
\begin{enumerate}
  \item[\textit{i.}] $H^1(X_0,\Theta_{X_0/k})=0$이면 모든 $A$에 대하여
    $F(A)$는 많아야 한 원소를 갖는다.
  \item[\textit{ii.}] $H^2(X_0,\Theta_{X_0/k})=0$이면 모든 $A$에 대하여
    $F(A)$는 적어도 한 원소를 갖는다.
\end{enumerate}

\src{206}
\result{따름정리}{1} $X_0$가 $k$ 위에서 고유라고 가정하자. 조건
\textit{(i)} 아래에서 모든 $A$에 대하여, $X_0$ 위에서 항등사상을
유도하는 동형을 제외하면, $X\otimes_A k=X_0$인
$A$ 위의 고유 평탄 스킴 $X$는 많아야 하나 존재한다.

여기서는 정리 3의 따름정리 2를 사용한다. 예를 들면 다음과 같다.

\result{따름정리}{2} $X$가 $A$ 위의 고유 평탄 스킴이고,
$X\otimes_A k$가 $k$ 위의 차원 $r$인 표준 사영공간
$\mathbf P_k^r$와 동형이면, $X$는 $\mathbf P_A^r$와 동형이다.

(이 결과는 명제 3의 따름정리 1에서도 도출할 수 있다.)

\result{따름정리}{3} $X_0$를 $k$ 위의 비특이 사영 스킴이라 하고,
다음을 가정하자.
\[
  H^2(X_0,\Osh_{X_0})=H^2(X_0,\Theta_{X_0/k})=0.
\]
그러면 모든 $A$에 대하여, $X\otimes_A k=X_0$인 사영
평탄 $A$-스킴 $X$가 존재한다.

정리 9\textit{(ii)}와 정리 4를 결합한다. 특히 다음을 얻는다.

\result{따름정리}{4} $X_0$를 $k$ 위의 완비 비특이 대수곡선의
스킴이라 하자. 그러면 잉여체가 $k$인 모든 완비 뇌터 국소환 $A$에 대하여,
$X\otimes_A k=X_0$인 ``비특이 곡선 스킴'' $X$가 $A$ 위에 존재한다.

\result{비고}{}

1) 따름정리 3과 4는 특히 $k$의 표수가
$p\ne0$일 때 흥미롭다. 이때 $A$로 표수가 $0$이고 잉여체가 $k$인
이산 값매김환을 취한다. 예를 들면 ``가능한 가장 작은 $A$'',
즉 $p$가 극대 아이디얼을 생성하는 환을 취한다. (실제로
Cohen의 정리들에 따라, 그러한 환 $A$ 하나에 대하여 따름정리 3과 4를
아는 것으로 충분하다.) 이와 관련하여 전문가들의 말에 따르면,
체 $k$ 위에 정의된 스킴 가운데 그러한 환 $A$ 위에 정의된 평탄 스킴의
mod $p$ 환원이 아닌 것이 존재하는지는 알려져 있지 않다. 적어도 이 절의 결과들은
이 문제를 체계적으로 조사할 방법을 준다. 먼저
$H^2(X_0,\Theta_{X_0/k})$에 놓이는 첫 번째 장애가
반드시 영인지 살펴보아야 할 것이다.

\emph{1962년 보충.} --- J.-P. Serre는 표수가 $p>0$인 대수적으로
닫힌 체 $k$ 위에 3차원 비특이 사영 다양체를 구성했는데, 이 다양체는
잉여체가 $k$이고 분수체의 표수가 영인 국소 정역 위의 고유 스킴을
환원하여 얻어지는 것이 아니다. Mumford는 비특이 사영 곡면에 관한
이와 비슷한 예를 찾았다고 한다.

2) 정리 3과 이에 대응하는 기법은 밑환을 (생각을 구체화하기 위해 국소환으로
잡을 때) 완비라고 해야만 유효함을 유의하자.

국소환의 완비화에 대하여 알려진 결과에서 그 국소환 자체에 대한
대응하는 결과로 넘어가려면, 그 최종적인 명제가 아직 찾아지지 않은
네 번째 ``기본정리''가 필요할 것이다.

\src{207}
3) 이 절의 결과들(특히 앞의 따름정리 1과 2)과 다음 절의 결과들을,
복소구조의 변형에 관한 Kodaira--Spencer의 결과들과 비교해야 한다.
방금 언급한 추측적 정리를 사용하면, 따름정리 1의 조건 아래에서
$A$가 더 이상 완비라고 가정되지 않더라도, $A$를 포함하면서
$A$ 위에서 유한이고 비분기인 환 $A'$이 존재하여
$X\otimes_A A'$과 $X'\otimes_A A'$이 $A'$-동형임을 결론지을 수 있어야 한다
($X,X'$은 $X\otimes_A k=X'\otimes_A k=X_0$을 만족하는, 주어진 두
고유 평탄 $A$-스킴이다). 적어도
$X_0=\mathbf P_k^r$일 때에는 명제 2의 따름정리를 사용하여 이를 증명할 수 있다. 어떤
경우든 $H^1(X_0,\Theta_{X_0/k})=0$이면, $Y=\operatorname{Spec}(A)$의
모든 점 $y$ 위에서 $X,X'$의 섬유들은
적어도 잉여체 $\kappa(y)$의 대수적 폐포로 넘어간 뒤에는 동형임을 증명할 수 있다.
(더 이상 $A$가 반드시 국소환이라고 가정하지 않는, 겉보기에는
더 강한 국소 결과도 있다.) 사영공간의 ``구조의 변형''과 관련하여,
Kodaira--Spencer의 대응하는 문제에서 제안된 다음 질문도 언급하자. $X$를
분수체 $K$와 잉여체 $k$를 갖는 국소 정역 $A$ 위의
고유 평탄 스킴이라 하자. $X\otimes_A K$가
$\mathbf P_K^r$와 동형이라고 가정할 때, $X\otimes_A k=X_0$는
$\mathbf P_k^r$와 동형인가(적어도 $k$의 대수적 폐포 위에서)? 이 문제에서는
$A$를 완비 이산 값매김환이라고 가정해도 된다. $X_0$가
아벨 다양체일 때에도 비슷한 문제가 있다.

\emph{1962년 보충.} --- 강연자는 이때 위에서 추측한 결과들에 대하여
덜 낙관적이다. 그러나 사영공간의 구조의 변형에 관한 문제는 Hironaka가
긍정적으로 해결했고, 아벨 다양체에 관한 비슷한 문제는 Koizumi가 해결했다.
\fsection{7}{``모듈라이 이론''에의 응용}

강연자는 이 이론을 이제 막 직접 다루기 시작했으므로, 여기서는 몇 가지
지적에 한정하지 않을 수 없다. 단순화를 위해 체 $k$ 위에 놓는다. 즉 우리는
동일 표수에서 작업한다. 정리 8을 사용하면 본질적인 변경 없이 일반적인
경우도 다룰 수 있을 듯하다. 현재 우리는 아직 ``형식적'' 단계를 넘어서지
못했다. 그러나 강연자는 여기에서 출발하여 몇몇 경우에는 실제 모듈라이
스킴을 구성하고, 특히 모든 정수 $g$에 대하여 종수 $g$인 비특이 곡선들의
보편 모듈라이 스킴 역할을 하는 정수환 위의 스킴을 구성할 수 있으리라
기대한다.

\emph{1962년 보충.} --- Mumford는 종수 $g$인 곡선들의 모듈라이 스킴을
구성하였다(cf. SHMI). 또한 정리 10은 모듈라이 이론의 ``레벨 $n$의 야코비
스킴''이 비특이이고, 더 나아가 $\mathbf Z$ 위에서 매끄럽다는 것을 증명한다.

\src{208}
이제 정리 9의 상황과 표기로 돌아가되, $A$가 $k$ 위의 유한 차원 국소대수라고
가정하자. 단순화를 위해 $k$는 대수적으로 닫혔다고 가정한다. 그러면 $F(A)$를
$A$에 관하여 집합의 범주에 값을 갖는 공변 함자로 볼 수 있다. 실제로
$k$-대수의 준동형 $A\longrightarrow B$는 사상
$F(A)\longrightarrow F(B)$를 정의한다. 왜냐하면
$X\otimes_A k=X_0$을 만족하는 모든 평탄 $A$-스킴 $X$는 같은 성질들을 갖는
$B$-스킴 $X\otimes_A B$를 낳기 때문이다. 완비 뇌터 국소 $k$-대수
$\mathcal O$와 함자적 동형사상
\begin{equation}
  \tag{*}
  \operatorname{Hom}(\mathcal O,A)\xrightarrow{\sim}F(A)
\end{equation}
을 찾을 수 있다고 가정하자(여기서 왼쪽 항은 $k$-대수의 준동형들을 뜻한다).
그러한 $\mathcal O$는 정준 동형까지 정해짐을 쉽게 알 수 있다. 이때
$\mathcal O$의 형식적 스펙트럼 $\mathfrak Y$, 즉 한 점으로 이루어진
위상공간에 $\mathcal O$ 하나로 이루어진 위상환의 층을 준 것을 $X_0$의
\emph{형식적 모듈라이 스킴}이라 한다(이것이 반드시 존재하는 것은 아니다).
$\mathfrak r$을 $\mathcal O$의 극대 아이디얼이라 하고, 모든 $n$에 대하여
$\mathcal O_n=\mathcal O/\mathfrak r^{n+1}$이라 하자(따라서
$\mathcal O_0=k$이다). 그러면 표준 준동형
$\mathcal O\longrightarrow\mathcal O_n$은
$\operatorname{Hom}(\mathcal O,\mathcal O_n)$의 원소이므로
$F(\mathcal O_n)$의 원소, 즉 $\mathfrak r$로 나눈 환원이 $X_0$인 평탄
$\mathcal O_n$-스킴 $X_n$을 정의한다. 이 $X_n$들은 스칼라 확대(즉 여기서는
환원)에 의해 서로에게서 얻어진다. 따라서 이들은 형식적 모듈라이 스킴
$\mathfrak Y$ 위의 잘 정해진 형식적 스킴 $\mathfrak X$에서 온다.
$\mathfrak X$는 $\mathfrak Y$ 위에서 평탄이고 $\mathfrak X_0=X_0$이다.
그러면 동형사상 (*)은, 바로 알 수 있듯이, 모든 $k$-대수의 준동형
$\mathcal O\longrightarrow A$에 $A$-스킴
$\mathfrak X\otimes_{\mathcal O}A$의 동치류를 대응시켜 주어진다. 즉 모든
$k$-스킴의 사상 $Y'=\operatorname{Spec}(A)\longrightarrow\mathfrak Y$에
밑준스킴의 변경으로 얻은 $Y'$-스킴
$\mathfrak X\otimes_{\mathfrak Y}Y'$을 대응시킨다. 더 나아가 $A$가 단지
완비 뇌터 국소 $k$-대수라고만 가정해도(반드시 아르틴일 필요는 없다)
동형사상 (*)과 방금의 설명은 여전히 성립함을 알 수 있다. 물론 통상 그렇듯
$\mathcal O$는 \textit{a priori} 멱영원을 가질 수 있으며,
$\mathcal O$ 자체가 $k$와 같지 않은 아르틴 환인 경우도 존재할 법하다. 이는
$A$를 정칙환으로 한정하는 Kodaira--Spencer의 관점이 일반적인 경우에는
\textit{a priori} 얼마나 부적절한지를 보여 준다.

\src{209}
$k$ 위에서 고유라고 가정한 $X_0$에 형식적 모듈라이 스킴이 존재하기 위한
충분조건을 제시하는 일이 남았다. 일반적으로, 유한 차원 국소 $k$-대수에서
집합으로 가는 함자 $A\longmapsto F(A)$가 적당한 $\mathcal O$에 대하여
$\operatorname{Hom}(\mathcal O,A)$ 꼴로 나타날 수 있기 위한 간단한
필요충분조건을 주기는 쉽다. 여기서는 그 조건들을 자세히 말하지 않겠다.
다만 지금의 경우 이 조건들은 $X_0$에 비자명한 코호몰로지적 조건들을
부과하며, 강연자가 반례를 구성하지는 않았지만 그 조건들이 언제나
성립할 가능성은 낮아 보인다. 한편
$H^0(X_0,\Theta_{X_0/k})=0$이라는 조건은 형식적 모듈라이 스킴의 존재를
위한 충분조건일 듯하다(다만 결코 필요조건은 아니다). 여기서는 특히
단순한 경우의 정리 하나만 진술한다. 해석공간 이론에서 이에 대응하는
결과는 잘 알려져 있으며(cf. Kodaira--Spencer), 이 정리는 앞 절의
결과들을 사용하여 어렵지 않게 증명된다.

\result{정리}{10} $X_0$를 체 $k$ 위에서 고유이고 비특이인 스킴이라 하고,
다음을 가정하자.
\[
  H^0(X_0,\Theta_{X_0/k})=H^2(X_0,\Theta_{X_0/k})=0.
\]
그러면 $X_0$의 형식적 모듈라이 스킴이 존재하며, 이는 정칙 국소환
$\mathcal O$, 즉 $k$ 위의 형식적 멱급수 대수에 대응한다.

이미 지적했듯이, 일반적으로 $\mathcal O$ 위의 형식적 스킴 $\mathfrak X$가
대수화 가능한 것은 아니다. 그러나 $X_0$가 사영적이고
$H^2(X_0,\Osh_{X_0})=0$이면 대수화 가능함이 알려져 있다(정리 4). 예를
들어 $X_0$의 차원이 1이면 그러하다. 이는 주어진 종수의 곡선들을 위한
정수환 위의 모듈라이 스킴을 구성할 수 있으리라는 기대를 어느 정도
뒷받침한다\ldots{}

또한 이 절에서 설명한 것과 같은 방법들은 피카르 다양체의 구성과 연구뿐
아니라 그 밖의 여러 구성에도 적용할 수 있다. 이에 관해서는 곧 다시
돌아오겠다.
\fsection{8}{기본군에의 응용}

앞에서 설명한 기법들을 사용하면 위상 이론을 본보기로 삼아 기본군을
체계적으로 연구할 수 있다. 이 절에서 서술하는 첫 두 정리는 랑--세르의
최근 논문에 있는 결과들을 일반화한 것이다.

$X$를 스킴이라 하자. $X$-스킴 $X'$이 다음 조건들을 만족하면 이를
$X$의 비분기 피복이라 한다.
\begin{enumerate}
  \item[\textit{i.}] $X'$은 $X$ 위에서 유한이다. 즉 $X$ 위의 연접인
    대수층 $\mathcal A=\mathcal A(X')$에 의해 정의된다.
  \item[\textit{ii.}] $\mathcal A$는 $X$ 위에서 국소 자유층이다.
  \item[\textit{iii.}] 모든 $x\in X$에 대하여 몫
    \[
      \mathcal A_x/\mathfrak m_x\mathcal A_x
      =\mathcal A_x\otimes_{\Osh_x}\kappa(x)
    \]
    은 $\kappa(x)$ 위의 분리 대수이다.
\end{enumerate}

\src{210}
이 비분기 피복의 개념(세르와 강연자에 의한 것이다)은 합리적으로 기대할
수 있는 모든 초보적 성질을 가지며, 여기서는 그 목록을 제시하지 않겠다.
고전적 갈루아 이론과 위상 피복의 갈루아 이론을 본보기로 삼은 갈루아
이론이 여기서 생긴다는 것만 말해 두자. 이 이론은 고전적 갈루아 이론을
포함하며, 그 증명들은 후자에 관하여 통상 받아들여지는 증명들보다 오히려
더 간단하다. 정확히 말해, 스킴 $X$의 \emph{기하적 점}이란 대수적으로
닫힌 체 $\Omega$의 스펙트럼 $\xi$에서 $X$로 가는 사상 $a$를 말한다.
이는 $X$의 점 $x=|a|$의 잉여체 $\kappa(x)$의 대수적으로 닫힌 확대체를
주는 것과 같다(이 $x$를 기하적 점 $a$의 \emph{소재점}이라 한다).
$X'$이 $X$의 비분기 피복이면 이에 ``$a$ 위에 있는 $X'$의 기하적
점들''의 집합 $E_a(X')$를 대응시킬 수 있다. 즉 이는 $x$ 위에 있는
$x'\in X'$과 $\kappa(x)$-준동형
$\kappa(x')\longrightarrow\Omega$로 이루어진 순서쌍들의 집합이다.
따라서 $(X,a)$를 고정하면 $X$의 비분기 피복 $X'$들의 범주 $R(X)$에서
유한집합의 범주로 가는 함자 $F(X,a)$를 얻는다. $X$가 연결이면 $R(X)$와
$F(X,a)$로 이루어진 쌍은 적당한 콤팩트 완전 불연결 위상군 $\pi$(즉
유한군들의 사영극한)가 정의하는 유사한 쌍과 동형이 되기 위해 필요한
형식적 성질들을 갖는다. 곧 $\pi$가 연속적으로 작용하는 유한집합 $E$들의
범주 $C(\pi)$와, 이 범주에서 유한집합의 범주로 가는 망각함자
$F(\pi)(E)=E$를 취한다. 더욱이 $(C(\pi),F(\pi))$가 주어진 쌍과
동형이라는 조건은 $\pi$를 정준 동형까지 결정한다. 이 경우 $\pi$를
기하적 점 $a$에서의 연결 스킴 $X$의 \emph{기본군}이라 하며
$\pi_1(X,a)$로 쓴다. $X$가 연결이 아니면 $X$를 $x=|a|$를 포함하는
연결 성분으로 바꾼다. 한편 $X$가 연결이면 $X$의 두 기하적 점
$a',a''$에 대응하는 군 $\pi_1(X,a)$들은 동형이며(그 동형은 내부
자기동형까지 정해진다), 따라서 통상처럼 목적에 가장 알맞게 $a$를
고를 수 있다. 예를 들어 $X$가 기약이라고 가정하고 그 일반점에서
$a$를 고를 수 있다. 물론 $\pi_1(X,a)$는 기점이 지정된 스킴 $(X,a)$에
관한 공변 함자이다. 비분기 피복의 분류에 관한 모든 명제\footnote{인쇄된
본문에는 ``비분리 피복''이라고 되어 있다. 그러나 이 문단은 비분기
피복들의 범주 $R(X)$와 연속적인 $\pi_1(X,a)$-작용을 갖는 유한집합들의
범주 사이의 동치를 다룬다.}는 잘 알려진 사전에 따라 군론의 언어로
옮길 수 있다(다만 여기서는 군들이 위상군이라는 사실을 고려해야 한다).

\src{211}
우리의 목적은 고유 사상 $f:X\longrightarrow Y$에 관하여 섬유공간의
호모토피 완전열에 대응하는 것을 전개하는 데 있다. 고차 호모토피군이
무엇인지 알지 못하므로, 얻는 결과가 불완전할 수밖에 없음은 분명하다.
제3절의 기본정리들을 적용하려면 먼저 체 또는 아르틴 환 위의 스킴에
관한 몇 가지 초보적 보조정리를 명시해야 한다(일반적인 절차에 따라서!).

\result{보조정리}{1} $(X',a')$를, 표시점 $e$를 갖춘 유한집합 $E$
위에서의 $\pi_1(X,a)$의 표현에 대응하는 기점이 지정된 비분기 피복이라
하자. 그러면 정준 준동형
$\pi_1(X',a')\longrightarrow\pi_1(X,a)$은 첫째 군을 $\pi_1(X,a)$ 안의
$e$의 안정자 부분군과 동일시한다(따라서 이 준동형은 단사이다).

\result{보조정리}{2} $X$를 체 $k$ 위의 대수적 스킴, $k'$을 $k$의
순수 비분리 확대체라 하자. 그러면 $X\otimes_k k'$의 모든 비분기 피복은
$X$의 어떤 비분기 피복의 역상(즉 스칼라 확대)으로 얻어지며, 그 피복은
정준 동형까지 결정된다.

특히 이 두 보조정리로부터, $k$의 모든 대수적 확대체 $K$와
$X$의 기하적 점 $a$로 사영되는 $X'=X\otimes_k K$의 모든 기하적 점
$a'$에 대하여 함자적 준동형
$\pi_1(X',a')\longrightarrow\pi_1(X,a)$이 단사임이 따른다.

\result{보조정리}{3} $A$를 아르틴 국소환이라 하고, $X$를 $A$ 위의
완전 스킴으로서 $H^0(X,\Osh_X)=A$인 것이라 하자. $X'$을 $X$의 비분기
피복이라 하고 $A'=H^0(X',\Osh_{X'})$로 놓자. 그러면 $A'$은
$A$-가군으로 유한 생성인 환이다(선험적으로는 $A$ 위에서 분기될 수도
있다). $X_0,X'_0$를 각각 $X,X'$에 딸린 축소 부분스킴이라 하자(이는
$\Osh_X,\Osh_{X'}$ 안의 멱영원들로 이루어진 층으로 나누어 얻는다).
$A/\mathfrak r(A)$의 부분체 $k$로서 $A/\mathfrak r(A)$가 $k$ 위에서
유한인 것을 택하자(따라서 $X_0$는 $k$ 위의 완전 대수적 스킴이고
$X'_0$는 그 비분기 피복이다). 끝으로 $\Omega$를 $k$의 대수적으로 닫힌
확대체라 하고, $X_0\otimes_k\Omega$의 비분기 피복
$X'_0\otimes_k\Omega$를 생각하자. 그러면 다음 두 조건은 동치이다.
\begin{enumerate}
  \item[\textit{i.}] $X'_0\otimes_k\Omega$는
    $X_0\otimes_k\Omega$ 위에서 완전히 분해되어 있다.
  \item[\textit{ii.}] 자연스러운 사상
    $X'\longrightarrow X\otimes_A A'$은 동형이다.
\end{enumerate}
이 조건들이 성립하면 $A'$은 $A$의 비분기 확대이다. 끝으로 $X'$이
연결이면 조건 \textit{(i)}는 겉보기에는 더 약한 다음 조건과 동치이다.
\begin{enumerate}
  \item[\textit{i bis.}] $X'_0\otimes_k\Omega$는
    $X_0\otimes_k\Omega$ 위에서 정칙 단면을 갖는다.
\end{enumerate}
조건 \textit{(ii)}가 성립할 때 $X$의 비분기 피복 $X'$을
\emph{기하적으로 자명하다}고 하겠다.

\src{212}
\result{보조정리}{4} $f:X\longrightarrow Y$를
$f_*(\Osh_X)=\Osh_Y$를 만족하는 고유 사상이라 하자. $a$를 $X$의
기하적 점, $b$를 $Y$ 위의 그 상이라 하자. 그러면
$\pi_1(X,a)\longrightarrow\pi_1(Y,b)$는 전사이다.

실제로 다음을 보이면 충분하다. $Y$의 비분기 피복 $Y'$이 국소 자유인
대수의 층 $\mathcal A$에 대응하고 $X\otimes_Y Y'$이 비연결이면,
$Y'$도 비연결이다. 실제로 $f^*(\mathcal A)$는 서로 영이 아닌 두 환의
층의 직합이고, 따라서 그 직접상도 그러하다. 그런데 이 직접상은 다름
아닌
$\mathcal A\otimes f_*(\Osh_X)=\mathcal A$이다.

\result{보조정리}{5} $X$를 체 $k$ 위의 완전 스킴이라 하자.
$H^0(X,\Osh_X)$가 국소환 $A$이고 $A/\mathfrak r(A)$가 $k$의 순수
비분리 확대체라고 가정하자. $\Omega$를 $k$의 대수적 폐포라 하고
$\overline X=X\otimes_k\Omega$라 하자(이는 연결이다). $\overline X$의
기하적 점 $\overline a$로서 $X$의 기하적 점 $a$로 사상되는 것을 하나
택하면 다음 완전열을 얻는다.
\[
  e\longrightarrow\pi_1(\overline X,\overline a)
  \longrightarrow\pi_1(X,a)
  \longrightarrow\pi_1(k,b)\longrightarrow e
\]
(여기서 $\pi_1(k,b)$는 $\Omega$의 $k$ 위의 갈루아 군이다.)

첫 번째 준동형이 단사라는 사실은 보조정리 1과 2에서 보았고, 가운데에서의
완전성은 보조정리 3에서 따르며, 마지막 준동형의 전사성은 보조정리 4에서
따른다. 이 마지막 사실만이 $A/\mathfrak r(A)$가 $k$의 순수 비분리
확대체라는 가정을 사용한다.

\result{명제}{4} $f:X\longrightarrow Y$를 고유 평탄 사상이라 하고,
모든 $y\in Y$에 대하여
$H^0(f^{-1}(y),\Osh_{f^{-1}(y)})$가 잉여체 $\kappa(y)$ 위의 분리가능
대수라고 하자. 예를 들어 $f^{-1}(y)$가 $\kappa(y)$ 위의 분리가능
스킴, 즉 축약이고 그 기약 성분들에 대응하는 체들이 $\kappa(y)$의
분리가능 확대체이면 이 조건이 성립한다. 그러면 $f_*(\Osh_X)$에
대응하는 $Y$의 피복 $Y'$은 비분기이다.

증명은 정리 2로부터 쉽게 따른다.

\src{213}
이 명제와 보조정리 1을 함께 쓰면, 분리가능 섬유를 갖는 고유 평탄
사상들의 호모토피론적 연구는 사실상 $f_*(\Osh_X)=\Osh_Y$인 경우로
환원된다. 실제로 스타인 인수분해를 사용하여 $Y$를 $Y'$으로 바꾸면 된다.

\result{비고}{} 유한형 평탄 사상으로서 그 섬유들이 분리가능 스킴
(각각 비특이 스킴)인 것을 \emph{분리가능}(각각 \emph{매끄러운})
사상이라 한다. $f$가 평탄하고 $f^{-1}(y)$가 분리가능(각각 비특이)이면,
$f^{-1}(y)$의 어떤 근방 위에서 $f$가 분리가능(각각 매끄러운) 사상이
됨을 보일 수 있다. 같은 결과가 “절대 정규”에 대해서도 성립한다
(“베르티니 정리”).

$f:X\longrightarrow Y$를 다음 조건을 만족하는 고유 사상이라 하자.
\begin{equation}
  \tag{i}
  f_*(\Osh_X)=\Osh_Y.
\end{equation}
또 $X'$을 $X$ 위에서 유한인 스킴이라 하자. $X'\longrightarrow Y$의
스타인 인수분해에 대응하는 $Y$의 피복을 $Y'$이라 하자(정리 5 참조).
$y\in Y$라 하자. 그러면 $y$ 위의 $X'$의 섬유 $F'$의 연결 성분들의
집합은 $y$ 위에 놓인 $Y'$의 점들의 집합과 동일시된다(정리 5). 자연스러운
사상 $X'\longrightarrow X$와
$X'\longrightarrow Y'$로부터 얻어지는 자명한 사상
\begin{equation}
  \tag{*}
  X'\longrightarrow X\times_Y Y'
\end{equation}
을 생각하자.\footnote{인쇄본에서는 두 번째 $X'$의 프라임을 빠뜨렸다.
$Y'$로 가는 사상은 $X'\to Y$의 스타인 인수분해에서 오는 사상이다.}
$Y$의 비분기 피복 $Y''$에 대하여 $X'=X\times_Y Y''$ 꼴이면 이 사상은
언제나 동형이고, 이때 $Y'=Y''$이며 (*)는 항등사상이다. 우리가 구하려는
것은 바로 $X'$이 이와 같은 꼴, 즉 $Y'$이 비분기이고 (*)가 동형이 되기
위한 조건이다. 이를 위하여 $y$에서 $X$의 섬유를 $F$라 하자. 이는
$\kappa(y)$ 위의 고유 스킴이며 $F'$은 $F$의 피복이다($X'$이 $X$ 위에서
비분기이면 이 피복도 비분기이다). $F'_1$을 $y$ 위에 놓인 $Y'$의 한 점
$y'_1$에 대응하는 $F'$의 연결 성분이라 하자. 다음 조건들을 더 가정한다.
\begin{enumerate}
  \item[\textit{(ii)}] $X'\longrightarrow X$는 $F'_1$의 점들에서
    비분기이다. 따라서 $F'_1$은 $F$의 비분기 피복이다.
  \item[\textit{(iii)}] $F'_1$은 $F$의 기하적으로 자명한 피복이다
    (보조정리 3 참조).
\end{enumerate}

\result{정리}{11} 이 조건들 아래에서 $Y'$ 안의 $y'_1$의 열린 근방
$U'$이 존재하여 (*)는 $U'$ 위에서 동형이다. 더 나아가
$Y'\longrightarrow Y$는 $y'_1$에서 비분기이다(그러나 $y$ 위에 놓인
$Y'$의 다른 점 $y'$들에서는 분기될 수 있다).

물론 조건 \textit{(ii)}와 \textit{(iii)}는 정리의 결론이 성립하기
위한 필요조건이기도 하다. 증명은 보조정리 3과 정리 2를 쓰면 쉽게 따른다.

\src{214}
\result{따름정리}{1} 여전히 \textit{(i)}가 성립한다고 가정하자. $X$ 위의
비분기 피복 $X'$이 $Y$의 비분기 피복 $Y'$의 역상과 동형일 필요충분조건은
$X'$이 각 섬유 $f^{-1}(y)$ 위에 기하적으로 자명한 피복을 유도하는 것이다.

정리 11에 따라 이 조건을 만족하는 $Y$의 점들의 집합은 열려 있으므로,
닫힌점 $y$들에 대해서만 확인하면 충분하다\ldots{} 따름정리 1과 동치인
다음 명제도 기록해 두자. (보조정리 4에 따라 전사인) 준동형
$\pi_1(X)\longrightarrow\pi_1(Y)$의 핵은
$\pi_1(\overline{f^{-1}(y)})$들의 $\pi_1(X)$ 안에서의 상들로 생성되는
닫힌 정규부분군이다. 여기서
\[
  \overline{f^{-1}(y)}
  =f^{-1}(y)\otimes_{\kappa(y)}\overline{\kappa(y)},
\]
$\overline{\kappa(y)}$는 $\kappa(y)$의 대수적 폐포를 나타낸다. 모든
섬유에 같은 기점을 택할 수 없으므로,
$\pi_1(\overline{f^{-1}(y)})\longrightarrow\pi_1(X)$인 준동형들은
(일단 $X$의 기점을, 따라서 $Y$의 기점도 택하고 나면) 어쨌든
$\pi_1(X)$의 내부 자기동형사상을 뒤에 합성하는 차이까지만 정해진다는
점에 주의하자.

\result{따름정리}{2} 정리 11의 일반 가정 아래에서, 더 나아가
$Y,X,X'$이 정역 스킴이라고 가정하고 그 유리함수체를 각각 $K,L,L'$이라
하자. 그러면 $L'$ 안에 $K$의 분리 중간확대체 $K'$가 존재하여, $K'$과
$L$은 $K$ 위에서 선형적으로 서로소이고 $L'=LK'$이다(따라서
$L'=L\otimes_K K'$이다).

(보조정리 3의 마지막 부분을 $X$의 기하적 섬유에 적용한다.) 정리 11의
가장 흥미로운 적용 경우는 $f$가 분리 가능한 사상일 때 얻어진다. 이때
$X'$도 $Y$ 위에서 분리 가능하다. 따라서 명제 4에 의해 $Y'$은 $Y$ 위에서
비분기이고, 그러므로 (*)의 우변 $X\times_Y Y'$은 $X$ 위에서 비분기이다.
이에 따라 쉽게 다음을 얻는다.

\result{따름정리}{3} \textit{(i)}에 더하여 $f$가 분리 가능하다고 가정하자.
$X'$을 $X$의 연결인 비분기 피복이라 하자. $X'$\footnote{인쇄본에는
여기서 ``$X$''라고 되어 있으나, 이 판정의 주어는 앞 문장에서 도입한
피복 $X'$이다.}이 $Y$의 비분기 피복 $Y'$의 역상일 필요충분조건은, 한
기하적 섬유 $\overline F=f^{-1}(y)$ 위에 유도된 피복 $\overline F'$이
정칙 단면을 갖는 것이다.

$\overline F'$이 $\overline F$ 위에서 기하적으로 자명하다고 가정할 필요는
없었다는 점에 주의하자(이것은 \textit{a posteriori} 참이지만,
\textit{a priori}에는 훨씬 더 강한 조건이다). 실제로 따름정리 3은 다음
명제와 동치이다.

\src{215}
\result{따름정리}{4} $f:X\longrightarrow Y$를
$f_*(\Osh_X)=\Osh_Y$를 만족하는 고유이고 분리 가능한 사상이라 하자.
$y\in Y$의 기하적 섬유를 $\overline F$라 하고, $\overline F$ 안에서
기하적 점 하나를 택하자. 그러면 사상들
$\overline F\longrightarrow X\longrightarrow Y$에 의해 $X,Y$에서도
기하적 점들이 정해지며, 이들을 $\overline F,X,Y$의 기본군의 기점으로
삼는다. 이 조건들 아래에서 다음 완전열을 얻는다.
\[
  \boxed{
    \pi_1(\overline F)\longrightarrow\pi_1(X)
    \longrightarrow\pi_1(Y)\longrightarrow e
  }.
\]

이로부터 정규성에 관한 어떤 가정도 필요 없는 다음 두 세르--랭 명제를
쉽게 얻는다.

\result{따름정리}{5} $X,Y$를 체 $k$ 위의 연결인 스킴들이라 하자. 축소
스킴 $X_{\mathrm{réd}}$가 $k$ 위에서 분리 가능하다고 가정한다($k$가
완전체이면 이는 자동으로 참이다). 또한 $X$ 또는 $Y$가 $k$ 위에서
고유라고 가정한다.\footnote{1962년 정오표에서 추가된 조건이다.} $X$(각각
$Y$) 안의 기하적 점 $a$(각각 $b$)를 택하면 $X\times_kY$ 안의 기하적 점
$c=(a,b)$와 자연스러운 사상
\[
  \pi_1(X\times_kY,c)\longrightarrow
  \pi_1(X,a)\times\pi_1(Y,b)
\]
을 얻는다(이는 $\pi_1(X\times Y,c)$에서 $\pi_1(X,a)$와
$\pi_1(Y,b)$로 가는 $\pi_1$의 함자적 사상들로부터 유도된다). 이 사상은
단사이고, $k$가 대수적으로 닫혀 있으면 전단사이기까지 하다.

(마지막 경우의 전사성은 거의 자명하다.) 이로부터 세르--랭과 함께 다음을
얻는다.

\result{따름정리}{6} $X$를 대수적으로 닫힌 체 $k$ 위의 연결이고 고유인
대수적 스킴이라 하고, $K$를 $k$의 대수적으로 닫힌 확대체라 하자. 그러면
$X$와 $X\otimes_kK$의 기본군은 서로 같다. 즉 뒤의 스킴의 모든 비분기
피복은 $X$의 한 비분기 피복으로부터 스칼라를 확대하여 얻어지며, 그 피복은
동형을 제외하고 유일하다.

\result{비고}{}

1) 명제 4를 사용하면 따름정리 4의 가정 $f_*(\Osh_X)=\Osh_Y$가 본질적이지
않음을 알 수 있다. 일반적인 경우에는 $\pi_1(Y)$ 뒤에 자명군 $e$를 놓는
대신 다음 열을 이어야 한다.
\[
  \pi_0(\overline F)\longrightarrow\pi_0(X)
  \longrightarrow\pi_0(Y)\longrightarrow e.
\]
이는 대수적 위상수학에서와 같다.

\src{216}
2) 일반적으로는 현재로서
$\pi_1(\overline F)\longrightarrow\pi_1(X)$의 핵에 관해 아무것도 말할 수
없으며, 여기에는 $\pi_2(Y)$가 관여해야 할 것이다. 그러나 $Y$가 국소환
$A$의 스펙트럼이면 아래의 정리 12($A$가 완비일 때 그렇다고 주장하는
정리)를 이용하여 $\pi_1(\overline F)\longrightarrow\pi_1(X)$가 단사임을
증명할 수 있을 것으로 보인다.

정리 11에서는 정리 1과 정리 2만을 사용하였다. 이제 다음의 초등적인
보조정리를 이용하여 정리 3을 적용하겠다.

\result{보조정리}{6} $X$를 스킴, $X_0$를 이에 대응하는 축소 스킴
(즉 멱영원들을 0으로 만든 스킴)이라 하자. 그러면 $X_0$의 모든 비분기
피복 $X'_0$은 $X$의 비분기 피복 $X'$에서 유도되며, $X'$은 표준 동형을
제외하고 결정된다.

순전히 국소적인 성질의 이 보조정리는 여기서 모듈라이 이론의 정리 8과
비슷한 구실을 한다. 이를 존재 정리(정리 3)와 결합하면 다음을 얻는다.

\result{정리}{12} $A$를 잉여체 $k$를 갖는 완비 뇌터 국소환이라 하고,
$X$를 $A$ 위에서 고유인 스킴이라 하자. 그러면
$X_0=X\otimes_A k$의 모든 비분기 피복 $X'_0$은 $X$의 비분기 피복
$X'$에서 유도되며, $X'$은 동형을 제외하고 유일하다.

다시 말하면 다음과 같다.

\result{따름정리}{1} $X_0$의 한 기하적 점을 $X_0$와 $X$의 기본군을 위한
기점으로 택하자. 그러면 표준 준동형
$\pi_1(X_0)\longrightarrow\pi_1(X)$는 동형이다.

이제 $X_0$에 보조정리 5를 적용하자(간단히 하기 위하여
$H^0(X_0,\Osh_{X_0})=k$라고 가정한다). 또한 $A$가 완비이므로 그 비분기
확대들은 잉여체의 비분기 확대들과 대응한다. 즉
$Y=\operatorname{Spec}(A)$라 할 때 $\pi_1(Y)=\pi_1(k)$이다. 이에 따라
다음 완전열을 얻는다.
\[
  e\longrightarrow\pi_1(\overline X_0)
  \longrightarrow\pi_1(X)\longrightarrow\pi_1(Y)\longrightarrow e.
\]

\result{따름정리}{2} $f:X\longrightarrow Y$를 고유 평탄 사상이라 하고,
$y_1$을 $Y$의 한 점, $y_0$를 $y_1$의 한 특수화라 하자. 이에 대응하는
``기하적'' 섬유 $\overline X_1$과 $\overline X_0$를 생각하고,
$\overline X_0$가 분리 가능하고 연결이라고 가정하자(그러면
$\overline X_1$도 같은 조건들을 만족한다). 이때 내부 자기동형을 제외하고
정의되는 군 준동형
$\pi_1(\overline X_1)\longrightarrow\pi_1(\overline X_0)$을 택할 수 있으며,
이 준동형은 전사이다.

\src{217}
이 준동형이 언제나 전단사이기를 기대할 수 있겠지만, 안타깝게도
$\kappa(y_0)$의 표수가 $>0$이면 일반적으로 그렇지 않다. 그러나 아래에서
(적어도 $\overline X_0$가 비특이일 때) 이 준동형의 핵에 대한 상계를 얻을
것이며, 이로부터 $\kappa(y_0)$의 표수가 $0$이면 위 준동형이 전단사임이
따른다(이는 초월적 방법으로도 증명할 수 있는 결과이다). 어쨌든 이미
일반 섬유의 $\pi_1$을 이용하여 특수 섬유의 $\pi_1$에 대한 상계를 얻는다.
예를 들어 표수 $p$의 대수곡선을 표수 $0$의 곡선으로 들어 올릴 수 있다는
사실(정리 9, 따름정리 4)을 이용하면, 초월적 방법으로 다음을 얻는다.

\result{따름정리}{3} $X_0$를 임의의 표수를 갖는 대수적으로 닫힌 체 위의
완비 비특이 곡선의 스킴이라 하고, $g$를 $X_0$의 종수라 하자. 그러면
$\pi_1(X_0)$는 잘 알려진 관계를 만족하는 $2g$개의 위상적 생성원을 갖는다.

잘 알려진 초평면 절단 기법으로부터 다음을 얻는다.

\result{따름정리}{4} $X$를 임의의 표수를 갖는 대수적으로 닫힌 체 위의
비특이 사영 스킴이라 하자. 그러면 $\pi_1(X)$는 유한 개의 위상적 생성원을
갖는다.

준동형 $\pi_1(\overline X_1)\longrightarrow\pi_1(\overline X_0)$의 핵을
구해 보자. 이를 위해 $Y$를 완비 이산 값매김환 $V=\Osh_y$의 스펙트럼이라고
가정해도 된다($y=y_0$). 문제는 다음과 동치이다. $\overline X_1$의 비분기
피복 $\overline X'_1$(원한다면 갈루아 피복이라고 가정해도 된다)가 주어졌을
때, 그것이 어떤 조건 아래 $X_0$의 비분기 피복에서 유도되는지를 결정하는
것이다. \textit{A priori}, 주어진 피복은 $V$의 분수체 $K$의 대수적 폐포
$\overline K$ 안의 유한 부분확대 $K'$에 대하여
$X_1\otimes_K K'$의 비분기 피복 $X'_1$을 스칼라 확대하여 얻어진다.
$\overline X'_1$이 군 $G$를 갖는 갈루아 피복이면, $X'_1$도 군 $G$를 갖는
갈루아 피복으로 택할 수 있다. 그러면 다음이 성립한다.
$X'_1\otimes_{K'}\overline K=\overline X'_1$이 $X_0$의 비분기 피복에서
유도될 필요충분조건은 $\overline K$ 안에 $K'$의 유한 확대 $K''$이 존재하여
$X''_1=X'_1\otimes_{K'}K''$이 $X''\otimes_{V''}K''$ 꼴이 되는 것이다.
여기서 $V''$은 $K''$ 안에서 $V$의 정수적 폐포이고, $X''$은
$X\otimes_VV''$의 비분기 피복이다. 예를 들어 $X_0$가 절대적으로
정규라고 가정하자. 그러면 $X\otimes_VV''$은 정규이다($V''$ 위에서
평탄하고 특수 섬유가 정규이기 때문이다). 그 유리함수체는
\[
  X_1\otimes_KK''
  =(X\otimes_VK)\otimes_KK''
  =X\otimes_VK''
  =(X\otimes_VV'')\otimes_{V''}K''
\]
의 유리함수체 $K''(X_1)$와 같다.

$L''=K''(X'_1)$을 $X'_1\otimes_{K'}K''$의 유리함수체라 하자. 이는
$K''(X_1)$의 유한 분리 확대이고, 위 조건은 $L''$이
$X\otimes_VV''$의 유리함수체의 비분기 확대라는 뜻이기도 하다(즉
$L''$ 안에서 $X\otimes_VV''$을 정규화한 것은 $X\otimes_VV''$ 위에서
비분기이다). 사실 $L''$이 $X\otimes_VV''$의 특수 섬유의 점들에서
비분기임을 확인하는 것으로 충분하다(일반 섬유
$X_1\otimes_KK''$ 위에서는 이미 비분기이기 때문이다). 이제 $X_0$가
비특이라면 Nagata--Zariski의 ``순수성 정리''에 따라, $L''$이
$X\otimes_VV''$의 특수 섬유의 일반점에서의 국소환 $O''$ 위에서
비분기임을 확인하는 것만으로도 충분하다. 이 환은
$K''(X_1)$ 안에서 $X$의 특수 섬유의 일반점에서의 국소환
$O\subset K(X_1)$의 정수적 폐포인 이산 값매김환이다. 따라서 다음을
얻었다.

\src{218}
\result{따름정리}{5} 앞의 조건과 표기 아래에서,
$\overline X_1=X_1\otimes_K\overline K$의 비분기 피복
$X'_1\otimes_{K'}\overline K$가 $\overline X_0$의 비분기 피복에서
유도될 필요충분조건은, $K''(X'_1)$이
$K''(X_1)$의 이산 값매김환 $O''\subset K''(X_1)$ 위에서 비분기가
되게 하는 $\overline K/K'$의 유한 중간 확대체 $K''$가 존재하는 것이다.

이제 $O''$는 $K''(X_1)$ 안에서 이산 값매김환
$O'\subset K'(X_1)$의 정수적 폐포이고($O'$은 $K'(X_1)$ 안에서
$O$의 정수적 폐포이다), $O'$은 $K'$ 안에서 $V$의 정수적 폐포인
$V'$을 포함하며, $V'$의 균등화원 $u$는 $O'$의 균등화원이기도 함을
주목하자. 이제 $X'_1$이 갈루아군 $G$를 갖는 갈루아 피복이고 그 위수
$n$이 $\kappa(y_0)$의 표수 $p$와 서로소라고 하자(이 표수는 $O'$의
잉여체의 표수이기도 하다). 그러면 $K'(X'_1)$은 $O'$에 관하여
``순분기''이고, 이로부터 쉽게 다음 사실이 따른다(``아브얀카르
보조정리''). $O'$의 한 균등화원의 $n$제곱근 $v$를 첨가하면 이 체는
$K'(X_1)(v)$ 안에서 $O'$의 정수적 폐포 위에서 비분기가 된다. 그런데
$v$는 $V'$의 한 균등화원의 $n$제곱근으로 택할 수 있으므로,
따름정리 5의 조건이 성립한다. (아브얀카르 보조정리와 순수성 정리를
사용할 수 있다는 이 가능성은 세르가 내게 팔았다.) 얻은 결과를
표현하기 위하여, 모든 완전 비연결 콤팩트군 $\pi$에 대하여
$\widetilde\pi$를 $\pi$의 $p$-실로우 부분군들로 생성되는 닫힌 부분군에
의한 몫군, 즉 위수가 $p$와 서로소인 $\pi$의 이산 몫군들의 사영극한이라
하자. 이 표기 아래에서 다음을 얻는다.

\result{정리}{13} $f:X\longrightarrow Y$를 고유하고 평탄한 사상이라
하고, $y_1$을 $Y$의 한 점, $y_0$을 $y_1$의 한 특수화라 하자.
$\overline X_0$가 연결이고 비특이라고 가정한다. 그러면 정리 12의
따름정리 2에 나오는 전사 준동형에서 유도되는 준동형
$\widetilde\pi_1(\overline X_1)\longrightarrow
\widetilde\pi_1(\overline X_0)$은 동형이다.

다시 말해 다음과 같다.

\src{219}
\result{따름정리}{1} 갈루아군의 위수가 $\kappa(y_0)$의 표수 $p$와
서로소인 비분기 갈루아 피복들의 분류는 $\overline X_0$와
$\overline X_1$에 대하여 동일하다.

특히 $\kappa(y_0)$의 표수가 0이면, 대수적 방법으로
$\pi_1(\overline X_1)\longrightarrow\pi_1(\overline X_0)$이
전단사임을 얻는다.

끝으로 사용한 기법들은 정리 13보다 더 일반적인 다음 결과도 준다.

\result{정리}{14} $f:X\longrightarrow Y$를 고유하고 매끄러운 사상이라
하고, $D$를 $Y$ 위에서 매끄럽고 모든 점에서 여차원이 1인 $X$의 닫힌
부분스킴이라 하자. $f$의 섬유 $Z=f^{-1}(z)$에 대하여
$Z'=Z-Z\cap D$라 놓고, $\pi_1^t(Z')$를 $\overline{Z\cap D}$ 위에서
순분기인 $Z'$의 비분기 피복들을 분류하는 기본군 $\pi_1(Z')$의 몫군이라
하자. $y_0,y_1$을 정리 13에서와 같이 택하자. 그러면 내부 자기동형까지
정의되는 전사 준동형
\[
  \pi_1^t(\overline X'_1)\longrightarrow\pi_1^t(\overline X'_0)
\]
이 존재하고, 이에 대응하는 준동형
\[
  \widetilde\pi_1^t(\overline X'_1)
  \longrightarrow\widetilde\pi_1^t(\overline X'_0)
\]
은 동형이다.

이로부터 정리 13의 따름정리들과 정리 12의 따름정리 4에 대응하는
변형들이 따른다. 마찬가지로 정리 9의 따름정리 3을 사용하면 초월적
방법으로 다음을 얻는다.

\result{따름정리}{} $X_0$를 임의 표수의 대수적으로 닫힌 체 위의 완비
비특이 곡선의 스킴이라 하고,
$S=(s_i)_{1\le i\le n}$을 $X_0$의 $n$개 원소로 이루어진 유한
부분집합이라 하자. 그러면 $\pi_1^t(X_0-S)$는 $2g+n$개의 위상적 생성원
$x_i,y_i$ ($1\le i\le g$)와 $\sigma_j$ ($1\le j\le n$)를 가지며,
이들은 다음 관계를 만족한다.
\[
  \left(\prod_i x_i y_i x_i^{-1}y_i^{-1}\right)
  \sigma_1\cdots\sigma_n=1,
\]
여기서 $\sigma_j$들은 $s_j$에 대응하는 관성군의 생성원이다. 표수와
서로소인 위수를 갖고 앞의 관계를 만족하는 원소
$x_i,y_i,\sigma_j$들로 생성되는 모든 유한군 $G$에 대하여, 군이
$G$이고 점 $s_j$에서의 관성군이 $\sigma_j$로 생성되는
$X_0-S$의 비분기 갈루아 피복이 존재한다.

$X_0$의 종수가 0이고 $n=3$이면, 적어도 위수가 표수와 서로소인 갈루아
피복에 대해서는 ``세 점 문제''의 해를 얻는다. (여기서는 정리 9가
실제로 필요하지 않으며, 한편 앞의 따름정리를 세 점 문제에서 다루는
특수한 경우로부터 이끌어 낼 수 있을 것으로 보인다.)

\src{220}
\result{비고}{}

1) $X,X_0,X_1$의 일반화된 갈루아 피복(갈루아군이 무한소적일 수도
있는 피복)을 아마도 사용하게 될 더 완전한 연구를 통하여 정리 12의
따름정리 2에서 핵을 복원할 수 있을 것이다. 반면 순분기가 아닌 분기를
허용하는 피복을 연구하는 일은 훨씬 더 어려울 것으로 보인다.

2) 보조정리 6을, 비특이 사영 스킴을 한 초평면 단면을 따라 형식적으로
완비화한 것에 관한 그라우어트의 결과(또는 정리 11 뒤의 비고 2에서
언급했으나 아직 증명되지 않은 정리)와 함께 사용하면, ``추상''
대수기하학에서도 기본군에 관한 고전적인 렙셰츠 정리를 증명할 수 있다.

\begin{center}
  \textbf{참고문헌}
\end{center}

\begin{description}
  \item[{[1]}] Dieudonné (J.) 및 Grothendieck (A.). --- \emph{Éléments de
    géométrie algébrique}, Publications mathématiques de l'Institut des Hautes
    Études Scientifiques (출간 예정).
  \item[{[2]}] Grothendieck (Alexander). --- \emph{The cohomology theory of
    abstract algebraic varieties}. International Congress of Mathematicians
    [1958. Edinburgh] (출간 예정).
  \item[{[3]}] Serre (Jean-Pierre). --- \emph{Faisceaux algébriques
    cohérents}, Annals of Math., Series 2, 제61권, 1955, 197--278쪽.
  \item[{[4]}] Serre (Jean-Pierre). --- \emph{Géométrie algébrique et
    géométrie analytique}, Ann. Institut Fourier Grenoble, 제6권, 1955--1956,
    1--42쪽.
\end{description}

\par\smallskip
\begingroup\footnotesize
\noindent\textit{인쇄문에 관한 주. 따름정리6에는 $X$의 고유성 가정이 필요하다. 명제2와 그에 따른 자동사상 서술에는 $\operatorname{rang}(E)\geq2$라는 가정이 필요하다. 이 두 가정은 인쇄문에 명시되어 있지 않다.}\par
\endgroup

\par
\renewcommand{\sourceunit}{e182}
\clearpage
\chapter*{강연 182의 정오표}
\addcontentsline{toc}{chapter}{강연 182의 정오표}
\setcounter{footnote}{1}
\src{298}

\noindent\textbf{2. 형식기하학과 대수기하학.}

\begin{center}
[부르바키 세미나, 제11권, 1958/59, 제182호, 28쪽.]
\end{center}

\noindent\textbf{8쪽, 정리 6.} --- ``$f$가 준사영이다''라는 가정은
더 약한 ``$f$가 분리 사상이다''라는 가정으로 바꿀 수 있다. 이는 다음
결과에 따른 것이다(cf. SGA VIII, 6.2): 모든 준유한 분리 사상
$f:X\longrightarrow Y$는 준사영이다.

\src{299}
\noindent\textbf{12쪽 및 14쪽, 주 1.} --- J.-P.~세르는 표수가 $p>0$인
대수적으로 닫힌 체 $k$ 위의 차원 3인 비특이 사영 다양체를 구성했다. 이
다양체는 잉여체가 $k$이고 분수체의 표수가 0인 정역 국소환 위의 고유 스킴을
환원하여 얻어지는 것이 아니다.\footnote{Serre (Jean-Pierre). --- Exemples de
variétés projectives en caractéristique $p$ non relevables en caractéristique
zéro, \emph{Proc. Nat. Acad. Sc. U.S.A.}, t. 47, 1961, p. 108--109.} 멈퍼드도
비특이 사영 곡면인 이와 유사한 예를 찾았다고 한다.

\noindent\textbf{15쪽, 주 2와 3.} --- 나는 현재 여기서 추측한 결과들에
관해서는 덜 낙관적이다. 그러나 주 3의 말미에서 제기한 사영공간의 구조
변형에 관한 문제는 히로나카가 긍정적으로 해결했으며, 아벨 다양체에 관한
유사한 문제는 고이즈미가 해결했다.

\noindent\textbf{16쪽, 3행.} --- 멈퍼드는 최근 종수 $g$의 곡선들에 대한
모듈라이 스킴을 구성했다(cf. SHMI). 또한 정리 10은 모듈라이 이론의
``준위 $n$의 야코비 스킴''이 비특이임을 증명한다(더 나아가 이들은
$\mathbf Z$ 위에서 매끄럽기까지 하다).

\noindent\textbf{23쪽, 12행, 따름정리 5의 제1행.} --- 다음을 덧붙인다:
``$X$ 또는 $Y$가 $k$ 위에서 고유인 경우''.

§§ 1--5의 내용은 EGA III의 출판된 부분에 들어 있고, §§ 6과 7의 내용은
SGA III에 들어 있다. 기본군의 연구에 관해서는 SGA V, IX, X, XI 및
SGA 1962(강연 X, XII, XIII)를 보라. 여기에는 르프셰츠 유형의 정리들과
수많은 열린 문제가 다루어져 있다. 순하게 분기된 피복의 이론(cf. 정리 14)만
아직 정식으로 집필되지 않았다. 정리 14의 따름정리는 대수적으로 닫힌 체
위의 대수곡선에 대해, 그 표수와 서로소인 위수를 갖는 갈루아 피복을 완전히
결정하며, 다음 세 경우에 본질적으로 사용되었다.

1\textsuperscript{o} 임의의 표수에서 비특이 사영 곡면에 대한 피카르
부등식을 이구사가 증명할 때;

2\textsuperscript{o} 임의의 표수에서 일변수 함수체 위에 정의된 아벨
다양체를 구조군으로 하는 동차 주다발들의 군을 연구할 때(오그와
샤파레비치가 서로 독립적으로 발전시켰다);

3\textsuperscript{o} 대수다양체의 ``베유 코호몰로지''에 관한 몇몇 핵심
정리를 아르틴이 최근 증명할 때.

\par
\renewcommand{\sourceunit}{190}
\unit{강연 190 — 하강 기법 I}
\sid{190}
\src{299}
\seminar{제12년차, 1959/60, 제190호\\1959년 12월}
\paperhead{대수기하학에서의 하강 기법과 존재 정리\\
I. 일반론. 충실 평탄 사상에 의한 하강}%
  {알렉상드르 그로텐디크 지음}

기술적 관점에서 보면, 이 강연과 그 뒤를 이을 강연들은
힐베르트의 유명한 ``정리 90''을 변주한 것으로 여길 수 있다.
대수기하학에 하강 방법을 도입한 것은 ``기초체의 하강''이라는 이름 아래
A. 베유가 한 것으로 보인다. 다만 베유는 체의 유한 분리 확대인 경우로
한정한다. 높이 1인 라디시엘 확대의 경우는 뒤에 P. 카르티에가 다루었다. 스킴의
언어가 없었고, 특히 고려가 허용된 환들에 멱영원이 없었기 때문에,
이 두 경우 사이의 본질적인 동일성을 카르티에는 정식화할 수
없었다.

현재로서는, 여기서 설명할 일반적인 하강 기법은 (필요한 경우 ``형식기하학''의
기본 정리들, [3] 참조, 과 결합하면) 대수기하학의 존재 정리 대부분의
기초를 이룬다고 생각된다.\footnote{1962년의 정오표는 이 서술이
지나치다고 판단한다. 이 서술은 TDTE I--VI 계열의 처음 두 강연에서 다룬
비사영적 기법들에는 상당한 정도로 타당하지만, 사영적 기법들
(TDTE IV, V, VI)에는 그렇지 않다.} 더 나아가 이 기법은 확실히
``해석기하학''으로 옮길 수 있음에 유의해야 한다. 가까운 장래에
전문가들이 강연 II에서 제시할 형식기하학의 존재 정리들의
``해석적'' 대응물들을 증명할 수 있으리라 기대할 수 있다. 어쨌든
코다이라--스펜서의 최근 연구에서 쓰인 방법들은 특이점 근방에서
``모듈라이 다양체''를 정의하고 연구하기에는 부적합해 보인다. 이는 스킴
이론에 더 가까운 방법들(물론 초월적 기법들이 이를 보완해야 한다)이
필요함을 아주 분명히 보여 준다.

이 강연 I에서는 (제목에 명시한) 가장 초보적인 하강의 경우를
다루겠다. 그럼에도 아래 정리 1, 2, 3의 응용은 이미 매우 많다.
여기서는 가능한 최대의 일반성을 부여하려 하지 않고, 그 가운데 몇 가지만
예로 드는 데 그쳤다.

스킴의 언어는 자유롭게 사용하겠다. 이에 관해서는 앞서 인용한 강연과
[2]를 참조한다. 다만 이 강연에서 고려하는 준스킴은 반드시
뇌터일 필요가 없고, 사상도

\src{300}
반드시 유한형일 필요가 없음을 명시해 둔다. 따라서 $A$가 국소 뇌터환이고
그 완비화가 $\widehat A$이면, 비뇌터환 $\widehat A\otimes_A\widehat A$와
고려하는 환들 사이의 포함관계에 대응하는 아핀 스킴의 사상들을
고려해야 한다.

\fsection{A}{범주에 관한 예비사항}

\subsection*{1. 섬유화 범주, 하강 데이터, $\mathcal F$-하강
사상}
\addcontentsline{toc}{subsection}{1. 섬유화 범주와 하강 데이터}

\result{a. 정의}{1.1} 밑범주가 $\mathcal C$인 \emph{섬유화 범주}
$\mathcal F$는 범주 $\mathcal C$, 각
$X\in\mathcal C$에 대한 범주 $\mathcal F_X$, 각
$\mathcal C$-사상 $f:X\longrightarrow Y$에 대한 함자
$f^*:\mathcal F_Y\longrightarrow\mathcal F_X$로 이루어진다. 이 함자는
\[
  f^*(\xi)=\xi\times_Y X
\]
로도 쓴다. 여기서 $\xi\in\mathcal F_Y$이고, $X$는 ``$Y$ 위의
$\mathcal C$의 대상'', 즉 사상 $f$가 주어진 대상으로 여긴다. 끝으로
합성할 수 있는 두 사상
$X\xrightarrow{f}Y\xrightarrow{g}Z$에 대하여 함자의 동형사상
\[
  c_{f,g}:(gf)^*\longrightarrow f^*g^*,
\]
이 주어지며, 이 데이터는 다음 조건들을 만족한다.

\begin{enumerate}
\item[(i)] $\operatorname{id}_X^*=\operatorname{id}$ ;
\item[(ii)] $f$ 또는 $g$가 항등 동형사상이면 $c_{f,g}$는
항등 동형사상이다 ;
\item[(iii)] 합성할 수 있는 세 사상
$X\xrightarrow{f}Y\xrightarrow{g}Z\xrightarrow{h}T$가 주어지면,
$c_{u,v}$ 꼴의 동형사상들로 이루어진 다음 동형사상의 도식은
가환이다.
\[
\begin{array}{ccc}
  (h(gf))^*=((hg)f)^* & \longrightarrow & f^*(hg)^* \\
  \big\downarrow && \big\downarrow \\
  (gf)^*h^* & \longrightarrow & (f^*g^*)h^*=f^*(g^*h^*).
\end{array}
\]
\end{enumerate}

\result{예}{1} 섬유곱이 존재하는 범주 $\mathcal C$를 생각하자.
그러면 밑범주가 $\mathcal C$인 섬유화 범주 $\mathcal F$를 다음과 같이
정의한다. $\mathcal F_X$는 $X$ 위의 $\mathcal C$-대상들의 범주이고,
사상 $f:X\longrightarrow Y$에 대응하는 함자
$f^*:\mathcal F_Y\longrightarrow\mathcal F_X$는 섬유곱
$Z\longmapsto Z\times_YX$로 정의한다.

\result{예}{2} $\mathcal C$를 준스킴의 범주라 하고, 각
$X\in\mathcal C$에 대하여 $\mathcal F_X$를 $X$ 위의 준연접
가군층들의 범주라 하자. $f:X\longrightarrow Y$가 준스킴의 사상이면
$f^*:\mathcal F_Y\longrightarrow\mathcal F_X$는 가군층의
역상 함자이다.

\src{301}
이로써 밑범주가 $\mathcal C$인 섬유화 범주를 얻는다.

\result{b. 정의}{1.2} 집합의 사상들로 이루어진 도식
\[
  E\xrightarrow{u}E'
  \overset{v_1}{\underset{v_2}{\rightrightarrows}}E''
\]
이 \emph{완전}하다고 함은, $u$가 $E$에서 $E'$의 원소 $x'$ 가운데
$v_1(x')=v_2(x')$를 만족하는 것들로 이루어진 부분집합 위로의
전단사임을 뜻한다.

\result{정의}{1.3} $\mathcal F$를 밑범주가 $\mathcal C$인 섬유화 범주라
하고, $\mathcal C$ 안의 사상들로 이루어진 도식
\[
  \tag{+}
  S\xleftarrow{\alpha}S'
  \overset{\beta_1}{\underset{\beta_2}{\rightrightarrows}}S''
\]
을 생각하자. $\alpha\beta_1=\alpha\beta_2$라 하자. 이 도식이
$\mathcal F$-\emph{완전}하다고 함은, $\mathcal F_S$의 임의의 두 대상
$\xi,\eta$에 대하여 다음 집합의 사상들의 도식
\[
  \tag{+'}
  \operatorname{Hom}(\xi,\eta)
  \xrightarrow{\alpha^*}
  \operatorname{Hom}(\alpha^*\xi,\alpha^*\eta)
  \overset{\beta_1^*}{\underset{\beta_2^*}{\rightrightarrows}}
  \operatorname{Hom}(\gamma^*\xi,\gamma^*\eta),
\]
이 완전함을 뜻한다. 여기서 $\gamma=\alpha\beta_1=\alpha\beta_2$이다.
마지막 도식에서는 간단히 하기 위해 $c_{\beta_i,\alpha}$를 써서
$\beta_i^*\alpha^*$를 $(\alpha\beta_i)^*=\gamma^*$와 동일시하였다.

\result{정의}{1.4} $\mathcal F$를 밑범주가 $\mathcal C$인 섬유화 범주라
하고, $\mathcal C$ 안에서 $S''$에서 $S'$로 가는 두 사상
$\beta_1,\beta_2$를 생각하자. 또한
$\xi'\in\mathcal F_{S'}$라 하자. $\xi'$ 위의 \emph{붙이기 데이터}
(쌍 $(\beta_1,\beta_2)$에 관한)란 $\beta_1^*(\xi')$에서
$\beta_2^*(\xi')$로 가는 동형사상이다. $\xi',\eta'\in\mathcal F_{S'}$에
각각 붙이기 데이터가 주어졌다고 하자. $\mathcal F_{S'}$의 사상
$u:\xi'\longrightarrow\eta'$가 \emph{붙이기 데이터들과 양립한다}고 함은
다음 도식이 가환임을 뜻한다.
\[
\begin{array}{ccc}
  \beta_1^*(\xi') & \longrightarrow & \beta_2^*(\xi') \\
  \big\downarrow && \big\downarrow \\
  \beta_1^*(\eta') & \longrightarrow & \beta_2^*(\eta').
\end{array}
\]
이와 같이 붙이기 데이터가 주어진 $\mathcal F_{S'}$의 대상들은
(쌍 $(\beta_1,\beta_2)$에 관하여) 하나의 범주를 이룬다. 만일
$\alpha:S'\longrightarrow S$가 $\alpha\beta_1=\alpha\beta_2$를 만족하는
사상이면, 모든 $\xi\in\mathcal F_S$에 대하여 대상 $\xi'=\alpha^*(\xi)$에는

\src{302}
표준 붙이기 데이터가 주어진다. 실제로
\[
  \beta_i^*\alpha^*(\xi)\simeq(\alpha\beta_i)^*(\xi)=\gamma^*(\xi),
\]
이고, 여기서 다시 $\gamma=\alpha\beta_1=\alpha\beta_2$로 놓았다. 더 나아가
$u:\xi\longrightarrow\eta$가 $\mathcal F_S$의 사상이면
\[
  \alpha^*(u):\alpha^*(\xi)\longrightarrow\alpha^*(\eta)
\]
는 $\mathcal F_{S'}$ 안에서 표준 붙이기 데이터들과 양립하는
사상이다. 따라서 $\mathcal F_S$에서 쌍 $(\beta_1,\beta_2)$에 관한
붙이기 데이터가 주어진 $\mathcal F_{S'}$의 대상들의 범주로 가는
표준 함자를 얻는다. 이제 정의 1.3은 다음과 같이 바꾸어 말할 수도 있다.
그 정의의 도식 (+)가 $\mathcal F$-완전하다는 것은 앞의 함자가 ``충만하고
충실하다'', 즉 $\mathcal F_S$와 쌍 $(\beta_1,\beta_2)$에 관한 붙이기
데이터가 주어진 $\mathcal F_{S'}$의 대상들의 범주의 한 부분범주 사이의
동치를 정의한다는 뜻이다.

\result{정의}{1.5} $\xi'\in\mathcal F_{S'}$ 위의 붙이기 데이터가
($\alpha$에 관하여) \emph{유효}하다고 함은, 이 데이터가 주어진 $\xi'$이
어떤 $\xi\in\mathcal F_S$에 대한 $\alpha^*(\xi)$와 동형임을
뜻한다.

도식 (+)가 $\mathcal F$-완전한 경우, 앞의 정의에 나오는 대상 $\xi$는
유일한 동형사상을 제외하고 결정되고, 범주 $\mathcal F_S$는
유효한 붙이기 데이터가 주어진 $\mathcal F_{S'}$의 대상들의 범주와
동치이다.

\medskip\noindent
c. 가장 중요한 경우는
\[
  S''=S'\times_S S',
\]
이고 $\beta_i$들이 $S'\times_S S'$에서 그 두 인자로 가는 두 사영
$p_1,p_2$인 경우이다(이제 $\mathcal C$에 섬유곱이 존재한다고 가정한다).
그러면 $\xi'\in\mathcal F_{S'}$ 위의 붙이기 데이터
$\varphi:p_1^*(\xi')\longrightarrow p_2^*(\xi')$가 유효하기 위한
두 가지 즉각적인 필요조건이 있다. 먼저
\[
  \tag{i}
  \Delta^*(\varphi)=\operatorname{id}_{\xi'},
\]
이다. 여기서 $\Delta:S'\longrightarrow S'\times_S S'$는 대각 사상이고,
$\Delta^*p_i^*(\xi')$는 $(p_i\Delta)^*(\xi')=\xi'$와
동일시하였다. 또한
\[
  \tag{ii}
  p_{31}^*(\varphi)=p_{32}^*(\varphi)\,p_{21}^*(\varphi),
\]

\src{303}
여기서 $p_{ij}$는 $S'\times_S S'\times_S S'$에서 그 $i$번째와
$j$번째 인자의 부분곱으로 가는 표준 사영을 나타낸다.

\result{정의}{1.6} 사상 $\alpha:S'\longrightarrow S$에 관하여
$\xi'\in\mathcal F_{S'}$ 위의 \emph{하강 데이터}란, 표준 사영의 쌍
$(p_1,p_2):S'\times_S S'\longrightarrow S'$에 관한 $\xi'$ 위의
붙이기 데이터로서 위의 조건 (i), (ii)를 만족하는 것을 말한다.

\result{정의}{1.7} 사상 $\alpha:S'\longrightarrow S$를
\emph{$\mathcal F$-하강 사상}이라 함은 사상들의 도식
\[
  S\xleftarrow{\alpha}S'
  \overset{p_1}{\underset{p_2}{\rightrightarrows}}S'\times_S S'
\]
이 $\mathcal F$-완전임을 뜻한다(정의 1.3). 더 나아가
$\mathcal F_{S'}$의 대상 위의 모든 하강 데이터(정의 1.6)가 유효하면
$\alpha$를 \emph{엄밀 $\mathcal F$-하강 사상}이라 한다.

마지막 조건은 좀 더 그림을 그리듯 다음과 같이 말할 수도 있다.
$\mathcal F_S$의 대상을 주는 일과, 하강 데이터가 주어진
$\mathcal F_{S'}$의 대상을 주는 일은 \fquote{동일한 일이다}.

사상 $\alpha:S'\longrightarrow S$가 \fquote{단면}
$s:S\longrightarrow S'$, 즉 $\alpha s=\operatorname{id}_S$를 만족하는
사상 $s$를 가지면, $\alpha$는 엄밀 $\mathcal F$-하강 사상임에
유의하자.\footnote{1962년 정오표는 $\alpha$가 미리
$\mathcal F$-하강 사상이라고 가정할 필요가 없다고 명시한다.}
실제로 $\xi'\in\mathcal F_{S'}$에 $\alpha$에 관한 하강 데이터가 주어지면,
그것은 $\xi=s^*(\xi')$에서 \fquote{온다}.

\medskip\noindent
d. 앞의 개념들은 $S$ 위의 $\mathcal C$의 매개변수 대상 $T$에 대하여
집합
\[
  \operatorname{Hom}_S(T,S')=S'(T)
\]
을 도입하면 더 직관적으로 제시할 수 있다. 그 원소들은 $t,t'$ 등으로
나타내겠다. 대상 $\xi'\in\mathcal F_{S'}$가 주어졌을 때 각
$t\in S'(T)$에는 $\mathcal F_T$의 대상 $t^*(\xi')$가 대응하며, 이를
$\xi'_t$로도 나타내겠다. 그러면 $\xi'$ 위의 붙이기 데이터
($p_1,p_2$에 관한 것)는 $S$ 위의 모든 $T$와 모든 점의 쌍
$t,t'\in S'(T)$에 대하여 동형사상
\[
  \varphi_{t',t}:\xi'_t\longrightarrow\xi'_{t'}
\]
을 주는 것으로 정의된다($T$가 변할 때의 자명한 함자성 조건들을
만족해야 한다). 1.c의 조건 (i), (ii)는 이제 각각
\[
  \tag{i bis}
  \varphi_{t,t}=\operatorname{id}_{\xi'_t}
  \qquad\text{모든 $T$와 모든 $t\in S'(T)$에 대하여},
\]

\src{304}
\[
  \tag{ii bis}
  \varphi_{t,t''}=\varphi_{t,t'}\varphi_{t',t''}
  \qquad\text{모든 $T$와 모든 $t,t',t''\in S'(T)$에 대하여}
\]
로 쓰인다. 실제로 (ii bis)에서 $t=t'=t''$로 놓으면
$\varphi_{t,t}=\varphi_{t,t}^2$을 얻는다. 가정에 따라
$\varphi_{t,t}$가 동형사상이므로 여기서 (i bis)가 따라온다. 따라서
(i bis)는 사실 (ii bis)의 결과이다(그러므로 (i)는 (ii)의 결과이다).
그러나 $\varphi_{t,t'}$들이 동형사상이라고, 즉
$\varphi:p_1^*(\xi')\longrightarrow p_2^*(\xi')$가 동형사상이라고 더는
미리 가정하지 않으면, (ii bis)가 반드시 (i bis)를 함의하지는 않는다.
그럼에도 (ii bis)와 (i bis)를 함께 가정하면 $\varphi_{t,t'}$들은
동형사상이다. 실제로
$\varphi_{t,t'}\varphi_{t',t}=\varphi_{t,t}
=\operatorname{id}_{\xi'_t}$이기 때문이다.

\subsection*{2. 완전 도식과 엄밀 에피사상, 하강 사상. 예}
\addcontentsline{toc}{subsection}{2. 완전 도식과 엄밀 에피사상}

\result{a. 정의}{2.1} $\mathcal C$를 범주라 하자. $\mathcal C$ 안의
사상들로 이루어진 도식
\[
  T\xrightarrow{\alpha}T'
  \overset{\beta_1}{\underset{\beta_2}{\rightrightarrows}}T''
\]
이 \emph{완전}이라 함은 모든 $Z\in\mathcal C$에 대하여 이에 대응하는
집합 사상들의 도식
\[
  \operatorname{Hom}(Z,T)\longrightarrow
  \operatorname{Hom}(Z,T')
  \rightrightarrows\operatorname{Hom}(Z,T'')
\]
이 완전임을 뜻한다(정의 1.2). 이때 $(T,\alpha)$를, 또는 말을
남용하여 $T$를 사상의 쌍 $(\beta_1,\beta_2)$의 \emph{핵}이라 한다.

이 핵은 유일한 동형을 제외하면 명백히 결정된다. $\mathcal C$가
집합의 범주이면 앞의 정의는 정의 1.2와 양립한다. 쌍대적으로
$\mathcal C$ 안의 사상들로 이루어진 도식
\[
  S\xleftarrow{\alpha}S'
  \overset{\beta_1}{\underset{\beta_2}{\rightrightarrows}}S''
\]
의 완전성을 정의한다. 이때 $(S,\alpha)$를 사상의 쌍
$(\beta_1,\beta_2)$의 \emph{여핵}이라 한다.

\result{정의}{2.2} 사상 $\alpha:S'\longrightarrow S$를
\emph{엄밀 에피사상}이라 함은, 그것이 에피사상이고 모든 사상
$u:S'\longrightarrow Z$에 대하여 $u$가
$S'\longrightarrow S\longrightarrow Z$로 인수분해되기 위한 다음
필요조건이 충분조건이기도 함을 뜻한다. 모든 $S''\in\mathcal C$와
$\alpha\beta_1=\alpha\beta_2$를 만족하는 모든 사상의 쌍
$\beta_1,\beta_2:S''\longrightarrow S'$에 대하여
$u\beta_1=u\beta_2$이다.

섬유곱 $S'\times_S S'$이 존재하면, 이는 도식
\[
  S\xleftarrow{\alpha}S'
  \overset{p_1}{\underset{p_2}{\rightrightarrows}}S'\times_S S'
\]

\src{305}
이 완전이라고, 즉 $S$가 쌍 $(p_1,p_2)$의 여핵이라고 말하는 것과
같다. 어떤 경우든 여핵 사상은 엄밀 에피사상이다. 또한 단사사상이기도
한 엄밀 에피사상은 동형사상임에 유의하자. 쌍대 개념인
\emph{엄밀 단사사상}을 전개하는 일은 독자에게 맡긴다.

$\mathcal F$-하강 사상의 개념과 엄밀 에피사상의 개념 사이의 관계를
정확히 밝히기 위해 다음 정의들을 더 도입한다.

\result{정의}{2.3} 사상 $\alpha:S'\longrightarrow S$를
\emph{보편 에피사상}(각각 \emph{보편 엄밀 에피사상})이라 함은,
$S$ 위의 모든 $T$에 대하여 섬유곱 $T'=S'\times_ST$가 존재하고
사영 $T'\longrightarrow T$가 에피사상(각각 엄밀 에피사상)임을 뜻한다.

매우 좋은 범주들, 이를테면 집합의 범주, 위상공간 위 집합층의 범주,
아벨 범주 등에서는 이렇게 도입한 에피성에 관한 네 개념이 일치한다.
반면 준스킴의 범주나 주어진 공집합이 아닌 준스킴 $S$ 위의 준스킴의
범주 같은 범주에서는, $S$ 위에서 유한인 $S$-스킴들만 생각하더라도
네 개념이 모두 서로 다르다.

\result{정의}{2.4} 사상 $\alpha:S'\longrightarrow S$를
\emph{하강 사상}(각각 \emph{엄밀 하강 사상})이라 함은, 그것이
$\mathcal F$-하강 사상(각각 엄밀 $\mathcal F$-하강 사상)임을 뜻한다
(정의 1.7 참조). 여기서 $\mathcal F$는 $\mathcal C$를 밑범주로 하는
섬유화 범주로서, 그 대상은 $\mathcal C$의 대상 위에 놓인
$\mathcal C$의 대상들이다(1번, 예 1).

\result{명제}{2.1} $\mathcal C$에 유한곱과 (유한) 섬유곱이 존재하면,
$\mathcal C$에서 하강 사상과 보편 엄밀 에피사상은 일치한다.

\result{b. 예}{} $\mathcal C$를 준스킴의 범주라 하자.
$S\in\mathcal C$라 하고, $S'$과 $S''$을 $S$ 위에서 유한인 두
준스킴이라 하자. 즉 이들은 $S$ 위의 대수층 $\mathcal A'$,
$\mathcal A''$에 대응하며, 이 층들은 가군층으로서 준연접이고
유한 생성이다($S$가 국소 뇌터이면 연접이다).
$\alpha:S'\longrightarrow S$를 $S'$의 구조 사상이라 하고,
$\beta_1,\beta_2$를 $S''$에서 $S'$으로 가는 두 $S$-사상이라 하자.
이들은 대수 준동형 $\mathcal A'\longrightarrow\mathcal A''$으로
정의되며, 이 준동형들도 $\beta_1,\beta_2$로 나타내겠다. 유한 사상은
닫힌 사상이라는 사실(첫 번째 코언--자이덴베르크 정리)을 이용하면,
$\mathcal C$ 안의 도식
\[
  \tag{+}
  S\xleftarrow{\alpha}S'
  \overset{\beta_1}{\underset{\beta_2}{\rightrightarrows}}S''
\]
이 완전일 필요충분조건은 $S$ 위의 층들의 도식

\src{306}
\[
  \mathcal O_S=\mathcal A\longrightarrow\mathcal A'
  \overset{\beta_1}{\underset{\beta_2}{\rightrightarrows}}\mathcal A''
\]
이 완전인 것임을 쉽게 증명할 수 있다. 특히
$\alpha:S'\longrightarrow S$가 $S$ 위의 대수층 $\mathcal A'$에
대응하는 유한 사상이면, $\alpha$가 엄밀 에피사상일 필요충분조건은
층들의 도식
\[
  \mathcal O_S=\mathcal A\longrightarrow\mathcal A'
  \overset{p_1}{\underset{p_2}{\rightrightarrows}}
  \mathcal A'\otimes_{\mathcal A}\mathcal A'
\]
이 완전인 것이다($\mathcal A\longrightarrow\mathcal A'$가 단사일
필요충분조건은 $\alpha$가 에피사상인 것이다). $S$가 환 $A$를 갖는
아핀 준스킴이면 $S'$도 $A$ 위에서 유한인 환 $A'$를 갖는 아핀
준스킴이다. 이때 $S'\longrightarrow S$가 엄밀 에피사상일
필요충분조건은 $A\longrightarrow A'$가 $A$를 다음 조건을 만족하는
$x'\in A'$들로 이루어진 $A'$의 부분환 위로 보내는 동형사상인 것이다:
\[
  1_{A'}\otimes_Ax'-x'\otimes_A1_{A'}=0.
\]
($A\longrightarrow A'$가 단사일 필요충분조건은 이 사상이
에피사상인 것이다.) 이미 지적했듯이, $S$가 아르틴 국소환의
스킴이더라도 에피사상인 유한 사상 $S'\longrightarrow S$가 반드시
엄밀 에피사상인 것은 아니다. 그러나 $S$가 뇌터 준스킴이면,
에피사상인 모든 유한 사상 $S'\longrightarrow S$는 유한 개의
(역시 유한인) 엄밀 에피사상의 합성임을 증명할 수 있다. 이는 두
엄밀 에피사상의 합성이 반드시 엄밀 에피사상인 것은 아님도 보여 준다.

\result{c}{} $(+)$가 유한 사상들의 완전 도식이면, 준스킴의 모든 평탄
사상 $T\longrightarrow S$에 의한 $(+)$의 역상 도식도 여전히 완전이다.
따라서 $X,Y$가 두 $S$-준스킴이고 $X$가 $S$ 위에서 평탄하면 다음
집합 사상들의 도식이 완전이다. 여기서 $X',Y'$은 $S'$ 위의 $X,Y$의
역상이고 $X'',Y''$은 $S''$ 위의 그 역상이다:
\[
  \operatorname{Hom}_S(X,Y)\longrightarrow
  \operatorname{Hom}_{S'}(X',Y')
  \rightrightarrows
  \operatorname{Hom}_{S''}(X'',Y'').
\]
특히 $\mathcal F$가 준스킴의 범주 $\mathcal C$를 밑범주로 하는
섬유화 범주이고, 각 $X\in\mathcal C$에 대하여 $\mathcal F_X$가
$X$ 위에서 평탄한 $X$-준스킴들의 범주이면, 도식 $(+)$는
$\mathcal F$-완전이다. (평탄성 가정을 하지 않으면 이 결과는 거짓이
된다. 특히 유한 엄밀 에피사상이 반드시 하강 사상인 것은 아니다.)
마찬가지로, $\mathcal F_X$가 준스킴 $X$ 위의 준연접 평탄 가군층들의
범주인 섬유화 범주를 $\mathcal F$라 해도 $(+)$는
$\mathcal F$-완전이다(여기에서도 평탄성 가정은

\src{307}
여기서도 평탄성 가정은 본질적이다). 어느 경우에든 $S'$ 위에서 평탄한
대상에 주어진 붙이기 자료(특히 $S''=S'\times_SS'$일 때의 강하 자료)가
유효한가 하는 문제는 미묘하며(여러 특수한 경우에서 이에 답하는 것이 이
강연들의 주요 목적 가운데 하나이다), 강연자는 모든 유한 엄밀 에피사상
$S'\longrightarrow S$에 대하여 $S'$ 위의 평탄 준연접층에 주어진 모든 강하
자료가 유효한지 알지 못한다($S$가 아르틴 국소환의 스펙트럼이라고 가정하고
랭크 1의 국소 자유층만을 생각하더라도 그러하다). 더 일반적으로 $A$를 환,
$A'$을 $A$-대수(모든 환은 가환이다)라 하고, 사상들의 도식
\[
  A\longrightarrow A'\rightrightarrows A'\otimes_AA'
\]
이 완전하다고 하자. 이는 $A,A'$의 스펙트럼에 대응하는 사상
$S'\longrightarrow S$가 $\mathcal F$-강하 사상이라는 것과도 동치이다.
여기서 $\mathcal F$는 평탄 준연접층들의 섬유화된 범주이다. $M'$을 $A$로의
강하 자료, 즉 1.c의 조건 (i), (ii)를 만족하는 $A'\otimes_AA'$-가군의 동형
\[
  \varphi:M'\otimes_AA'\xrightarrow{\sim}A'\otimes_AM'
\]
을 갖춘 평탄 $A'$-가군이라 하자(이 조건들을 가군의 말로 명시하는 일은
독자에게 맡긴다). 이 자료는 (평탄 준연접층들의 섬유화된 범주에 관하여)
유효한가? $M'$의 원소 $x'$ 가운데
\[
  \varphi(x'\otimes_A1_{A'})=1_{A'}\otimes_Ax'
\]
을 만족하는 것들의 부분집합을 $M$이라 하자. 이는 $M'$의 부분
$A$-가군이다. 표준 단사사상 $M\longrightarrow M'$은 $A'$-가군의 준동형
$M\otimes_AA'\longrightarrow M'$을 정의한다. 그러면 $\varphi$가
유효하다는 것은 다음을 뜻한다. 즉 $M$은 평탄 $A$-가군이고, 앞의
준동형은 동형이다.

\result{비고}{} 앞의 고찰에서는 도식 $(+)$의 사상들에 아무런 평탄성
가정도 두지 않았다. 따라서 강하 기법을 얻기 위해서는 우리가 다루는
$S,S'$ 위의 대상들에 평탄성 가정을 두어야 했다. 제2절에서는
$\alpha:S'\longrightarrow S$에 평탄성 가정을 둘 것이며, 그러면 더 이상
아무 평탄성 조건도 부과되지 않은 $S,S'$ 위의 대상들에 대해서도 강하
기법을 사용할 수 있다. 어느 경우든 평탄성 가정이 개입한다. 이것이
평탄성 개념이 대수기하학에서 중요한 주요 이유 가운데 하나이다(기초체
위에서는 실제로 무엇이든 평탄하므로, 기초체에만 한정하는 동안에는 그
역할이 드러날 수 없었다!).

\src{308}
\subsection*{3. 에탈 사상에 대한 응용}
\addcontentsline{toc}{subsection}{3. 에탈 사상에 대한 응용}

$A$를 국소환이라 하고 $B$를 $A$ 위의 국소대수로서, 그 극대 아이디얼이
$A$의 극대 아이디얼을 유도한다고 하자. 다음 조건들을 만족할 때 $B$가
$A$ 위에서 \emph{에탈}이라고 하겠다(다른 곳에서 쓰이는 ``비분기''라는
말 대신 이 말을 쓴다).
\begin{enumerate}
\item[(i)] $B$는 $A$ 위에서 평탄하다.
\item[(ii)] $B/\mathfrak mB$는 $A/\mathfrak m=k$의 유한 분리 확대체이다
  (여기서 $\mathfrak m$은 $A$의 극대 아이디얼이다).
\end{enumerate}
$A,B$가 뇌터이고 $k$가 대수적으로 닫혀 있을 때, 이는
$A\longrightarrow B$를 연장하는 완비화 사이의 준동형
$\widehat A\longrightarrow\widehat B$가 동형이라는 뜻이다. 유한형 사상
$f:T\longrightarrow S$에 대하여 $\mathcal O_x$가
$\mathcal O_{f(x)}$ 위에서 에탈이면 $f$가 \emph{$x\in T$에서 에탈}이라고,
또는 $T$가 \emph{$x$에서 $S$ 위에 에탈}이라고 한다. 모든 $x\in T$에서
$f$가 에탈이면 $f$가 \emph{에탈}이라고 하거나 $f$를 \emph{에탈 사상}이라
하거나, 또는 $T$가 \emph{$S$ 위에 에탈}이라고 한다. 또한 $S$가 국소
뇌터이면 $f$가 에탈인 $T$의 점들의 집합은 열려 있고, 자리스키의
``주정리''를 쓰면 그러한 점의 근방에서 $T/S$의 구조를 (잘 알려진 꼴의
방정식으로) 더 구체적으로 기술할 수 있다.

$S$가 복소수체 위의 유한형 스킴이면 세르 [5]의 의미에서 이에 대응하는
해석공간 $\overline S$가 있다. 다만 $\overline S$의 구조층에는 멱영원이
있을 수 있으나, 이는 [5]에서 본질적인 것을 바꾸지 않는다. 이때 $f$가
에탈 사상일 필요충분조건은
$\overline f:\overline T\longrightarrow\overline S$가 에탈 사상인 것이다.
즉 $\overline T$의 모든 점이, 그 위에서 $\overline f$가 $\overline S$의 한
열린집합 위의 동형을 유도하는 근방을 가져야 한다. 특히 $S$의 모든
\emph{에탈 덮개} $T$(즉 유한 에탈 사상 $f:T\longrightarrow S$)에는
$\overline S$의 에탈 덮개 $\overline T$가 대응하며, 이것이 연결일
필요충분조건은 $T$가 연결인 것이다 [5]. 또한 $T,T'$이 $S$ 위에 에탈인
두 스킴이면 자연스러운 사상
\[
  \operatorname{Hom}_S(T,T')\longrightarrow
  \operatorname{Hom}_{\overline S}(\overline T,\overline{T'})
\]
이 전단사임을 쉽게 알 수 있다. 다시 말해 $S$ 위에 에탈인 스킴들의
범주에서 $\overline S$ 위에 에탈인 해석공간들의 범주로 가는 함자
$T\longmapsto\overline T$는 ``충실충만''하며, 따라서 첫째 범주와 둘째
범주의 한 부분범주 사이의 범주 동치를 정의한다. 그라우어트--렘메르의
정리 [2]에 따르면 $S$가 정규이면, 이 방법으로 $S$의 \emph{에탈 덮개}들의
범주와 $\overline S$의 (유한) 에탈 덮개들의 범주 사이의 범주 동치를
얻는다. 즉 $\overline S$의 모든 에탈 덮개는 $S$의 어떤 에탈 덮개 $T$에
대하여 $\overline T$와 $\overline S$-동형이다. 이제 \emph{그라우어트--렘메르
정리가 $S$의 정규성 가정 없이도 성립함}을 보이자. 먼저
$S'\longrightarrow S$를 유한 엄밀 에피사상이라 하고, 이 정리가 $S'$에
대하여 증명되었다고 가정하여 $S$에 대해서도 참임을 보이자.

\src{309}
$\mathfrak T$를 $\overline S$의 에탈 덮개라 하고, 그
$\overline{S'}$ 위의 역상 $\mathfrak T'$을 생각하자. 이는
$\mathfrak T$를 정의하는 $\overline S$ 위의 대수층 $\mathcal A$의
역상인, $\overline{S'}$ 위의 연접 해석적 대수층 $\mathcal A'$에
대응한다. 가정에 따라 $\overline{S'}$ 위의 $\mathfrak T'$은 $S'$의
에탈 덮개 $T'$에서 온다. 즉 $\mathcal A'$은 $S'$ 위의 연접 대수층
$A'$에서 온다. 한편 $\mathcal A'$에는
$\overline{S'}\longrightarrow\overline S$에 관한 표준 강하 자료, 즉
\[
  \overline{S'}\times_{\overline S}\overline{S'}
  =\overline{(S'\times_SS')}
\]
위의 두 역상 사이의 동형이 주어져 있다(조건 (i), (ii)를 만족한다).
앞에서 말한 바에 따라 이 동형은 대응하는 대수적 층들 사이의 동형에서
온다. 즉 $S'\longrightarrow S$에 관한 $A'$ 위의 강하 자료에서 온다.
이 마지막 자료가 유효함은 쉽게 확인된다($\mathcal A'$ 위의 자료가
유효하고, 앞 절에서 명시한 것과 같은 강하 자료의 유효성은 관여하는
가군들의 \emph{완비화}에서 국소적으로 판별할 수 있기 때문이다). 따라서
$S$ 위의 연접 대수층 $A$가 얻어지고, 이는 찾고 있던 $S$의 덮개 $T$를
정의한다. 앞의 결과는 $S'\longrightarrow S$가 유한 엄밀 에피사상들의
유한 합성일 때, 다시 말해 (제2절에서 언급한 분해정리에 따라) 임의의
유한 에피사상일 때에도 명백히 성립한다. 따라서 그라우어트--렘메르 정리는
$S$가 \emph{축소} 스킴일 때에도 성립한다. 즉 $\mathcal O_S$에 멱영원이
없다면, 그 정규화 $S'$을 도입하여 이를 알 수 있다. 여기서 일반적인
경우로 넘어가는 것은 쉽다.

유한 엄밀 에피사상에 관한 분해정리와 강하 자료 유효성의 ``형식적'' 성격을
다시 사용하는 전적으로 비슷한 증명으로 다음 결과를 보일 수 있다.
$S$를 국소 뇌터 준스킴이라 하고 $S'\longrightarrow S$를 \emph{유한이고
전사이며 라디시엘인} 사상이라 하자(이는 다음 조건과 같다. 즉 유한형
사상으로서, $S$ 위의 모든 $T$에 대하여
$T'=S'\times_ST\longrightarrow T$가 위상동형사상이어야 한다. 이는
$S'\longrightarrow S$가 \emph{보편 위상동형사상}이라고 말하는 것과도
같다). $S$ 위에 에탈인 모든 $T$에 대하여 그 역상
$T'=T\times_SS'$을 생각하자. 이는 $S'$ 위에 에탈이다. 그러면 함자
$T\longmapsto T'$은 \emph{$S$ 위에 에탈인 준스킴 $T$들의 범주에서
$S'$ 위에 에탈인 준스킴 $T'$들의 범주로 가는 범주 동치}이다. (이를
위해 $S$ 위에 에탈인 두 준스킴 $T_1,T_2$에 대한 사상
\[
  \operatorname{Hom}_S(T_1,T_2)\longrightarrow
  \operatorname{Hom}_{S'}(T'_1,T'_2)
\]
이 전단사라는 쉽게 직접 확인되는 사실과, $S'=(S,\mathcal O_S/\mathcal J)$인
경우에 앞의 정리가 참이라는 사실을 사용한다.

\src{310}
여기서 $\mathcal J$는 $\mathcal O_S$의 멱영 연접 아이디얼층이다
([4], 보조정리 6). 또한 여기서 생각하는 $T$가 $S$ 위에서 유한이라고는
가정하지 않음에 유의하자. 특히 이 결과로부터 사상
$S'\longrightarrow S$가 $S'$의 기본군 $[\,]$에서 $S$의 기본군으로 가는
동형을 유도함을 얻는다(``\emph{준스킴 기본군의 위상적 불변성}'').

\subsection*{4. 1차 코호몰로지와의 관계}
\addcontentsline{toc}{subsection}{4. 1차 코호몰로지와의 관계}

\medskip\noindent
a. 두 대상의 곱이 언제나 존재하는 범주를 $\mathcal C$라 하고
$T\in\mathcal C$라 하자. 공집합이 아닌 모든 유한집합 $I$에 대하여
$T^I$를 생각할 수 있다. $I$를 움직이면 공집합이 아닌 유한집합들의
범주에서 $\mathcal C$로 가는 공변 함자를 얻는다. 즉 $\mathcal C$의
\emph{단체 대상}이라 부를 수 있는 것을 얻으며, 이를 $K_T$로 나타낸다.
이는 $T$에 공변적으로 의존한다. 또한 $u,v$가 $T\longrightarrow T'$인 두
사상이면 이에 대응하는 사상들 $K_T\longrightarrow K_{T'}$은 서로
호모토픽이다. $\operatorname{Hom}(T,T')\ne\varnothing$일 때 $T$가 $T'$을
\emph{지배한다}고 하자. 이는 $\mathcal C$에서 상향 여과적인 준순서
관계이다. 앞에서 $T$가 $T'$을 지배하면 단체 대상의 사상
$K_T\longrightarrow K_{T'}$들의 표준적인 호모토피류가 존재함을 얻는다.
특히 $K_T$와 $K_{T'}$이 각각 서로를 지배하면 $K_T$와 $K_{T'}$은
호모토피 동치이다. 이제 $F$를 $\mathcal C$에서 아벨 범주
$\mathcal C'$로 가는 함자(구체적으로 반변 함자라 하자)라고 하면
\[
  C^*(T,F)=F(K_T)
\]
는 $\mathcal C'$의 코심플리시얼 대상이고, 따라서 잘 알려진 방식으로
$\mathcal C'$ 안의 (공사슬) 복합체를 정의한다. 그 코호몰로지는
\[
  H^*(T,F)=H^*(C^*(T,F))=H^*(F(K_T))
\]
이다(혼동의 염려가 있으면 $H^*$에 $\mathcal C$를 아래첨자로 붙일 수
있다). 이는 $F$에 관한 코호몰로지 함자이며, $T$를 움직일 때의 공변성은
$K_T$에 관하여 말한 것으로부터 따른다. 더 정확히 말해 $F$를 고정하고
$T$를 $\mathcal C$에서 움직이면(지배 관계로 준순서를 주었을 때)
$H^*(T,F)$들은 $\mathcal C'$의 등급 대상들의 귀납계를 이룬다. 특히
$T,T'$이 각각 서로를 지배하면 $H^*(T,F)$와 $H^*(T',F)$는 표준적으로
동형이다.

$\mathcal C$에 섬유곱이 존재한다고 하자. 그러면 고정된
$S\in\mathcal C$에 대하여 위의 것을 $S$ 위의 $\mathcal C$의 대상들의
범주 $\mathcal C_S$에 적용할 수 있다. $\mathcal C_S$에서 생각함을
명시하려면 $C^*(T,F)$와 $H^*(T,F)$ 대신 $C^*(T/S,F)$와
$H^*(T/S,F)$라고 쓴다. 따라서

\src{311}
$C^*(T/S,F)$는 $\mathcal C'$ 안의 공사슬 복합체이고, 그 $n$차 항은
\[
  F(T\times_ST\times_S\cdots\times_ST)
\]
이다(괄호 안에는 $n+1$개의 인자가 있다).

통상과 같이 $\mathcal C'$이 아벨 범주라고 가정하지 않고도
$H^0(T/S,F)$를 정의할 수 있음에 유의하자. 이것은 두 사영
$p_1,p_2:T\times_ST\rightrightarrows T$에 대응하는 두 사상
$F(p_i)$ ($i=1,2$)
\[
  F(T)\rightrightarrows F(T\times_ST)
\]
의 핵(정의 2.1)이 존재할 때 그 핵이다. 특히 자연스러운 사상(\emph{증대}
사상이라 한다)
\[
  F(S)\longrightarrow H^0(T/S,F)
\]
이 있으며, 이는 좋은 경우(특히 $T\longrightarrow S$가 엄밀
에피사상이고 $F$가 ``왼쪽 완전''일 때) 동형이다. 마찬가지로 $F$가 어떤
범주 $\mathcal C''$ 안의 군들의 범주에 값을 가지면 $H^1(T/S,F)$도 정의할
수 있다. $\mathcal C''$이 집합들의 범주인 경우(즉 $F$가 반드시 가환일
필요는 없는 보통 군들의 범주에 값을 갖는 경우), $H^1(T/S,F)$는
$C^1(T/S,F)=F(T\times_ST)$의 원소 $g$ 가운데
\[
  F(p_{31})(g)=F(p_{32})(g)\,F(p_{21})(g)
\]
을 만족하는 것들로 이루어진 부분군 $Z^1(T/S,F)$을 작용자군 $F(T)$으로
나눈 몫이다. 여기서 $F(T)$는 $C^1(T/S,F)$에, 특히 부분집합
$Z^1(T/S,F)$에 다음과 같이 작용한다.
\[
  \rho(g')\cdot g=F(p_2)(g')\,g\,F(p_1)(g')^{-1}.
\]

\medskip\noindent
b. 예를 들어 $\mathcal F$를 $\mathcal C$를 밑으로 하는 섬유화된 범주라
하자. $\xi,\eta\in\mathcal F_S$라 하고, $S$ 위의 모든 $S'$에 대하여
\[
  F_{\xi,\eta}(S')=
  \operatorname{Hom}(\xi\times_SS',\eta\times_SS')
\]
라 하자. 그러면 $F_{\xi,\eta}$는 $\mathcal C_S$에서 집합들의 범주로
가는 반변 함자이다. 이때 \emph{모든 두 원소
$\xi,\eta\in\mathcal F_S$에 대하여 증대사상}
\[
  F_{\xi,\eta}(S)\longrightarrow H^0(S'/S,F_{\xi,\eta})
\]
\emph{이 동형이라는 것은 $\alpha:S'\longrightarrow S$가
$\mathcal F$-강하 사상이라는 뜻이다}(정의 1.7).

\medskip\noindent
c. 마찬가지로 $\xi\in\mathcal F_S$와 $S$ 위의 $\mathcal C$의 모든 대상
$S'$에 대하여
\[
  G_\xi(S')=\operatorname{Aut}(\xi\times_SS')
\]
라 두자. 이로써 $\mathcal C_S$에서 군들의 범주로 가는 반변 함자
$G_\xi$를 정의했다.

\src{312}
그러면 $Z^1(S'/S,G_\xi)$는 \emph{$S'\longrightarrow S$에 관하여
$\xi'=\xi\times_SS'$ 위에 주어진 강하 자료들의 집합과 표준적으로
동일시}되고(정의 1.6), $H^1(S'/S,G_\xi)$는 \emph{$\alpha:S'\longrightarrow S$에
관한 강하 자료를 갖춘 $\mathcal F_{S'}$의 대상 가운데,
$\mathcal F_{S'}$의 대상으로서는 $\xi'=\xi\times_SS'$와 동형인 것들의
동형류 집합과 동일시}된다. 따라서 $\alpha:S'\longrightarrow S$가
$\mathcal F$-강하 사상이면 ((b) 참조), $H^1(S'/S,G_\xi)$는
\emph{$\eta\times_SS'$이 $\mathcal F_{S'}$에서 $\xi\times_SS'$와 동형인
$\mathcal F_S$의 대상 $\eta$들의 동형류 집합을 부분집합으로 포함한다.
그리고 이 포함이 항등일 필요충분조건은 $\alpha:S'\longrightarrow S$에
관하여 $\xi'=\xi\times_SS'$ 위에 주어진 모든 강하 자료가 유효한 것이다}.
(특히 $\alpha:S'\longrightarrow S$가 엄밀한 $\mathcal F$-강하 사상이면
그러하다.)

\result{비고}{} $C^*(T/S,F)$ 꼴의 공사슬 복합체들은 잘 알려진 표준
복합체들 대부분(체흐 코호몰로지, 군 코호몰로지 등)을 특수한 경우로
포함하며, 대수기하학에서, 특히 준스킴의 ``베유 코호몰로지''에서 중요한
역할을 한다.

\medskip\noindent
d. \result{예}{1} $S\in\mathcal C$ 위의 대상을 $S'$이라 하고,
$\Gamma$를 $S'$의 자기동형들의 군으로서, $S'$이 ``군 $\Gamma$에 관하여
$S$ 위에서 형식적으로 주동차적''이라고 하자. 즉 자연스러운 사상
\[
  \Gamma\times S'\longrightarrow S'\times_SS'
\]
이 동형이라고 하자. 여기서 $\Gamma\times S'$은 $S'$의 $\Gamma$개
복사본의 직합을 뜻한다. (여기서 관여하는 직합들이 $\mathcal C$에
존재한다고 가정한다.) $F$를 $\mathcal C$에서 아벨군들의 범주로 가는
반변 함자라 하자. 그러면 $C^*(S'/S,F)$는 표준 동차 공사슬들의 단체군
$C^*(\Gamma,F(S'))$과 표준적으로 동형이고, 따라서 $H^*(S'/S,F)$는
$H^*(\Gamma,F(S'))$과 표준적으로 동형이다.

\medskip\noindent
e. \result{예}{2} $\mathcal C$를 준스킴들의 범주라 하자. $G_a$(``가법군'')로
다음과 같이 정의되는, $\mathcal C$에서 아벨군들의 범주로 가는 반변 함자를
나타낸다.
\[
  G_a(X)=H^0(X,\mathcal O_X).
\]
마찬가지로 $G_m$(``곱셈군'')을
\[
  G_m(X)=H^0(X,\mathcal O_X)^*
\]
로 정의한다(환 $H^0(X,\mathcal O_X)$의 가역원들의 군). 더 일반적으로
$\operatorname{Gl}(n)$(``$n$차 일반선형군'')을 다음과 같이 정의한다.

\src{313}
\[
  \operatorname{Gl}(n)(X)=
  \operatorname{Gl}\bigl(n,H^0(X,\mathcal O_X)\bigr),
\]
이며, 이는 $\mathcal C$에서 군의 범주로 가는 함자이다($n>1$이면
반드시 가환일 필요는 없고, $n=1$이면 $G_m$을 다시 얻는다). 한편
$\operatorname{Gl}(n)$은 자기동형사상 함자로도 해석할 수 있다((c) 참조).
실제로 $\mathcal C$를 밑으로 하는 섬유화 범주 $\mathcal F$를, 각
$X\in\mathcal C$에 대하여 $\mathcal F_X$가 $X$ 위의 국소 자유층들의
범주가 되도록 잡으면
\[
  \operatorname{Gl}(n)(X)=
  \operatorname{Aut}_{\mathcal F_X}(\mathcal O_X^n)
\]
이다.
(b)에 따르면, $\alpha:S'\longrightarrow S$가 $\mathcal F$-강하 사상이면
(2.c절 참조), $H^1(S'/S,\operatorname{Gl}(n))$은 $S'$ 위에서의 역상이
$\mathcal O_{S'}^n$과 동형인 $S$ 위의 국소 자유층들의 동형류의 집합을
포함한다. 이 포함이 등식일 필요충분조건은
$\mathcal O_{S'}^n$ 위의 모든 강하 자료($\alpha:S'\longrightarrow S$에
관한)가 유효한 것이다. $S$가 국소환의 스펙트럼이면 $S$ 위의 모든
국소 자유층이 자명하므로, 이는
$H^1(S'/S,\operatorname{Gl}(n))=(e)$라는 뜻이다.

사상 $\alpha:S'\longrightarrow S$에 대한 다음 조건들은 서로 동치이다.
\begin{enumerate}
\item[(i)] 증대 준동형
\[
  H^0(S,\mathcal O_S)=G_a(S)\longrightarrow H^0(S'/S,G_a)
\]
은 동형사상이다.
\item[(ii)] $S'\longrightarrow S$는 $\mathcal F$-강하 사상이다
($\mathcal F$는 위에서 살펴본, $\mathcal C$를 밑으로 하는 섬유화
범주이다).
\end{enumerate}
$S'\longrightarrow S$가 유한 사상이면 이 조건들은 다음 조건과도
동치이다.
\begin{enumerate}
\item[(iii)] $S'\longrightarrow S$는 엄밀 에피사상이다(2.c절 참조).
\end{enumerate}

이제 $S=\operatorname{Spec}(A)$,
$S'=\operatorname{Spec}(A')$라고 하자. 그러면
\[
  C^n(S'/S,G_a)=C^n(A'/A,G_a)=\bigotimes_A^{n+1}A'
\]
이고, 공경계 연산자
$C^n(A'/A,G_a)\longrightarrow C^{n+1}(A'/A,G_a)$는 면 연산자
\[
  \partial_i(x_0\otimes x_1\otimes\cdots\otimes x_n)
  =x_0\otimes\cdots\otimes x_{i-1}\otimes1_{A'}\otimes
   x_i\otimes\cdots\otimes x_n
\]
들의 교대합이다. 마찬가지로
\[
  C^n(S'/S,G_m)=C^n(A'/A,G_m)
  \simeq\left(\bigotimes_A^{n+1}A'\right)^*
\]
이고, $C^*(A'/A,G_m)$의 단체 연산들은
$C^*(S'/S,G_a)$의 단체 연산들로부터 유도된다.
\linebreak $C^*(A'/A,\operatorname{Gl}(n))$의 단체 연산들도 같은 방식으로
구체적으로 적을 수 있다. \emph{강연자가 아는 모든 경우에는}
\[
  H^i(A'/A,G_a)=0\qquad\text{($i>0$)}
\]
\emph{이고, $A$가 국소환이면}
\[
  H^1(A'/A,G_m)=0
  \quad\text{이고 더 일반적으로}\quad
  H^1(A'/A,\operatorname{Gl}(n))=(e)
\]
이다($S'\longrightarrow S$가 $\mathcal F$-강하 사상인 경우, 즉
$A\longrightarrow A'\rightrightarrows A'\otimes_AA'$가 완전 도식인
경우이며, 2.c절과 비교하라). 또한 \emph{힐베르트의 “정리 90”은 바로
관계식}
\src{314}
\[
  H^1(S'/S,G_m)=0
\]
\emph{이다. 여기서 $A$는 체이고 $A'$은 그 체의 유한 갈루아 확대체이다
(예 1 참조). 이는 위의 경우에 $S'\longrightarrow S$가 계수 1인
국소 자유층들의 섬유화 범주에 대한 엄밀 강하 사상이라고 말하는 것으로도
표현할 수 있다.} 바로 이 마지막 형태로 힐베르트의 정리를 일반화해야 하며,
그때에는 $S'\longrightarrow S$라는 사상에 대한 가정과 고려하는
준연접층에 대한 가정을 모두 바꾸어야 한다.

마지막으로, $A$가 극대 아이디얼 $\mathfrak m$을 갖는 아르틴 국소환이고
$A'$이 $A$-대수일 때 다음 성질들은 서로 동치이다(모든 정수 $k>0$에
대하여 $A_k$와 $A'_k$는 각각
$A/\mathfrak m^{k+1}$과 $A'/\mathfrak m^{k+1}A'$을 나타낸다).
\begin{enumerate}
\item[(i)] 모든 $k$에 대하여 $H^1(A'_k/A_k,G_a)=0$이다.
\item[(ii)] 모든 $k$에 대하여 $H^1(A'_k/A_k,G_m)=0$이다.
\item[(iii)] 모든 $k$와 모든 $n$에 대하여
$H^1(A'_k/A_k,\operatorname{Gl}(n))=(e)$이다.
\end{enumerate}
$S'\longrightarrow S$가 엄밀 에피사상이면, 앞의 조건들은 더 나아가 이
사상이 $A'$ 위의 자유 가군들(유한 생성이든 아니든)에 대한 엄밀 강하
사상임을 함의한다.

\result{주의}{} $S,S'$가 각각 체 $A,A'$의 스킴인 경우의 군
$H^i(S'/S,G_m)$의 정의는 아미츠르에게서 비롯되었다. 군
$H^2(S'/S,G_m)$은 브라우어 군의 “전역적” 변형으로서 특히 흥미롭고,
이 변형에 관해서는 [1] 제VII장을 참조할 수 있다.

\fsection{B}{충실 평탄 사상에 의한 강하}

\subsection*{1. 강하 정리들의 진술}
\addcontentsline{toc}{subsection}{1. 강하 정리들의 진술}

\result{정의}{1.1} 준스킴의 사상
$\alpha:S'\longrightarrow S$가 \emph{평탄}하다는 것은 모든
$x'\in S'$에 대하여 $\mathcal O_{x'}$가
$\mathcal O_{\alpha(x')}$-환 위의 평탄 가군이라는 뜻이다(즉
$\mathcal O_{x'}\otimes_{\mathcal O_{\alpha(x')}}M$이
$\mathcal O_{\alpha(x')}$-가군 $M$에 관한 완전 함자라는 뜻이다).
사상이 \emph{충실 평탄}하다는 것은 그것이 평탄하고 전사라는 뜻이다.

예를 들어 $S=\operatorname{Spec}(A)$,
$S'=\operatorname{Spec}(A')$이면, $S'$가 $S$ 위에서 평탄일
필요충분조건은 $A'$이 평탄 $A$-가군인 것이고, $S'$가 $S$ 위에서
충실 평탄일 필요충분조건은 $A'$이 충실 평탄 $A$-가군인 것이다(즉
$A'\otimes_A M$이 $A$-가군 $M$에 관한 완전하고 충실한 함자인 것이다).
이는 세르 [5]의 용어로 쌍 $(A,A')$가 평탄하다는 뜻이기도 하다.
$S'$가 $S$ 위에서 충실 평탄이면 $S$ 위의 준연접층에 대한 역상 함자는
완전하고 충실하다. 다시 말해, $S$ 위의 층들의 준동형들로 이루어진 열이

\src{315}
준연접층들의 열일 때, 그 열이 완전일 필요충분조건은 $S'$ 위에서 그
역상이 완전인 것이다
(특히 $S$ 위의 준연접층들의 준동형이 단사사상, 각각 에피사상, 각각
동형사상일 필요충분조건은 $S'$ 위에서 그 역상이 그러한 것이다). 이
성질은 $S$의 임의의 열린집합 위로 제한해도 성립하며, 이 형태로
충실 평탄 사상을 특징짓는다.

\result{정의}{1.2} 사상 $\alpha:S'\longrightarrow S$가
\emph{준콤팩트}하다는 것은 $S$의 모든 준콤팩트 열린 부분집합 $U$의
역상이 준콤팩트하다는 뜻이다(즉 유한개의 아핀 열린집합의 합집합이라는
뜻이다).

이 성질은 $S$의 아핀 열린집합들에 대해서만 확인하면 충분함이 자명하다.
예를 들어 아핀 사상(즉 아핀 열린집합의 역상이 아핀이 되는 사상)은
준콤팩트하다.

평탄 사상, 각각 충실 평탄 사상, 각각 준콤팩트 사상의 부류는 합성 및
“밑준스킴의 확장”에 대하여 안정하고, 물론 동형사상들을 포함한다.

\result{정리}{1} $\alpha:S'\longrightarrow S$를 충실 평탄이고
준콤팩트인 준스킴의 사상이라 하자. 그러면 $\alpha$는 준연접층들의
섬유화 범주 $\mathcal F$에 대한 \emph{엄밀 강하 사상}이다
(A, 정의 1.7 및 A, 1절, 예 2).

이 진술은 두 가지를 뜻한다.
\begin{enumerate}
\item[(i)] $F,G$가 $S$ 위의 두 준연접층이고 $F',G'$이 $S'$ 위에서의
각각의 역상이면, 자연스러운 준동형
\[
  \operatorname{Hom}(F,G)\longrightarrow\operatorname{Hom}(F',G')
\]
은 왼쪽 항에서 오른쪽 항의 다음 부분군으로 가는 전단사이다. 이 부분군은
이 층들 위의 표준 강하 자료와 양립하는 준동형
$F'\longrightarrow G'$들, 즉 $S''=S'\times_SS'$에서 $S'$로 가는 두
사영에 의한 역상들이 같은 준동형 $F''\longrightarrow G''$을 주는
준동형들로 이루어진다.
\item[(ii)] $S'$ 위의 준연접층 $F'$에 사상
$\alpha:S'\longrightarrow S$에 관한 강하 자료가 주어지면
(A, 정의 1.6), $F'$은 그 강하 자료까지 포함하여 $S$ 위의 어떤
준연접층 $F$의 역상과 동형이다.
\end{enumerate}
(i)에서 $F=\mathcal O_S$로 두면 다음을 얻는다.

\result{따름정리}{1} $G$를 $S$ 위의 준연접층이라 하고, $G'$과 $G''$을
각각 $S'$과 $S''=S'\times_SS'$ 위에서의 역상이라 하며, $p_1,p_2$를
$S''$에서 $S'$로 가는 두 사영이라 하자. 그러면 다음 집합 사상들의
도식

\src{316}
\[
  \Gamma(G)\xrightarrow{\alpha^*}\Gamma(G')
  \underset{p_2^*}{\overset{p_1^*}{\rightrightarrows}}\Gamma(G'')
\]
은 완전이다(A, 정의 1.(a)).

한편 (i), (ii)와 정의 1.1을 함께 적용하면 다음을 얻는다.

\result{따름정리}{2} $G$를 따름정리 1에서와 같이 잡자. 그러면 $G$의
준연접 부분층들과, 두 사영 $p_1,p_2$에 의한 $S''$ 위의 역상들이
$G''$의 같은 부분층을 주는 $G'$의 준연접 부분층들 사이에는
일대일 대응이 있다.

물론 몫층의 용어로도 이와 동치인 진술이 성립한다. 앞에서 본 바와 같이
(A, 4절, (e)), 정리 1은 힐베르트의 “정리 90”의 일반화로 보아야 하며,
1차 코호몰로지로 표현되는 여러 진술을 특수한 경우로 포함한다. 증명을
위해서는 $S=\operatorname{Spec}(A)$,
$S'=\operatorname{Spec}(A')$인 경우로 쉽게 환원된다. (i)에 대해서는
다시 따름정리 1, 즉 모든 $A$-가군 $M$에 대하여 도식
\[
  M=A\otimes_A M\longrightarrow A'\otimes_A M
  \rightrightarrows A'\otimes_AA'\otimes_A M
\]
이 완전임을 증명하는 것으로 쉽게 환원된다. 이는 다음의 더 일반적인
보조정리에서 따른다.

\result{보조정리}{1.1} $A'$을 충실 평탄 $A$-대수라 하자. 그러면 모든
$A$-가군 $M$에 대하여 $M$로 증대된 복합체
$C^*(A'/A,G_a)\otimes_A M$(A, 4절, (e) 참조)는 $M$의
\emph{분해}이다.

앞의 복합체에 $A$에서 $A'$로 스칼라를 확장하여 얻은 증대 복합체가
같은 결론을 만족함을 보이면 충분하다. 이는 $A$를 $A'$으로, $A'$을
$A'\otimes_AA'$으로 바꾸었을 때의 진술을 확인하는 일이므로,
$A$-대수의 준동형 $A'\longrightarrow A$가 존재하는 경우(기하학적
용어로는 $S'$이 $S$ 위에서 단면을 갖는 경우)로 환원된다. 이 경우의
주장은 A, 4절, (a)의 일반론에서 따른다. 그 과정에서 갈루아
코호몰로지의 잘 알려진 진술을 일반화하는 다음 따름정리를 함께 적어 둔다
(A, 4절, (e)와 비교하라).

\result{따름정리}{} $A'$이 $A$ 위에서 충실 평탄이면
$H^0(A'/A,G_a)=A$이고
\[
  H^i(A'/A,G_a)=0\qquad\text{($i\geq1$)}
\]
이다.

정리 1의 (ii)를 증명할 때에도 (i)에서와 같이 진행하여, $S'$이 $S$
위에서 단면을 갖는 경우로 환원한다. 그 경우에는 (i)에서 결론이 따른다
(A, 1절, (c) 참조).

고려하는 준연접층(또는 층들의 계)에 여러 추가 구조를 도입함으로써
정리 1과 그 따름정리들을 여러 방식으로 자유롭게 변형할 수 있음은
자명하다. 예를 들어 $S$ 위에서 하나의 층을 주는 것은
\src{317}
가환 대수의 준연접층을 주는 것은 $S'$ 위에서 그러한 층과
$\alpha:S'\longrightarrow S$에 관한 하강 자료를 함께 주는 것과
“동치”이다. $S$ 위의 그러한 준연접층과 $S$ 위에서 아핀인
준스킴 사이의 함자적 대응을 고려하면, 다음 정리의 두 번째 명제를
얻는다.

\result{정리}{2} $\alpha:S'\longrightarrow S$가 정리 1에서와 같다고
하자. 그러면 $\alpha$는 (일반적으로 엄밀하지는 않은)
\emph{하강 사상}이고(A, 정의 2.4), 준스킴 위에서 아핀인 스킴들의
섬유화 범주에 대해서는 \emph{엄밀 하강 사상}이다(A, 정의 1.7).

정리의 첫 번째 명제는 다음을 뜻한다. $X,Y$를 $S$ 위의 두 준스킴,
$X',Y'$을 $S'$ 위에서의 그 역상, $X'',Y''$을
$S''=S'\times_SS'$ 위에서의 그 역상이라 하자. 그러면 자연스러운
사상들의 다음 그림은
\[
  \operatorname{Hom}_S(X,Y)\xrightarrow{\alpha^*}
  \operatorname{Hom}_{S'}(X',Y')
  \underset{p_2^*}{\overset{p_1^*}{\rightrightarrows}}
  \operatorname{Hom}_{S''}(X'',Y'')
\]
완전하다. 즉 $\alpha^*$는 $\operatorname{Hom}_S(X,Y)$에서
$X',Y'$ 위의 표준 하강 자료와 양립하는 사상들(다시 말해 $S''$에서
$S'$로 가는 두 사영에 의한 그 역상들이 서로 같은 사상들)로 이루어진
$\operatorname{Hom}_{S'}(X',Y')$의 부분집합으로 가는 전단사이다.
$Y$를 $S$ 위에서 아핀인 경우로 한정하면 이는 정리 1의 따름정리 1에서
쉽게 따른다. 일반적인 경우에는 정리 1을 다음 결과와 함께 써야 한다.

\result{보조정리}{1.2} $\alpha:S'\longrightarrow S$를 충실 평탄이고
준콤팩트인 사상이라 하자. 그러면 $S$는 $S'$의 \emph{몫 위상공간}과
동일시된다. 즉 $\alpha^{-1}(U)$가 열리도록 하는 $S$의 모든 부분집합
$U$는 열린집합이다.

정리 2를 완결하려면 $S'$-준스킴 $X'$ 위의 하강 자료에 대한 유효성
판정법을 주어야 한다($X'$이 $S'$ 위에서 아핀이라고 가정하지 않는
경우이다). 먼저 \emph{그러한 하강 자료가 반드시 유효한 것은
아니다}. 이는 $S$가 체 $k$의 스펙트럼이고, $S'$이 그 이차 확대체
$k'$의 스펙트럼이며, $X'$이 $S'$ 위의 2차원 고유 대수적 스킴인
경우에도 마찬가지이다(세르를 따라 나가타의 비사영 곡면을 이용하면
이를 알 수 있다). 충실 평탄이고 준콤팩트인
$\alpha:S'\longrightarrow S$에 관한 $X'/S'$ 위의 하강 자료가
유효하기 위한 필요충분조건은 $X'$이 $S'$ 위에서 아핀이고 $X'$의
하강 자료에 관하여 “안정”인 열린집합 $X'_i$들의 합집합인
것이다. $\alpha:S'\longrightarrow S$가 \emph{라디시엘}인 경우,
즉 단사이고 잉여체 확대들이

\src{318}
순수 비분리인 경우에는 $X'/S'$와 $X'$ 위의 하강 자료가 무엇이든
이 조건이 성립한다. 또한 $\alpha:S'\longrightarrow S$가
\emph{유한}이고, $S$ 위의 $X'$의 한 섬유에 들어 있는 $X'$의 모든
유한 부분집합이 $S$ 위에서 아핀인 $X'$의 한 열린집합에 들어 있으면
여전히 이 조건이 성립함을 보일 수 있다(이것이 \emph{베유 판정법}이다).
특히 $X'/S'$이 준사영이면 이 조건이 성립하며, 이 경우 “하강된”
준스킴 $X/S$도 준사영임을 보일 수 있다($X'/S'$이 사영이면 $X/S$도
사영이다). 요약하면 다음과 같다.

\result{정리}{3} $\alpha:S'\longrightarrow S$를 준스킴의 충실 평탄
준콤팩트 사상이라 하자. $\alpha$가 \emph{라디시엘}이면 이는
\emph{엄밀 하강 사상}이다. $\alpha$가 유한이면, 준스킴 위의 준사영
(또는 사영) 준스킴들로 이루어진 섬유화 범주에 관하여 엄밀 하강
사상이다.

\result{비고}{} 위의 두 번째 명제에서 $\alpha$가 \emph{유한} 사상이라는
가정이 정말 필요한지는 알지 못한다. 어쨌든 “형식적으로” 확인하면,
이를 겉보기에는 더 약한 다음 가정으로 바꿀 수 있다. \emph{$S$의 모든
점에 대하여 열린 근방 $U$, $U$ 위에서 유한이고 충실 평탄인 $U'$,
그리고 $U'$에서 $S'$으로 가는 $S$-사상이 존재한다.} 앞의 경우에
들어맞지 않는 전형적인 예는
$S=\operatorname{Spec}(A)$,
$S'=\operatorname{Spec}(\widehat A)$이고, $A$가 뇌터 국소환이며
$\widehat A$가 그 완비화인 경우이다. 또 다른 예는 $S'$이 $S$ 위에서
준유한이지만(즉 국소적으로 어떤 유한 $S$-스킴의 열린집합과
동형이지만) 유한이지는 않은 경우이다. 이 두 경우에 발표자는 다음
질문의 답도 알지 못한다. $X$가 $S$-스킴이고
$X'=X\times_SS'$이 $S'$ 위에서 사영이면, $X$도 $S$ 위에서
사영인가?\footnote{1962년의 정오표는 부정적인 답을 준다.
준유한이고 에탈이며 전사인 사상 $S'\longrightarrow S$, 또는
$\operatorname{Spec}(\widehat A)\longrightarrow\operatorname{Spec}(A)$
꼴의 사상은 언제나 엄밀 하강 사상인 것은 아니다. 이는 $A$가
대수적으로 닫힌 체 $k$ 위의 대수곡선의 국소환이고
$S=\operatorname{Spec}(A)$인 경우에도 마찬가지이다. 따라서
$X$를 밑이 $S$이고 타원곡선 $E$를 구조군으로 하는 주다발로 만드는 고유하고
매끄러운 사상 $f:X\longrightarrow S$로서,
$f':X'\longrightarrow S'$은 사영이지만 $f$는 사영이 아닌 것을
찾을 수 있다. 그러므로 이는 아벨 스킴을 구조군으로 하는 등상적이지
않은 동차 주다발의 예이기도 하다.}

\subsection*{2. 일부 사상 성질의 하강에 대한 응용}
\addcontentsline{toc}{subsection}{2. 일부 사상 성질의 하강에 대한 응용}

$P$를 준스킴 사상들의 한 부류라 하자.
$\alpha:S'\longrightarrow S$를 준스킴의 사상,
$f:X\longrightarrow Y$를 $S$-준스킴의 사상이라 하고,
$f':X'\longrightarrow Y'$을 $\alpha$에 의한 $f$의 역상이라 하자.
그러면 “$f'\in P$”가 “$f\in P$”를 함의하는지 물을 수 있다.
$\alpha$가 \emph{충실 평탄이고 준콤팩트}라고 가정하면 많은 중요한
경우에 답이 긍정적임을 알 수 있다(뒤의 가정은 때로 불필요하다).
$P$가 전사 사상, 라디시엘 사상(이 두 경우는 $\alpha$의 전사성에서
따른다), 평탄 사상, 충실 평탄 사상, 매끄러운 사상(이 세 경우는
$\alpha$의 충실 평탄성에서 따른다), 또는 유한형 사상의 부류인
경우에는 이를 어려움 없이 직접 알 수 있다. 정리 1, 정리 2와
보조정리 1.2를 이용하면, $P$가 다음 부류 가운데 하나인 경우에도
마찬가지임을 알 수 있다. 즉 동형사상, 열린 몰입 사상, 닫힌 몰입
사상, 몰입 사상($f$가 유한형이고 $Y$가 국소 뇌터인 경우), 아핀
사상, 유한 사상, 준유한 사상, 열린 사상, 닫힌 사상, 위상동형사상,
분리 사상,

\src{319}
고유 사상의 부류이다. 해명되지 않은 유일한 중요한 경우는
제1절의 비고에서 이미 지적한 사영 사상과 준사영 사상의 경우이다.

\subsection*{3. 유한 충실 평탄 사상에 의한 하강}
\addcontentsline{toc}{subsection}{3. 유한 충실 평탄 사상에 의한 하강}

$\alpha:S'\longrightarrow S$를 유한 사상이라 하자. 이는 $S$ 위의
대수의 층 $\mathcal A'$에 대응하며, $\mathcal A'$은 가군의 층으로서
유한 계수로 국소 자유이고 모든 점에서 $0$이 아니다. 그러면
$\alpha$는 충실 평탄이고 준콤팩트인 사상이므로 앞의 결과들을 적용할
수 있다. $S'$ 위의 준연접층 $F'$을 주는 것은 $S$ 위의 준연접층
$\alpha_*(F')$과 그 $\mathcal A'$-가군 구조를 주는 것과 동치이다
($\mathcal A'=\alpha_*(\mathcal O_{S'})$임에 유의하자). 간단히 하기
위해 $S$ 위의 이 층도 $F'$으로 나타내겠다. $S'\times_SS'$ 위의
$F'$의 두 역상 $p_i^*(F')$도 마찬가지로 다음
$(\mathcal A'\otimes_{\mathcal O_S}\mathcal A')$-가군의 준연접층들에
대응한다.
\[
  F'\otimes_{\mathcal O_S}\mathcal A'
  \quad\text{및}\quad
  \mathcal A'\otimes_{\mathcal O_S}F'.
\]
첫 번째 것에서 두 번째 것으로 가는
$(S'\times_SS')$-준동형을 주는 것은
$(\mathcal A'\otimes_{\mathcal O_S}\mathcal A')$-가군의 준동형을
주는 것과 동치이다. 또한 $\mathcal A'$이 국소 자유라는 점을
고려하면, 이는 다음 $(\mathcal A'\otimes\mathcal A')$-가군의
준동형을 주는 것과 동치이다.
\[
  \mathcal U=
  \operatorname{Hom}_{\mathcal O_S}(\mathcal A',\mathcal A')
  =\mathcal A'\otimes\check{\mathcal A}'
  \longrightarrow
  \operatorname{Hom}_{\mathcal O_S}(F',F').
\]
즉 열린집합 $V$ 위의 $\mathcal U$의 모든 단면 $\xi$마다 다음
조건들을 만족하는 $\mathcal O_S$-가군의 준동형
$T_\xi:F'|_V\longrightarrow F'|_V$를 주는 것과 동치이다.
\begin{equation}
  T_{f\xi}(x)=fT_\xi(x),
  \qquad T_{\xi f}(x)=T_\xi(fx),
  \tag{3.1}
\end{equation}
여기서 $f,x$는 각각 $V$에 들어 있는 $S$의 한 열린집합 위의
$\mathcal A',F'$의 단면이다. 하강 자료의 조건 (i), (ii)
(A, 제1절, (c))는 각각 다음과 같이 쓰인다.
\begin{equation}
  T_{1_{\mathcal U}}x=x,
  \qquad\text{즉}\qquad
  T_{1_{\mathcal U}}=\operatorname{id}_{F'},
  \tag{3.2}
\end{equation}
\begin{equation}
  T_{\xi\eta}=T_\xi T_\eta.
  \tag{3.3}
\end{equation}
다시 말해, \emph{$F'$ 위의 하강 자료를 주는 것은
$\mathcal O_S$-대수의 층
$\mathcal U=\operatorname{Hom}_{\mathcal O_S}(\mathcal A',\mathcal A')$
에서 $\mathcal O_S$-대수의 층
$\operatorname{Hom}_{\mathcal O_S}(F',F')$로 가는, 두 선형성 조건
(3.1)을 만족하는 표현을 주는 것과 동치이다.} $S'$ 위의 준연접층들
사이에 짝지음
\[
  F'_1\times F'_2\longrightarrow F'_3
\]
이 주어졌다고 하자(이는 $S$ 위의 $\mathcal A'$-가군의 층들 사이의
짝지음으로 해석할 수 있다).

\src{320}
또한 $F'_i$ 위에 다음 준동형들로 정의되는 붙이기 자료가 주어졌다고
하자.
\[
  T_i:\mathcal U\longrightarrow
  \operatorname{Hom}_{\mathcal O_S}(F'_i,F'_i)
  \qquad (i=1,2,3).
\]
그러면 이 자료들이 명백한 의미에서 주어진 짝지음과 양립할
필요충분조건은 다음 조건이 성립하는 것이다.

한 열린집합 위의 $\mathcal U$의 모든 단면 $\xi$에 대하여
\[
  \Delta\xi=\sum_i\xi'_i\otimes_{\mathcal A'}\xi''_i
\]
로 쓰자. 이는 $\mathcal U\otimes_{\mathcal A'}\mathcal U$의
단면인데($\mathcal U$는 그 왼쪽 구조에 의하여 $\mathcal A'$-가군으로
본다), 다음 식으로 정의된다.
\[
  \xi\mathbin{\cdot}(fg)=\sum_i\xi'_i(f)\xi''_i(g)
\]
(여기서 $f,g$는 더 작은 열린집합 위의 $\mathcal A'$의 두 단면이다).
그러면 더 작은 열린집합 위의 $F'_1,F'_2$의 단면 $x,y$의 모든
쌍에 대하여 다음 식이 성립한다.
\begin{equation}
  T^{(3)}_\xi(x\mathbin{\cdot}y)
  =\sum_iT^{(1)}_{\xi'_i}x\mathbin{\cdot}T^{(2)}_{\xi''_i}y
  \tag{3.4}
\end{equation}
\footnote{인쇄본에는 “$\mathcal A'$의”라고 되어 있다.
$T^{(1)}$과 $T^{(2)}$의 정의역 및 짝지음
$F'_1\times F'_2\longrightarrow F'_3$에 따르면, $x,y$는 각각
$F'_1,F'_2$의 단면이어야 한다.} (이 성질은 준동형 $T^{(i)}$들이
주어진 짝지음에 관하여 $\mathcal U$의 대각 사상과 양립한다고
표현할 수 있다.) 특히 식 (3.1)--(3.4)를 쓰면 $S'$ 위의 대수의
준연접층에 대한 하강 자료를 대각 사상이 주어진 대수의 표현으로
해석할 수 있다. 따라서 가환대수로 한정하면 아핀 $S'$-스킴에 대한
하강 자료도 그렇게 해석할 수 있다.

이로부터 임의의 $S'$-준스킴 $X'$ 위의 하강 자료에 대한 비슷한
해석으로 넘어간다. 그러한 $X'$을 주는 것은 $S$ 위의 준스킴
$X'$과 $\mathcal O_S$-대수의 준동형
\[
  \mathcal A'\longrightarrow\mathcal O_{X'}
\]
을 함께 주는 것과 동치이며, $X'$ 위의 하강 자료를 주는 것은
다음 층의 준동형($h:X'\longrightarrow S'$과 양립하는 것)을
주는 것과 동치이다.
\[
  \mathcal U\longrightarrow
  \operatorname{Hom}_{h^{-1}(\mathcal O_S)}
  (\mathcal O_{X'},\mathcal O_{X'}).
\]
이 준동형은 위의 조건 (3.1)--(3.4)와 비슷한 조건들을 만족한다.

\textbf{예 1 (“베유”).} -- $S'/S$가 갈루아 군 $\Gamma$를 갖는
갈루아 에탈 피복이라고 하자(A의 제3절과 제4절 d 참조). 그러면
$S'$ 위의 준연접층 $F'$ 위의 하강 자료(각각 $S'$-준스킴 $X'$ 위의
하강 자료)를 주는 것은 $S'$ 위의 $\Gamma$의 작용과 양립하는
$(S',F')$(각각 $(S',X')$)의 자기동형들에 의한 $\Gamma$의 표현을
주는 것과 동치이다. 이 결과는

\src{321}
는 “형식적”이다. 즉, 범주의 언어로 증명할 수 있다. 그러나 이 절의
관점에서도 이 결과는 다음 결과 덕분에 완전히 알려진 $\mathcal U$의 명시적
구조($\mathcal U$의 환 구조, 환 준동형
$\mathcal A'\longrightarrow\mathcal U$, 그리고 대각 사상이 주어진 구조)로부터
따라 나온다. 즉, $\mathcal U$는 왼쪽 $\mathcal A'$-가군으로서
$\Gamma$의 원소들에 대응하는 $\mathcal U$의 단면들로 이루어진 기저를 갖는다.

\textbf{예 2(“카르티에”).} -- $p$를 소수라 하고,
$p\mathcal O_S=0$(즉, $\mathcal O_S$의 표수가 $p$이다),
$(\mathcal A')^p\subset\mathcal O_S=\mathcal A$(즉, $S'/S$는 높이 1의
라디시엘 사상이다)라고 가정하자. 또한 $\mathcal A$ 위의 대수층
$\mathcal A'$가 국소적으로 $p$-기저를 갖는다고 가정하자. 즉, 단면들의
족 $(x_i)$가 있어서 $\mathcal A'$는 $x_i$들에 의해 대수로서 생성되고,
이들이 만족하는 관계는 $x_i^p=0$뿐이다. 지표 $i$들의 집합은 유한하며
그 크기가 $n$이라고 가정한다. $\mathfrak g$를 $\mathcal A'$의
$\mathcal A$-미분들의 층이라 하자. 이는 계수(階數) $n$인 국소 자유
$\mathcal A'$-가군층이고, 더 나아가 $\mathcal A$ 위의 $p$-리 대수의
층이지만 $\mathcal A'$ 위의 것은 아니며, 다음 조건을 만족한다.
\begin{equation}
  [X,fY]=X(f)Y+f[X,Y].
  \tag{3.5}
\end{equation}

\textbf{보조정리.} -- $\mathcal U=
\operatorname{Hom}_{\mathcal O_S}(\mathcal A',\mathcal A')$는 대수 준동형
$\mathcal A'\longrightarrow\mathcal U$가 주어진 $\mathcal O_S$-대수로서,
왼쪽 $\mathcal A'$-부분가군 $\mathfrak g$에 의해 생성되며 다음 관계들이
추가로 성립한다.
\begin{equation}
  \left\{
  \begin{aligned}
    Xf-fX&=X(f),\\
    XY-YX&=[X,Y],\\
    X^p&=X^{(p)}.
  \end{aligned}
  \right.
  \tag{3.6}
\end{equation}

앞의 보조정리에 따라 $S'$ 위의 준연접층 $F'$에 강하 자료를 주는 것은
각 $X\in\mathfrak g$에 대하여 다음 조건들을 만족하는 $F'$의
$\mathcal O_S$-자기준동형 $\overline X$를 주는 것과 동치이다.
\begin{align}
  \overline{fX}&=f\overline X, \tag{3.7}\\
  \overline X(fx)&=X(f)x+f\overline X(x), \tag{3.8}\\
  \overline{[X,Y]}&=[\overline X,\overline Y], \tag{3.9}\\
  \overline{X^{(p)}}&=\overline X^p. \tag{3.10}
\end{align}
(이를 곡률이 없고 $p$제곱과 양립하는 $F'$ 위의 선형 접속이라고 부를 수
있다.) $S'$-준스킴 $X'$ 위의 강하 자료라는 개념도 같은 방식으로 명시할 수
있다. 여기서는 관계 (3.4)를 $\overline X$들이 $\mathcal O_{X'}$의
미분이라는 조건으로 바꾼다. 사상 $S'\longrightarrow S$가 라디시엘이므로,
정리 3에 따라 이러한 모든 강하 자료는

\src{322}
유효하며, 따라서 어떤 $S$-준스킴 $X$를 정의한다.

$F'$ 또는 $X'$에 평탄성, 비특이성, 또는 어떠한 유한성 가정도 둘 필요가
없었다는 점에 유의하자.

\subsection*{4. 유리성 판정법에 대한 응용}
\addcontentsline{toc}{subsection}{4. 유리성 판정법에 대한 응용}

$X$를 $\mathcal O_X$의 $S$로의 직상이 $\mathcal O_S$인
$S$-준스킴이라 하자. 이 성질은 밑 $S$의 임의의 평탄 확대
$S'\longrightarrow S$에 의해서도 그대로 유지된다. $F$가 $X$ 위의
가역층(즉, 계수(階數) 1인 국소 자유층)이면, $F$의 자기동형사상들은
$\mathcal O_X$의 가역 단면들과 동일시되므로 $\mathcal O_S$의 가역
단면들과 전단사 대응한다. 이제 $s$를 $S$ 위의 $X$의 단면이라 하자.
비유적으로, $s$ 위의 $F$의 단면이란 $S$ 위의 가역층 $s^*(F)$의 단면을
뜻한다고 하겠다. 앞의 사실에 따라 $F_i$($i=1,2$)가 $X$ 위의 두
가역층이고 각각 $s$ 위의 단면을 갖추고 있으며 $F_1$과 $F_2$가
동형이면, 문제의 단면들과 양립하는(즉, 첫째 단면을 둘째 단면으로 보내는)
$F_1$에서 $F_2$로의 동형사상이 유일하게 존재한다. 한편 단면 $s$와
무관하게, $S$의 모든 점이 다음 성질을 갖는 열린 근방 $U$를 가질 때
$X$ 위의 두 가역층 $F_1$과 $F_2$를 동치라고 약속하자. 그 성질이란
$F_1$과 $F_2$를 $X|U$에 제한한 층들이 동형이라는 것이다. 그러면
$X$ 위의 모든 가역층 $F$는 $s$ 위의 지정 단면을 갖춘 가역층 $F_1$과
동치이고($F_1=F\,s^*(F)^{-1}$로 잡는다), $F_1$은 동형까지 결정된다.
다시 말해, $X$ 위의 가역층들을 동치까지 분류하는 것은 지정 단면을 갖춘
가역층들을 동형까지 분류하는 것과 같다.

이 성질들은 밑의 평탄 확대 $\alpha:S'\longrightarrow S$에 의해서도
그대로 유지되므로(단면 $s$는 $\alpha$에 의한 그 역상 $s'$로 바꾼다),
정리 1을 고려하면 다음 결론을 얻는다.

\emph{$X/S$가 위와 같은 준스킴이고 단면 $s$를 갖는다고 하자.
$\alpha:S'\longrightarrow S$를 충실 평탄 준콤팩트 사상이라 하고,
$F'$을 $X'=X\times_SS'$ 위의 가역층이라 하자. $F'$이 $X$ 위의 어떤
가역층 $F$의 $X'$로의 역상과 동치이기 위한 필요충분조건은
$X'\times_XX'=X\times_S(S'\times_SS')$ 위의 두 역상
$p_1^*(F')$와 $p_2^*(F')$가 동치인 것이다. 이 조건이 성립하면 $F$는
동치까지 결정된다.} (이때 $F'$을 $S$ 위에서 유리적이라고 한다.)

앞 절의 예 1 또는 예 2에서와 같은 $\alpha:S'\longrightarrow S$의
경우에 이 원리를 적용하면 베유 또는 카르티에의 유리성 판정법을 얻는다.
(이 저자들은 $S$와 $S'$가 체들의 스펙트럼인 경우만 다룬다는 점에
유의하자. 더구나 이때 $S$는 국소환의 스펙트럼이므로, 위에서 도입한

\src{323}
동치관계는 다름 아닌 동형관계이다.) 첫째 경우에 $F'$이 $S$ 위에서
유리적이기 위한 필요충분조건은 $\Gamma$에 의한 그 변환들이 $F'$과
동치인 것이다. 둘째 경우의 유리성 판정법을 나타내기 위하여, 일반적으로
$X'/X$의 대각 사상
\[
  X'\longrightarrow X''=X'\times_XX'
\]
과 이에 대응하는 $X'\times_XX'$ 위의 아이디얼층 $\mathcal I$, 그리고
그 역상인 $X'$ 위의 $\Omega^1_{X'/X}$와 동일시되는 층
$\mathcal I/\mathcal I^2$을 생각하자. $F''_i=p_i^*(F')$($i=1,2$)의
대각선으로의 제한들은 동형이므로(각각 $F'$과 동형이므로), 즉
$F''_1(F''_2)^{-1}=F''$의 대각선으로의 제한은 자명하므로, $F''$을
$(X'',\mathcal O_{X''}/\mathcal I^2)$에 제한한 층은 동형까지 다음 군의
잘 결정된 원소 $\xi$로 주어진다.
\[
  H^1(X'',\mathcal I/\mathcal I^2)
  =H^1(X',\Omega^1_{X'/X}).
\]
한편 이 경우에는
\[
  \Omega^1_{X'/X}
  =\Omega^1_{S'/S}\otimes_{\mathcal O_S}\mathcal O_X
\]
이고, 따라서 $\Omega^1_{S'/S}$가 $S$ 위에서 국소 자유이면(카르티에의
경우처럼), $\xi$는 $S'$ 위의
\[
  R^1f'_*(\mathcal O_{X'})\otimes\Omega^1_{S'/S}
\]
의 단면을 정의한다. 이를 $X'/S$ 위의 가역층 $F'$의
아티야--카르티에 류라고 부른다. 이 류가 영인 것은
\[
  (X'',\mathcal O_{X''}/\mathcal I^2)
  =X\times_S(S'',\mathcal O_{S''}/\mathcal J^2)
\]
의 두 사영에 의한 $F'$의 $X'$ 위로의 역상들이 동치이기 위한
필요충분조건이다. 여기서 $\mathcal J$는 대각 사상
$S'\longrightarrow S'\times_SS'$가 $S''=S'\times_SS'$ 위에 정의하는
아이디얼층이다. 따라서 이 류가 영인 것은 $X''=X\times_SS''$ 자체 위의
$F'$의 역상들이 동치이기 위하여 명백히 필요하고, 따라서 $F'$이 $X$ 위의
어떤 가역층 $F$의 역상과 동치이기 위하여도 필요하다. 한편
아티야--카르티에 류는 $S'$ 위에서 국소적으로, $X'/X$의 미분들에
상대적인 $F'$의 접속이 존재하는 데 대한 장애로도 해석할 수 있다. 더
나아가 그러한 접속은 $S'/S$의 미분들을 $X'$로 자연스럽게 연장한 것들에
대응하는 $F'$의 미분들을 알면 결정된다. 이 사실과 앞 절의 전개들로부터,
그 절의 예 2의 경우이고 $X/S$가 단면을 가지면 아티야--카르티에 류가
영이라는 것은 $F'$이 $S$ 위에서 유리적이기 위하여 충분하기도 하다는
결론을 쉽게 얻는다.

\subsection*{5. 아벨 스킴에서 기저 스킴의 제한에 대한 응용}
\addcontentsline{toc}{subsection}{5. 아벨 스킴에서 기저 스킴의 제한}

$S$를 준스킴이라 하자. $S$ 위의 아벨 스킴이란 $S$ 위에서 매끄럽고
고유하며

\src{324}
$S$의 각 점 $x\in S$에서의 섬유가 $\kappa(x)$ 위의 아벨 다양체인
스킴 $X$를 말한다. $S$가 뇌터이고 정칙(즉 그 국소환들이 정칙)이라고
가정하자. 그러면 Murre의 연결성 정리 [4]를 써서(적어도 인용한 정리가
현재 증명되어 있는 ``동표수''의 경우에는) $S$ 위의 $X$의 모든 유리
절단이 모든 곳에서 정의됨(즉 절단임)을 보일 수 있다. 이는 Weil의
고전적 정리를 일반화한다. 더 일반적으로, $X'$이 $S$ 위에서 매끄러운
스킴이면 $X'$에서 $X$로 가는 모든 유리 $S$-사상이 모든 곳에서
정의된다는 것이 따라온다. 따라서 Chow--Iang의 결과를 일반화하는 다음
사실을 얻는다.\footnote{인쇄본에는 ``Igusa--Lang''이라고 되어 있으나,
1962년 정오표는 ``Chow--Iang''으로 고치도록 한다.} $S$가 뇌터이고
정칙이며, $K$가 그 유리함수환(체들의 직접곱)이라고 하자. $X$를 $K$
위의 아벨 스킴이라 하자. $X$가 어떤 $S$ 위의 아벨 스킴 $X_0$에 대해
$X_0\times_S\operatorname{Spec}(K)$ 꼴의 $K$-스킴과 동형이면, $X_0$는
유일한 동형사상까지 결정된다.

앞의 유일성 결과를 쓰면 $X$에서 기저를 제한하는 문제는 $S$ 위에서
국소적임을 알 수 있다. 따라서 각 $x\in S$에 대해
$\operatorname{Spec}(\mathcal O_x)$로 제한할 수 있음만 알면 충분하다.
같은 방법으로, $\alpha:S'\longrightarrow S$가 유한형 매끄러운 사상이고,
$K'$이 $S'$의 유리함수환이며, $X\otimes_KK'$이
$X'_0\times_{S'}\operatorname{Spec}(K')$ 꼴이면, $X'_0$에는
$\alpha$에 관한 표준 하강 데이터가 주어진다. 정리 3을 고려하면 다음이
따라온다.

\result{명제}{5.1} $S$를 유리함수체가 $K$인 기약 뇌터 정칙 준스킴이라
하고, $K'$을 $S$ 위에서 비분기인 $K$의 유한 확대라 하며, $S'$을
$K'$ 안에서의 $S$의 정규화라 하자(따라서 $S'$은 $S$의 에탈 피복이다).
$X$를 $K$ 위의 아벨 스킴이라 하고, $X\otimes_KK'$이 어떤 $S'$ 위의
사영 아벨 스킴 $X'_0$에 대해
$X'_0\times_{S'}\operatorname{Spec}(K')$ 꼴이라고 하자. 그러면 $X$은
어떤 $S$ 위의 사영 아벨 스킴 $X_0$에 대해
$X_0\times_S\operatorname{Spec}(K)$ 꼴이다.

\result{비고}{} 강연자는 $S'\longrightarrow S$가 전사 에탈 피복이라는
가정(이 가정으로 정리 3을 쓸 수 있다)을 전사 유한형 매끄러운 사상이라는
가정으로 바꿀 수 있는지(심지어 그것이 에탈 사상이라고 가정하더라도),
또는 $X'_0$가 $S'$ 위에서 사영이라고 가정하지 않아도 이 명제가
성립하는지 알지 못한다. 후자의 조건은 어쩌면 자동으로 성립할지도 모른다.

\subsection*{6. 국소 자명성과 등자명성 판정기준에 대한 응용}
\addcontentsline{toc}{subsection}{6. 국소 자명성과 등자명성 판정기준}

$S$를 준스킴이라 하고, $G$를 $S$ 위의 ``군 준스킴'', $P$를 $G$가
오른쪽에서 ``작용하는'' $S$ 위의 준스킴이라 하자. $G$의 $P$ 위 작용에서
유도되는 익숙한 사상
\[
  G\times_SP\longrightarrow P\times_SP
\]

\src{325}
이 동형이면 $P$가 $G$ 아래에서 \emph{형식적으로 주동차적}이라고 한다.
이제부터 $G$가 $S$ 위에서 평탄(따라서 충실 평탄)이라고 가정하며,
$G$ 아래의 \emph{동차 주다발}이라는 이름은 충실 평탄이고 준콤팩트인
형식적 주동차 다발 $P$에만 쓰겠다. 이는 어떤 충실 평탄 준콤팩트 기저
확대 $S'\longrightarrow S$가 존재하여 $G'=G\times_SS'$ 아래의 형식적
주동차 다발
\[
  P'=P\times_SS'
\]
이 자명해지는 것, 즉 $G'$과 동형인 것(즉 절단을 갖는 것)과 동치이다.
특히 $S'=P$로 잡을 수 있다. 또한 $S$가 국소 뇌터이면, $P$의 충실
평탄성은 모든 $s\in S$에 대해
\[
  \widehat P_s=P\times_S\operatorname{Spec}(\widehat{\mathcal O}_s)
\]
가 $\widehat{\mathcal O}_s$ 위에서 충실 평탄하다는 조건과 동치임을
유의하라. 여기서 $\widehat{\mathcal O}_s$는 국소환 $\mathcal O_s$의
완비화이다. 이는 $\widehat{\mathcal O}_s$가 $\mathcal O_s$ 위에서 충실
평탄하다는 사실에서 따라온다. 더 나아가 $P$가 국소 뇌터 $S$ 위에서
유한형이면 위 조건을 만족하는 점 $s$들의 집합은 구성가능이다. 따라서
$S$가 ``야콥슨 준스킴''이면(예를 들어 체 위의 유한형 스킴이거나, 더
일반적으로 야콥슨 환의 스펙트럼이면), 문제의 조건을 $S$의 닫힌점들에서만
확인하면 충분하다. 이로써 기저가 완비 국소환 $A$의 스펙트럼인 경우로
환원된다. $S=\operatorname{Spec}(A)$이고($A$는 완비 뇌터 국소환), $P$가
$S$ 위에서 유한형이면, $P/S$의 충실 평탄성은 $S$ 위에서 유한이고
평탄인 어떤 $S'$이 존재하여 $P'$이 자명해지는 것과도 동치이다. 더욱이
$G$가 $S$ 위에서 매끄러우면 $S'$을 $S$ 위의 에탈 사상으로 잡을 수
있다. 따라서 $A$의 잉여체가 대수적으로 닫혀 있기까지 하면(``기하학적
경우''), $P$가 $A$ 위에서 충실 평탄일 필요충분조건은 $P$가 자명한
것이다. 그러므로 $S$가 대수적으로 닫힌 체 위의 대수적 준스킴이고
$G$가 $S$ 위에서 매끄러운 유한형 스킴이면, $S$ 위의 충실 평탄성 조건은
Serre의 해석적 자명성 조건(SLF)과 동치이다([6], 1--12쪽).

$P$에 대해 ``국소 자명성''의 성격을 지닌 더 강한 다른 조건들도 도입할
수 있다. 특히 모든 $s\in S$에 대해 $s$의 열린 근방 $U$와 유한 충실
평탄 사상(각각 전사 에탈 피복) $U'\longrightarrow U$가 존재하여
$P'=P\times_SU'$이 자명하면, $P$가 \emph{등자명}(각각 \emph{엄밀
등자명})이라고 하겠다. (Serre [6]가 국소 등자명이라고 부르는 것을
여기서는 엄밀 등자명이라고 부르므로, 우리는 Serre의 용어법과 다르다.)
엄밀 등자명성은 $G$가 $S$ 위에서 매끄러울 때 특히 유용하지만, 다른
경우에는 오히려 부적절한 개념이다.

$G$가 $S$ 위에서 아핀이면, 2절에 따라 $G$ 아래의 모든 동차 주다발
$P$는 아핀이다. 따라서 정리 2를 써서 충실 평탄 준콤팩트 사상에 의해
그러한 다발들을 ``하강''시킬 수 있다.

\src{326}
특히 $G=\operatorname{Gl}(n)_S$로 잡자. 이는 $S$-준스킴의 함자
\[
  S'\longmapsto\operatorname{Hom}_S(S',G)
\]
(군들의 범주에 값을 갖는다)가 A의 4.e절의 함자
\[
  \operatorname{Gl}(n)(S')
  =\operatorname{Gl}\bigl(n,H^0(S',\mathcal O_{S'})\bigr)
\]
와 동일시된다는 조건으로 정의된다. 다음 사실들을 사용하자. (i) $G$
아래의 모든 동차 주다발(각각 $S$ 위의 계수 $n$인 모든 국소 자유층)은
$S$의 적당한 충실 평탄 준콤팩트 확대 뒤에 ``자명한'' 대상 $G$(각각
$\mathcal O_S^n$)와 동형이 된다. (ii) 여기서 다루는 유형의 대상들,
즉 $G$ 아래의 동차 주다발(각각 계수 $n$인 국소 자유층)은 그러한
사상에 의해 하강시킬 수 있다. (iii) $S'/S$ 위의 자명한 다발의
자기동형군은 $S'$ 위의 계수 $n$인 자명한 국소 자유층의 자기동형군과
함자적으로 동형이다. 이로부터 $S$ 위에(또는 $S'/S$ 위에) 군이 $G$인
동차 주다발을 주는 것과 그 위에 계수 $n$인 국소 자유층을 주는 것이
``같은 것''임이 ``형식적으로'' 따라온다. (더 정확히는 섬유화 범주들의
동치가 있다.) 특히 다음을 얻는다.

\result{명제}{6.1} 군이 $\operatorname{Gl}(n)_S$인 모든 동차 주다발은
국소 자명하다.

잘 알려진 논법으로부터 $\operatorname{Sl}(n)_S$,
$\operatorname{Sp}(n)_S$ 및 그러한 군들의 곱과 같은 구조군들에 대해서도
같은 결과가 따라온다. 또한 $F$가 $G=\operatorname{Gl}(n)_S$의 닫힌
부분군이고 $S$ 위에서 평탄하며 몫 $G/F$가 존재한다고 하자. $G$가
$G/F$ 위에서, 군이 $F\times_S(G/F)$인 등자명(각각 엄밀 등자명) 동차
주다발이면, 군이 $F$인 모든 동차 주다발은 등자명(각각 엄밀 등자명)이다.
이는 지금까지 사용된 $S$ 위의 모든 ``선형군''에 적용된다. 특히 $S$가
체 $k$ 위의 준스킴이고 $\Gamma$가 $k$ 위의 고전적 의미의 선형군,
특히 매끄러운 군일 때의 $G=S\times_k\Gamma$에 적용된다. 따라서 이러한
군들에 대해서는 Serre의 질문 하나가 해결된다(앞에서 인용한 곳).

또한 알려진 $S$ 위의 매끄러운 군들 가운데 대부분(선형이든 아니든)에
대해, 그리고 어쨌든 위와 같은 $S\times_k\Gamma$ 꼴의 군들에 대해,
모든 등자명 동차 주다발이 엄밀 등자명임을 보일 수 있음을 언급해 둔다.
특히 카르티에 방식의 하강으로 얻은 동차 주다발(3절 예 2 참조)은 명백히
등자명이므로, 이를 고려하면 Serre의 또 다른 질문(앞에서 인용한 곳,
1--14쪽)이 해결된다.

\result{비고}{} 이 문제들에서 본질적인 난점 가운데 하나는(몫스킴
$G/F$의 존재 문제는 별도로 하고) 유한이 아닌 충실 평탄 사상에 관한
하강 데이터의 유효성 판정기준이 부족하다는 것이다.

\src{327}
\begin{center}
  \textbf{참고문헌}
\end{center}

\begin{description}
  \item[{[1]}] Dieudonn\'e (J.) 및 Grothendieck (Alexander). --- \emph{Éléments
    de géométrie algébrique}, Publications mathématiques de l'Institut des
    Hautes Études Scientifiques에 게재 예정.
  \item[{[2]}] Grauert (H.) und Remmert (R.). --- \emph{Komplexe Räume}, Math.
    Annalen, t. 136, 1958, p. 245--318.
  \item[{[3]}] Grothendieck (Alexander). --- \emph{Géométrie formelle et
    géométrie algébrique}. Séminaire Bourbaki, t. 11, 1958/59, no 182.
  \item[{[4]}] Murre (J. P.). --- \emph{On a connectedness theorem for a
    birational transformation at a simple point}, Amer. J. Math., t. 80, 1958,
    p. 3--15.
  \item[{[5]}] Serre (Jean-Pierre). --- \emph{Géométrie algébrique et géométrie
    analytique}, Ann. Inst. Fourier Grenoble, t. 6, 1955--56, p. 1--42.
  \item[{[6]}] Serre (Jean-Pierre). --- \emph{Espaces fibrés algébriques},
    Séminaire Chevalley, t. 2, 1958 : Anneaux de Chow et applications, no 1.
\end{description}


\par
\renewcommand{\sourceunit}{e190}
\clearpage
\chapter*{강연 190의 정오표}
\addcontentsline{toc}{chapter}{강연 190의 정오표}
\setcounter{footnote}{0}
\src{300}

\noindent\textbf{3. 대수기하학에서의 하강 기법과 존재 정리
[TDTE I]. I: 일반론. 충실 평탄 사상에 의한 하강.}

\begin{center}
[부르바키 세미나, 제12권, 1959/60, 제190호, 29쪽.]
\end{center}

\noindent\textbf{1쪽, 10행.} --- 이제 하강 기법이 ``대수기하학의
존재 정리 대부분의 기초''라고 말하는 것은 지나쳐 보인다. 이는 TDTE
I--VI 계열의 처음 두 강연에서 다루는 비사영적 기법들에는 상당한
정도로 타당하지만, 사영적 기법들(TDTE IV, V, VI)에는 그렇지 않다.

\noindent\textbf{5쪽, 16행.} --- $\alpha$가 $\mathcal F$-하강 사상이라고
가정할 필요는 없다.

\noindent\textbf{20쪽, 비고.} --- 준유한이고 에탈이며 전사인 사상
$S'\longrightarrow S$, 또는
$\operatorname{Spec}(\widehat A)\longrightarrow\operatorname{Spec}(A)$
꼴의 사상은 언제나 엄밀 하강 사상인 것은 아니다. 이는 $A$가
대수적으로 닫힌 체 $k$ 위의 대수곡선의 국소환이고
$S=\operatorname{Spec}(A)$인 경우에도 마찬가지이다. 따라서 $X$가 밑이
$S$이고 타원곡선 $E$를 구조군으로 하는 주다발이 되게 하는 고유하고
매끄러운 사상 $f:X\longrightarrow S$로서,
$f':X'\longrightarrow S'$은 사영이지만 $f$는 사영이 아닌 것을 찾을
수 있다. 그러므로 이는 아벨 스킴을 구조군으로 하는 등상적이지 않은
동차 주다발의 예이기도 하다.

\noindent\textbf{26쪽, 9행.} --- ``Igusa--Lang'' 대신 ``Chow--Lang''으로
읽어야 한다.

하강 이론의 여러 세부사항에 대해서는 SGA VI, VII, VIII를 보라.


\par
\renewcommand{\sourceunit}{195}
\unit{강연 195 — 하강 기법 II}
\sid{195}
\src{369}
\seminar{제12년차, 1959/60, 제195호\\1960년 2월}
\paperhead{대수기하학에서의 하강 기법과 존재 정리\\
II. 형식적 모듈라이 이론에서의 존재 정리}%
  {알렉상드르 그로텐디크 지음}

\fsection{A}{표현 가능한 함자와 프로-표현 가능한 함자}

\subsection*{1. 표현 가능한 함자}
\addcontentsline{toc}{subsection}{1. 표현 가능한 함자}

$\mathcal C$를 범주라 하자. 각 $X\in\mathcal C$에 대하여 $h_X$를
$\mathcal C$에서 집합의 범주 $(\operatorname{Ens})$로 가는 반변 함자
\[
  h_X:\mathcal C\longrightarrow(\operatorname{Ens})
\]
로서, 다음 식으로 정의하자:
\[
  h_X(Y)=\operatorname{Hom}(Y,X).
\]
$\mathcal C$에 사상 $X\longrightarrow X'$가 있으면, 이로부터 자명한
방식으로 함자 준동형 $h_X\longrightarrow h_{X'}$가 나온다. 즉 $h_X$는
$X$에 관하여 공변 함자이며, 다시 말해 표준 공변 함자
\[
  h:\mathcal C\longrightarrow
  \operatorname{Hom}\bigl(\mathcal C^o,(\operatorname{Ens})\bigr)
\]
를 정의하였다. 이는 $\mathcal C$에서, $\mathcal C$의 쌍대범주
$\mathcal C^o$로부터 집합의 범주로 가는 공변 함자들의 범주로 가는
함자이다. 다음 사실을 상기하자.

\result{명제}{1.1} 이 함자 $h$는 \emph{충실충만하다}. 다시 말해
$\mathcal C$의 임의의 두 대상 $X,X'$에 대하여 자연스러운 사상
\[
  \operatorname{Hom}(X,X')\longrightarrow
  \operatorname{Hom}(h_X,h_{X'})
\]
은 전단사이다.

특히 함자 $F\in\operatorname{Hom}(\mathcal C^o,(\operatorname{Ens}))$가
$h_X$ 꼴의 함자와 동형이면, $X$는 유일한 동형을 제외하고 결정된다.
이때 함자 $F$가 \emph{표현 가능하다}고 한다. 앞의 명제는 표준 함자 $h$가
범주 $\mathcal C$와
$\operatorname{Hom}(\mathcal C^o,(\operatorname{Ens}))$의 표현 가능한
함자들로 이루어진 충만 부분범주 사이의 동치를 정의한다는 뜻이다. 이
사실은 ``보편 문제의 해''라는 개념의 기초가 된다. 그러한 문제는 언제나
$\mathcal C$에서 $(\operatorname{Ens})$로 가는 주어진 함자(여기서처럼
반변이거나, 쌍대적인 경우에는 공변)가 표현 가능한지를 조사하는 문제이다.

\src{370}
더 나아가 범주에서 곱이라는 개념 자체의 정의에 따라 [1], 함자
$h:X\longmapsto h_X$는 곱이 존재할 때마다 곱과 가환한다. 더 일반적으로
유한 또는 무한 사영극한, 특히 섬유곱이나 ``핵'' [2] 등의 형성과도,
그것들이 존재할 때마다 가환한다. 즉 $X\times X'$가 존재할 때마다
함자의 동형사상
\[
  h_{X\times X'}\xrightarrow{\sim}h_X\times h_{X'}
\]
이 있으며, 다시 말해 $Y$에 관하여 함자적인 전단사
\[
  h_{X\times X'}(Y)\xrightarrow{\sim}h_X(Y)\times h_{X'}(Y)
\]
가 있다. 특히 $\mathcal C$에서 사상
\[
  X\times X'\longrightarrow X''
\]
을 주는 것(즉 $\mathcal C$에서 $X,X',X''$ 사이의 ``합성법''을 주는 것)은
사상
\[
  h_{X\times X'}=h_X\times h_{X'}\longrightarrow h_{X''}
\]
을 주는 것과 동치이다. 다시 말해 각 $Y\in\mathcal C$에 대하여 집합의
합성법
\[
  h_X(Y)\times h_{X'}(Y)\longrightarrow h_{X''}(Y)
\]
을 주되, $\mathcal C$의 모든 사상 $Y\longrightarrow Y'$에 대하여 집합
사상들의 계
\[
  h_{X^{(i)}}(Y)\longleftarrow h_{X^{(i)}}(Y')
  \qquad(\text{$i=0,1,2$에 대하여})
\]
가 $Y$와 $Y'$에 대응하는 두 합성법 사이의 사상이 되도록 하는 것과
동치이다. 이로부터 $\mathcal C$의 대상 $X$ 위의 ``$\mathcal C$-군'',
``$\mathcal C$-환'' 등의 구조는 다음과 같이 나타내는 것이 이론적으로나
실제적으로 가장 편리함을 알 수 있다. 각 $Y\in\mathcal C$에 대하여 집합
$h_X(Y)$ 위에 통상적인 의미의 군의 법칙(각각 환의 법칙 등)이 있고,
사상 $Y\longrightarrow Y'$에 대응하는 사상
$h_X(Y)\longleftarrow h_X(Y')$는 군의 준동형(각각 환의 준동형 등)이어야
한다. 예컨대 이 방법은 임의의 밑준스킴 $S$ 위에서 여러 고전군
$G_a$, $G_m$, $\operatorname{Gl}(n)$ 등을 정의하고 이 군들 사이의
고전적인 관계들을 쓰는 데, 준연접층 $\mathcal F$로 정의되는 $S$ 위의
아핀 스킴 $V(\mathcal F)$에 ``벡터 다발''의 구조를 주는 데, 그리고
이에 대응하는 여러 깃발 다양체(그라스만 다양체, 사영 다발)를 정의하고
연구하는 데 가장 직관적이고 편리하다. 일반적 방법의 요지는 표준 함자
$h$를 이용하여 $\mathcal C$의 대상들을 $\mathcal C$에서 집합의 범주로
가는 특별한 반변 함자들, 곧 표현 가능한 함자들과 그대로 동일시하는
것이다.

\src{371}

예를 들어 아핀 스킴의 경우 기하학의 언어에서 가환대수의 언어로
옮겨가기 위해 필요한 통상적인 화살표 뒤집기 절차는 앞의 고찰을
쌍대화하게 한다. 특히 $\mathcal C\longrightarrow(\operatorname{Ens})$인
표현 가능한 공변 함자, 즉
\[
  Y\longmapsto\operatorname{Hom}(X,Y)=h'_X(Y)
\]
꼴의 함자도 도입하게 한다.

\subsection*{2. 프로-표현 가능한 함자와 프로대상}
\addcontentsline{toc}{subsection}{2. 프로-표현 가능한 함자와 프로대상}

$\widetilde X=(X_i)_{i\in I}$를 $\mathcal C$의 대상들의 사영계라 하자.
이에 대응하는 공변 함자는
\[
  h'_{\widetilde X}=\varinjlim_i h'_{X_i}
\]
이고, 더 구체적으로는
\[
  h'_{\widetilde X}(Y)=\varinjlim_i h'_{X_i}(Y)
  =\varinjlim_i\operatorname{Hom}(X_i,Y)
\]
로 주어지는 $\mathcal C$에서 $(\operatorname{Ens})$로 가는 함자이다.
$I$가 여과될 때 이 꼴의 함자와 동형인 $\mathcal C$에서
$(\operatorname{Ens})$로 가는 함자를 \emph{프로-표현 가능하다}고 한다.
앞 절에 따르면 이들은 정확히 표현 가능한 함자들의 여과 귀납극한과
동형인 함자들이다. $\widetilde X'=(X'_j)_{j\in J}$를 $\mathcal C$의 두
번째 여과 사영계라 하자. 여기서 $J$는 두 번째 여과된 사전순서
지표집합이다. 그러면 명제 1.1을 일반화하는 표준 전단사
\[
  \operatorname{Hom}(h'_{\widetilde X'},h'_{\widetilde X})
  =\varprojlim_j\varinjlim_i\operatorname{Hom}(X_i,X'_j)
\]
가 있음을 쉽게 확인할 수 있다.

이로써 $\mathcal C$의 프로대상들의 범주
$\operatorname{Pro}(\mathcal C)$를 도입하게 된다. 그 대상들은
$\mathcal C$의 대상들의 사영계들이다. 이때 지표집합은 여과된 사전순서
집합들 사이에서 달라질 수 있다. 이러한 두 대상
$\widetilde X=(X_i)_{i\in I}$와
$\widetilde X'=(X'_j)_{j\in J}$에 대하여
\[
  \operatorname{ProHom}(\widetilde X,\widetilde X')
  =\varprojlim_j\varinjlim_i\operatorname{Hom}(X_i,X'_j)
\]
로 놓으며, 프로-사상들의 합성은 명백하다. 구성 자체에 따라
$\widetilde X\longmapsto h'_{\widetilde X}$는 $\widetilde X$에 관한 반변
함자로 볼 수 있고, $\mathcal C$의 프로대상들의 범주
$\operatorname{Pro}(\mathcal C)$의 쌍대범주와, $\mathcal C$에서
$(\operatorname{Ens})$로 가는 프로-표현 가능한 공변 함자들의 범주 사이의
동치를 준다. 물론 $\mathcal C$의 대상 $X$는 표준적으로 하나의 프로대상을
정의하며, 이 프로대상도 $X$로 나타낸다.

\src{372}
따라서 $\mathcal C$는 $\operatorname{Pro}(\mathcal C)$의 충만 부분범주와
동치이다. 이제 $\widetilde X=(X_i)_{i\in I}$가 $\mathcal C$의 임의의
프로대상이면, 앞의 동일시 아래에서
\[
  \widetilde X=\varprojlim_i X_i
\]
이다. 여기서 사영극한은 $\operatorname{Pro}(\mathcal C)$에서 취한다
($h'_{\widetilde X}=\varinjlim_i h'_{X_i}$이기 때문이다). $X_i$들의
사영극한이 $\mathcal C$에서 존재하더라도, 일반적으로 이것은
$\operatorname{Pro}(\mathcal C)$에서의 사영극한 $\widetilde X$와 동형이
아님에 주의해야 한다. $\mathcal C$가 집합의 범주일 때 이미 이 사실은
명백하다. 정의 자체에 따라 $\varprojlim_{\mathcal C}X_i=L$이라는 것은
$Y\in\mathcal C$에 관한 $(\operatorname{Ens})$-값 함자
\[
  \varprojlim_i\operatorname{Hom}_{\mathcal C}(Y,X_i)
  =\operatorname{Hom}_{\operatorname{Pro}(\mathcal C)}(Y,\widetilde X)
\]
가 $L$에 의해 표현 가능하다는 조건, 즉
$\operatorname{Hom}_{\mathcal C}(Y,L)$과 동형이라는 조건으로 정의된다.
따라서 $\varprojlim_{\mathcal C}X_i$는 이미 프로대상 $\widetilde X$를
이용하여 정의되며, 정의될 때에는 정확히 말해 $\widetilde X$에
함자적으로 의존한다. 그러므로 이를
$\varprojlim_{\mathcal C}(\widetilde X)$로 나타내도 무방하다.
$\mathcal C$에서 사영극한이 언제나 존재하면
$\varprojlim_{\mathcal C}(\widetilde X)$는
$\operatorname{Pro}(\mathcal C)$에서 $\mathcal C$로 가는 함자이고,
함자들의 표준 준동형
$\varprojlim_{\mathcal C}(\widetilde X)\longrightarrow\widetilde X$가 있다.

생각을 고정하기 위해 공변이라고 한 임의의 함자
\[
  F:\mathcal C\longrightarrow\mathcal C'
\]
는 자명한 방식으로 함자
\[
  \operatorname{Pro}(F):\operatorname{Pro}(\mathcal C)
  \longrightarrow\operatorname{Pro}(\mathcal C')
\]
로 연장된다. 따라서 $\mathcal C'$에서 사영극한이 언제나 존재하면,
$F$는 표준적으로 합성함자
\[
  \overline F:\operatorname{Pro}(\mathcal C)\longrightarrow\mathcal C',
  \qquad
  \overline F=\varprojlim_{\mathcal C'}\circ\operatorname{Pro}(F)
\]
도 정의한다. 이 함자는 $\widetilde X=(X_i)_{i\in I}$를
$\varprojlim_{\mathcal C'}F(X_i)$로 보낸다.

프로대상 $\widetilde X$가, 전이사상 $X_i\longrightarrow X_j$가
에피사상인 프로대상 $(X_i)_{i\in I}$와 동형이면 이를
\emph{엄밀한 프로대상}이라 한다. 이러한 대상이 정의하는 함자를
\emph{엄밀하게 프로-표현 가능하다}고 한다. 이때 $I$가 여과된
순서집합이라고 더 요구할 수 있고, 모든 에피사상
$X_i\longrightarrow X'$은 적당한 $j\in I$에 대한 에피사상

\src{373}
$X_i\longrightarrow X_j$와 동치라고 더 요구할 수 있다. 이 조건으로
$j$는 유일하게 결정된다. 이 조건들 아래에서 사영계
$(X_i)_{i\in I}$는 사영계의 통상적인 동형이라는 의미에서 유일한 동형을
제외하고 결정된다. 따라서 $\operatorname{Pro}(\mathcal C)$에서 엄밀한
프로대상들 $\widetilde X^{(\alpha)}$의 사영계의 사영극한은 언제나
존재하며, 위의 $F$, $\overline F$라는 표기 아래에서 다음이 성립한다:
\[
  \overline F\left(\varprojlim_\alpha\widetilde X^{(\alpha)}\right)
  =\varprojlim_{\alpha,\mathcal C'}
  \overline F\left(\widetilde X^{(\alpha)}\right).
\]
특히 $\mathcal C$의 모든 프로대상이 엄밀하면(다음 절 참조), 연장된 함자
$\overline F$는 사영극한과 가환한다.
\subsection*{3. 프로표현 가능 함자의 특징짓기}
\addcontentsline{toc}{subsection}{3. 프로표현 가능 함자의 특징짓기}

유한 사영극한(즉 반드시 여과될 필요는 없는 유한 준순서집합에 관한
사영극한)이 존재하는 두 범주 $\mathcal C$, $\mathcal C'$을 생각하자. 이는
유한곱과 유한 섬유곱이 존재한다고 해도 같다(특히 이 조건은 모든 $X$에
대하여 $\operatorname{Hom}(X,e)$가 한 원소로만 이루어지는 ``오른쪽 단위원
대상'' $e$가 존재함을 뜻한다). $F$를 $\mathcal C$에서 $\mathcal C'$으로 가는
공변 함자라 하자. 다음 조건들은 서로 동치이다.

\begin{enumerate}
\item[(i)] $F$는 유한 사영극한과 교환한다.
\item[(ii)] $F$는 유한곱 및 유한 섬유곱과 교환한다.
\item[(iii)] $F$는 유한곱과 교환하고, $\mathcal C$의 모든 완전 도식
\[
  X\longrightarrow X'\rightrightarrows X''
\]
에 대하여([3], A, 정의 2.1), $F$로 변환한 도식
\[
  F(X)\longrightarrow F(X')\rightrightarrows F(X'')
\]
도 완전하다.
\end{enumerate}

이때 $F$를 \emph{왼쪽 완전}이라고 한다.

이제부터는 $\mathcal C$에 유한 사영극한이 존재한다고 가정한다. 그러면
정의에서 바로, 표현 가능 함자는 왼쪽 완전이고 사영극한으로 넘기면
프로표현 가능 함자도 왼쪽 완전임을 알 수 있다.

역을 얻기 위하여
\[
  F:\mathcal C\longrightarrow(\operatorname{Ens})
\]
를 공변 함자라 하고, $X\in\mathcal C$, $\xi\in F(X)$라 하자. 모든
$X'\in\mathcal C$, $\xi'\in F(X')$ 및 $\xi=F(u)(\xi')$를 만족하는 모든
엄밀 단사사상([3], A, 2) $u:X'\longrightarrow X$에 대하여 $u$가
동형사상이면, $\xi$(또는 쌍 \src{374}$(X,\xi)$)를 \emph{극소}라고 한다.
한편 $\xi\in F(X)$, $\xi''\in F(X'')$일 때, $\xi''=F(v)(\xi)$를 만족하는
사상 $v:X\longrightarrow X''$가 존재하면 쌍 $(X,\xi)$가 $(X'',\xi'')$을
\emph{지배한다}고 한다. $\xi$가 극소이고 $F$가 왼쪽 완전이면 이 사상
$v$는 유일하며, $\xi''$가 극소이면 $v$는 전사이다. 이로부터 다음 명제를
쉽게 얻는다.

\result{명제}{3.1} $F$가 엄밀하게 프로표현 가능하기 위한 필요충분조건은
다음 두 조건을 만족하는 것이다.
\begin{enumerate}
\item[(i)] $F$는 왼쪽 완전이다.
\item[(ii)] $\xi\in F(X)$인 모든 쌍 $(X,\xi)$는 어떤 극소 쌍의 지배를
받는다.
\end{enumerate}

$\mathcal C$의 모든 대상이 아르틴이면 마지막 조건은 자동으로 성립한다.
실제로 $\xi$가 어떤 $\xi'\in F(X')$의 상이 되게 하는 $X$의 부분대상
$X'$ 가운데 극소인 것을 택하면 된다. 따라서 다음을 얻는다.

\result{따름정리}{} 모든 대상이 아르틴이고 유한 사영극한이 존재하는 범주
$\mathcal C$를 생각하자. 이때 $\mathcal C$에서 $(\operatorname{Ens})$로 가는
프로표현 가능 함자들은 정확히 왼쪽 완전 함자들이며, 더 나아가 모두
엄밀하게 프로표현 가능하다.

마지막 사실은 \emph{$\mathcal C$의 모든 프로대상이 엄밀하다}는 뜻이기도
하다.

\subsection*{4. 예: 갈루아형 군과 프로대수군}
\addcontentsline{toc}{subsection}{4. 갈루아형 군과 프로대수군}

$\mathcal C$가 통상적인 유한군들의 범주이면 $\operatorname{Pro}(\mathcal C)$는
콤팩트 완전불연결 위상군들의 범주와 동치이다. 이러한 군들과, 통상적인
유한군을 주어진 밑준스킴 위의 유한 군 스킴(예를 들어 체 $k$ 위의 유한
대수군)으로 바꾸어 얻는 일반화들은 준스킴에 대한 절대 및 상대 기본군,
호모토피군, 호몰로지군의 역할을 하게 될 것이다. 이 모든 예에는 명제
3.1의 따름정리를 적용할 수 있으며, $\pi_1$은 바로 대응하는 함자를 통해
정의해야 한다 [2]. 체 위의 대수군 또는 준대수군의 범주(더 일반적으로는
뇌터 준스킴 위의 그러한 범주)에서 출발해도 마찬가지이며, 이때 Serre의
``프로대수군''을 얻는다 [4].

\subsection*{5. 예: ``형식적 다양체''}
\addcontentsline{toc}{subsection}{5. 예: 형식적 다양체}

$\Lambda$를 뇌터 환이라 하고, $\mathcal C$를 유한 생성이고 아르틴인
$\Lambda$-가군인 $\Lambda$-대수 $A$들의 범주(더 짧게 말하면,

\src{375}
아르틴 $\Lambda$-대수들의 범주)라 하자. 명제 3.1의 따름정리의 조건들이
성립한다. 여기서 범주 $\operatorname{Pro}(\mathcal C)$는 $\Lambda$ 위의 위상대수
$O$ 가운데 $O_i\in\mathcal C$인 대수들의 위상적 사영극한
\[
  O=\varprojlim_i O_i
\]
과 동형인 것들의 범주와 동치이다. 다시 말해 그 위상은 선형이고 분리이며
완비이고, $O$의 모든 열린 아이디얼 $J_i$에 대하여 $O/J_i$는 $\Lambda$ 위의
아르틴 대수이다. 그러한 대수에 대응하는 함자
$\mathcal C\longrightarrow(\operatorname{Ens})$는 다름 아닌
\begin{align*}
  F(A)=h'_O(A)
  &=\{\text{$O\longrightarrow A$인 위상적 $\Lambda$-대수의 연속 준동형}\}\\
  &=\varinjlim_i
    \operatorname{Hom}_{\Lambda\text{-대수}}(O_i,A)
\end{align*}
이다.

또한 범주 $\mathcal C$는 본질적으로, $\Lambda$의 극대 아이디얼
$\mathfrak m$들에 대한 $\Lambda_{\mathfrak m}$의 완비 국소환에 대응하는
같은 종류의 범주들의 곱이다. 따라서 원한다면 $\Lambda$가 그러한 완비
국소환인 경우로 한정할 수 있다. 어느 경우든 $O$는 $O$가 정의하는 형식적
스킴 [2]의 ``점''들에 대응하는 국소 성분들의 위상적 곱으로 정준적으로
분해된다. 그러한 점은 $K\in\mathcal C$인 어떤 체 $K$에 대한 $F(K)$의 원소
$\xi$로 정의된다(예를 들어 $K$는 생각하는 국소 성분의 잉여체이다). 두 쌍
$(\xi,K)$와 $(\xi',K')$가 같은 점을 정의하기 위한 필요충분조건은 동일한
$(\xi'',K'')$이 두 쌍을 지배하는 것이며, 이는 두 쌍이 동일한
$(\xi''',K''')$을 지배하는 것과도 같다. (모든 $\Lambda/\mathfrak m$이
대수적으로 닫혀 있으면 $F(\Lambda/\mathfrak m)$들의 서로소 합집합을
취하는 것으로 충분하다.)

$\xi\in F(K)$에 대응하는 $O$의 국소 성분 $O_\xi$가 뇌터 환이 되기 위한
조건을 주는 것이 중요하다. $\Lambda$가 완비 국소환(다시 말하지만
뇌터 환)일 때, 이는 $O_\xi$가 $\Lambda$ 위의 형식적 멱급수환
$\Lambda[[t_1,\ldots,t_n]]$의 어떤 몫환과 동형이라고 하는 것과 같다.
그러한 판정조건을 주기 위하여, 모든 환 $A$에 대해 $A$의 ``쌍대수''의
$A$-대수 $I_A$를
\[
  I_A=A[t]/t^2A[t]
\]
로 도입하자. $\varepsilon:I_A\longrightarrow A$를 증대 준동형이라 하자.
이는 $A\in\mathcal C$일 때 사상
\[
  F(\varepsilon):F(I_A)\longrightarrow F(A)
\]

\src{376}
을 정의한다. $F$가 왼쪽 완전이라는 사실을 이용하면, $F(I_A)$에서
``$\xi$로 환원되는'' $\xi'\in F(I_A)$들로 이루어진 부분집합
\[
  F(I_A,\xi)=F(\varepsilon)^{-1}(\xi)
\]
에 $A$-가군 구조를 내재적으로 정의할 수 있다. $O$로 나타낸 $F$의
구체적인 꼴을 이용하면 이 $A$-가군은
\[
  \operatorname{Hom}_{\Lambda}(\mathfrak m_\xi/\mathfrak m_\xi^2,A)
\]
와 동일시된다. 여기서 $\mathfrak m_\xi$는 준동형
$\xi:O\longrightarrow A$의 핵이며, $A$가 체이면 $O$의 국소 성분
$O_\xi$의 극대 아이디얼이다. 이로부터 곧바로 다음 명제를 얻는다.

\result{명제}{5.1} $K\in\mathcal C$가 체이고 $\xi\in F(K)$라 하자.
$\xi$에 대응하는 $O$의 국소 성분 $O_\xi$가 뇌터 환이기 위한
필요충분조건은 $\xi$로 환원되는 $F(I_K)$의 원소들의 집합
$F(I_K,\xi)$가 $K$ 위의 유한 차원 벡터공간인 것이다. 이 조건 아래에는
정준 동형
\[
  F(I_K,\xi)=
  \operatorname{Hom}\!\left(
    \frac{\mathfrak m_\xi/\mathfrak m_\xi^2}
         {\operatorname{im}(\mathfrak n_\xi/\mathfrak n_\xi^2)},K
  \right)
\]
이 있다. 여기서 $\mathfrak n_\xi$는 $\Lambda$의 극대 아이디얼, 즉
$\Lambda\longrightarrow K$의 핵이다.
\footnote{인쇄본에는 여기서 $\Lambda$가 체일 때에만 옳은 공식이 실려
있다. 1962년 정오표는 일반적인 경우 $\mathfrak m_\xi/\mathfrak m_\xi^2$를
$\mathfrak n_\xi/\mathfrak n_\xi^2$의 상으로 나눈 몫을 쓰도록 정정한다.}
특히 $K$-벡터공간 $F(I_K,\xi)$의 차원은 이렇게 고친 벡터공간의 체
$O_\xi/\mathfrak m_\xi=K(\xi)$ 위에서의 차원과 같다.

$O_\xi$가 뇌터라고 가정하고, 표기를 간단히 하기 위하여 $\Lambda$가 완비
국소환이고 $O=O_\xi$라고 가정하자. $O$가 $\Lambda[t_1,\ldots,t_n]$을
이 환의 한 극대 아이디얼에서 국소화한 것의 완비화 위에서 유한 에탈
대수이고, 이 극대 아이디얼이 $\Lambda$의 극대 아이디얼을 유도할 때,
$O$가 \emph{$\Lambda$ 위에서 매끄럽다}고 한다. $O$의 $\Lambda$ 위의 잉여
확대가 자명하면(예를 들어 $\Lambda$의 잉여체가 대수적으로 닫혀 있으면),
이는 $O$ 자체가 그러한 형식적 멱급수 대수와 동형이라고 하는 것과 같다.
\footnote{1962년 정오표는 ``$\Lambda$ 위에서 매끄럽다''는 이 정의가 잉여
확대 $k'/k$가 분리일 때에만 옳다고 명시하며, 일반적인 경우에는 SGA III,
1.1을 참조한다.}
마지막으로 $O$가 반드시 뇌터라고 가정하지 않을 때에도, $O$가 앞의
의미에서 뇌터이고 $\Lambda$ 위에서 매끄러운 몫 $\Lambda$-대수들의 위상적
사영극한과 동형이면 $O$가 $\Lambda$ 위에서 매끄럽다고 한다. $\Lambda$와
$O$가 더 이상 국소라고 가정하지 않는 경우로도 곧바로 일반화할 수 있다.
이제 다음 명제가 성립한다.

\result{명제}{5.2} $O$가 $\Lambda$ 위에서 매끄럽기 위한 필요충분조건은
대응하는 함자 $F$가 에피사상을 에피사상으로 보내는 것이다.

이는 $\mathcal C$의 모든 전사 준동형 $A\longrightarrow A'$에 대하여
$F(A)\longrightarrow F(A')$가 전사라는 뜻이다. 물론 $A$가 국소인 경우,
그리고 (단계적으로 진행하여) $A\longrightarrow A'$의 핵인 $A$의
아이디얼이 $A$의 극대 아이디얼에 의해 소멸되는 경우에 이 조건을
확인하는 것으로 충분하다. 실제로는 함자 $F$를 낳는 상황의 무한소
불변량들과 관련된 어떤 장애가 0임을 확인하는 문제로 환원되며, 이는
코호몰로지적 성격의 문제이다.

마지막으로 앞의 맥락에서 \emph{정의환}에 관하여 몇 마디 하자. 계속해서
$F$를 $\mathcal C$에서 $(\operatorname{Ens})$로 가는 함자라 하고, 위상적
$\Lambda$-대수 $O$가 $F$를 프로표현한다고 하자. 그러면 모든
$A\in\mathcal C$와 모든 $\xi\in F(A)$에 대하여, $\xi$가 어떤
$\xi'\in F(A')$의 상이 되게 하는 $A$의 부분환 $A'$ 가운데 가장 작은 것이
존재하며, 이때 $\xi'$도 유일하게 정해진다. 실제로 $O$에서 $A$로 가는
준동형으로 $\xi$를 해석하고, 이 준동형에 의한 $O$의 상을 $A'$로 취하면
된다. 이때 $A'$를 대상 $\xi\in F(A)$의 정의환이라고 한다.
$u:A\longrightarrow B$가 대수 준동형이고 $\eta=F(u)(\xi)$이면 $\eta$의
정의환은 $\xi$의 정의환을 $u$로 보낸 상이다. $\mathcal C$에서
$(\operatorname{Ens})$로 가는 함자 $F$에서 출발할 때, 정의환의 존재와
방금 말한 성질들은 약간의 차이를 제외하면 $F$가 프로표현 가능하다는
사실과 동치이다. 따라서 이 성질들은 대개 결코 자명하지 않다.

\fsection{B}{두 존재정리}

A의 5절의 표기를 유지하고 공변 함자
\[
  F:\mathcal C\longrightarrow(\operatorname{Ens})
\]
에서 출발하자. $F$가 프로표현 가능하기 위한, 즉 위와 같이 어떤
$\Lambda$-대수 $O$로 나타낼 수 있기 위한 다루기 쉬운 판정조건을 찾고자
한다. A, 명제 3.1의 따름정리에 따라 이를 위한 필요충분조건은 $F$가 왼쪽
완전인 것이다. 현재의 하강 기법으로는([3], p.~9의 질문들 참조) 가장
중요한 경우에 이 판정조건을 이 꼴로 직접 확인할 수 없으므로, 겉보기에는
덜 강한 판정조건들이 필요하다.

\result{정리}{1} 함자 $F$가 프로표현 가능하기 위한 필요충분조건은 다음
두 조건을 만족하는 것이다.
\begin{enumerate}
\item[(i)] $F$는 유한곱과 교환한다.
\item[(ii)] 모든 대수 $A\in\mathcal C$와 모든 준동형
$A\longrightarrow A'$에 대하여, 도식
\[
  A\longrightarrow A'\rightrightarrows A'\otimes_A A'
\]
이 완전하면([3], A, 정의 1.2), 변환한 도식
\[
  F(A)\longrightarrow F(A')\rightrightarrows F(A'\otimes_A A')
\]

\src{378}
도 완전하다.
\end{enumerate}
또한 $A$가 국소이고 다음 두 경우 가운데 하나인 때에 (ii)를 확인하는
것으로 충분하다.
\begin{enumerate}
\item[(a)] $A'$는 $A$ 위의 자유 가군이다.
\item[(b)] 몫가군 $A'/A$는 길이 1인 $A$-가군이다.
\end{enumerate}

이 정리의 증명은 상당히 섬세하여 여기서는 개요를 제시할 수 없다. 다만
그 증명이 아르틴 대수의 스펙트럼 안의 동치관계(범주론적 의미에서)에 대한
연구에 본질적으로 의존함을 지적해 둔다. 이 연구에는 아직도 여러 문제가
남아 있으며, 이 이론을 더 발전시키려면 그 해결이 꼭 필요한 듯하다.

응용에서 조건 (i)의 확인은 언제나 자명하다. 조건 (ii)의 확인은 두
부분으로 나뉜다. $A'$가 자유 $A$-가군인 경우는 평탄 사상에 의한 하강
기법([3], 정리 1, 2, 3)에 속하며 어려움이 없다. 경우 (b)를 다루기
위해서는 다음 결과를 사용한다.

\result{정리}{2} $A$를 극대 아이디얼 $\mathfrak m$을 갖는 아르틴
국소환이라 하고, $A'$을 $A$를 포함하는 $A$-대수라 하자. 다음을 가정한다.
$\mathfrak mA'\subset A$이고 도식
\[
  A\longrightarrow A'\rightrightarrows A'\otimes_A A'
\]
은 완전하다(특히 $A'/A$가 길이 1인 $A$-가군이면 그러하다).
$\mathcal F$를 변하는 준스킴 위의 준연접이고 평탄한 층들로 이루어진
섬유화 범주([3], A, 정의 1.1)라 하자. 그러면 사상
$\operatorname{Spec}(A')\longrightarrow\operatorname{Spec}(A)$은 엄밀한
$\mathcal F$-하강 사상이다([3], A, 정의 1.7).

다시 말해 평탄 $A$-가군 $M$을 주는 것은, 평탄 $A'$-가군 $M'$과
$M'\otimes_AA'$에서 $A'\otimes_AM'$로 가는
$A'\otimes_AA'$-가군 동형사상 가운데 하강 자료의 통상적인 추이성 조건을
만족하는 것을 주는 것과 완전히 동치이다(앞에 인용한 곳).

정리 2는 먼저
\[
  H^i(A'/A,G_a)=0\qquad\text{$i\geq 1$에 대하여}
\]
임을 증명하여 얻는다([3], A, 4(e)). 가정 $\mathfrak mA'\subset A$를
이용하면 $A$가 체 $k$, 즉 $A/\mathfrak m$인 경우로 쉽게 환원할 수 있다.
그런 다음 [3], 16쪽에서 언급한 동치들을 적용한다.

\fsection{C}{몇 가지 특수한 경우에의 응용}
\subsection*{1. 준스킴으로 표현되는 함자에 관한 일반적 고찰}
\addcontentsline{toc}{subsection}{1. 준스킴으로 표현되는 함자에 관한 일반적 고찰}

$S$를 국소 뇌터 준스킴이라 하자. $S$ 위의 준스킴 $X$가 $S$ 위에서
\src{379}국소 유한형이라는 것은, $y\in S$로 사영되는 모든 $x\in X$에
대하여,\footnote{인쇄본에는 여기에서 ``$y\in Y$''라고 되어 있고 이 문단
끝에서도 다시 $Y$를 쓴다. 고려하는 사상이 $X\longrightarrow S$이므로 두
곳 모두 $S$로 읽는다.} 환이 $A$인 $y$의 아핀 근방과, 그 위에 놓이며
환이 $B$인 $x$의 아핀 근방이 존재하여 $B$가 유한형 $A$-대수가 되는
것을 말한다. 고전적 보편 문제의 해로서 나타나는 중요한 예들 가운데에는
$S$ 위에서 국소 유한형이지만 $S$ 위에서 유한형은 아닌 준스킴이 많이
있다. 따라서 곡선의 피카르 스킴을 무한히 많은 연결 성분의 합집합으로
다룰 수 있어야 한다. 이 연결 성분들을 항등원의 연결 성분, 즉
``피카르 다양체''와 혼동해서는 안 된다. 그러므로 때로는 $S$ 위에서
국소 유한형인 준스킴들의 범주 $\mathcal C$ 안에서 반변 함자 $F$의 표현
가능성을 검토하는 것이 편리하다. 이 강연들의 주된 목적은 이러한 함자
$F$가 표현 가능한지를 판별하고, 이에 대응하는 $S$-준스킴 $X$의 성질을
$F$의 성질을 통해 연구하게 해 주는 일반적 기법을 전개하는 데 있다.
덧붙여 말하면, 이 연구에서는 $S$ 위에서 분리되지 않은 준스킴의
비병리적인 예들이, 특히 excellent인 $S$-스킴의 ``피카르 준스킴''으로서
나타난다. 그러므로 스킴이 아닌 준스킴을 대수기하학에서 배제해서는 안 된다.

$X$를 $S$ 위에서 국소 유한형인 준스킴이라 하고,
\[
  F(Y)=\operatorname{Hom}_S(Y,X)
\]
를 이에 결부된 반변 함자라 하자. $F$를, $S$ 위에서 아르틴이고 유한인
준스킴 $Y$들로 이루어진 $\mathcal C$의 부분범주 $\mathcal C_0$에
제한한 함자 $F_0$를 생각할 수 있다. 따라서
$S=\operatorname{Spec}(\Lambda)$일 때 $\mathcal C_0$는 B에서 다룬
아르틴 $\Lambda$-대수의 범주의 쌍대범주이다.
$Y=\operatorname{Spec}(A)$이고 $A$가 아르틴 국소환이면, $Y$는 $S$의
닫힌점 $s$ 위에 놓인 한 점 $y$로 이루어진다. 이때 $Y$에서 $X$로 가는
$S$-사상(즉 $F(Y)$의 한 원소)은 $s$ 위에 놓인 점 $x\in X$와
$\mathcal O_{x,X}$에서 $A$로 가는 $\mathcal O_{s,S}$-준동형을 주는
것으로 정해진다. 그러한 준동형이 존재하면 $x$는 반드시 $X$의 닫힌점이다
(그 잉여체가 $s$의 잉여체 위에서 대수적이기 때문이다). 따라서 $F$를
``아르틴 $\Lambda$-대수''에 제한한 $F_0$는 pro-표현가능하며, 닫힌점
$s\in S$ 위에 놓인 닫힌점 $x\in X$에서의 $X$의 국소환의 완비화
$\widehat{\mathcal O}_{x,X}$들을 국소 성분으로 갖는 위상
$\Lambda$-대수로 표현된다. 이는 $F_0$만 알아도 $X$의 구조, 즉 표시한
점들에서 그 국소환들의 완비화의 구조에 관한 정확한 정보를 얻는다는 것을
보여 준다. $S$가

\src{380}
대수적으로 닫힌 체의 스펙트럼인 경우에도, $\mathcal O_Y$가 멱영원을 가질
수 있는 ``다양체'' $Y$를 체계적으로 고려하고, 특히 아르틴 국소환의
스펙트럼을 사용해야만 고전적 보편 문제의 ``올바른 정식화''에 도달하고
그 ``무한소적'' 측면을 파악할 수 있다는 점에 유의하자.

미리 주어진 함자 $F$가 표현 가능한지를 판정하려 할 때, 정리 1과 2를
이용하여 $F_0$를 연구하면 거의 완전한 지침을 얻는다. 예를 들어
$F(I_K,\xi)$라는 집합들의 성질과 그 함자적 거동을 간단히 조사하여
(A를 보라), 흔히 $F_0$조차 pro-표현가능하지 않음을 알게 되는 경우가
있다. 이는 지금까지 계수가 $>1$인 벡터다발을 분류하는 모듈라이 다양체를
다소나마 자연스럽게 정의하려 했던 시도들이 실패한 까닭을 설명한다.
또는 $F_0$가 분명 표현 가능하지만 벡터 공간 $F(I_K,\xi)$들이 유한
차원이 아니어서 ``형식적'' 해로 만족해야 하는 경우도 있다. 마지막으로
$F_0$가 완비 뇌터 국소환들의 곱으로 잘 표현되는 경우가 있다. 이는 $F$
자체도 표현 가능하리라는 매우 강한 근거를 제공하며, 우리가 뒤에 전개할
더 대역적인 성격의 유사한 성질들과 결합하면 실제로 그러함을 함의하기에
충분할 것이다. 끝으로, ``대역적'' 함자 $F$에서 유래하지 않는 $F_0$만
주어지는 흥미로운 기하 문제들도 나타나며(아래의 문단 4와 5를 보라),
그 경우에는 $F_0$에 ``형식적 모듈라이 스킴''을 결부시킬 수만 있어도
만족할 만하다.

이 일반적 고찰을 마치면서, 스킴 이론이 겉보기의 이상 현상을 어떻게
설명하는지 보자. 예를 들어 이구사 곡면 $V$의 ``피카르 다양체'' $P$는
한 점으로 이루어져 있지만
\[
  H^1(V,\mathcal O_V)\ne 0 ;
\]
이다. 이 경우 $P$는 항등군으로 축소되지 않는 ``순수하게 무한소적인''
군이다. 즉 $P$는 기초체 $k$ 위에서 유한 계수인 국소대수 $O$로 정의되고,
$P$의 곱셈 구조에 대응하는 대각사상을 갖는다. $\mathfrak m$이 $O$의
극대 아이디얼이면 $\mathfrak m/\mathfrak m^2$의 쌍대는
$H^1(V,\mathcal O_V)$와 표준적으로 동형이다(아래의 문단 3을 보라).
피카르 군이 고전적 의미의 대수군, 즉 기초체 $k$ 위에서 매끄러운 군일
때에만 $H^1(V,\mathcal O_V)$의 차원(이는 언제나
$\mathfrak m/\mathfrak m^2$의 차원과 같다)이 피카르 군의 차원과 같다.

\subsection*{2. 스킴 $\underline{\operatorname{Hom}}_S(X,Y)$,
$\Pi_{X/S}Z$, $\underline{\operatorname{Aut}}_S(X)$ 등}
\addcontentsline{toc}{subsection}{2. Hom, Pi 및 Aut 스킴}

$X,Y$를 $S$ 위의 두 준스킴이라 하자. $S$ 위의 모든 준스킴 $T$에 대하여

\src{381}
$X_T=X\times_ST$, $Y_T=Y\times_ST$라 하고, 집합
\[
  F(T)=\operatorname{Hom}_T(X_T,Y_T)
      =\operatorname{Hom}_S(X_T,Y)
      =\operatorname{Hom}_S(X\times_ST,Y)
\]
을 $T$의 반변 함자로 생각하자. 이 함자가 표현 가능하면 이를 표현하는
$S$ 위의 준스킴을 $\underline{\operatorname{Hom}}_S(X,Y)$로 나타낸다.
따라서 함자적 동형사상
\[
  \operatorname{Hom}_S
  \bigl(T,\underline{\operatorname{Hom}}_S(X,Y)\bigr)
  \xrightarrow{\sim}
  \operatorname{Hom}_S(T\times_SX,Y)
\]
이 존재한다. 이 보편 문제로 환원되는 변형은 많다. $S$-준스킴 $X$의
$S$-자기동형사상들의 준스킴은 군 준스킴이 되고, 한 $S$-군 준스킴에서
다른 $S$-군 준스킴으로 가는 $S$-준동형사상들의 준스킴은 두 번째 군
스킴이 가환이면 가환 군 준스킴이 되는 것 등이 그 예이다.
$S$ 위의 준스킴 $X$ 위에 준스킴 $Z$가 주어졌을 때 함자
\[
  F(T)=\operatorname{Hom}_{X_T}(X_T,Z_T)
\]
(즉 $X_T$ 위에서 $Z_T$의 ``단면''들의 집합)를 생각함으로써
$\underline{\operatorname{Hom}}_S(X,Y)$의 정의를 일반화할 수도 있다.
이 함자가 표현 가능하면 이를 표현하는 $S$-준스킴을 $\Pi_{X/S}Z$로
나타낸다. 따라서 정의에 의해 함자적 동형사상
\[
  \operatorname{Hom}_S(T,\Pi_{X/S}Z)
  =\operatorname{Hom}_{X_T}(X_T,Z_T)
\]
이 존재한다. $Z=Y\times_SX$로 두면
$\underline{\operatorname{Hom}}_S(X,Y)$를 다시 얻는다. 이 정의들에서
새로 도입한 준스킴에 관한 자명하지만 유용한 공식 모음이 따라 나온다.
그 공식들은 곱과 섬유곱이 존재하는 모든 범주에서 성립하므로 여기서는
제시하지 않겠다. 더 본질적인 것은
$\underline{\operatorname{Hom}}_S(X,Y)$ 꼴의 스킴이 존재하는가 하는
문제이다. 먼저 $X$를 고정했을 때, 모든 $S$ 위의 $Y$에 대하여
$\underline{\operatorname{Hom}}_S(X,Y)$가 존재하려면(또는 모든 $X$
위의 $Z$에 대하여 $\Pi_{X/S}Z$가 존재하려면) $X$가 $S$ 위에서 평탄해야
함을 알 수 있다. 더 일반적인 $Y$에 대해 해가 존재하리라고 합리적으로
기대하려면 $X$가 또한 $S$ 위에서 고유해야 한다고 여겨진다. 더 나아가
필요한 경우 $Y/S$, 각각 $Z/X$에 ``준사영''과 같은 가정을 두면, 이
조건들은 $\underline{\operatorname{Hom}}_S(X,Y)$와 $\Pi_{X/S}Z$의
존재에 충분한 듯하다. 적어도 상당히 많은 경우에는 이를 확인할 수 있다
(예를 들어 $Y$가 $S$ 위에서 아핀인 경우, 또는 직접적인 초등적 구성에
의해 $X$가 $S$ 위에서 유한인 경우). 정리 1과 2로부터 다음을 얻는다.

\src{382}
\result{명제}{2.1} $\Lambda$를 뇌터환, $X$와 $Y$를 $\Lambda$ 위의
임의의 두 준스킴이라 하고, 아르틴 $\Lambda$-대수들의 범주
$\mathcal C_0$ 위의 함자
\[
  F(A)=\operatorname{Hom}_A(X_A,Y_A)
\]
를 생각하자. $X$가 $\Lambda$ 위에서 평탄하면 이 함자는
pro-표현가능하다.

또한 모든 $A\in\mathcal C_0$와 $\xi\in F(A)$에 대하여 표준 동형사상
\[
  F(I_A,\xi)=H^0\!\left(
    X_A,
    \underline{\operatorname{Hom}}_{\mathcal O_{X_A}}
    \bigl(\xi^*(\Omega^1_{Y_A/A}),\mathcal O_{X_A}\bigr)
  \right)
\]
이 존재한다. 여기서 $\Omega^1_{Y_A/A}$는 $A$에 대한 $Y_A$의 카엘러
$1$-미분들의 층이다. $A$를 체로 두고 A.5.1과 [2]의 유한성 정리를
사용하면 다음 따름정리를 얻는다.

\result{따름정리}{} $X$가 $S$ 위에서 평탄하고 고유이며, $Y$가 $S$
위에서 국소 유한형이라고 하자. 그러면 $F$는 pro-표현가능하고, 이에
대응하는 위상 $\Lambda$-대수의 국소 성분들은 뇌터환이다.

\result{주석}{} 이 절에서 다룬 문제들과 그 밖의 많은 문제들은 고전
대수기하학에서 흔히 사영공간 안의 순환의 ``차우 좌표''를 이용하여
연구되었다. 차우 좌표를 쓰면 이 순환들을 적절한 사영 다양체의 점으로
볼 수 있었다. 그러나 이 방법과 일반적으로 차우 좌표를 사용하는 방법은
스킴의 관점에서는 돌이킬 수 없을 만큼 불충분해 보인다. 매개변수
다양체에서 멱영원을 소거하므로 특히 순환의 무한소 변동을 제대로
연구하는 데 맞지 않으며, 사영공간에 의존하는 비내재적 성격도 지니기
때문이다. 유감스럽게도 많은 대수기하학자가 다양체의 족이나 순환의 족을
연구할 때 차우 좌표라는 언어만 사용했다. 분명한 기술적 가치가 있었지만
(아마 일시적인 가치였겠지만), 이는 이러한 개념들을 명료하게 하는 데
심각한 장애가 되었던 듯하다. 스킴 이론에서 차우 다양체에 대응하는
대상을 얻으려면 다음 보편 문제에 이른다. $X$를 $S$ 위의 준스킴이라
하자. 모든 $S$ 위의 준스킴 $T$에 대하여, $T$ 위에서 평탄한
$X_T=X\times_ST$의 닫힌 부분준스킴들의 집합을 $F(T)$라 한다. 이
$T$의 함자를 $S$ 위의 준스킴으로 표현하려 한다. 더 일반적으로 $X$ 위의
준연접층 $G$를 주고, $T$ 위에서 평탄한 $G_T$의 몫층들의 집합을
$F(T)$로 둘 수 있다.

\src{383}
$S$가 국소 뇌터이고 $X$가 $S$ 위에서 고유이며 $G$가 연접이면, 이
문제에는 $S$ 위에서 국소 유한형인 스킴 $C$로 주어지는 해가 존재하는
듯하다. 아무튼 $S$가 국소 뇌터라는 가정만으로도 $F$를 ``아르틴
$S$-대수''에 제한한 함자는 pro-표현가능하다. 더 나아가 $X$가 $S$
위에서 고유이고 $G$가 연접이면, 이에 대응하는 위상환 $O$의 국소 성분은
뇌터환이다. 물론 $S$ 위에서 $X$의 ``차우 스킴''이 존재함을 증명한
뒤에도, 이를 서로소 열린집합 $C_i$들로 분해해야 한다. 이 $C_i$들은
변동하는 순환의 차수와 차원 같은 연속 불변량을 고정하는 데 대응하며,
$S$ 위에서 단지 국소 유한형일 뿐 아니라 유한형이어야 한다. 또한 이
스킴과 고전적 차우 다양체 사이의 관계를 밝히고, 언제 $C_i$가 $S$ 위에서
사영이거나 적어도 준사영인지도 밝혀야 한다.\footnote{1962년의 정오표에
따르면 사영인 경우는 힐베르트 스킴(TDTE IV)에 의해 완전히 해결된다.
반면 나가타와 히로나카의 예들은 사영성 가정이 없으면 비특이 완비
3차원 다양체의 0차원 부분다양체에 대해서조차 이러한 함자들이 표현
가능하지 않을 수 있음을 보여 준다. 이 현상은 대칭제곱이 존재하지 않을
가능성과 관련된다.}
\subsection*{3. 피카르 스킴}
\addcontentsline{toc}{subsection}{3. 피카르 스킴}

$f:X\longrightarrow S$를 $S$-준스킴이라 하자.\footnote{더 완전한 연구에
대해서는 1962년 정오표가 TDTE V를 참조하도록 한다.} $X$의 구조층의
가역원들로 이루어진 곱셈층 $\mathcal O_X^*$과 군
\[
  P(X/S)=H^0\bigl(S,R^1f_*(\mathcal O_X^*)\bigr)
\]
을 생각하자. 이를 $X/S$의 상대 피카르 군이라 한다. 따라서 이 군의 한
원소는 $S$의 열린 덮개 $(U_i)$와 각 $f^{-1}(U_i)$ 위의 가역층 $L_i$를
주되, 모든 쌍 $(i,j)$에 대하여
$L_i|f^{-1}(U_i\cap U_j)$와 $L_j|f^{-1}(U_i\cap U_j)$가 적어도
$U_i\cap U_j$ 위에서 국소적으로 동형이 되도록 함으로써 정의된다(즉 이
두 층은 [3], B, 4의 의미에서 ``동치''이다). $X/S$가 단면을 가지면
$P(X/S)$는 다름 아닌 $X$ 위의 가역층들을 이 ``동치''로 나눈 동치류들의
집합이다(
\textit{loco citato}). 이제 $S$ 위의 모든 $T$에 대하여
\[
  F(T)=P(X_T/T)
\]
로 놓자. 그러면 $T$에 대한 반변 함자를 얻으며, 이를 $X/S$의 피카르
함자라 부를 수 있다. 이 함자가 표현 가능하면, 이를 표현하는 $S$ 위의
준스킴을 $X/S$의 피카르 준스킴이라 부르고 $\underline P(X/S)$로 나타낸다.
따라서 함자들의 동형사상
\[
  \operatorname{Hom}_S\bigl(T,\underline P(X/S)\bigr)
  \xrightarrow{\sim}P(X_T/T)
\]
을 갖는다. 피카르 준스킴을 만드는 과정은 밑변경과 양립한다. 특히 $S$
위의 $X$의 섬유들에 대한 피카르 준스킴들(이는 $s\in S$의 잉여체
$K(s)$ 위의 준스킴이다)은 $\underline P(X/S)$의 섬유들이다. 물론
$P(X_T/T)=F(T)$가 가환군이므로 피카르 준스킴들은 군 준스킴이다. 또한
특이점을 가질 수 있는 완비 곡선의 피카르 스킴에서 항등원을 포함하는 연결
성분은 다름 아닌 로젠리히트의 일반화 야코비안이다. 따라서 존재가 증명되고
나면 그 성질 대부분은 명백해질 것이다.

\result{비고}{} 여기서 채택한 정의는 $S$의 모든 점\footnote{인쇄본에는
``$Y$''라고 되어 있으나, 여기서는 밑 $Y$가 정의되지 않았고 이 조건은 밑
$S$에 관한 것이다.}이 열린 근방 $U$를 가져 그 위에서 $X$가 단면을 갖는
경우에만 타당하다. 일반적인 경우에도 존재정리를 얻으려면 피카르 함자의
정의를 조금 수정해야 한다.

여기서 피카르 준스킴이 존재할 법한 조건들은 다음과 같다. $X$는 $S$
위에서 고유이고 평탄하며, $f_*(\mathcal O_X)=\mathcal O_S$이고, $X$는
$S$ 위에서 국소적으로 단면을 갖는다. 마지막 조건은 하강 기법을 적용할 때
자연스럽게 나타난다. 즉 $X$ 위의 가역층 $L$에 단면 $s$ 위에서 표시된
단면을 갖추어 그 자기동형들을 제거한다([3], B, 4). 특히 다음을 얻는다.

\src{384}
\result{명제}{3.1} $S=\operatorname{Spec}(\Lambda)$이고 $\Lambda$가
뇌터 환이라고 하자. $X$가 $S$ 위에서 평탄이며, $S$ 위의 모든 유한형
$T$에 대하여
\[
  f_{T*}(\mathcal O_{X_T})=\mathcal O_T
\]
라고 가정하자($f$가 고유이고 분리가능이거나,\footnote{1962년 정오표는
``$f$가 고유이고 분리가능이다''라는 가정을 명시적으로 요구한다.} $S$가
한 체의 스펙트럼이면, K\"unneth 정리에 의해 마지막 조건은
$f_*(\mathcal O_X)=\mathcal O_S$와 동치이다). 그러면 아르틴
$\Lambda$-대수들의 범주 위에서 $X/S$의 피카르 함자는 pro-표현 가능하다.

더 나아가 여기서는
\[
  F(I_A,\xi)=H^1(X_A,\mathcal O_{X_A})
\]
이고, 특히 다음이 성립한다.

\result{따름정리}{} $X$가 $S$ 위에서 고유이면, 피카르 함자에 대응하는
위상적 $\Lambda$-대수 $O$의 국소 성분들은 뇌터이다.

\result{비고}{} 이 번호의 정의와 결과는 $X$ 위의 주다발들의 분류로
일반화할 수 있다. 여기서 구조군은 $S$ 위에서 아핀이고 평탄인 가환 군
스킴 $G$이다. $G$가 가환이 아니면 주다발에 딸린 수반 군다발(그 단면들은
주다발의 자기동형들이다)이 더는 자명하지 않으므로, 명제 3.1은 진술된
그대로는 성립하지 않는다. 그러나 보편 문제를 수정하여 여전히 해를 얻을
수 있다(적어도 현재로서는 형식기하에서). 이 번호와 뒤따르는 번호들의
맥락에서, 그리고 동형까지밖에 정의되지 않는 대상들의 부류에 대한
``모듈라이 스킴''을 찾을 때마다 기억해야 할 황금률은 언제나 다음과 같다.
분류하려는 대상에 있을 수 있는 자기동형들을, 필요하다면 보조 구조(표시된
점이나 표시된 단면의 원소, 미분형식의 고정 등)를 도입하여 제거하되, 그
구조는 원래 문제를 본질적으로 바꾸지 않을 만큼 무해하게 택한다.

\subsection*{4. 다양체의 형식적 모듈라이}
\addcontentsline{toc}{subsection}{4. 다양체의 형식적 모듈라이}

$\Lambda$를 잉여체가 $k$인 뇌터 국소환이라 하고(대부분의 경우
$\Lambda$는 $k$ 또는 Cohen의 $p$-환이다), $X_0$를 $k$ 위의 준스킴이라
하자. 모든 아르틴 국소 $\Lambda$-대수 $A$에 대하여, $A$ 위에서 평탄인
$A$-준스킴 $X$ 가운데 다음 동형사상을 갖춘 것들의 동형류 집합을 $F(A)$라
하자.
\[
  (*)\qquad
  X\otimes_A\kappa(A)\xleftarrow{\sim}
  X_0\otimes_k\kappa(A),
\]
여기서 $\kappa(A)$는 $A$의 잉여체이다. 물론 이러한 평탄
$A$-준스킴들 사이의 동형사상은 구조의 일부로 주어진 위 동형사상을
보존해야 한다. $A$가 반드시 국소일 필요는 없는 아르틴
$\Lambda$-대수이고 그 국소 성분들이 $A_i$이면, $F(A)$를 $F(A_i)$들의
곱으로 잡는다. 이로써 $F$는 $A$에 대한 곱을 보존하는 함자가 되며, 이를
$X_0$(및 $\Lambda$)의 모듈라이 함자라 부를 수 있다. 이 함자가 표현
가능하면 잉여체가 $k$인 국소 위상적 $\Lambda$-대수 $O$가 이에 대응하며,
$O$의 형식적 스펙트럼을 $X_0$(및 $\Lambda$)의 형식적 모듈라이 스킴이라
한다. 이 개념에 관한 몇 가지 자세한 내용은 [2]를 보라.

여기서 하강 기법을 적용할 때 $X_0$의 ``유한'' 자기동형들은 무해하다.
이들은 $A$-준스킴 $X$가 앞에서 정한 의미의 자기동형을 갖는지 여부에
영향을 주지 않기 때문이다. $A$가 체인 경우가 아니면, $X$가 비자명한
$A$-자기동형을 갖지 않기 위한 필요충분조건은
\[
  H^0(X_0,\Theta_{X_0/k})=0
\]
이다. 여기서 $\Theta_{X_0/k}$는 $X_0$의 $k$-미분들로 이루어진 층(즉
접다발)을 나타낸다. 또한 적어도 $X_0$가 $k$ 위에서 매끄러우면 다음을
쉽게 얻는다.
\[
  F(I_A,\xi)=H^1(X_A,\Theta_{X_A/A}).
\]
그러므로 평소와 같이 다음을 결론 내린다.

\src{386}
\result{명제}{4.1} $H^0(X_0,\Theta_{X_0/k})=0$이라고 가정하자. 그러면
$X_0$의 형식적 모듈라이 스킴이 존재한다. 더 나아가 $X_0$가 $k$ 위에서
고유이면 이 형식적 모듈라이 스킴은 뇌터이다.

\result{비고}{} 1\textsuperscript{o} $X_0$가 $k$ 위에서 매끄럽다고
가정하지 않으면, $F(I_A,\xi)$는 다음 $A$-가군의 한 부분가군과 동일시된다.
\[
  \operatorname{Ext}^1_{\mathcal O_{P_A}}
  \bigl(P_A;\mathcal I_{X_A},\mathcal O_{X_A}\bigr),
\]
여기서 $P_A=X_A\times_\Lambda X_A$로 놓고, 대각 사상
$X_A\longrightarrow P_A$를 통해 $\mathcal O_{X_A}$를 $P_A$ 위의
연접층으로 보며, $\mathcal I_{X_A}$는 이 대각 사상이 정의하는 $P_A$
위의 연접 아이디얼층을 나타낸다. 더 정확히 말하면, 호흐실트 이론을
간단히 전역화하여 위의 $\operatorname{Ext}^1$을 다음 데이터와 함께 주어진
$X_A$ 위의 평탄한 $I_A$-대수층 $\mathcal O$들의 동형류 집합과 동일시할
수 있다.
\[
  \mathcal O\otimes_{I_A}A\longrightarrow\mathcal O_{X_A}
\]
이 사상은 증대 동형사상이다(여기서 $I_A=A[t]/(t^2)$로 놓음을 상기하자).
부분가군 $F(I_A,\xi)$는 가환 대수층들에 대응하는 것이다. 따라서 [2]가
시사했던 것과 달리 매끄러움 가정들은 모듈라이 이론에서 본질적이지 않다.

2\textsuperscript{o} 특히 $k$ 위에서 매끄럽고 고유인 모든 대수곡선
$X_0$는 $\Lambda$ 위에서 매끄러운 형식적 모듈라이 스킴을 갖는다는 것을
상기하자(
\textit{loco citato}). 그 상대 차원은 종수 $g$가 $\geq2$이면 $3g-3$이고,
$g=0$ 또는 $1$이면 $g$이다. 뒤의 두 경우는 더는 명제 4.1에 직접
포함되지 않는다. 그러나 타원곡선($g=1$)의 경우에는 뒤따르는 고찰을
이용하여 이 명제로 환원할 수 있을 것이다.

물론 여러 구조를 갖춘 $k$ 위의 스킴들의 계를 생각하면 명제 4.1을 마음껏
변형할 수 있다. 예를 들어 $X_0$가 표시된 원점을 갖는 $k$ 위의 아벨
스킴이라고 하자(즉 $X_0$를 $k$ 위의 군 스킴으로 본다). $A$ 위의 아벨
스킴들(즉 $A$ 위에서 고유이고 매끄러운 군 스킴들) 가운데 아벨 스킴의
동형사상 $(*)$를 갖춘 것들의 동형류 집합을 $F(A)$라 하자. 곱셈 구조를
부과하면(심지어 ``단위원 단면''만 부과해도) 무한소 자기동형들이 제거되고,
따라서 완비 뇌터 국소환 $O$에 대응하는 형식적 모듈라이 스킴이 존재함을
확인할 수 있다. 더 나아가 $S$를 국소 뇌터 준스킴이라 하고, $X$를
``절대 연결''인 섬유들을 갖는 $S$ 위의 고유 매끄러운 스킴이라 하자.
$X$ 위의 곱셈 구조가 단위원 단면을 가지면, 적어도 한 섬유 위에서 결합적이고
가환이며 $S$가 연결인 한, 그 구조는 반드시 결합적이고 가환이다. 또한 그
구조는

\src{387}
단위원 단면만 알면 유일하게 결정된다. 더 나아가 $S$가 잉여체 $k$를 갖는
아르틴 국소환 $A$의 스펙트럼이고, $X$가 $A$ 위에서 고유이며 단면 $s$를
갖는다고 하자. 끝으로 $X\otimes_Ak$가 $k$ 위의 아벨 스킴 구조를 가지며,
그 항등원은 $s$에 대응하는 $X\otimes_Ak$의 점이라고 하자. 간단한 장애
계산과 세르에게서 비롯된 논증을 함께 쓰면, $s$를 단위원 단면으로 갖는
곱셈 구조가 실제로 $X$ 위에 존재함을 증명할 수 있다. (여기서 출발하여,
$A$가 완비 뇌터 국소환인 경우로 넘어가기 위해 [2]의 ``존재정리''를 쓰고
일반적인 경우에는 [3]의 하강 기법을 쓰면, 연결인 모든 국소 뇌터 $S$에
대해 같은 명제를 증명할 수 있다.) 이는 여기서 생각한 함자 $F(A)$가 이
번호의 처음에서 $X_0$ 위의 곱셈 구조를 잊고 정의한 대응하는 함자와
동형임을 증명한다. 특히 $\mathfrak m$이 $O$의 극대 아이디얼이면
$\mathfrak m/\mathfrak m^2$은 $H^1(X_0,\Theta_{X_0/k})$의 쌍대와
표준적으로 동형이며, 따라서 그 차원은 $n^2$이다. 여기서
$n=\dim X_0$이다. $O$가 실제로 $\Lambda$ 위에서 매끄러운지, 즉
$\Lambda$ 위의 $n^2$개 미정원에 대한 형식적 멱급수 대수와 동형인지
밝히는 것은 매우 흥미로울 것이다. 제1절 명제 5.2를 이용하면 이 문제를
주어진 아벨 스킴으로 환원되는 아벨 스킴들의 존재 문제와 동치인 형태로
진술할 수 있다. 어쨌든 초월적 방법을 쓰면 $k$의 표수가 0일 때 답이
긍정적임을 알 수 있다. 표수가 $p\ne0$이면 $\Lambda$가 대수적으로 닫힌
체 $k$의 비트 벡터 환인 경우만 다루면 충분함은 명백하다. 이제야말로
``그린버그 함자''가 그 진가를 보일 때인지도 모른다\ldots\footnote{1962년
정오표는 한 체 위의 아벨 다양체의 형식적 모듈라이 스킴이 비트 벡터 환 위에서
매끄럽다고 밝힌다. 즉 아르틴 국소 몫환 위의 모든 아벨 스킴은 그 몫환을
낳는 환으로 올려진다. 증명은 강연 182, p.~12에서 도입한 올림의 장애류가
변하는 방식을 사용한다. 또한 종수 $g$인 곡선과 편극된 아벨 스킴에 대한
Mumford의 구성들도 상기시킨다(cf. SHMT).}
\subsection*{5. 피복의 연장}
\addcontentsline{toc}{subsection}{5. 피복의 연장}

$\mathfrak X$를 뇌터 형식적 준스킴 [2]라 하고, $U$를 $\mathfrak X$의
열린 부분이라 하자. $U$는 $\mathcal O_{\mathfrak X}$의 한 단면이
``소멸하지 않음''으로 국소적으로 정의되며, 이 단면은 $\mathfrak X$ 위에서뿐
아니라 모든 $X_n$ 위에서도 영인자가 아니다.\footnote{이 명확화는 1962년의
정오표가 요구한 것이다.} 따라서 $U$는 충분히 커서, 열린집합 $V$ 위의
$\mathcal O_{\mathfrak X}$의 단면이 $U\cap V$에서 영이면 그 단면 자체가
영이다. $\mathcal J$를 $\mathfrak X$의 ``정의 아이디얼층''이라 하고
\[
  X_n=(\mathfrak X,\mathcal O_{\mathfrak X}/\mathcal J^{n+1})
\]
이라 놓자. 그러면 이는 통상적인 뇌터 준스킴이다. 이제
$\mathfrak X',\mathfrak X''$을 $\mathfrak X$의 두 평탄 피복, 즉 연접이고
가군층으로서 국소 자유인 대수층으로 정의되는 $\mathfrak X$ 위의 두
준스킴이라 하고, 이들이 $U$ 위에서 비분기라고 하자. 그러면 자명한 사상
\[
  \operatorname{Hom}_{\mathfrak X}(\mathfrak X',\mathfrak X'')
  \longrightarrow
  \operatorname{Hom}_{X_0}(X'_0,X''_0)
\]
은 단사이다. 특히 $X'_0$ 위에 항등을 유도하는 $\mathfrak X'$의
자동사상은 항등사상이다. 이로써 이 상황에 하강 기법을 적용할 수 있다.
특히 $U_0$ 위에서 비분기인 $X_0$의 평탄 피복 $X'_0$에서 출발하자.

\src{388}
$G(\mathfrak X)$를 $X_0$ 위에 $X'_0$를 유도하는(따라서 필연적으로
$U$ 위에서 비분기인) $\mathfrak X$의 평탄 피복 $\mathfrak X'$의 동형류
전체로 두되, 여기서 동형은 $X'_0$ 위에 항등을 유도하는 것으로 한다.
$\mathfrak X$의 모든 열린 부분 $V$에 대하여 $G(V)$를 같은 방식으로
정의하고, 더 일반적으로 $\mathfrak X$ 위의 모든 형식적 준스킴
$\mathfrak Y$에 대하여 $G(\mathfrak Y)$를 정의한다. 그러면 [2]와 [3]의
결과들로부터 우선 다음 결과들이 따른다.

\noindent a. $V$가 $\mathfrak X$의 열린집합들을 달릴 때 $G(V)$들은
$\mathfrak X$ 위의 층을 이루며, 이를 $G_{\mathfrak X}=G$라 쓴다. 이 층을
$U$에 제한하면 각 섬유가 한 원소로만 이루어진 상수층이다.

더 일반적으로 $G_{\mathfrak X}$의 섬유를 결정하는 문제는 완비 국소환의
문제이다. 정확히 말하면 다음과 같다.

\noindent b. 모든 $x\in\mathfrak X$에 대하여
\[
  G_x=G\bigl(\operatorname{Spec}(\mathcal O_{\mathfrak X,x})\bigr)
  \subset
  G\bigl(\operatorname{Spec}(\widehat{\mathcal O}_{\mathfrak X,x})\bigr),
\]
이다. (오른쪽 항은 $\widehat{\mathcal O}_{\mathfrak X,x}$ 위의 유한 자유
대수 $B$ 가운데
$B\otimes_{\widehat{\mathcal O}_{\mathfrak X,x}}\mathcal O_{X_0,x}$와
$\widehat{\mathcal O}'_{0x}$ 사이의 동형사상이 주어진 것들의 동형류이며,
여기서 $\mathcal O'_0$는 $X'_0$를 정의하는 $X_0$ 위의 대수층이다.)
\footnote{인쇄된 식은 $G_x$를 완비된 항과 직접 동일시하고
$\mathcal O'_{0x}$ 위의 모자를 빠뜨렸다. 이 포함관계와 완비화는 1962년
정오표가 규정한 것이다.}

\noindent c. 다음 정준 동형사상이 있다.
\[
  G_{\mathfrak X}=\varprojlim_nG_{X_n} ;
\]
다시 말해 $\mathfrak X$의 모든 열린집합 $V$에 대하여
\[
  G(V)=\varprojlim_nG(V_n).
\]

\noindent d. $\mathfrak X$가 정의 아이디얼 $\mathfrak m$을 갖는 완비
뇌터 국소환 $\Lambda$ 위의 통상적인 고유 스킴 $X$로부터 얻어진다고 하자.
구체적으로 $\mathcal J=\mathfrak m\mathcal O_X$라 놓고
$\mathcal O_X$의 $\mathcal J$-아딕 완비화를 취한다고 하자. 그러면
$G(\mathfrak X)$는 ``$X'_0$로 환원되는'' 통상적인 스킴 $X$의 평탄
피복들의 동형류 전체와 정준적으로 동형이다.

비유적으로 말하면, (a)와 (b)는 이 문제의 국소적 측면과 대역적 측면 사이의
근본 관계를 확립하고, (c)는 ``유한한'' 측면과 ``무한소적'' 측면 사이의
관계를 주며, 마지막으로 (d)는 명시된 조건 아래에서 ``형식적'' 측면과
``대수적'' 측면이 동일함을 상기시킨다.

이제 $\mathfrak X$가 완비 뇌터 국소환 $\Lambda$ 위에서 정의되고
$\mathcal J=\mathfrak m\mathcal O_{\mathfrak X}$라고 하자. 따라서 $X_0$는
$\Lambda/\mathfrak m$ 위의 준스킴이다. $A$가 $\Lambda$-가군으로 유한
생성인 모든 $\Lambda$-대수일 때
\[
  F(A)=G(\mathfrak X\times_\Lambda A)
\]
라 놓자. 여기서 표시는 문자 그대로 $A/\mathfrak mA$ 위에서가 아니라,
축소 대수 $A_{\mathrm{red}}$로 밑변경한 뒤에 취한다. 이는 $A$에 관한
공변 함자로서 집합의 범주에 값을 가진다. (c)에 따르면 이 함자는 아르틴
$A$들에 대해서만 알면 완전히 알 수 있다. 또한 이 함자가 pro-표현 가능하다고
말하는 것, 즉
\[
  F(A)=\operatorname{Hom}_{\Lambda\text{-algèbres top.}}(O,A)
\]

\src{389}
꼴이라고 말하는 것은, $O$가 A의 5절에서 고찰한 유형의 위상
$\Lambda$-대수인 경우나 $A$를 아르틴 대수들로 제한한 경우나 마찬가지이다.
따라서 정리 1과 정리 2를 함께 적용하면 실제로 다음을 얻는다.

\result{명제}{5.1} 이 함자가 분류하는 모든 아르틴 족이 그 계수 대수 위에서
평탄하고, 선택한 열린집합이 관련 밑변경 뒤에도 스킴론적으로 조밀하다고 하자.
그러면 앞의 함자는 pro-표현 가능하다.

물론 (a)에 따라 $U=\mathfrak X$이면 $\mathfrak X$ 위의 모든
$\mathfrak Y$에 대하여 $G(\mathfrak Y)$는 한 원소로만 이루어지고, 이때
함자 $F$는 별다른 흥미가 없다($O=\Lambda$이다). 사실상 그 밖의 거의 모든
경우에는 위상 국소환 $O$가 뇌터가 아닌 듯하다. 그럼에도 그 존재는 해들의
집합 $G(\mathfrak X)$가 갖는 ``연속적'' 성격을 두드러지게 드러낸다. 이는
직관적으로 $X'_0$를 연장할 때 분기가 전파되는 방식을 ``연속적으로'' 선택할
수 있다는 사실에 대응한다. 이 결과를 국소 유체론에서 J.-P. Serre [4]가
취한 관점과 비교해 볼 수 있다. 그 관점 역시 ``기하적'' 국소체의 극대 아벨
확대가 갖는 위상 갈루아군의 연속적 성격을 드러내며, 그 쌍대군은
(Pontrjagin의 의미에서) 대수군들(또는 적어도 준대수군들)의 귀납극한으로
나타난다. 여기에서도 확대들의 분류는 무한 차원의 ``다양체들''로 주어진다.
더욱이 위에서 $\mathfrak X$를 완비 국소환의 형식적 스펙트럼으로 취할 수
있으며, 이때 $\Lambda$는 예를 들어 그 환의 Cohen 부분환으로 취할 수 있다.
그리고 이 절의 결과들을 차원이 $>1$인 완비 국소환의 확대를 연구하는 데
사용할 수 있으리라 기대할 수 있다. 국소적인 경우와 대역적인 경우 모두에서,
이 결과들은 고차 분기 현상과 표수 영에서의 현상(초월적 방법으로 다룰 수
있는 현상) 사이의 정확한 관계를 정식화하게 해 줄지도 모른다. 어쨌든 [2]에서
기본군을 연구하기 위해 설명한 방법을 ``순분기''인 경우로 확장하고, 초월적
방법으로 ``세 점 문제''를 해결할 수 있게 하는 것은 바로 명제 5.1에 앞선
분석이다.

끝으로 $X_0$의 차원이 1이면 상황이 단순해짐을 지적하자. 이때 (a)를 점을
뺀 완비 국소 스킴에서의 형식적 붙이기로 보충하면
\[
  G(\mathfrak X)\simeq
  \prod_iG\bigl(\operatorname{Spec}(\widehat{\mathcal O}_{\mathfrak X,x_i})\bigr),
\]
이며, 여기서 $x_i$들은 $X_0-U$의 점들이다. 즉 분기점들에서의 ``국소''
확대들은 임의로 주어도 된다. $X_0$의 정규성만으로는 5.1의 형식적 모듈라이
스킴이 $\operatorname{Spec}(\Lambda)$ 위에서 매끄럽게 되지 않는다. 인쇄된
마지막 명제는 일반적으로 거짓이다.

\src{390}
\begin{center}
  \textbf{참고문헌}
\end{center}

\begin{description}
  \item[{[1]}] Grothendieck (Alexander). --- \emph{Sur quelques points
    d'algèbre homologique}, Tohoku Math. J., t. 9, 1957, p. 119--221.
  \item[{[2]}] Grothendieck (Alexander). --- \emph{Géométrie formelle et
    géométrie algébrique}, Séminaire Bourbaki, t. 11, 1958/59, fasc. 3,
    n\textsuperscript{o} 182.
  \item[{[3]}] Grothendieck (Alexander). --- \emph{Technique de descente et
    théorèmes d'existence en géométrie algébrique, I : Généralités, Descente
    par morphismes fidèlement plats}, Séminaire Bourbaki, t. 12, 1959/60,
    fasc. 1, n\textsuperscript{o} 190.
  \item[{[4]}] Serre (Jean-Pierre). --- \emph{Corps locaux et isogénies},
    Séminaire Bourbaki, t. 11, 1958/59, fasc. 3, n\textsuperscript{o} 185.
\end{description}

\par\smallskip
\begingroup\footnotesize


\noindent\textit{2.1의 인쇄본은 접공간 공식에 $H^1$을 쓰지만, 고정된 사상의 일차 변동은 전역 도함수들, 즉 $H^0$로 주어진다. 5.1에는 계수 대수 위의 평탄성과 관련 밑변경 뒤의 스킴론적 조밀성이 필요하다. 함자의 표시는 $A/\mathfrak mA$가 아니라 축소 계수 대수 $A_{\mathrm{red}}$로 밑변경한 뒤에 취한다.}\par
\noindent\textit{$X_0$의 정규성만으로 형식적 모듈라이 스킴의 매끄러움이 따르지는 않는다. 일차원 경우의 완비 국소 곱 분해에는 층 성질 (a)와 점을 뺀 완비 국소 스킴에서의 형식적 붙이기가 필요하다. 정오표에서 (b)의 완비화 등식은 포함관계로 약화되므로, (a)와 (b)만으로는 그 전사성을 얻을 수 없다.}\par
\endgroup

\par
\renewcommand{\sourceunit}{e195}
\clearpage
\chapter*{강연 195의 정오표}
\addcontentsline{toc}{chapter}{강연 195의 정오표}
\setcounter{footnote}{0}
\src{300}

\noindent\textbf{4. TDTE II: 형식적 모듈라이 이론의 존재 정리.}

\begin{center}
[부르바키 세미나, 제12권, 1959/60, 제195호, 22쪽.]
\end{center}

\noindent\textbf{8쪽.} --- 명제 5.1에 쓰인 공식은 $\Lambda$가 체인
경우에만 옳다. 일반적인 경우에는
$\mathfrak m_\xi/\mathfrak m_\xi^2$을 이 공간을
$\mathfrak n_\xi/\mathfrak n_\xi^2$의 상으로 나눈 몫으로 바꾸어야 한다.
여기서 $\mathfrak n$은 $\Lambda$의 극대 아이디얼이다. 또한 $O$가
$\Lambda$ 위에서 매끄럽다는 것에 관해 주어진 정의는 잉여체 확대 $k'/k$가
분리 가능할 때에만 옳다. 일반적인 경우에는 SGA III, 1.1을 보라.

\noindent\textbf{14쪽, 주.} --- 여기서 제기한 문제들은 사영의 경우
``힐베르트 스킴''에 의해 완전히 해결된다(cf. TDTE IV). 그러나 Nagata와
Hironaka의 예들은 사영성 가정을 두지 않으면 고려한 함자들이 반드시 표현
가능한 것은 아님을 보여 준다. 이는 차원 3인 비특이 완비 다양체의 0차원
부분다양체들을 분류하는 데에만 국한하더라도 마찬가지이다. 이 사실은 그러한
다양체의 2차 대칭곱이 존재하지 않을 수 있다는 점과 관련된다.

\src{301}
\noindent\textbf{15쪽, 제3항.} --- 더 자세한 연구는 TDTE V를 보라.

\noindent\textbf{16쪽, 명제 3.1.} --- ``$f$가 고유이면'' 대신
``$f$가 고유이고 분리 가능하면''으로 읽어야 한다.

\noindent\textbf{18쪽과 19쪽, 주, 제2항.} --- 나는 최근 체 위의 아벨
다양체에 대한 형식적 모듈라이 스킴이 실제로 비트 환 위에서 매끄럽다는 것을
증명했다. 다시 말해, 어떤 다른 환의 몫인 아르틴 국소환 위의 모든 아벨
스킴은 그 다른 환 위의 아벨 스킴을 환원하여 얻어진다. 증명은 강연 182,
12쪽에서 도입한 올림 장애류의 함자적 성질만을 사용한다. 또한 종수 $g$인
곡선들의 모듈라이 스킴과 편극된 아벨 스킴들의 모듈라이 스킴을 Mumford가
구성했음을 상기하자(cf. SHMT).

\noindent\textbf{19쪽, 제5문단.} --- $U$를 정의하는 데 쓰는 단면은
$\mathfrak X$ 위에서뿐 아니라 모든 $X_n$ 위에서도 영인자가 아니라고
가정해야 한다.

\noindent\textbf{20쪽, 12--13행.} ---
\[
  G_x=G\bigl(\operatorname{Spec}(\widehat{\mathcal O}_{\mathfrak X,x})\bigr),
\]
대신
\[
  G_x=G\bigl(\operatorname{Spec}(\mathcal O_{\mathfrak X,x})\bigr)
  \subset
  G\bigl(\operatorname{Spec}(\widehat{\mathcal O}_{\mathfrak X,x})\bigr),
\]
으로 읽고, ``$\mathcal O'_{0x}$'' 대신
``$\widehat{\mathcal O}'_{0x}$''으로 읽어야 한다.

\par
\renewcommand{\sourceunit}{212}
\unit{강연 212 — 몫 준스킴}
\sid{212}
\src{99}
\seminar{제13년차, 1960/61, 제212호\\1961년 2월}
\paperhead{대수기하학에서의 구성 기법과 존재 정리들\\
III. 몫 준스킴}%
  {알렉상드르 그로텐디크 지음}

\subsection*{서론}
\addcontentsline{toc}{subsection}{서론}

이 강연에서 다루는 문제들은 가변 스킴에 대한 어떤 공변 함자들을
나타내려 한다는 점에서, 반변 함자들을 나타내려 했던 앞의 두 강연과
다르다. 그렇지만 몫을 취하는 과정은 강연 I과 II의 문제들을 포함하여
대수기하학의 많은 구성 문제에서 본질적이다([1], [2]). 이를테면 충실
평탄이고 준콤팩트인 사상 $T\longrightarrow S$에 관하여 $T$-준스킴
$X$ 위의 \emph{하강 데이터가 유효한가} 하는 문제는, 그 하강 데이터가
$X$ 안에 정의하는 평탄 동치관계에 대한 $X$의 몫이 존재하는가 하는
문제와 동치이다. 여기서 그 몫은 아래에서 살펴볼 합당한 성질들을
만족해야 한다. [1], A, \S~2 c에서 제기한 문제들도 아마 이 강연의
제2절에서 제기하는 문제들과 함께 해결될 것이다. 마찬가지로
$S$-스킴 $X$의 \emph{피카르 스킴}(정의는 [2], C, \S~3 참조)은
몇몇 다른 스킴들(유효 약수들의 스킴 또는 사영공간으로의 몰입 사상들의
스킴)을 평탄 동치관계로 나눈 몫으로 여러 방식으로 정의할 수 있다.
그런데 이 보조 스킴들의 정의와 구성은 기술적으로 더 간단하다. 실제로
이들은 [2], C, \S~2에서 정의한
$\underline{\operatorname{Hom}}_S(X,Y)$ 꼴의 스킴들과 그 변형들이며,
적절한 사영성 가정 아래에서 이들의 구성이 다음 강연의 주제가 될
것이다. 따라서 이 강연과 다음 강연의 결과들을 결합하면 적절한 가정
아래에서 피카르 스킴들을 구성하게 된다.

준스킴에서 몫을 취하는 문제에는 여전히 해결되지 않은 문제가 여럿
있다. 가장 중요한 문제는 제8절에서 언급한다. 현재 이 문제는
\emph{임의 종수의 곡선과 극화된 아벨 다양체 등에 대한 정수 위의
모듈라이 스킴}을 구성하는 데 남은 유일한 장애물이다. 그러므로 이
문제의 해결은 대수적 군 전문가들이 힘쓸 만한 가치가 있다.

\src{100}
\subsection*{1. 동치관계와 유효 동치관계}
\addcontentsline{toc}{subsection}{1. 동치관계와 유효 동치관계}

$\mathcal C$를 범주라 하고 $X$를 $\mathcal C$의 대상이라 하자. 사상쌍
\[
  p_1,p_2:R\rightrightarrows X
\]
을 $\mathcal C$에서의 \emph{동치쌍}이라고 하며, 그 \emph{공역}은
$X$, \emph{정의역}은 $R$이라고 한다. 이는 $\mathcal C$의 모든 대상
$T$에 대하여 대응하는 두 사상
\[
  p_1(T),p_2(T):R(T)\rightrightarrows X(T)
\]
(여기서 $\mathcal C$의 모든 대상 $Y$에 대하여
$Y(T)=\operatorname{Hom}(T,Y)$로 둔다)이 사상
\[
  R(T)\longrightarrow X(T)\times X(T)
\]
을 정의하고, 이 사상이 $R(T)$에서 집합 $X(T)$ 위의 한 동치관계의
그래프로 가는 전단사를 유도한다는 뜻이다. 공역이 $X$인 동치쌍들
사이에는 자명한 동치관계가 있다. 이 후자의 동치관계에 대한 하나의
동치류를 \emph{$X$ 안의 $\mathcal C$-동치관계}라 하며, 혼동할
염려가 없으면 간단히 동치관계라 한다.

$X\times X$가 존재하면, $X$ 안의 동치관계를 주는 것은 $X\times X$의
부분대상 $R$을 주되, $\mathcal C$의 모든 대상 $T$에 대하여
\[
  (X\times X)(T)=X(T)\times X(T)
\]
에서 $R(T)$에 대응하는 부분집합이 $X(T)$ 위의 한 동치관계의
그래프가 되도록 하는 것과 동치이다. 사영
$\operatorname{pr}_1,\operatorname{pr}_2$가 유도하는 $R$에서 $X$로
가는 사상들을 $p_1,p_2$라 하면, 앞의 조건은 $(p_1,p_2)$가
동치쌍이라는 뜻이다. $X\times X$와 섬유곱
$(R,p_2)\times_X(R,p_1)$이 존재한다는 전제 아래에서는, [2], A,
\S~1에서 설명한 일반 원리에 따라 $X(T)$ 안의 $R(T)$에 대한 집합론적
동치관계의 공리들을 $\mathcal C$ 안에서 도식으로 나타낼 수도 있다.
여기서는 이것이 필요하지 않을 것이다.

같은 정의역 $R$과 같은 공역 $X$를 갖는 사상쌍 $(p_1,p_2)$가 주어질
때마다, 그 쌍의 \emph{여핵}을 $Z$를 변수로 하는 공변 함자
\[
  \operatorname{Hom}_{p_1,p_2}(X,Z),
\]
즉 $up_1=up_2$를 만족하는 $X$에서 $Z$로 가는 사상 $u$들의 집합을
나타내는 $\mathcal C$의 대상 $Y$로 정의할 수 있다. $Y$가 존재하면
유일한 동형사상을 제외하고 결정된다. 이를 $X/(p_1,p_2)$로 쓰거나,
기호를 남용하여 $X/R$로 쓴다. 후자의 기호는 특히
\src{101}$(p_1,p_2)$가 동치쌍일 때 사용한다. 이 경우 그 쌍이 정의하는
동치관계와 $R$을 기호에서 동일시하는 것이 관례이다. $Y$를 $X$의
몫으로 본다면, 실제로 $Y$는 그 쌍 $(p_1,p_2)$가 정의하는 동치관계에만
의존함을 유의하자.

이제 사상
\[
  f:X\longrightarrow Y
\]
에서 출발하자. 이 사상으로 $X$를 ``$Y$ 위의 대상''으로 볼 수 있다.
섬유곱
\[
  R(f)=X\times_Y X
\]
가 존재한다고 가정하고, 그 두 사영을 $p_1,p_2$라 하자. 그러면
$(p_1,p_2)$는 동치쌍이며, 이를 사상 $f$에 \emph{결부된} 동치쌍이라
한다. 따라서 이 동치쌍은 하나의 동치관계를 정의하며, 이를 $f$에
결부된 동치관계라 한다.

공역이 $X$이고 공통 정의역이 $R$인 사상쌍 $(p_1,p_2)$가 다음을
만족하면 이를 \emph{유효 동치쌍}이라 한다.
\begin{enumerate}
  \item[(i)] 여핵 $X/(p_1,p_2)=Y$가 존재한다.
  \item[(ii)] 섬유곱 $X\times_Y X$가 존재한다.
  \item[(iii)] 성분이 $p_1,p_2$인 사상
  $R\longrightarrow X\times_Y X$가 동형사상이다.
\end{enumerate}
그러면 $(p_1,p_2)$는 실제로 동치쌍이다. 이 쌍이 정의하는 동치관계도
\emph{유효 동치관계}라 한다.

사상 $f:X\longrightarrow Y$가 다음을 만족하면 이를 \emph{유효
에피사상}이라 한다.
\begin{enumerate}
  \item[(i)] 섬유곱 $X\times_Y X=R$이 존재한다.
  \item[(ii)] $p_1,p_2$를 $R$에서 $X$로 가는 두 사영이라 할 때,
  몫 $X/(p_1,p_2)$가 존재한다.
  \item[(iii)] $f$가 유도하는 사상
  $X/(p_1,p_2)\longrightarrow Y$가 동형사상이다.
\end{enumerate}
그러면 $f$는 실제로 에피사상이며, 더 나아가 엄밀 에피사상이다
([1], A, \S~2.3 참조). 섬유곱 $X\times_Y X$가 존재하면 그 역도
성립한다. 에피사상 $f$가 정의하는 $X$의 몫대상도 $X$의
\emph{유효 몫}이라 한다.

앞의 정의들은 다음의 ``갈루아 대응''을 함의한다.

\src{102}
\result{명제}{1.1} $X$ 안의 유효 동치관계 $R$들의 집합과 $X$의 유효
몫 $Y$들의 집합 사이에는 자연스러운 순서들을 보존하는 일대일 대응이
있다. $R$에는 유효 몫 $X/R$이 대응하고, $Y$에는 표준 사영
$X\longrightarrow Y$가 정의하는 유효 동치관계가 대응한다. 이 관계는
두 사영을 갖춘 섬유곱 $X\times_Y X$로 정의된다.

매우 좋은 범주들(집합의 범주, 집합의 층의 범주 등)에서는 모든 몫이
유효이고 모든 동치관계가 유효이다. 주어진 준스킴 $S$ 위의 준스킴
범주와 같은 범주에서는 더 이상 그렇지 않다. 이는 $S$가 체의
스펙트럼인 경우에도, 심지어 $S$ 위에서 유한인 스킴들로 한정한
경우에도 마찬가지이다. 유효성 문제, 그리고 $S$ 위에서 유한이지 않은
준스킴의 경우에는 몫의 존재 문제조차 대부분 까다롭다.

\subsection*{2. 예: $S$ 위에서 유한인 준스킴}
\addcontentsline{toc}{subsection}{2. 예: S 위에서 유한인 준스킴}

$S$를 국소 뇌터 준스킴이라 하고, $S$ 위에서 유한인 준스킴들의
범주를 $\mathcal C$라 하자. 이 범주는 $S$ 위의 연접인 가환 대수의
층들의 범주의 반대범주와 동치이다. 또 $S$가 환 $A$의 아핀
준스킴이면, 이 범주는 $A$ 위에서 유한인 $A$-대수들(즉 $A$-가군으로서
유한 생성인 대수들)의 범주의 반대범주와도 동치이다. 따라서
$\mathcal C$에는 유한 사영극한과 유한 귀납극한이 존재함이 곧바로
따른다. 전자는 (아무런 유한성 가정 없이도\ldots) 잘 알려져 있다.
예를 들어 $S$ 위의 준스킴 $X,Y$의 섬유곱은 대응하는 대수들의
텐서곱 $B\otimes_A C$에 대응한다. 두 $A$-대수 준동형
$u,v:C\rightrightarrows B$가 정의하는 두 사상
$X\rightrightarrows Y$의 핵은 $u(c)-v(c)$들이 생성하는 아이디얼로
$B$를 나눈 몫에 대응한다. 유한 귀납극한에 관해서는 한편으로는
$A$-대수들의 보통 곱에 대응하는 유한 합들을, 다른 한편으로는
사상쌍 $X\rightrightarrows Y$의 여핵들을 생각하면 충분하다. 실제로
후자는 (곧바로 알 수 있듯이) 준동형
$u,v:C\rightrightarrows B$가 일치하는 $C$의 원소들로 이루어진
부분환에 대응한다. 이 부분환은 뇌터 가정에 의해 $A$ 위에서 유한이다.
더욱이 뇌터 가정을 사용하면, $S$ 위에서 유한인 준스킴들의 범주
$\mathcal C$에서 이렇게 구성한 유한 귀납극한들, 특히 몫들은 사실
\emph{모든} 준스킴의 범주에서도 몫임을 보일 수 있다.

\src{103}
[1]에서 말했듯이 $\mathcal C$에는 유효가 아닌 에피사상들이 있다
(또는 엄밀하지 않은 에피사상들이 있다\footnote{인쇄본에는
``엄밀하지 않은''에서 ``않은''이 빠져 있다. 지금 생각하는 범주에는
섬유곱이 존재하며, 바로 앞 문단은 유효성과 엄밀 에피사상임 사이의
동치를 확립한다.}; 섬유곱이 존재하므로 이는 같은 말이다).
평탄성 가정을 두지 않을 때 \emph{동치관계가 언제나 유효인지는 알지
못한다}. 이 방향에서 얻은 것은 가군의 형식 이론의 기본 정리를
증명하는 데 꼭 필요한 매우 부분적인 몇몇 긍정적 결과뿐이다([2], B,
정리 1 참조). 제기된 문제에서는 $S$가 대수적으로 닫힌 잉여체를 갖는
국소 아르틴환 $A$의 스펙트럼인 경우로 쉽게 귀착시킬 수 있음을
언급해 둔다. 그러나 $A$가 체인 경우에도 답은 알려져 있지 않다.

$S$ 위의 준스킴 $X$가 더 이상 $S$ 위에서 유한이라고 가정하지 않되,
$X$ 안의 동치관계 $R$에 대하여 $p_1:R\longrightarrow X$가 유한
사상인 경우도 생각할 수 있다. 이때 $R$을 유한 동치관계라 한다.
간단히 하기 위해 $S$와 $X$가 아핀이라고 가정하자. 그러면 $R$도
아핀이므로 상황은 순수한 가환대수의 상황으로 귀착된다. \emph{이
경우에도 몫 $X/R=Y$가 존재하는지, 그리고 표준 사상
$X\longrightarrow Y$가 유한인지조차 알려져 있지 않다}. (가장 간단한
경우는 $S$가 체 $k$의 스펙트럼이고 $X$가 $k[t]$의 스펙트럼, 즉 아핀
직선인 경우이다.) 물론 앞의 두 문제가 긍정적으로 해결되면, 지금의
상황에서 $R$이 유효임을 결론내릴 수 있다. 동치관계의 그래프 대신
제4절의 의미에서 $X$ 안의 동치 전그래프만 주어진 경우에도, 몫 $Y$의
\emph{존재}와 $f:X\longrightarrow Y$의 유한성 문제는 정확히 같은
방식으로 제기됨을 유의하자.

다소 임의적인 유한 동치관계로 몫을 취하는 문제는 주어진 준스킴
$X_i$들을 어떤 닫힌 부분준스킴들을 따라 ``붙여서'' 준스킴을 구성할
때 생긴다. 붙이기 법칙은 $X_i$들의 합인 준스킴 $X$ 안의 유한
동치관계로 정확히 표현된다. 또한 여기에서 제기한 문제들과 그 여러
변형을 해결하는 것은 [2]에서 시작한 방향으로 비사영적 구성의 일반
기법을 확립하기 위한 예비 조건이 될 것이라고 예상해야 한다.

필자가 아는 유일한 일반적 긍정 결과는 다음과 같다.

\result{명제}{2.1} $S$를 국소 뇌터 준스킴, $s$를 $S$의 점,
$\Omega$를 $k(s)$의 대수적으로 닫힌 확대체라 하자. 대응하는
\src{104}``섬유 함자'' $F$, 즉 $S$ 위에서 유한인 모든 $S$-스킴
$X$에 $X_s$의 $\Omega$-값점들의 집합을 대응시키는 함자를 생각하자.
이 함자는 (자명하게 왼쪽 완전이며) \emph{오른쪽 완전}이다. 즉 유한
귀납극한, 특히 사상쌍의 여핵과 가환한다.

이 결과를 $S$의 모든 ``기하적 점''에 사용하면 다음이 따른다.
전사 라디시엘 사상들을 법으로 하여 추론함으로써(즉 해당 사상들의
역을 형식적으로 덧붙여서) $\mathcal C$에서 얻는 ``몫'' 범주
$\mathcal C'$은 ``기하적'' 범주, 즉 집합의 범주와 같은 ``유한적인
성격''의 성질들을 만족하는 범주이다. 특히 그 안의 모든 동치관계는
유효이다. 따라서 $R$이 $S$ 위에서 유한인 $X$ 안의 동치관계이면
표준 사상
\[
  R\longrightarrow X\times_Y X\qquad (\text{여기서 }Y=X/R)
\]
은 \emph{라디시엘이고 전사}이다. 사실 이 사상은 단사사상이므로
전사 닫힌 몰입 사상이다.

\subsection*{3. 작용군의 경우}
\addcontentsline{toc}{subsection}{3. 작용군의 경우}

다시 $\mathcal C$가 임의의 범주라고 가정하자. $G,X$를 $\mathcal C$의
대상이라 하고, $G$가 $\mathcal C$의 대상 $X$에 작용하는
$\mathcal C$-작용군이라고 가정하자. 이는([2], A, \S~1 참조) $\mathcal
C$의 모든 대상 $T$에 대하여 $G(T)$에는 군 구조가, $X(T)$에는
$G(T)$-집합 구조가 주어져 있고, $T$가 변할 때 이 구조들이 $T$에 따라
``함자적으로 변한다''는 뜻이다. $\mathcal C$에서 곱 $G\times G$와
$G\times X$가 존재하면, 그러한 구조를 다음 두 사상의 쌍으로 정의할
수도 있다.
\[
  G\times G\longrightarrow G,
  \qquad \pi:G\times X\longrightarrow X,
\]
여기에는 $\mathcal C$의 모든 대상 $T$에 대하여 집합 $G(T)$와 $X(T)$에
대응하는 연산들이 $G(T)$를 $X(T)$에 작용하는 군으로 만들어야 한다는
조건이 붙는다. 이 공리를 $\mathcal C$의 몇몇 도식의 가환성으로 옮기는
일은 쉽지만 번거로우며, 실제로 내가 아는 모든 경우에 전혀 쓸모가 없다.

$G\times X$가 존재한다고 가정하고 다음 두 사상을 생각하자.
\[
  p_1,p_2:G\times X\rightrightarrows X
\]
여기서
\src{105}
\[
  p_1=\operatorname{pr}_2,\qquad p_2=\pi.
\]
쌍 $(p_1,p_2)$가 동치쌍일 필요충분조건은 $\mathcal C$의 모든 대상
$T$에 대하여 이 쌍이 정하는 사상
\[
  G(T)\times X(T)\simeq(G\times X)(T)
  \longrightarrow X(T)\times X(T)
\]
이 단사인 것이다. 다시 말해, 군 $G(T)$가 집합 $X(T)$에
\emph{자유롭게} 작용하는 것이다. 즉 $g\in G(T)$, $x\in X(T)$에 대하여
$g\cdot x=x$이면 $g$는 군 $G(T)$의 항등원이어야 한다. 이때 $G$가
\emph{$X$에 자유롭게 작용한다}고 한다(또는 $X$가 $G$에 관한
\emph{$\mathcal C$-주공간}이라고 한다). 그러면 쌍 $(p_1,p_2)$에 결부된
동치관계를 \emph{$X$에 자유롭게 작용하는 군 $G$가 정하는 동치관계}라고
한다. $X\times X$도 존재하고 쌍 $(p_1,p_2)$가 정하는 사상
\[
  p:G\times X\longrightarrow X\times X
\]
을 생각할 때, $G$가 자유롭게 작용한다는 조건은 $p$가
\emph{단사사상}이라는 뜻이다.

물론 $G$가 $X$에 자유롭게 작용하지 않을 때에도 $G$에 의한 $X$의 몫,
즉 앞의 쌍 $(p_1,p_2)$의 여핵이 존재하는 기준을 얻고자 한다.

이 여핵은 흔히 $X/G$로 쓰며, $G$가 왼쪽에서 작용한다면 $G\backslash
X$로 쓰는 편이 낫다(앞의 표기는 $G$가 오른쪽에서 작용하는 경우를 위해
남겨 둔다). $p$에 의한 $G\times X$의 ``상''이 존재하더라도(예를 들어
이 상을 $p$가 그것을 거쳐 인수분해되는 $X\times X$의 가장 작은
부분대상으로 정의하더라도), 이를 $R$이라 하면 $R$은 대개 $X$의
동치관계가 아니다. 이때 $R$에 의한 몫(더 정확히는 두 사영
$\operatorname{pr}_i$가 유도하는 $R$에서 $X$로 가는 사상들의 쌍에 의한
몫)으로 곧바로 넘어가려 하면, 처음 쌍 $(p_1,p_2)$의 고유한 성질들을
돌이킬 수 없이 잃는다. 따라서 작용하는 $\mathcal C$-군이 정하는 쌍에
직접 적용되는 동치관계 개념의 일반화를 찾아야 한다.

\subsection*{4. 전동치관계}
\addcontentsline{toc}{subsection}{4. 전동치관계}

모든 사상이 동형사상인 범주를 \emph{준군}이라고 부름을 상기하자. 한편
범주는 두 바탕집합 $(X,R)$, 곧 \emph{대상}들의 집합과 \emph{화살표}들의
집합에 다음 구조들을 준 것으로 정의해야 한다.
\src{106}
\begin{enumerate}
\item[(i)] 다음 두 사상
\[
  p_1,p_2:R\rightrightarrows X
\]
을 각각 \emph{시점 사상}과 \emph{종점 사상}이라고 한다.

\item[(ii)] 다음 사상
\[
  \pi:(R,p_2)\times_X(R,p_1)\longrightarrow R
\]
을 \emph{합성 사상}이라고 한다.
\end{enumerate}
이 데이터는 여기서 되풀이하지 않을 잘 알려진 공리들을 만족해야 한다.
그 공리들은 몇몇 도식의 가환성과, 다른 두 도식을 가환하게 하는 사상
$D:X\to R$의 존재(그러한 사상은 반드시 유일하다)로 명시할 수 있다.
$D$는 각 대상을 그에 대응하는 항등화살표로 보내며 다음을 만족한다.
\[
  p_1D=p_2D=\operatorname{id}_X.
\]
이 범주가 준군이라는 말은 각 화살표를 역화살표로 보내는, $R$의
\emph{대칭사상}이라 불리는 사상
\[
  s:R\longrightarrow R
\]
이 존재한다는 뜻이다(그러한 사상은 반드시 유일하다). 이는 $s$,
$\Delta$와 앞의 데이터를 사용해 만든 다른 네 도식의 가환성으로 명시할
수 있으며, 그 가운데 처음 두 식은 다음과 같다.
\[
  p_1s=p_2,\qquad p_2s=p_1.
\]

이 점들을 상기하면 ([2], A, \S~1)의 일반 정의는 특히 임의의 범주
$\mathcal C$의 두 대상 $X,R$에 $\mathcal C$-\emph{범주} 구조, 각각
$\mathcal C$-\emph{준군} 구조를 준다는 말의 뜻을 보여 준다. 정의상 이는
$\mathcal C$의 모든 대상 $T$에 대하여, 대상들의 집합이 $X(T)$이고
화살표들의 집합이 $R(T)$인 집합론적 의미의 범주 구조, 각각 준군 구조를
주되, $T$가 변할 때 이 구조들이 $T$에 따라
\guillemotleft 함자적으로 변하게\guillemotright{} 하는 것이다. 따라서
\emph{시점 사상}과 \emph{종점 사상}이라 불리는 다음 두 사상이 정해지고,
\[
  p_1,p_2:R\rightrightarrows X
\]
필요한 섬유곱이 존재할 때에는 다음 사상도 정해진다.
\src{107}
\[
  \pi:(R,p_2)\times_X(R,p_1)\longrightarrow R
\]
이를 \emph{합성 사상}이라고 한다. 이 세 데이터만으로 $X,R$ 위의 범주
구조가 결정되며, 이 데이터에 부과할 공리는 다음과 같다. 모든 $T$에
대하여 $X(T),R(T)$에 대응하는 세 데이터가 이 두 집합 위에 범주 구조
(각각 준군 구조)를 정해야 한다. 필요하다면 이는 잘 정해진 사상
\[
  D:X\longrightarrow R
\]
을 포함하는 몇몇 도식의 가환성으로 나타낼 수 있고, 준군의 경우에는 잘
정해진 사상
\[
  s:R\longrightarrow R
\]
도 포함한다. 이 도식들은 ``집합론적'' 경우와 마찬가지로 명시된다. 이처럼
공리를 번거롭게 해석하는 일은 다행히 실제 계산에는 쓸모가 없다. 데이터와
공리들을 몇몇 섬유곱 사이의 사상들과 사상들의 등식으로 나타낼 수 있다는
사실의 이론적 의의는 오직 다음과 같다. 왼쪽 완전 함자
$F:\mathcal C\to\mathcal C'$가 있으면(즉 유한곱 및 섬유곱과 가환하는
함자가 있으면), 이 함자는 $\mathcal C$-범주(각각 $\mathcal C$-준군)를
$\mathcal C'$-범주(각각 $\mathcal C'$-준군)로 보낸다(단, $\mathcal
C$에서 유한곱과 섬유곱이 존재한다고 가정한다).

실제로는 사상 $p_1,p_2,\pi,D,s$를 $\mathcal C$의 알맞은 준단체 대상
또는 단체 대상에서의 \emph{단체 연산}으로 해석하는 법을 아는 것이
중요하다. 적어도 $\mathcal C$에서 섬유곱들이 존재할 때에는 그렇다.
용어를 고정하기 위해 \emph{표준 단체들의 범주 $\mathcal S$}를 다음과
같이 도입하자. 그 대상들은
\[
  \Delta_n=[0,n],\qquad\text{단 }n\geq 0
\]
꼴의 유한집합($0$부터 $n$까지의 정수 구간)이고, 사상 또는 화살표는 이
유한집합들 사이의 \emph{임의의 사상}이다. 범주 $\mathcal S$는 유한집합
사상을 사상으로 삼는 \emph{공집합 아닌} 유한집합들의 범주와 동치임을
유의하자. $\mathcal S$에서는 \emph{공집합 아닌} 유한 대상족의 쌍대곱이
명백히 존재하고, 한 대상 아래의 두 대상의 밂(섬유곱의 쌍대 연산)도
존재한다. $\mathcal S'$은 $\mathcal S$와 대상은 같지만 사상이
$\Delta_n$들 사이의 \emph{증가 사상}인 부분범주를 뜻한다. 이 범주는
전순서가 주어진 공집합 아닌 유한집합들의 범주와 동치이다. 이 범주에서는
두 대상의
\src{108}
쌍대곱이 결코 존재하지 않고, 한 대상 $C$ 아래의 두 대상 $A,B$의 밂도
일반적으로 존재하지 않는다. 예를 들어 $C=\Delta_0$, $A=B=\Delta_1$로
두고 두 구조사상 $u:C\to A$와 $v:C\to B$를 서로 같게 잡으면 된다.
그러나 다음과 같은 몇몇 경우에는 존재한다.
\[
  C=\Delta_0,\qquad A=\Delta_m,\qquad B=\Delta_n,\qquad
  u(0)=m,\qquad v(0)=0,
\]
이때
\[
  \Delta_m\amalg_{\Delta_0}\Delta_n=\Delta_{m+n}
\]
이다.

범주 $\mathcal C$의 \emph{단체 대상}(각각 \emph{준단체 대상})은 정의상
$\mathcal S$(각각 $\mathcal S'$)에서 $\mathcal C$로 가는 반변 함자
$K$이다. 따라서 단체 대상을 제한하면 준단체 대상이 되지만, 첫 개념과
둘째 개념의 본질적 차이는 $K_n=K(\Delta_n)$에 \emph{대칭 연산}이 있다는
점이다. 이 연산들은 $\Delta_n$의 $\mathcal S$에서의 자기동형군으로 보는
$n+1$개 문자 위의 대칭군 원소들을 함자 $K$로 보낸 것에 대응한다.

이제 모든 $n$에 대하여 $\Delta'_n$(각각 $\Delta''_n$)을 다음 유한범주로
두자. 그 대상들의 집합은 $\Delta_n$이고, $\Delta_n$ 위의 무차별
순서관계(각각 자연스러운 전순서)가 이 범주를 정한다(화살표들의 집합은 그
순서관계의 그래프이다). $\Delta'_n$(각각 $\Delta''_n$)은 $\mathcal
S$(각각 $\mathcal S'$)의 대상 $\Delta_n$에 함자적으로 의존함이 명백하다.
이제 $Z$가 범주이면, $\Delta_n$이 변할 때
$\operatorname{Hom}(\Delta'_n,Z)$(각각
$\operatorname{Hom}(\Delta''_n,Z)$)는 범주 $\mathcal S$(각각 $\mathcal
S'$)에서 집합들의 범주로 가는 함자, 즉 $Z$라는 범주에 \emph{결부된}
단체 집합(각각 준단체 집합)이다. 이를 각각 $Z'$와 $Z''$라 하자. 더구나
$Z'$에서 유도되는 준단체 집합에서 $Z''$로 가는 자명한 자연 준동형이
있으며, 이 준동형이 동형일 필요충분조건은 $Z$가 준군인 것이다. 이제
다음이 성립한다.

\result{PROPOSITION}{4.1} 범주들의 범주에서 준단체 집합들의 범주로 가는
함자 $Z\longmapsto Z''$는 충만충실하며, \emph{범주}들의 범주와 다음
준단체 집합들의 범주 사이의 동치를 정한다. 곧 $\mathcal S'$에서
$(\mathrm{Ens})$로 가는 반변 함자 $K$ 가운데 위에서 지정한 꼴의 밂
$A\amalg_C B$를 집합들의 섬유곱으로 보내는 것들의 범주이다. 마찬가지로
준군들의 범주에서 단체 집합들의 범주로 가는 함자 $Z\longmapsto Z'$는
충만충실하며, \emph{준군}들의 범주와
\src{109}
단체 집합들의 범주 사이의 동치를 정한다. 곧 $\mathcal S$에서
$(\mathrm{Ens})$로 가는 반변 함자 $K$ 가운데 밂을 섬유곱으로 보내는
것들의 범주이다.

따라서 어떤 구조를 다른 구조의 말로 해석할 때 요구되는 대로 ``동형을
제외하고'' 논한다면, 범주를 특수한 준단체 집합으로, 준군을 특수한 단체
집합으로 볼 수 있다. 집합론적 경우로 환원하는 통상적인 방법으로부터
다음이 따른다.

\result{COROLLAIRE}{4.2} 앞의 명제는 다음을 바꾸어도 여전히 성립한다:
여전히 성립한다. 즉 범주, 준군, 단체 집합 대신 각각
$\mathcal C$-범주, $\mathcal C$-준군, $\mathcal C$의 단체 대상을 쓰되,
$\mathcal C$에서 섬유곱이 존재한다고 가정한다.

$\mathcal C$의 범주 $(X,R,\ldots)$에 결부된 $\mathcal C$의 반단체 대상
$K$는 다음과 같이 명시할 수 있다. $K$의 성분
$K_n=K(\Delta_n)$을 $(R,p_1)$의 $X$ 위에서의 $n$번째 섬유
거듭제곱으로 해석하거나\footnote{인쇄본에는
“$(n+1)$번째 섬유 거듭제곱”이라 되어 있고, 뒤이은 공식에는
$K_0=R$이라 되어 있다. 신경의 대상 성분은 $K_0=X$이고 화살표
성분은 $K_1=R$이므로, 두 첨자 수정은 서로 연동된다.}, 더 나아가
다음 점화식으로 정의한다.
\[
  K_0=X,\qquad
  K_n=(K_{n-1},p_{n-1}^{(n-1)})\times_X(R,p_1),
\]
여기서 $p_i^{(n-1)}$ ($0\leq i\leq n-1$)는 $K_{n-1}$에서 $X$로 가는
자연 사영들이며, 이들 역시 점화적으로 정의된다. 이렇게 하면
$p_1,p_2,\pi,D,s$는 각각 $\mathcal S$의 다음 사상에 대응하는 단체
연산으로 해석된다. 곧 $\Delta_1$의 $0$번 면, $\Delta_1$의 $1$번
면, $\Delta_2$의 $(0,2)$번 면, 퇴화 $\Delta_1\to\Delta_0$,
$\Delta_1$의 대칭이다. 나머지 모든 반단체(각각 단체) 연산은 앞의
네 개(각각 다섯 개)로부터 합성과 섬유곱에 의해 형식적으로 따른다.

이제 범주의 대상 $X$에서 일반적으로 \emph{전동치관계}란, 대상들의
대상(이렇게 말해도 된다면)이 $X$인 준군을 주는 자료를 말한다. 따라서
이런 자료에는 특히 하나의 대상 $R$과 두 사상
\[
  p_1,p_2:R\rightrightarrows X
\]
이 포함된다. 그러나 동치쌍의 경우와 달리, 이 자료들만으로는 지금
고려하는 구조가 결정되지 않는다. 이 강연에서는 몫으로 넘어갈 수
있는 판정기준, 즉 쌍 $(p_1,p_2)$의 여핵을 구성할 수 있는 판정기준을
얻기 위해 이 개념을 고려한다. 따라서 문제의 진술 자체는 준군에 내재한
추가 자료를 사용하지 않는다. 그러나 뒤이을 결과들의 증명에서는
\src{110}
이 추가 자료, 특히 차원 $3$까지의 단체 연산(대칭 연산 포함)을
사용한다. 여기에는 $R$의 $X$ 위에서의 네 겹 섬유 거듭제곱이 관여한다.

$\mathcal C$의 대상 $X$에서의 동치관계는 전동치관계를 정한다. 이는
집합의 경우만 보면 충분하다. 즉 집합 $X$의 동치관계에, 대상의
집합이 $X$이고 화살표의 집합이 그 동치관계의 그래프인 준군을 대응시킨다.

$\mathcal C$-모노이드 $G$가 $\mathcal C$의 대상 $X$에 작용하면, 바탕
대상이 $R=G\times X$와 $X$인 $\mathcal C$-범주를 정한다(단,
$G\times X$가 존재한다고 가정한다). 이 범주가 $\mathcal C$-준군일
필요충분조건은 $G$가 군인 것이다. 이 역시 집합의 경우에 확인하면
충분하다. 화살표 $(g,a)$와 $(g',g\cdot a)$의 합성은 다음으로
정의한다.
\[
  (g',g\cdot a)(g,a)=(g'g,a),
\]
즉 $a,b\in X$이면 $\operatorname{Hom}(a,b)$는 정의상 $a$를 $b$로 보내는
운반자이고, 사상들은 $G$의 원소들을 곱하여 합성한다.

\medskip
\noindent\textbf{비고. ---} (4.1)과 같은 진술이 일으키는 논리적 어려움은,
고려하는 모든 대상이 하나의 고정된 “우주”(이 우주 자체도 집합이다) 안에
있다고 암묵적으로 가정하면 피할 수 있다.

\fsection{5}{유한 평탄 동치관계로 나눈 몫}

\result{정리}{5.1} $X=\operatorname{Spec}(B)$를 아핀 스킴, $R$을 $X$의
전동치관계라 하고, 그 성분 $R_1$이 아핀이라 하자. 곧
$R_1=\operatorname{Spec}(C)$라 하자. 제1사영 $p_1:R_1\to X$가 유한이고
국소 자유인 사상이라고 가정하자. 다시 말해 대응하는 환 준동형사
$p_1^*:B\to C$에 의해 $C$가 유한 생성 사영 $B$-가군이 된다고
가정한다. 두 준동형사
\[
  p_1^*,p_2^*:B\rightrightarrows C
\]
의 쌍의 핵인 $B$의 부분환을 $A$라 하자($p_1^*(b)=p_2^*(b)$를 만족하는
원소 $b$들의 집합이다). $Y=\operatorname{Spec}(A)$라 하고, $A$에서 $B$로의
포함에 의해 정의되는 사상을 $f:X\to Y$라 하자. 이 조건들 아래에서
다음이 성립한다.
\begin{enumerate}
\item[(i)] $B$는 $A$ 위에서 정수적이다. 즉 $f$는 정수적 사상이다.

\item[(ii)] 사상 $f$는 전사이고, 그 섬유들은 $X$에서 $R$로 나눈
집합론적 동치류
\src{111}
$p_2\allowbreak(p_1^{-1}(x))$이다. 또 $Y$의 위상은 $X$의 몫위상이다.

\item[(iii)] $Y$는 준스킴의 범주에서 $X$를 $R$로 나눈 몫이다.

\item[(iv)] $R$이 \emph{동치관계}에서 나왔다면 사상 $f:X\to Y$는
유한 국소 자유이고(즉 $B$는 유한 생성 사영 $A$-가군이고), 그
동치관계는 유효하다. 즉 $R_1\to X\times_Y X$는 동형사이다.
\end{enumerate}

이 정리는 유한군 $G$가 환 $B$에 자기동형사들로 작용하는 경우와 불변
부분환 $A$에 관한 잘 알려진 정리를 일반화하며, 증명도 알려진
증명과 유사하다. (iii)은 다음과 같이 더 정밀히 진술할 수 있다.

\result{따름정리}{5.2} 표준 사상 $R_1\to X\times_Y X$는 \emph{전사}이다.

$R$을 여전히 $X$의 “유한이고 국소 자유인” 전동치관계라 하되, 이제
$X$는 임의의 준스킴이라 하자. $fp_1=fp_2$를 만족하는 준스킴 $Y$와
사상 $f:X\to Y$를 찾을 수 있고, 더나아가 $Y$ 위의 환의 층의
준동형사열
\[
  \mathcal O_Y\longrightarrow f_*(\mathcal O_X)
  \rightrightarrows g_*(\mathcal O_R)
\]
이 완전하다고 가정하자(여기서 $g=fp_i$). 그러면 이 정리에서 (i)--(iv)와
같은 결론을 얻는다. 특히 (iii)에 의해 $Y$는 $X$를 $R$로 나눈
몫이고, 따라서 유일한 동형사까지 결정된다. 이 조건들 아래에서
$X$의 전동치관계 $R$을 \emph{허용 가능}하다고 한다. 이 정의로부터
다음을 얻는다.

\result{정리}{5.3} $X$를 준스킴, $R$을 $X$의 전동치관계라 하고,
$p_1:R_1\to X$가 유한 국소 자유 사상이라 하자. $R$이 허용 가능할
필요충분조건은 $X$에서 $R$로 나눈 모든 집합론적 동치류
$p_2(p_1^{-1}(x))$가 아핀 열린집합에 포함되는 것이다. 이 조건은
$X$의 모든 유한 부분집합이 아핀 열린집합에 포함되면 항상 만족하며,
예를 들어 $X$가 아핀 스킴 위의 준사영이면 그렇다.

실제로 이때 $X$에서 $R$로 나눈 모든 동치류가 \emph{$R$-불변인}
아핀 열린집합에 포함됨을 쉽게 보일 수 있고, 정리 5.1에서 얻은
조각들을 붙여 몫 $Y$를 구성한다.
\src{112}

\result{따름정리}{5.4} 이 조건이 만족하고 더나아가 $R$이 동치관계에서
나왔다고 하자. 그러면 그 동치관계는 유효하다. 즉
$R_1\to X\times_Y X$는 동형사이고, $f:X\to Y$는 유한 국소 자유
사상이다.

“강하”에 의해 곧바로 다음을 얻는다.

\result{따름정리}{5.5} (5.4)의 조건 아래서, $X$가 $Y$ 위에서 모든
곳에서 계수 $n$일 필요충분조건은 $(R_1,p_1)$이 $X$ 위에서 모든 곳에서
계수 $n$인 것이다. $X,R_1$이 $Z$-준스킴이고 $p_1,p_2$가
$Z$-사상이라 하자. 그러면 $Y$도 $Z$-준스킴이며, $X$가 $Z$
위에서 평탄할 필요충분조건은 $Y$가 $Z$ 위에서 평탄한 것이다.

요약하면 다음과 같다.

\result{스콜리움}{} 준스킴의 유한이고 국소 자유이며 전사인 사상
$f:X\to Y$를 주는 것은, 유한 국소 자유 사상 $p_1:R\to X$를
갖고 모든 류 $p_2(p_1^{-1}(x))$가 아핀 열린집합에 포함되는 동치관계
$R$를 갖춘 준스킴 $X$를 주는 것과 동치이다.

\medskip
\noindent\textbf{비고 5.6.}

\smallskip
\noindent $1^\circ$ 뇌터 가정은 필요하지 않았다.

\smallskip
\noindent $2^\circ$ 이 몫 구성 개념은 카르티에의 “비분리 강하”를 특수한
경우로 포함한다. 이는 $f_*(\mathcal O_X)$가 $\mathcal O_Y$에 대해
$p$-기저를 갖는 유한 국소 자유 사상 $f:X\to Y$를 결정하는 문제에
해당한다. 여기서 $X$는 소수 $p>0$에 의해 구조층 $\mathcal O_X$가
소멸되는 주어진 준스킴이다. 이 결과는 국소환에 어떤 정칙성 가정도
두지 않고, $X$가 체 위의 대수적 스킴이라고 가정하지 않은 채로 쉽게
표현할 수 있음에 유의하라. $X$를 특성 $p$인 체의 스펙트럼으로 취하면
야코브슨--부르바키 이론을 얻는다.

\smallskip
\noindent $3^\circ$ 가브리엘은 이전에 체 $k$ 위의 유한 가환군에 대한
몫 구성 이론에서 정리 (5.3)의 특수한 경우를 얻었다. [(7.4)와 비교하라.]

\fsection{6}{고유 평탄 동치관계로 나눈 몫}

\result{정리}{6.1} $S$를 국소 뇌터 준스킴이라 하고, $X$를 준사영
$S$-스킴이라 하며, $R$을 $X$ 안의 전동치관계라 하자. 다음을 가정한다.
\src{113}
\begin{enumerate}
\item[(a)] $p_1:R_1\to X$는 고유하고 평탄하다.
\item[(b)] $R_1\to X\times_S X$는 유한 사상이다. 또는 (a)에 따라 이와
동치이게, 모든 섬유가 유한하다(이 조건은 $R$이 동치관계에서 오는 경우
자동으로 성립한다).
\end{enumerate}
이 조건들 아래에서 다음이 성립한다.
\begin{enumerate}
\item[(i)] $X/R=Y$가 존재하고, ($S$가 뇌터이고 $R$이 동치관계에서 오는
경우) $S$ 위에서 준사영이다.\footnote{인쇄본에는 두 번째 가정이 없다.
강연 212에 딸린 정오표는 이때 준사영성이 $R$이 동치관계에서 오는 경우에만
증명된 것으로 보인다고 명시한다.}

\item[(ii)] 표준 사상 $f:X\to Y$는 전사이고 고유이며 열린 사상이다. 그
섬유들은 $R$에 대한 $X$의 동치류 $p_2(p_1^{-1}(x))$이므로, $Y$는
$R$이 정하는 집합론적 동치관계에 의한 $X$의 몫 위상공간과 동일시된다.
끝으로 $R_1\to X\times_Y X$는 전사이다.

\item[(iii)] $R$이 동치관계에서 오는 경우, 그 동치관계는 유효하다. 즉,
$R_1\to X\times_Y X$는 동형사상이다. 또한 $f:X\to Y$는 평탄하다
(따라서 충실 평탄하다).
\end{enumerate}

증명에서는 $R$에 대한 $X$의 적절한 준단면들을 사용하여 (5.1)로
귀착한다. 증명은 슈발레 세미나에서 몫 대수군을 구성하는 것과 유사하다.

요약하면 다음과 같다.

\result{스콜리움}{} $S$가 국소 뇌터이고 $X$가 $S$ 위에서 준사영이라고
하자. $X$에서 $S$-준스킴 $Y$로 가는 고유하고 충실 평탄이며 전사인 사상
$f:X\to Y$를 주는 것은, $p_1:R\to X$가 고유하고 평탄한 $X$ 위의
동치관계 $R$을 주는 것과 동치이다.

같은 방법으로 다음 결과를 얻는다.

\result{정리}{6.2} $S$를 뇌터 준스킴, $X$를 $S$ 위에서 유한형인
준스킴, $R$을 $S$-준스킴 $X$ 안의 전동치관계라 하자. 다음을 가정한다.
\begin{enumerate}
\item[(a)] $p_1:R_1\to X$는 평탄하고 유한형이다.
\item[(b)] 사상 $R_1\to X\times_S X$는 준유한이다(즉, 모든 섬유가
유한하다).
\end{enumerate}
그러면 $R$에 대해 \emph{포화된} $X$의 조밀한 열린집합 $U$가 존재하여
다음을 만족한다.
\begin{enumerate}
\item[(i)] $R_U$가 $R$이 $U$에 유도하는 전동치관계이면, $U/R_U$가
존재하고 $S$ 위에서 유한형이다.

\item[(ii)] 표준 사상 $U\to U/R_U$는 전사이고 열린 사상이며, 그 섬유들은
$R_U$에 대한 집합론적 동치류이다(따라서 $U/R_U$는 $R_U$가 정하는
집합론적 동치관계에 의한 $U$의 몫 위상공간이다). 끝으로 사상
\[
  (R_U)_1\longrightarrow U\times_{U/R_U}U
\]
은 전사이다.
\src{114}

\item[(iii)] $R$이 동치관계에서 오는 경우에는, $U\to U/R_U$가 충실
평탄이고 $R_U$가 유효하도록 $U$를 택할 수 있다.
\end{enumerate}

이는 본질적으로 ``쌍유리적'' 성격의 결과이다.

\medskip
\noindent\textbf{비고 6.3.}

\smallskip
\noindent $1^\circ$ (6.1)과 (6.2)에서 가정 (b)가 불필요한지는 알지
못한다. 이 가정 때문에 군으로 몫을 취할 때 실제로는 안정자들이 모두
유한군인 경우에 한정해야 한다.

\smallskip
\noindent $2^\circ$ 평탄성 가정 없이도 (6.1), (6.2)와 비슷한 결과가
있는지 물을 수 있다. 이 방향의 반례는 하나도 알지 못한다. 반면 평탄성
가정을 유지하고 $p_1:R\to X$가 평탄하고 준유한이지만 유한하지 않은
동치관계로 한정하더라도, $X$가 아핀일 때 $R$이 유효하지 않은 일이 생길
수 있다. 나가타 다양체를 덮는 아핀 열린집합들 위에 유도되는 동치관계들을
취하면 된다(그 다양체에서는 원소가 두 개인 군이 ``허용 가능하지 않게''
작용한다).

\fsection{7}{응용}

서론에서 말했듯이, (6.1)의 가장 중요한 응용은 피카르 스킴을 구성하는
것과 여러 다른 ``모듈라이'' 문제의 해를 구성하는 것이다. 이 문제들은
나중에 다시 다루겠다.

시무라의 다음 결과에 대한 간단한 증명을 얻는다.

\result{명제}{7.1} $V$를 분수체가 $K$인 이산 값매김환이라 하고,
$A$를 $V$ 위에 정의된 아벨 스킴이라 하자. 그러면
$A\otimes_V K$의 한 몫과 아이소제니로 연결되는 모든 $K$ 위의 아벨 스킴
$B'$은 ``$V$에 대해 좋은 환원을 갖는다.'' 다시 말해 $V$ 위의 어떤 아벨
스킴 $B$에 대하여 $B\otimes_V K$와 동형이다.\footnote{인쇄본에는
$B\otimes_V A$라고 되어 있다. $B$는 $V$-스킴이고 이 주장은 그 일반섬유에
관한 것이므로, 바탕 인자는 반드시 $K$여야 한다.} (상기하듯 이러한 $B$는
본질적으로 유일하다.)

$B'$이 $A_K$를 한 부분군 스킴 $C'$로 나눈 몫이라고 가정해도 된다.
(주의: 일반적으로 $C'$은 ``축소''가 아니다. 즉 그 국소환들은 멱영원을
가질 것이다.) $A$ 안에서 $C'$의 ``폐포''인 닫힌 부분스킴 $C$, 다시 말해
$C_K$가 $C'$을 포함하도록 하는 $A$의 가장 작은 닫힌 부분스킴을 생각하자.
그러면 $C_K=C'$이고, $V$가 이산 값매김환이라는 사실로부터 $C$가 $V$ 위의
$A$의 부분군 스킴임을 쉽게 얻는다. $A$가
$\operatorname{Spec}(V)=S$ 위에서 고유하므로 $C$도 그러하고, 한편 $A$는
더 나아가 $S$ 위에서 사영이다. 따라서 (6.1)을 적용하여 $A/C=B$를
구성할 수 있고,
\src{115}
이것이 구하는 $B$이다.

마지막으로 본질적으로 잘 알려진 논증을 쓰면 (6.2)로부터 다음 결과를
얻을 수 있다.

\result{정리}{7.2} $S$를 아르틴 환의 스펙트럼이라 하고, $F$와 $G$를
$S$ 위의 유한형 군 스킴이라 하며, $u:F\to G$를 $S$ 위의 군 스킴의
준동형이라 하자. 다음을 가정한다.
\begin{enumerate}
\item[(i)] $F$는 $S$ 위에서 평탄하다.
\item[(ii)] $u$의 핵은 유한하다.
\end{enumerate}
이 조건들 아래에서 몫 스킴 $G/F$가 존재하고, 표준 사상 $G\to G/F$는
전사인 열린 사상이며, 그 섬유들은 $G$ 위의 $F$의 오른쪽 작용이 정하는
집합론적 동치류들이다. 끝으로 $u$가 단사사상이면 사상 $G\to G/F$는
평탄하고, 사상
\[
  G\times F\longrightarrow G\times_{G/F}G
\]
은 동형이다. 다시 말해 $G$는 $G/F$ 위의 주동차 공간이고 그 구조군은
오른쪽에서 작용하는 $F$, 더 정확히는 $G/F$ 위의 군 스킴으로 보는
$F\times_S(G/F)$이다(cf. [1], B, \S~6).

\result{따름정리}{7.3} 이 조건들 아래에서 $G$가 $S$ 위에서 평탄일
필요충분조건은 $G/F$가 $S$ 위에서 평탄인 것이다. 이 조건이 성립하면
$F$로 나눈 몫을 취하는 일은 $S$의 임의의 기저확대와 가환하며, $F$가
$G$의 정규 부분군 스킴이면 $G/F$에 $G$를 $F$로 나눈
\emph{몫군}의 구조를 줄 수 있다.

$S$가 체의 스펙트럼이면 모든 $S$-준스킴이 자동으로 $S$ 위에서
평탄하므로 상황은 특히 간단하다. 다음을 얻는다.

\result{따름정리}{7.4} $F,G$를 체 $k$ 위의 유한형 군 스킴이라 하고,
$u:F\to G$를 $k$-군의 준동형이라 하자. 그러면 $u$는
$F\to F'\to G$로 인수분해된다. 여기서 $F\to F'$은 $F$의 닫힌 부분군
$\operatorname{Ker}u$로 나누어 몫을 취하는 준동형이고, $F'\to G$는
닫힌 몰입 사상인 군 준동형이다. 몫 $G/F=G/F'$이 존재한다. 통상적인
형식 체계(가령 뇌터 정리류와 같은 것)가 $k$ 위의 대수군들 사이에서
성립한다.

이 결과를 쓰면 고전적 의미의 대수군(즉 $k$ 위에서 기약이고 $k$ 위에서
매끄러운 대수군)에서 몫을 취하는 일과 카르티에가 고찰한 ``무한소''
부분군으로 몫을 취하는 일을 통일적으로 다룰 수 있다.
\src{116}
뒤도네의 형식군 연구에 이어 카르티에가 도입한 ``하이퍼대수''를 $k$ 위의
형식적 스킴의 범주 안의 군으로 보고, 해당하는 경우(그것들이 $k$ 위에서
유한 계수인 하이퍼대수에 대응하는 경우)에는 $k$ 위의 유한 대수군으로
보는 것이 유리하다.

\fsection{8}{한 추측}

이는 ``힐베르트 스킴''의 몇몇 부분스킴에 작용하는 사영군으로 몫을
취할 필요에 답한다. (스킴 이론에서 힐베르트 스킴은 차우 다양체를
대신한다.)

$S$를 준스킴이라 하고 $n$을 정수라 하자. $S$ 위의 각 준스킴 $S'$에
$\mathcal O_{S'}$의 단면환에 값을 갖는 $n$행 $n$열 가역행렬들의 군
$\operatorname{Gl}(n,\Gamma(S',\mathcal O_{S'}))$를 대응시키자. 이로써
$S'$에 대한 반변 함자를 얻으며, 이 함자가 표현 가능함은 쉽게 보인다.
따라서 이 함자는 $S$ 위의 군 스킴에 대응하고, 그 군 스킴은 실제로
$S$ 위에서 아핀이며 $\operatorname{Gl}(n)_S$로 나타낸다. 이 구성은
기저변환과 양립하므로, 실제로는 모든 것이 $\operatorname{Gl}(n)$으로
나타내는 $\mathbf Z$ 위의 한 군 스킴에서 나온다.
$\operatorname{Gl}(1)$은 곱셈군이라 하며 흔히 $\mathbf G_m$으로 나타내고,
함자
$S\longmapsto\Gamma(S,\mathcal O_S)^*$, 즉 $S$ 위의 ``가역원''들의 군에
대응한다. 명백한 준동형
$\operatorname{Gl}(1)\to\operatorname{Gl}(n)$이 있고, 몫군을 쉽게
구성할 수 있다. 이를 $\operatorname{GP}(n-1)$로 나타내며
\emph{$\mathbf Z$ 위의 $n-1$차 사영군}이라 한다. 이는 $S$에 층
$\widetilde{\operatorname{Gl}(n)}_S/
\widetilde{\operatorname{Gl}(1)}_S$의 단면군을 대응시키는 함자를
표현한다. 여기서 $\widetilde{\operatorname{Gl}(n)}_S$는 $S$ 위의
$\operatorname{Gl}(n)_S$의 단면들의 싹의 층을 뜻한다. ($S$ 위의
$\operatorname{GP}(n-1)_S$의 단면들이 일반적으로 $S$ 위의
$\operatorname{Gl}(n)_S$의 단면들에서 오지는 않는다는 점에 주의해야
한다!) 또한 적어도 $S$가 뇌터일 때 $\operatorname{GP}(n-1)$은 함자
$S\longmapsto\operatorname{Aut}_S(\mathbf P_S^{n-1})$도 표현함을 증명할
수 있다. 여기서 $\mathbf P_S^{n-1}$은 $S$ 위의 상대 차원 $n-1$인 표준
사영 스킴이다. 이것이 사영군이 모듈라이 이론에 들어오는 방식이다.

$S$를 뇌터 스킴이라 하자. 원하면 아핀이라고 가정해도 된다. $X$를
준사영 $S$-준스킴이라 하고, $S$에 대해 매우 풍부한 가역층
$\mathcal L$을 갖추었다고 하자. 군 $G=\operatorname{GP}(n)_S$가 $X$에
작용하는 동시에 $\mathcal L$에도 작용하고(그 $X$ 위의 작용과 양립하는
방식으로), $X$에 \emph{자유롭게} 작용한다고 가정하자.\footnote{인쇄본에는
``$S$에''라고 되어 있다. 여기서 생각하는 작용과 몫 조건은 $X$에 관한
것이다.}
\src{117}

\result{추측}{8.1}\footnote{강연 212와 관련된 정오표는 이 추측이
``결국 거짓''이라고 지적한다. 또한 비고 8.1의 긍정적 사실이 그 뒤
증명되었다고 밝히고 멈포드의 연구와 그 정련들을 참조한다.} 앞의 조건들
아래에서 다음이 성립한다.

\smallskip
\noindent $1^\circ$ \emph{$G$가 정하는 동치관계는 유효하고, 몫
$Y=X/G$는 $S$ 위에서 유한형이며, 표준 사상 $f:X\to Y$는 평탄이고
전사이다.} (따라서 $X$는 $Y$ 위에서
$G\times_S Y=\operatorname{GP}(n)_Y$를 구조군으로 하는 동차 주다발이
된다.)

\smallskip
\noindent $2^\circ$ $\mathcal L'$을 $f$에 의한 ``충실 평탄 강하''로
$\mathcal L$에서 얻는 $Y$ 위의 가역층이라 하자(cf. [1], B,
\S~th. 1). 그러면 $\mathcal L'$은 \emph{$S$에 대해 $Y$ 위에서
예비풍부}이다. 즉 어떤 정수 $m$과 $Y$에서 알맞은 표준 사영 스킴
$\mathbf P_S^N$으로 가는 준유한 사상이 존재하여,
$(\mathcal L')^{\otimes m}$은 $\mathcal O_{\mathbf P_S^N}(1)$의 역상과
동형이다.

$X$가 $S$ 위에서 분리되어 있더라도 $Y$는 $S$ 위에서 분리되어 있지 않을
수 있음에 유의하라. (이는 전혀 병적이지 않은 ``모듈라이 문제''에서
나타나는 상황이다.) $1^\circ$가 성립하면 $Y$가 분리일 필요충분조건은
$G$가 정하는 동치관계의 그래프가 닫혀 있는 것, 즉
$G\times X\to X\times X$의 상이 닫혀 있는 것이다. (이때 이 사상은
닫힌 몰입 사상이다.) $Y$가 분리이면 $\mathcal L'$이 $S$에 대해 $Y$
위에서 예비풍부일 필요충분조건은 그것이 풍부한 것이다. 다시 말해
알맞은 텐서 거듭제곱이 사영 몰입 사상을 정의하는 것이다. 서론에서
언급한 모듈라이 문제들에서는 이렇게 얻는 동치관계의 그래프가 실제로
닫혀 있음을 보일 수 있다.

\medskip
\noindent\textbf{비고 8.1. ---} 논의를 구체화하기 위하여, 그리고 현재
실제에서 가장 중요한 경우이므로 $G=\operatorname{GP}(n)_S$라고
가정하였다. $G$에 둘 만한 합리적인 가정은 오히려 $G$가 도호쿠 군들
가운데 어느 하나의 $S$ 위의 ``꼴''이라는 것일 듯하다. (이 군들의
정수 위의 구성은 슈발레도 수행하였다.) 앞의 추측을 뒷받침하는 방향에서
강연자가 아는 유일한 긍정적 사실은 다음과 같다. $X$를 \emph{표수 $0$인
체 $k$ 위의 아핀 스킴이라 하고, 군 $\operatorname{Gl}(n)_k$ 또는
$\operatorname{GP}(n-1)_k$가 $X$에 자유롭게 작용한다고 하자. 그러면
$G$가 정하는 동치관계는 유효하고, 몫 $X/G$도 아핀이며, 사상
$X\to X/G$는 평탄이고 전사이다.} 그 증명은 다음 사실을 사용한다.
(현재로서는 표수 $0$에서만 증명되어 있다.) $G$가 $G$의 아핀 좌표환
$A$에 작용하게 하고 이를 $k$-벡터공간으로 보면, $G$의 자명 표현은
$G$-안정인 유한차원 $k$-벡터 부분공간의 한 합성열 안에서 단 한 번만
나타난다. $G$의 선형표현론을 체계적으로 사용하면
\src{118}
적어도 바탕체 위에서는 결국 이 추측의 증명을 얻을 수 있을 듯하다.
바탕이 더 이상 체가 아니면 강연자는 아무것도 알지 못한다.

\section*{참고문헌}
\addcontentsline{toc}{section}{참고문헌}

\begin{description}
\item[{[1]}] Grothendieck (Alexander). --- \emph{Techniques de descente et
théorèmes d'existence en géométrie algébrique, I: Généralités}, Séminaire
Bourbaki, 1959/60, 제190호, 29쪽.

\item[{[2]}] Grothendieck (Alexander). --- \emph{Techniques de descente et
théorèmes d'existence en géométrie algébrique, II: Le théorème d'existence en
théorie formelle des modules}, Séminaire Bourbaki, 1959/60, 제195호, 22쪽.

\item[{[3]}] Grothendieck (A.) 및 Dieudonné (J.). --- \emph{Éléments de
géométrie algébrique, I: Le langage des schémas}. --- Paris, Presses
universitaires de France, 1960 (Institut des hautes Études scientifiques,
Publications mathématiques, 4).
\end{description}


\par
\renewcommand{\sourceunit}{e212}
\clearpage
\chapter*{강연 212의 정오표}
\addcontentsline{toc}{chapter}{강연 212의 정오표}
\setcounter{footnote}{2}
\src{301}

\noindent\textbf{5. TDTE III: 몫 준스킴.}

\begin{center}
[부르바키 세미나, 제13권, 1960/61, 제212호, 20쪽.]
\end{center}

\noindent\textbf{15쪽, 주장 (i).} --- $Y=X/R$가 $S$ 위에서
준사영이라는 사실은, 현재로서는 $R$이 동치관계에서 유래하는 경우에만
증명된 것으로 보인다.

\noindent\textbf{18쪽과 19쪽.} --- 다음 강연의 말미에서 지적하듯이,
문제의 추측은 분명히 거짓이다. 주 8.1에서 언급한 ``긍정적인 사실''은
여러 저자(나가타, 로젠리히트, 그로텐디크, \ldots)가 동시에 증명한 것으로
보인다.

TDTE V에서 피카르 스킴을 구성하기 위해 여기서 전개한 몫을 취하는 이론을
적용한 부분은, 힐베르트 스킴을 적절히 이용하는 것으로도 대체할 수 있음을
지적해 둔다(cf. SHMT). § 8에서 말했듯이, 여기서 제시된
\src{302}
이론의 가장 큰 공백은 고유하지 않은 동치관계로 나눈 몫의 존재 판정법이
없다는 점이다. 그러한 동치관계에는 사영군의 어떤 작용들에서 유래하는
것들이 있다. 이 방향에서 중요한 정리 하나를 멈퍼드가 얻었다\footnote{MUMFORD
(David). --- An elementary theorem in geometric invariant theory, \emph{Bull.
Amer. Math. Soc.}, t. 67, 1961, p. 483--486.}. 그의 결과를 더 정교하게 만든
형태와 이 이론에 대한 여러 응용에 대해서는 SHMT를 보라.\par

\par
\renewcommand{\sourceunit}{221}
\unit{강연 221 — 힐베르트 스킴}
\sid{221}
\src{249}
\seminar{제13년차, 1960/61, 제221호\\1961년 5월}
\paperhead{대수기하학에서의 구성 기법과 존재 정리\\
IV. 힐베르트 스킴}%
  {알렉상드르 그로텐디크 지음}

\subsection*{서론}
\addcontentsline{toc}{subsection}{서론}

[2]의 I 및 II에서 제시한 기법들은 본질적으로 고려하는 스킴에 관한
어떠한 사영성 가정에도 무관하였다. 유감스럽게도 현재로서는 이 기법들로
[2]의 강연 II에서 제기한 존재 문제들을 해결할 수 없다. 이번 강연과
다음 강연에서는 사영성 가정 아래 이 문제들을 해결하겠다. 사용하는
기법들은 전형적으로 사영적이며, [2]의 I 및 II의 결과들은 사실상 사용하지
않는다. 여기서는 [2]의 II, 제2절에서 말한 바와 같이 저우 좌표의 사용을
대신할 ``힐베르트 스킴''을 구성하겠다. 다음 강연에서는 III에서 전개한
스킴에서 몫을 취하는 이론과 힐베르트 스킴 이론을 함께 써서, 예컨대
상당히 일반적인 조건 아래 피카르 스킴([2]의 II, 제3절에서 정의하였다)을
구성할 수 있게 된다.

요약하면, 이제 사영적 구성을 위한 기법은 대체로 만족스러운 수준에
이르렀다고 할 수 있다. 다만 사영군과 같은 군이 ``고정점 없이'' 작용할 때
그 군에 의한 몫으로 이행하는 정리가 아직\footnote{강연 끝의 부록을 보라.}
빠져 있다([2]의 III, 제8절 참조). 해석기하학에서는 상황이 오히려 조금 더
나아 보인다(주어진 해석공간 위에서 ``사영적인'' 해석공간만 연구할 때).
해석공간의 경우에는 좋은 방식으로 작용하는 군에 의한 몫으로 이행할 때의
어려움이 사라지기 때문이다. 한편 대수기하학에서와 마찬가지로 해석기하학에서도
사영성 가정 없이 쓸 수 있는 구성 기법은 여전히 확립해야 한다.

\fsection{1}{층의 유계인 족: 안정성}

$k$를 체, $X$를 $k$-준스킴이라 하자. 간단히 하기 위해 $X$는 유한형이라고
가정한다. 모든 확대체 $K/k$에 대하여 $K$-준스킴
$X_K=X\otimes_k K$를 얻는다. $\mathcal F$가 $X_K$ 위의 연접층이고
$K'$이 $K$의 확대체이면,
$\mathcal F\otimes_K K'=\mathcal F_{K'}$은
$X_K\otimes_K K'=X_{K'}$ 위의 준연접층이다. 이제 $K$와 $K'$을 $k$의
임의의 두 확대체, $\mathcal F$를 $X_K$ 위의 준연접층, $\mathcal F'$을
\src{250}
$X_{K'}$ 위의 준연접층이라 하자. $K$와 $K'$에서 $k$의 한 공통 확대체
$K''$로 가는 $k$-준동형들이 존재하고 $X_{K''}$ 위에서
$\mathcal F_{K''}$과 $\mathcal F'_{K''}$이 동형이면, $\mathcal F$와
$\mathcal F'$은 \emph{동치}라고 한다. 이는 동치관계이며, 이하에서는 이
관계에 관한 층의 동치류와 그러한 동치류들의 집합을 다룬다. $X$가 $k$
위의 유한형이면,
\footnote{인쇄본에는 여기서 $X_0$라고 되어 있으나, 이 문단에서 정의된 것은
$X$뿐이다.} 연접층의 모든 동치류는 $k$의 어떤 유한형 확대체 $K$에 대한
$X_K$ 위의 한 연접층으로 정의될 수 있음에 유의하자. 따라서 연접층의
동치류를 정의할 때 $k$의 대수적으로 닫힌 확대체들만으로 한정할 수 있다.
또 초월차수가 무한인 $k$의 고정된 대수적으로 닫힌 확대체 $\Omega$ 하나로
한정할 수도 있다. 이 경우 $X_\Omega$ 위의 두 연접층 $\mathcal F$와
$\mathcal F'$이 동치일 필요충분조건은 $\Omega$의 $k$-자동동형
$\sigma$가 존재하여 $\mathcal F\otimes_\Omega(\Omega,\sigma)$가
$\mathcal F'$과 동형인 것이다. 두 정의에서 얻는 연접층의 동치류들은
서로 일대일로 대응한다.

$E$, $E'$을 $X$ 위의 연접층의 동치류들로 이루어진 두 집합이라 하자.
같은 $X_K$ 위의 연접층 $\mathcal F$, $\mathcal F'$ 가운데
$\mathcal F$의 동치류가 $E$에 속하고 $\mathcal F'$의 동치류가 $E'$에
속하는 것들에 대하여, $\mathcal F\otimes\mathcal F'$ 꼴인 모든 층의
동치류를 생각하자. 이렇게 얻는 연접층의 동치류들의 집합을
$E\otimes E'$으로 나타낸다. $\operatorname{Tor}_i(E,E')$ 등도 같은
방식으로 정의한다. 더 일반적으로, 같은 $X_K$ 위의 $n$개 연접층으로
이루어진 모든 열 $\mathcal F_1,\ldots,\mathcal F_n$에 $X_K$ 위의
연접층들의 집합 $U(\mathcal F_1,\ldots,\mathcal F_n)$을 대응시키며,
층의 동형 및 밑변환에 의한 역상과 자명한 방식으로 양립하는 함수 $U$가
주어졌다고 하자. 그러면 같은 기호 $U$로 나타내는 함수를 얻는다. 이 함수는
연접층의 동치류들로 이루어진 $n$개 집합의 모든 열
$E_1,\ldots,E_n$에 연접층의 동치류들의 집합
$U(E_1,\ldots,E_n)$을 대응시킨다.

이 절의 목적은 층의 동치류들의 어떤 집합들을 \emph{유계}라고 정의하고,
가장 흔히 쓰는 연산 $U$를 유계인 집합들에 적용하면 다시 유계인 집합을
얻음을 보이는 데 있다.

$X$를 뇌터 준스킴 $S$ 위의 유한형 준스킴이라 하자. 모든 $s\in S$에
대하여 섬유 $X_s$는 $k(s)$ 위의 유한형 준스킴이고, 앞의 의미에서
$X_s$ 위의 연접층의 동치류들을 생각한다. 이로써 ``$X/S$의 한 섬유 위의
연접층의 동치류''라는 표현과 그와 비슷한 표현들이 의미를 갖는다. 마찬가지로
각 섬유에서 따로 진행하면, $X/S$의 섬유들 위의 연접층의 동치류들의
집합들로 이루어진 계에 연산 $E\otimes E'$ 등을 적용하여 다른
\src{251}
$X/S$의 섬유들 위의 연접층의 동치류들의 집합을 얻을 수 있다.

\result{정의}{1.1} $E$를 $X/S$의 섬유들 위의 연접층의 동치류들로
이루어진 집합이라 하자. $S$ 위의 유한형 준스킴 $S'$과
$X'=X\times_S S'$ 위의 연접층 $\mathcal F'$이 존재하여, $E$가
$\mathcal F'$으로 정의되는 $X/S$의 섬유들 위의 층의 동치류들의 집합에
포함되면 $E$를 \emph{유계}라고 한다.

마지막 집합은 정의에 따라 각 $s\in S$에 다음 층들의 동치류를 대응시킨다:
$s$ 위에 놓인 $S'$의 점 $s'$이 움직일 때의
$\mathcal F'\otimes_{S'}k(s')$. 여기서 $k(s')$은 $k(s)$의 확대체이고,
$X_s\otimes_{k(s)}k(s')$은 $s'$에서의 $X'$의 섬유
$X'\otimes_{S'}k(s')=X'_{s'}$과 동일시된다. 따라서 유계인 족이란
$S$ 위의 유한형 준스킴 $S'$으로 매개되는 연접층의 한 \emph{대수적 족}에
포함되는 족이라고 말할 수 있다.

유계인 족들의 유한 합집합은 유계이다(그 족들을 포함하는 대수적 족들을
정의하는 매개변수 준스킴 $S_i$들의 서로소 합인 준스킴을 취한다).
밑변환 $T\longrightarrow S$는 $X/S$에 상대적인 유계인 족을
$X_T/T$에 상대적인 유계인 족으로 바꾼다. $T\longrightarrow S$가
전사이면 그 역도 성립한다(더 일반적으로는, 그 상이 $X/S$에 대하여 주어진
족 $E$에 실제로 나타나는 $s$들을 포함하면 된다). 이로써 유계인 족을
결정하는 문제는 이론상 $S$가 정수환 $\mathbf Z$ 위의 유한형 대수의
스펙트럼인 경우로 환원된다.

$E$와 $E'$이 $X/S$에 상대적인 층의 동치류들의 유계인 족이면
$E\otimes E'$도 유계이다. 실제로 $E$와 $E'$을 각각 포함하는 대수적
족들이 $T\longrightarrow S$와 $X_T$ 위의 $\mathcal F$, 그리고
$T'\longrightarrow S$와 $X_{T'}$ 위의 $\mathcal F'$으로 정의된다고
하자. 그러면 $E\otimes E'$은 $T''=T\times_S T'\longrightarrow S$와
$X_{T''}$ 위의 층 $\mathcal F''$으로 정의되는 대수적 족에 포함된다.
여기서 $\mathcal F''$은 $X_{T''}$ 위에서 $\mathcal F$와 $\mathcal F'$의
역상들을 텐서곱한 층이다. 이 논증이 옳은 것은 $\mathcal F$, $\mathcal F'$에
관한 함자 $\mathcal F\otimes\mathcal F'$이 오른쪽 완전하여 밑변환과
가환하기 때문이다(특히 섬유로의 이행과 가환한다). 이 논증을
$\operatorname{Tor}_i(E,E')$, $\operatorname{Hom}(E,E')$,
$\operatorname{Ext}^i(E,E')$ 같은 국소 연산에 그대로 적용할 수는 없다.
그러나 $E\otimes E'$의 경우와 같이 진행하면서 다음 유형의 결과들을
(모두 [3]의 IV, 6.11에 들어 있다) 더 사용하면, 이 연산들도 유계인
집합들을\footnote{인쇄본에는 한정사 ``des''가
빠져 있다.} 다시 유계인 집합들로 바꿈을 보일 수 있다. 즉 유계인 족 $E$는
항상 $X_T$ 위의 연접층 $\mathcal F$로 정의되는 대수적 족에 포함되도록
할 수 있고($T$는 $S$ 위의 유한형이다), 이때 $\mathcal F$는 $T$에
상대적으로 평탄하다. (처음의 매개변수 공간을 ``조각으로
나눈다''.
\footnote{인쇄본에는 잘못된 구절 ``coupe au
morceaux''가 쓰여 있다.}) 알맞은 층들에 관한 이러한
\src{252}
평탄성은 $\operatorname{Tor}_i(\mathcal F,\mathcal F')$ 같은 연산이
임의의 밑변환과 가환하도록 보장한다. 같은 방법은 전역적인 성질의 연산에도
적용된다. 고유 사상에 의한 연접층의 직상과 유도 직상, 고유 사상에 상대적인
전역 $\operatorname{Ext}$([5]의 III, \S~6 참조) 등이 그 예이다. 이 모든
연산은 층의 유계인 족들을 층의 유계인 족들로 바꾼다. (주의 --- 여기서
여러 층을 취하는 준스킴들은 고려하는 유형의 연산을 거치면서 달라질 수 있다.)

다음 두 명제는 본질적으로 같은 평탄성 기법으로 증명한다. 유한형 사상의
섬유들 위에서의 준소분해에 관해서는 특히 [5]의 IV를 보라.

\result{명제}{1.2} $E$, $E'$을 $X/S$의 섬유들 위의 층의 동치류들로
이루어진 유계인 집합이라 하고, $X$는 $S$ 위에서 고유라고 가정하자. 그러면
다음이 성립한다.

\medskip
\noindent (i) $\mathcal F$의 동치류가 $E$에 속하고 $\mathcal F'$의
동치류가 $E'$에 속할 때, 준동형
$\mathcal F\longrightarrow\mathcal F'$들의 핵, 여핵, 상으로 이루어진
족은 유계이다.

\smallskip
\noindent (ii) $\mathcal F$의 동치류가 $E$에 속하고 $\mathcal F'$의
동치류가 $E'$에 속할 때, $\mathcal F'$에 의한 $\mathcal F$의 확장인
층 $\mathcal F''$들로 이루어진 족은 유계이다.

알맞은 밑변환을 한 뒤, $E$와 $E'$이 각각 $X_T/T$ 위의 연접층
$\mathcal G$와 $\mathcal G'$으로 정의된다고 가정할 수 있다. 여기서
$T$는 $S$ 위의 유한형이다. 더 나아가 어떤 평탄성 가정들이 성립한다고
할 수 있다. 이 가정들에 따라 층
$f_{T*}(\operatorname{Hom}_{X_T}(\mathcal G,\mathcal G'))$과
$\operatorname{Ext}^{1}_{f_T}(\mathcal G,\mathcal G')$의 형성은
임의의 사상 $T'\longrightarrow T$에 의한 밑변환과 가환한다. 또한 $T$
위의 앞의 연접층들이 국소 자유라고 가정할 수 있다. $T_0$와 $T_1$을
단면의 싹들의 층이 각각 앞의 두 층인 $T$ 위의 벡터 다발이라 하자. 그러면
$X_{T_0}$ 위의 연접층들의 준동형
$\mathcal G_{T_0}\longrightarrow\mathcal G'_{T_0}$와, $X_{T_1}$ 위의
연접층들의 확장\footnote{인쇄본은 보편 확장의 두
끝 항을 서로 뒤바꾸어 놓았다. 고려하는 류는
$\operatorname{Ext}^1(\mathcal G,\mathcal G')$에 속한다.}
\[
0\longrightarrow\mathcal G'_{T_1}\longrightarrow\mathcal G''
\longrightarrow\mathcal G_{T_1}\longrightarrow0
\]
을 표준적으로 정의할 수 있고, 후자는 자명한 보편 성질을 갖는다. 이 두 번째
층은 (ii)에서 고려한 족을 포함하는 대수적 족을 정의한다. 앞의 준동형의
\src{253}
핵, 여핵, 상도 마찬가지이며, 이들을 생각하면 (i)가 증명된다(다만 여핵이
$T_0$에 상대적으로 평탄하다고 가정한다. 이 경우에도 $T_0$를 조각으로
나누어 환원한다\ldots).

\result{명제}{1.3} $E$를 $X/S$의 섬유들 위의 층의
동치류들로\footnote{인쇄본에는 단수 ``classe''가
쓰여 있다.} 이루어진 유계인 족이라 하자. $K$가 대수적으로 닫혀 있고
$\mathcal F$가 $X_K$ 위의 연접층이며 그 동치류가 $E$에 속할 때,
$(\operatorname{supp}\mathcal F)_{\mathrm{réd}}$의 구조층들의 동치류는
유계인 족을 이룬다.

여기서 $(\operatorname{supp}\mathcal F)_{\mathrm{réd}}$은
$\mathcal F$의 지지집합에 유도된 축소 구조를 준 것이다. 즉 그 구조층은
$\operatorname{supp}\mathcal F$를 정의하는 가장 큰 아이디얼층으로
$\mathcal O_{X_K}$를 나눈 몫이다. 준소분해 이론으로 $\mathcal F$에서
표준적으로 얻는 층들에 대해서도 (1.3)과 같은 결과를 증명할 수 있다. 예를
들어 $\mathcal F_i$들이 $\mathcal F$의 지지집합의 성분들에 대한 준소
부분층이고 이 성질에 관하여 극소일 때의 $\mathcal F/\mathcal F_i$들,
$\mathfrak p$가 $\mathcal F$에 연관된 소 아이디얼층일 때의
$\mathcal O_{X_K}/\mathfrak p$들, 또는 $\mathfrak q$가 $\mathcal F$의
지지집합의 한 성분에 연관된 준소 아이디얼층일 때의
$\mathcal O_{X_K}/\mathfrak q$들이 그러하다(기초체는 대수적으로 닫혀 있다고
한다).

\fsection{2}{유계인 족과 힐베르트 다항식}

이하에서는 $X$가 $S$ 위에서 사영이고 $\mathcal O_X(1)$로 나타내는 매우
풍부한 층이 주어졌다고 가정한다. $S$의 점 $s$의 잉여체 $k(s)$의 모든
확대체 $K$에 대하여, $X_K$ 위에서 이에 대응하는 층
$\mathcal O_{X_K}(1)$을 생각한다. 이 층도 매우 풍부하다.

$X_K$ 위의 모든 연접층 $\mathcal F$에 함수
\[
P_{\mathcal F}(n)=\text{$X_K$ 위의 $\mathcal F(n)$의
오일러--푸앵카레 특성}
\]
을 대응시킨다. 이는 정수 $n$에 관한 다항식이며 $\mathcal F$의
\emph{힐베르트 다항식}이라고 한다. $n$이 충분히 크면 $i>0$에 대하여
$H^i(X_K,\mathcal F(n))$이 영이므로, $P(n)$은 바로
$H^0(X_K,\mathcal F(n))$의 $K$ 위의 차원이다.

이제 $\mathcal F$가 $X$ 위의 연접층이고 $S$에 상대적으로 평탄하다고 하자.
그러면 $S$의 같은 연결 성분에 속하는 섬유 $X_s$들 위에 유도된 층
$\mathcal F_s$들의 힐베르트 다항식은 모두 같다([5]의 III, \S~7).
따라서 평탄성 가정 없이도, $X$ 위의 모든 연접층 $\mathcal F$에 대하여
$s\in S$일 때의 층 $\mathcal F_s$들의 힐베르트 다항식들로 이루어진
집합은 유한하다.
\src{254}
한편 $\mathcal F$가 $X$ 위의 연접층이면, 충분히 큰 $n$, $N$에 대하여
$\mathcal F$는 $\mathcal O_X(-n)^N$의 한 몫층과 동형임을 상기하자.
따라서 섬유들 위에 유도된 층 $\mathcal F_s$들도 그 섬유 위의 층
$\mathcal O(-n)^N$의 몫이다.

이 두 비고로부터 다음 정리의 ``필요조건'' 부분을 얻는다.
\result{정리}{2.1} $X$를 뇌터 준스킴 $S$ 위에서 사영인 준스킴이라 하고,
$\mathcal O_X(1)$이 $S$에 관해 $X$ 위에서 매우 풍부하다고 하자. $E$를
$X/S$의 섬유들 위의 층들의 동형류로 이루어진 집합이라 하자. $E$가
유계일 필요충분조건은\footnote{인쇄본에는 ``필요하다''만 적혀 있다.
Exposé 221에 대한 1962년의 정오표는 ``필요충분하다''로 고치도록 한다.}
다음 조건들을 만족하는 것이다.

\medskip
\noindent (a) $X$ 위의 연접층 $\mathcal L$이 존재하여($\mathcal
O_X(-n)^N$ 꼴이라고 가정해도 된다), $E$가 $\mathcal L_K$ 꼴의 층들의
연접 몫층의 동형류로 이루어진 족에 포함된다.

\smallskip
\noindent (b) 동형류가 $E$에 속하는 층 $\mathcal F$들의 힐베르트 다항식
$P_{\mathcal F}$들이 하나의 같은 유한한 다항식 집합에 속한다.

이제 ``충분하다''는 부분만 증명하면 되며, 이는 더 정밀한 결과의 특수한
경우가 될 것이다. 체 $K$ 위의 유한형 준스킴 위의 임의의 연접 가군
$\mathcal F$와 임의의 정수 $r$에 대하여, 열린집합 위의 단면들이 그
열린집합 위에서 지지집합의 차원이 $<r$인 $\mathcal F$의 단면들인
$\mathcal F$의 부분가군을 $N_r$이라 하자. 그러면
$r>\dim\operatorname{supp}\mathcal F$이면 $N_r=\mathcal F$이고,
$r\leq0$이면 $N_r=0$이다. 따라서 $\mathcal F$의 유한한 증가 여과를
얻으며, 그 인자 $N_{r+1}/N_r$\footnote{인쇄본은 몫을
$N_r/N_{r+1}$로 뒤집어 적었는데, 이는 방금 정의한 증가 여과와 양립하지
않는다.}의 연관 소순환은 정확히 $\mathcal F$의 연관 소순환 가운데
차원이 $r$인 것들이다. 다음과 같이 놓겠다.
\[
\mathcal F_{(r)}=\mathcal F/N_r.
\]
그러면 $\mathcal F_{(r)}$의 연관 소순환은 정확히 $\mathcal F$의 연관
소순환 가운데 차원이 $\geq r$인 것들이다. 특히 $\mathcal F_{(r)}$가
$\mathcal F$와 같을 필요충분조건은 $\mathcal F$의 연관 소순환들의
차원이 $\geq r$인 것이다. 이제 다음이 성립한다.

\result{정리}{2.2} (2.1)의 조건 아래에서 $s$를 정수라 하고, $E$가 조건
(a)와 다음과 같이 약화한 (b)를 만족한다고 가정하자.

\medskip
\noindent $(b_s)$ 동형류가 $E$에 속하는 층 $\mathcal F$들의 푸앵카레
다항식 $P_{\mathcal F}$에서 차수가 $\geq s-1$인 계수들이 유계이다.

이 조건들 아래에서 $\mathcal F$의 동형류가 $E$를 달릴 때 층
$\mathcal F_{(s)}$들은 유계인 족을 이룬다. 더 나아가
$P_{\mathcal F}$의 $s-2$차 계수들은 아래로 유계이다.
\src{255}
따라서 다음을 얻는다.

\result{따름정리}{2.3} 동형류가 $E$에 속하는 층 $\mathcal F$들의 모든
연관 소순환의 차원 $d$가 $s\leq d\leq r$을 만족한다고 가정하자. 그러면
(2.1)의 (b)에서는 $P_{\mathcal F}$의 $s-1$차부터 $r$차까지의 계수들만
고려하면 된다.

이 번호의 나머지 부분에서는 (2.2)의 증명을 개략적으로 설명한다. 핵심은
다음 두 보조정리이며, 첫째 것은 잘 알려져 있다(그리고 Chow 좌표의 유용한
수학적 내용을 요약한다).

\result{보조정리 2.4 (CHOW)}{} $S$의 잉여체들의 대수적으로 닫힌 확대체
$K$에 대한 섬유 $X_K$의 부분스킴 $Y$들의 구조층을 생각하자. 여기서
$Y$는 축소되어 있고 그 모든 성분의 차원이 같은 $r$이며
($\mathcal O_Y$는 $\mathcal O_X$의 몫층으로 본다), $Y$들의 차수가
유계이면 $Y$들은 유계인 족을 이룬다.

여기서 $Y$의 차수 $a$는 다음 식의 최고차항의 계수로 정의하는 것이 가장
편리하다.
\[
P_{\mathcal O_Y}=a\frac{n^r}{r!}+\cdots.
\]

\result{보조정리}{2.5} $\mathcal L$을 $X$ 위의 연접층, $E$를
$\mathcal L_K$ 꼴의 층($K$는 $S$의 잉여체의 확대체)의 몫층
$\mathcal F$들의 동형류로 이루어진 집합이라 하자. $S$ 위의 $X$의
섬유들의 차원이 $\leq r$이라고 가정하고 다음과 같이 놓자.
\[
P_{\mathcal F}(n)=a_{\mathcal F}\frac{n^r}{r!}
+b_{\mathcal F}\frac{n^{r-1}}{(r-1)!}
+\text{차수가 $<r-1$인 항들}.
\]
그러면 계수 $a_{\mathcal F}$는 유계이고 $b_{\mathcal F}$는 아래로
유계이다. $b_{\mathcal F}$가 유계이면 $\mathcal F_{(r)}$들의 족은
유계이다.

\emph{증명.} -- $S$를 덮는 $S$의 부분스킴들의 합집합으로 $S$를
바꾸어도 무방하므로, 유한 사상
$f:X\longrightarrow\mathbf P^r_S$가 존재하고 $\mathcal O_X(1)$이
$\mathcal O_{\mathbf P^r_S}(1)$의 역상과 동형이라고 가정할 수 있다.
따라서 $X$ 위의 모든 연접층 $\mathcal F$에 대하여
$P_{\mathcal F}=P_{f_*(\mathcal F)}$이다. 한편 (앞 번호의 방법으로)
$X$ 위의 층 $\mathcal F$들로 이루어진 집합이 유계일 필요충분조건은
$f_*(\mathcal F)$들로 이루어진 집합이 유계인 것임을 쉽게 확인할 수
있다. 끝으로
\[
f_*(\mathcal F)_{(r)}=f_*(\mathcal F_{(r)})
\]
이다.
\src{256}
따라서 $X=\mathbf P^r_S$인 경우로 환원된다. 더 나아가 적당한 $k,s$에
대하여 $\mathcal L=\mathcal O_{\mathbf P^r_S}(k)^s$라고 가정할 수 있다.
계수 $a_{\mathcal F}$는
\[
0\leq a_{\mathcal F}\leq s
\]
를 만족하므로 유계이다. 이제 $P_{\mathcal F}(n)$의 $n^{r-1}$의
계수가 아래로 유계(각각 유계)라고 말하는 것과
$P_{\mathcal F}(n-k)=P_{\mathcal F(-k)}(n)$에 대하여 그렇게 말하는 것은
같은 뜻이다. 따라서 다음 경우로 환원된다.
\[
\mathcal L=\mathcal O_{\mathbf P^r_S}^{s}.
\]

완전열
\[
0\longrightarrow N_r\longrightarrow\mathcal F
\longrightarrow\mathcal F_{(r)}\longrightarrow0
\]
을 생각하자. 이로부터
\[
P_{\mathcal F}=P_{\mathcal F_{(r)}}+P_{N_r}
\]
이고, $P_{N_r}$에서 $n^{r-1}$의 계수는 양수이므로
($\dim\operatorname{supp}N_r\leq r-1$이기 때문이다),
\[
b_{\mathcal F_{(r)}}\leq b_{\mathcal F}
\]
이다. 따라서 보조정리의 주장들을 증명할 때 $\mathcal F$를
$\mathcal F_{(r)}$로 바꾸어도 된다. 즉, 고려하는 $\mathcal L$의 몫층
$\mathcal F$들이 꼬임이 없다고 가정할 수 있다.

$\mathbf P^r_K$가 정규이므로 $\mathcal F$는
$U=\mathbf P^r_K-Y$ 위에서 계수(階數) $a=a_{\mathcal F}$인 국소
자유층이며, 여기서 $Y$의 여차원은 $\geq2$이다. 그러므로
$\bigwedge^a\mathcal F$는 $\mathbf P^r_K$ 위의 층으로서 $U$로의 제한이
가역이다. 따라서 ($\mathbf P^r_K$가 정칙이고 $Y$의 여차원이
$\geq2$이므로) 이는 $\mathbf P^r_K$ 위의 어떤 가역층의 제한과
동형이고, 그 가역층은 동형을 제외하고 결정된다. 이 가역층은 잘 결정된
어떤 정수 $d$에 대하여 $\mathcal O_{\mathbf P^r_K}(d)$ 꼴이다.
$\bigwedge^a\mathcal F$는
$\bigwedge^a\mathcal O_{\mathbf P^r_K}^{s}\simeq
\mathcal O_{\mathbf P^r_K}^{N}$\footnote{인쇄본은 여기서 지수 $s$를
정의되지 않은 $n$으로 바꾸어 적었다. $s$는 $\mathcal F$로 전사하는
자유층의 계수(階數)이다.}의 몫이고 $N=\binom{s}{a}$이므로, 이 층은
$N$개의 표준 단면을 갖는다. 따라서 이 단면들은 $U$ 위의
\src{257}
$\mathcal O_{\mathbf P^r_K}(d)$의 단면들을 정의하고, 이들은
$\mathcal O_{\mathbf P^r_K}(d)$의 단면 $s_i$ ($1\leq i\leq N$)들의
제한이다($\mathbf P^r_K$가 정규이고 $Y$의 여차원이 $\geq2$이기
때문이다). 이 단면들은 $U$의 점들에서
$\mathcal O_{\mathbf P^r_K}(d)$를 생성하므로 모두 영일 수 없고, 따라서
$d\geq0$이다. 한편 간단한 계산으로
\[
b_{\mathcal F}=a_{\mathcal F}(r+1)/2+d
\]
를 얻는다. 특히 이로부터 $b_{\mathcal F}>0$임을 알 수 있으므로
$b_{\mathcal F}$는 아래로 유계이다. 이것이 유계일 필요충분조건은 $d$가
유계인 것이며, 이때 $\mathcal F$가 유계인 족에 속함을 보이자.
$a_{\mathcal F}$와 $b_{\mathcal F}$가 각각 고정된 $a,b$라고 가정할 수
있고, 따라서 $d$도 고정되어 있다.
$\mathcal O_{\mathbf P^r_K}(d)$의 $N$개 단면 $s_i$, 즉 준동형
\[
s:\bigwedge^a\mathcal L_K\longrightarrow\mathcal O_{\mathbf P^r_K}(d)
\]
를 알면, 이에 대응하는 합성 준동형
\[
\mathcal L_K\longrightarrow
\operatorname{Hom}\!\left(\bigwedge^{a-1}\mathcal L_K,
\bigwedge^a\mathcal L_K\right)\longrightarrow
\operatorname{Hom}\!\left(\bigwedge^{a-1}\mathcal L_K,
\mathcal O_{\mathbf P^r_K}(d)\right)
\]
의 여상으로 $\mathcal F$를 복원할 수 있다. 여기서 첫째 준동형은
외승곱에서 나오는 표준 준동형이고, 둘째 준동형은 $s$로부터 얻어진다.
이제 (1.2) (i)에 따라 결론이 나온다.

앞의 두 보조정리를 결합하면 다음을 증명할 수 있다.

\result{보조정리}{2.6} (2.1)의 예비 조건 아래에서 모든 $\mathcal F$에
대하여
\[
P_{\mathcal F}(n)=a_{\mathcal F}\frac{n^r}{r!}
+b_{\mathcal F}\frac{n^{r-1}}{(r-1)!}
+\text{차수가 $<r-1$인 항들}
\]
이고 계수 $a_{\mathcal F}$들이 유계라고 가정하자. 그러면 계수
$b_{\mathcal F}$들은 아래로 유계이다. $b_{\mathcal F}$들이 유계이면
$\mathcal F_{(r)}$들은 유계인 족을 이룬다.

층 $\mathcal F$들의 밑체 $K$는 대수적으로 닫혀 있다고 가정할 수 있다.
각 $\operatorname{supp}\mathcal F_{(r)}$, 즉
$\operatorname{supp}\mathcal F$의 $r$차원 성분들의 합집합에 유도된
축소 구조를 주자. 그러면 $\operatorname{supp}\mathcal F_{(r)}$들의
차수는 $a$를 위에서 유계로 가지므로 (2.4)에 따라
$\operatorname{supp}\mathcal F_{(r)}$들은 유계인 집합을 이룬다. 더
나아가 $\operatorname{supp}\mathcal F_{(r)}$의 각 성분에 대한
$\mathcal F_{(r)}$의 길이는 $\leq a$이다. 따라서
$\mathcal J_{\mathcal F}$가 $\operatorname{supp}\mathcal F_{(r)}$를
정의하는 아이디얼이면 $\mathcal F_{(r)}$는
\src{258}
$\mathcal J_{\mathcal F}^{a}$로 정의되는 $X$의 부분스킴
$Y_{\mathcal F}$ 위의 가군으로 볼 수 있다. 앞 보조정리에서와 같이
$\mathcal F=\mathcal F_{(r)}$인 경우로도 환원할 수 있으므로,
$\mathcal F$는 $Y_{\mathcal F}$ 위의 가군으로부터 온다.
$Y_{\mathcal F}$들은 $\mathcal O_{X_K}$의 몫가군들의 유계인 족에
대응하므로, 어떤 스킴 $X\times_S T$의 닫힌 부분스킴 $Y$로부터 온다.
이제 $Y/T$와 $\mathcal L\otimes_XY$에 (2.5)를 적용하면 결론을 얻는다.

이제 $\dim\operatorname{supp}\mathcal F$의 상계 $r$에 대한 귀납법으로
(2.2)를 증명할 수 있다. $r<0$이면 주장은 자명하다. 따라서 $r\geq0$이고
$r'<r$에 대해서는 주장이 증명되었다고 가정하자. (2.6)에 따라
$\mathcal F_{(r)}$들은 유계인 집합을 이루므로, (1.2) (i)에 따라
준동형 $\mathcal L_K\longrightarrow\mathcal F_{(r)}$들의 핵들도
유계인 집합을 이룬다. 따라서 문제의 핵들, 그러므로 또한
\[
N_r(\mathcal F)=\operatorname{Ker}(\mathcal F\longrightarrow\mathcal F_{(r)})
\]
이 $\mathcal L'_K$들의 몫이 되게 하는 $X$ 위의 연접 가군
$\mathcal L'$이 존재한다. $\mathcal F_{(r)}$들이 유계이므로
$P_{\mathcal F_{(r)}}$들도 유계이고, 공식
\[
P_{\mathcal F}=P_{\mathcal F_{(r)}}+P_{N_r}
\]
에 따라 $P_{N_r}$들은 $P_{\mathcal F}$들과 같은 조건 $(b_s)$를
만족한다. 그러므로 $N_r$들은 조건 (a)와 $(b_s)$를 만족하고, 귀납
가정에 따라 $(N_r)_{(s)}$들은 유계이다. 그런데 $\mathcal F_{(s)}$는
$(N_r)_{(s)}$에 의한 $\mathcal F_{(r)}$의 확대이므로 (1.2) (ii)에 따라
$\mathcal F_{(s)}$들은 유계이다. (2.2)의 마지막 주장에 대해서는
유계성과 그 다음 계수의 아래로 유계임을 동시에 귀납한다. 고정된
원천에서 $\mathcal F_{(s)}$로 가는 사상들의 유계인 핵들은 각 $N_s$로
전사하는 하나의 같은 연접층을 제공하지만, (1.2) (i)만으로는 이
몫들의 족이 유계임을 얻기에 충분하지 않다. 이 몫들에 계수에 대한
주장을 귀납적으로 다시 적용하면, 힐베르트 다항식의 가법성과 보조정리
2.6에 따라 원하는 다음 계수의 하계를 얻는다.
이것으로 증명이 끝난다.

\fsection{3}{힐베르트 스킴: 정의와 존재 정리}

$X$를 다른 준스킴 $S$ 위의 준스킴이라 하고, $\mathcal F$를 $X$ 위의
준연접 가군이라 하자. 다음 기호로
\[
\operatorname{Quot}(\mathcal F/X/S)
\]
$\mathcal F$의 준연접 몫가군 가운데 $S$ 위에서 평탄인 것들의 집합을
나타내자. 이제 $S'\longrightarrow S$를 밑준스킴의 변경 사상이라 하고,
$X'=X\times_S S'$ 및
$\mathcal F'=\mathcal F\otimes_{\mathcal O_S}\mathcal O_{S'}$로 놓자.
그러면 $X'$은 준연접 가군
\src{259}
$\mathcal F'$이 주어진 $S'$ 위의 준스킴이고,
$\operatorname{Quot}(\mathcal F'/X'/S')$를 생각할 수 있다. 다음과 같이
놓는다.
\[
\underline{\operatorname{Quot}}_{\mathcal F/X/S}(S')
=\operatorname{Quot}(\mathcal F'/X'/S')
\qquad\text{(여기서 $X'=X\times_SS'$)}.
\]
이제 $S''\longrightarrow S'$가 $S$-사상이면
$X''=X\times_SS''$은 $X'\times_{S'}S''$과 동형이고, $\mathcal F''$은
$\mathcal F'\otimes_{\mathcal O_{S'}}\mathcal O_{S''}$과 동형이다. 또한
$X'$ 위의 준연접 가군의 범주에서 $X''$ 위의 준연접 가군의 범주로 가는
역상 함자
\[
\mathcal G'\longmapsto
\mathcal G'\otimes_{\mathcal O_{S'}}\mathcal O_{S''}
\]
는 오른쪽 완전이고 $S'$-평탄 가군들을\footnote{인쇄된 프랑스어 본문은 여기서 직접목적어 앞의 관사 ``les''를
빠뜨렸다.} $S''$-평탄 가군들로 보내므로, 자연스러운 사상
\[
\operatorname{Quot}(\mathcal F'/X'/S')\longrightarrow
\operatorname{Quot}(\mathcal F''/X''/S'')
\]
을 얻는다. 따라서 $\underline{\operatorname{Quot}}_{\mathcal F/X/S}(S')$는
$S$ 위의 준스킴 $S'$에 대한 반변 함자로서 집합의 범주에 값을 갖는다.
이하에서는 $X$가 뇌터 준스킴 $S$ 위에서 사영이고 $\mathcal F$가
연접이라고 가정하며, 단순화를 위해 $S$ 위의 $S'$ 가운데 국소 뇌터인
것들로 한정하겠다.

\result{정리}{3.1} 이 조건들 아래에서 국소 뇌터 $S$-준스킴들의 범주
위의 반변 함자 $\underline{\operatorname{Quot}}_{\mathcal F/X/S}$는
$S$-준스킴 $\operatorname{Quot}_{\mathcal F/X/S}$에 의해 표현 가능하며,
이 준스킴은 사영 $S$-스킴들의 한 열의 합이다(따라서
$\operatorname{Quot}_{\mathcal F/X/S}$는 $S$ 위에서 국소 유한형이다).

이러한 분해는 다음과 같이 얻는다. $\mathcal O_X(1)$을 $X$ 위의
가역층으로서 $S$에 관해 매우 풍부한 것이라 하자. 유리수 계수를 갖는
각 다항식 $P(n)$에 대하여 $\operatorname{Quot}^P(\mathcal F/X/S)$를
$\operatorname{Quot}(\mathcal F/X/S)$ 가운데 $\mathcal F$의 연접
몫가군 $\mathcal G$로서 $S$ 위에서 평탄하고 모든 $s\in S$에서의
힐베르트 다항식이 $P$와 같은 것들로 이루어진 부분이라 하자. 그러면
다음과 같이 놓는다.
\[
\underline{\operatorname{Quot}}^P_{\mathcal F/X/S}(S')
=\operatorname{Quot}^P(\mathcal F'/X'/S')
\]
이로써 $\underline{\operatorname{Quot}}_{\mathcal F/X/S}$의 부분함자를
얻는다. 2절에서 상기한 힐베르트 다항식의 불변성으로부터 다음이
따른다. $\underline{\operatorname{Quot}}_{\mathcal F/X/S}$가 표현
가능할 필요충분조건은 모든
$\underline{\operatorname{Quot}}^P_{\mathcal F/X/S}$가 표현 가능한
것이다. 이때 전자를 표현하는 $S$-준스킴
$\operatorname{Quot}_{\mathcal F/X/S}$는 후자의 함자들을 표현하는
$\operatorname{Quot}^P_{\mathcal F/X/S}$들의 합인 준스킴과 동형이다.
따라서 (3.1)은 다음 정리의 귀결이다.
\src{260}
\result{정리}{3.2} 앞의 표기를 사용하면, 함자
$\underline{\operatorname{Quot}}^P_{\mathcal F/X/S}$는 사영
$S$-준스킴 $\operatorname{Quot}^P_{\mathcal F/X/S}$에 의해 표현
가능하다.

이 절의 나머지 부분에서는 (3.2)를 증명한다.

$\nu$를 정수라 하자. $S$ 위의 각 $S'$에 대하여 $A_\nu(S')$로
$\mathcal F'=\mathcal F\otimes_{\mathcal O_S}\mathcal O_{S'}$의
몫 $\mathcal G=\mathcal F'/\mathcal K$ 가운데 연접이고 $S'$ 위에서
평탄하며 다음 조건들을 만족하는 것들의 집합을 나타내자.
\begin{align*}
\text{a.}\quad &R^if'_*(\mathcal G(n))=0 &&\text{$i>0$, $n\geq\nu$일 때},\\
\text{b.}\quad &R^if'_*(\mathcal K(n))=0 &&\text{$i>0$, $n\geq\nu$일 때},\\
\text{c.}\quad &f'_*(\mathcal K(\nu+k))
=\mathcal S'_k f'_*(\mathcal K(\nu)) &&\text{$k\geq0$일 때}.
\end{align*}
마지막 관계에서는 $X$가 $S$ 위의 준연접 등급 대수 $\mathcal S_*$의
동차 소 스펙트럼으로 주어져 있고, 이 대수는 양의 차수로 등급이 주어져
$\mathcal S_1$에 의해 생성된다고 가정한다. 또한
\[
\mathcal S'_* =\mathcal S_*\otimes_{\mathcal O_S}\mathcal O_{S'}
\]
로 놓아서 $X'$이 $\mathcal S'_*$의 동차 소 스펙트럼이 되게 한다.
그런데 (3.2)를 증명할 때에는 $X=\mathbf P^r_S$인 경우로 쉽게
귀착된다. 실제로 $S$는 $X|U$가 $\mathbf P^r_U$의 닫힌 부분스킴이고
$\mathcal O_X(1)$이 $\mathcal O_{\mathbf P^r_U}(1)$에서 유도되는
열린집합 $U$들의 합이다. 또한 $\mathcal F$가
$\mathcal O_{\mathbf P^r_S}^{N}$ 꼴인 경우로 귀착되므로 $\mathcal F$는
$S$ 위에서 평탄하다. 그러면 앞에서의 층 $\mathcal K$들도 $S'$ 위에서
평탄하다. 이제 퀸네트 관계식([5], III, \S~7)에 따라 조건 (a), (b)는
밑준스킴의 변경에 대해 안정하며, $n\geq\nu$일 때
$f'_*(\mathcal G(n))$과 $f'_*(\mathcal K(n))$의 형성이 밑준스킴의
확장과 가환함을 함의한다(같은 곳). 따라서 (c)도 밑준스킴의 확장에
대해 안정하다.

다시 말해 $A_\nu(S')$는 $S'$에 관한 반변 함자이며, 더 정확히는
$A(S')=\underline{\operatorname{Quot}}^P_{\mathcal F/X/S}(S')$의
부분함자이다. $\nu$가 변할 때 이렇게 $A(S')$의 부분집합들로 이루어진
증가열 $A_\nu(S')$를 얻으며, 세르의 잘 알려진 정리 [5], III, \S~2에
따르면 그 합집합은 $A(S')$이다. 이제 $\mathcal G$가 $\mathcal F'$의
연접 몫으로서 $S'$ 위에서 평탄하고, $s$가 $S'$의 점이며, 밑준스킴의
변경 $\operatorname{Spec}(k(s))\longrightarrow S'$으로부터 조건 (a),
(b), (c)를 만족하는 $\mathcal F_s$의 몫 $\mathcal G_s$, 즉
$A_\nu(\operatorname{Spec}(k(s)))$의 원소가 얻어진다고 하자. 그러면
$s$의 열린 근방 $U$가 존재하여 $\mathcal G|f'^{-1}(U)$에 대해서도
같은 조건들이 성립한다. 즉 이 몫은 $A_\nu(U)$에 속한다. (a), (b)에
대해서는 이는 실제로 「정칙 함수」 정리 [5], III, \S~7에서 따르고,
(c)는 나카야마 보조정리와, 어쨌든 $n$이 충분히 크고 $k\geq0$일 때
\[
f'_*(\mathcal K(n+k))=\mathcal S'_k f'_*(\mathcal K(n))
\]
\src{261}
임을 알고 있다는 사실([5], III, \S~2 참조)에서 따른다. 이 고찰들로부터
다음과 같은 결론을 얻는다([4], IV와 비교하라). 함자 $A$가 표현 가능할
필요충분조건은 모든 함자 $A_\nu$가 표현 가능한 것이다. 이때 $A$를
표현하는 $S$-준스킴 $Q$는 $A_\nu$들을 표현하는 열린 부분스킴
$Q_\nu$들의 증가하는 합집합이다.

다음과 같이 놓자.
\[
M_* =\bigoplus_{n\geq0}f_*(\mathcal F(n))=\mathcal S_*^{N},
\]
그러면
\[
M'_* =M_*\otimes_{\mathcal O_S}\mathcal O_{S'}
=\bigoplus_{n\geq0}f'_*(\mathcal F'(n))=\mathcal S'_*{}^{N}.
\]
(a)에서 다음이 따른다.

\noindent (a') $n\geq\nu$일 때 $f'_*(\mathcal G(n))$은 계수 $P(n)$인
국소 자유 가군이다([5], III, \S~7). 또한 $i=1$에 대해 (b)를 적용하면
다음이 따른다.

\noindent (a'') $f'_*(\mathcal G(n))$은 $M'_n$의 몫가군이다.

더구나 $n=\nu$일 때 이 몫가군을 알면, (c)에 따라 $n\geq\nu$인
$M'_n$의 부분가군 $f'_*(\mathcal K(n))$들을 알 수 있고, 따라서
$\mathcal K$와 이어서 $\mathcal G$를 알 수 있다. 이로써 단사 사상
\[
A_\nu(S')\longrightarrow\operatorname{Grass}_{P(\nu)}(M'_\nu)
\]
을 얻는데, 이는 $A_\nu(S')$에서 $M'_\nu$의 계수 $P(\nu)$인 국소 자유
몫가군들의 집합으로 가는 사상이다. 따라서 함자 준동형
\[
i_\nu:A_\nu(S')\longrightarrow
\underline{\operatorname{Grass}}_{P(\nu)}(M_\nu)(S')
\]
을 얻는다. 오른쪽의 함자는 $S$ 위에서 사영적인 그라스만 스킴
$\operatorname{Grass}_{P(\nu)}(M_\nu)$에 의해 표현 가능하다
([4], V와 비교하라). 이제 다음을 주장한다.

\result{보조정리}{3.3} $A_\nu(S')$는 표현 가능한 함자이고, 준동형
$i_\nu$를 표현하는 사상
$Q_\nu\longrightarrow\operatorname{Grass}_{P(\nu)}(M_\nu)$는
몰입 사상이다(따라서 $Q_\nu$는 $S$ 위에서 준사영적이다).

이 주장은 다음 진술과 동치이다([4], IV와 비교하라). $M'_\nu$의
몫가군 $N$으로서 계수 $P(\nu)$인 국소 자유 가군이 주어졌다고 하자.
그러면 $S'$의 부분준스킴 $Z$가 존재하여, $S'$ 위의 모든 국소 뇌터
준스킴 $T'$에 대하여
\src{262}
$T'$ 위에서 $N$의 역상이 $\operatorname{Im}A_\nu(T')$에 속할
필요충분조건은 $T'\longrightarrow S'$가 부분준스킴 $Z$의 포함 사상을
통하여 인수분해되는 것이다. 표기를 바꾸어 $S'=S$라고 가정할 수 있다.
즉 부분가군 $R_\nu$에 의한 $M_\nu$의 몫 $N_\nu$가 주어졌다고 하자.
이것이 $A_\nu(S)$의 한 원소에서 유도될 필요충분조건은 다음 두 조건을
만족하는 것이다.

\medskip
\noindent (i) 모든 $k\geq0$에 대하여
$M_{\nu+k}/\mathcal S_kR_\nu$는 계수 $P(\nu+k)$인 국소 자유 가군이다.

\smallskip
\noindent (ii) $M_*$의 등급 부분가군
\[
R_* =\sum_{k\geq0}\mathcal S_kR_\nu
\]
([5], II, \S~3)에 의해 정의되는 $\mathcal F$의 부분층 $\mathcal K$와
몫 $\mathcal G=\mathcal F/\mathcal K$는 위의 조건 (a), (b)를 만족한다.

이 조건들이 필요함은 명백하다. 역으로 이 조건들이 성립한다고 하자.
그러면 (ii)에서 정의한 층 $\mathcal G$는
$M_{\nu+k}/\mathcal S_kR_\nu$들의 합인 등급 $\mathcal S_*$-가군
$N_*$에 대응하는 층과 동형이므로 $S$ 위에서 평탄하다. 실제로 그
섬유들은 $\mathcal S_*$의 동차 소아이디얼들에 관한 $N_*$의 국소화들의
직접인자이다\footnote{인쇄본에는 $N$이라고 되어
있다. 동차 국소화가 적용되는 것은 등급 $\mathcal S_*$-가군 $N_*$이다.}.
또한 (i)에 따라 $\mathcal G$의 힐베르트 다항식은 $P$이다. (ii)를
고려하면 (a'), (a'')도 성립한다. 따라서 $n\geq\nu$일 때 자연스러운
준동형
\[
N_n\longrightarrow f_*(\mathcal G(n))
\]
은 같은 계수의 국소 자유 가군들 사이의 전사 준동형이고, 그러므로
동형이다. 따라서
\[
f_*(\mathcal K(\nu+k))=\mathcal S_kR_\nu
\quad\text{모든 $k\geq0$에 대하여},
\]
이며, 이는 $\mathcal G\in A_\nu(S)$이고 $R_\nu$가 $\mathcal G$에 의해
정의되는 $\operatorname{Grass}_{P(\nu)}(M_\nu)$의 원소임을 증명한다.

이 판정법 (i), (ii)는 밑준스킴의 변경 $S'\longrightarrow S$ 뒤에 얻는
상황에도 적용된다. 먼저 밑준스킴의 변경 $S'\longrightarrow S$ 뒤에
조건 (i)가 성립한다는 사실을, $S'\longrightarrow S$가 $S$의 어떤
부분준스킴 $Z$의 포함 사상을 통하여 인수분해된다는 조건으로 나타낼 수
있음을 증명하겠다. 이 결과를 얻으면 $S$를 $Z$로 바꾸어 조건 (i)가 이미
$S$ 위에서 성립하는 경우로 귀착된다. 이 조건은 밑준스킴의 변경에 대해
안정하므로 조건 (ii)를 나타내는 일만 남는다. 이제 $U$를 다음과 같은
$s\in S$들의 집합이라 하자. $n\geq\nu$일 때 $\mathcal G(n)$과
$\mathcal K(n)$이 섬유 $X_s$ 위에 유도하는 층들의 양의 차수 코호몰로지가
소멸한다. $U$가 열려 있음은 이미 지적하였다. 또한 밑준스킴의 변경
$S'\longrightarrow S$ 뒤에 조건 (ii)가 성립할 필요충분조건은 이 사상이
$U$의 포함 사상을 통하여 인수분해되는 것이다. 이로써 보조정리 3.3을
증명하였다. 이제 다음 보조정리만 증명하면 된다.

\result{보조정리}{3.4} $S$를 국소 뇌터 준스킴이라 하고, $S$ 위에 양의
차수로 등급이 주어지고 $\mathcal S_1$에 의해 생성되는 준연접 등급 대수
$\mathcal S_*$가 주어졌다고 하자. $M_*$를 유한형인 준연접 등급
$\mathcal S_*$-가군, $P$를 유리수 계수의 다항식, $\nu$를 정수라 하자.
그러면 $S$의 부분준스킴 $Z$로서 다음 성질을 갖는 것이 존재하며 분명
유일하다.
\src{263}
$S$ 위의 모든 준스킴 $S'$에 대하여, 모든 $n\geq\nu$에 대해
$M_n\otimes_{\mathcal O_S}\mathcal O_{S'}$가 계수 $P(n)$인 국소 자유
가군일 필요충분조건은 $S'\longrightarrow S$가 $Z$의 포함 사상을 통하여
인수분해되는 것이다.

$S$가 아핀이고 따라서 뇌터라고 분명 가정할 수 있다. 이제 다음이
성립한다.

\result{보조정리}{3.5} 각 정수 $N\geq\nu$에 대하여 $U_N$을 다음을
만족하는 $s\in S$들로 이루어진 $S$의 열린집합이라 하자.
\[
\operatorname{rang}_{k(s)}
M_{ns}\otimes_{\mathcal O_{S,s}}k(s)\leq P(n)
\quad\text{모든 $\nu\leq n\leq N$에 대하여}.
\]
그러면 열린집합들의 감소열 $(U_N)$은 어느 항부터 일정하다.

[3], IV에 따르면 $S$에는 축소 부분스킴 $S_i$들로 이루어진 유한 분할이
존재하여, 각 $M\otimes_{\mathcal O_S}\mathcal O_{S_i}$는 $S_i$ 위에서
평탄하다. 따라서 $M$이 평탄하고, 그러므로 $M_n$들이 평탄하다고 가정할
수 있다. 끝으로 $S$가 연결이라고 분명 가정할 수 있다. 그러면 정수
$n_0$와 다항식 $Q$가 존재하여
\[
\operatorname{rang}_{k(s)}
M_{ns}\otimes_{\mathcal O_{S,s}}k(s)=Q(n)
\quad\text{$n\geq n_0$일 때},
\]
이다([5], III, \S~7). 먼저 $P<Q$라고 가정하자\footnote{인쇄본에는
$P\neq Q$라고 되어 있다. 1962년 정오표는 이를 $P<Q$로 고치고,
``$P=Q$''인 경우를 그 반대인 $n$이 클 때 $P(n)\geq Q(n)$인 경우로
바꾸며, (*)의 엄격한 부등식을 넓은 부등식으로 바꾸도록 한다.}. 그러면
$n$이 클 때 $P(n)<Q(n)$이므로, $N$이 크면 분명
$U_N=\varnothing$이다. 따라서 $(U_N)$은 어느 항부터 일정하다. 반대
경우에는 $n$이 클 때 $P(n)\geq Q(n)$이고, 따라서 $N\geq n_0$일 때
$U_N=U_{n_0}$이다. 그러므로 이 경우에도 $(U_N)$은 어느 항부터 일정하다.

특히 다음을 만족하는 $s\in S$들의 집합 $U_\infty$를 생각하자.
\[
\tag{*}
\operatorname{rang}_{k(s)}
M_{ns}\otimes_{\mathcal O_{S,s}}k(s)\leq P(n)
\quad\text{모든 $n\geq\nu$에 대하여}.
\]
이 집합은 $U_N$들의 교집합이므로 열린집합이다. 따라서 (3.4)를 증명할
때 $S$를 열린집합 $U_\infty$로 바꿀 수 있다\footnote{인쇄본에는 $U$라고 되어 있다. 방금 정의한 열린집합은 $U_\infty$이다.}.
이로써 모든 $s\in S$에서 부등식 (*)이 성립하는 경우로 귀착된다. 또한
다음이 성립한다.

\result{보조정리}{3.6} $M$을 국소 뇌터 준스킴 $S$ 위의 가군, $r$을
정수라 하자. 그러면 $S$의 부분준스킴 $Z$로서 다음 성질을 갖는 것이
존재하며 분명 유일하다. $S$ 위의 모든 $S'$에 대하여
$M\otimes_{\mathcal O_S}\mathcal O_{S'}$가 계수 $r$인 국소 자유 가군일
필요충분조건은 $S'\longrightarrow S$가 $Z$의 포함 사상을 통하여
인수분해되는 것이다. 만일
\[
\operatorname{rang}_{k(s)}
M_s\otimes_{\mathcal O_{S,s}}k(s)\leq r
\quad\text{모든 $s$에 대하여},
\]
이면 $Z$는 $S$의 닫힌 부분준스킴이다(여기서 $M$은 연접이라고
가정한다).
\src{264}
실제로 앞의 논증에서, 필요하면 $S$를 이 부등식이 성립하는 $s$들로
이루어진 열린 부분으로 바꿈으로써, 모든 $s\in S$에 대해 부등식 (*)이
성립하는 경우로 귀착된다. 이제 $M$이 완전열
\[
\mathcal O_S^q\longrightarrow\mathcal O_S^r\longrightarrow M\longrightarrow0
\]
에 들어간다고 가정할 수 있다. $S$ 위의 $S'$들에 관해 생각하는 조건은,
이에 대응하는 완전열
$\mathcal O_{S'}^q\longrightarrow\mathcal O_{S'}^r
\longrightarrow M'\longrightarrow0$에서 두 번째 화살표가 동형이라는 것,
즉 첫 번째 화살표가 영이라는 것이기도 하다. 따라서 준동형
$\mathcal O_S^q\longrightarrow\mathcal O_S^r$을 정의하는 행렬의 성분들이
생성하는 아이디얼로 정의되는 $S$의 닫힌 부분준스킴 $Z$가 요구되는
조건을 만족함을 알 수 있다.

이제 (3.4)의 증명을 멈춘 곳으로 돌아가자. $Z_n$을 (3.6)에 따라 가군
$M_n$과 정수 $r=P(n)$에 대응하는 $S$의 닫힌 부분스킴이라 하고,
$Z'_N$을 $\nu\leq n\leq N$인 $Z_n$들의 하한이라 하자. 그러면
$Z'_N$들은 $S$의 닫힌 부분스킴들로 이루어진 감소열을 이룬다\footnote{인쇄본에는 $Z$라고 되어 있으나, $Z$는 다음 문장에서야 정의된다.
$Z'_N$들은 $S$의 닫힌 부분스킴으로 구성된다.}. 따라서 이 열은 반드시
어느 항부터 일정하다. $N$이 클 때 $Z'_N$이 갖는 일정한 값을 $Z$라 하자.
이것이 (3.4)에서 요구한 $Z$이다. 이로써 (3.4)의 증명이 끝나고, 따라서
(3.3)의 증명도 끝난다.

따라서 $\underline{\operatorname{Quot}}^P_{\mathcal F/X/S}$가
$S$-준스킴 $Q$에 의해 표현 가능하고, $Q$는 $S$ 위에서 준사영적인 열린
부분스킴 $Q_\nu$들의 증가하는 합집합임을 증명하였다. 더 나아가려면
정리 2.1을 적용해야 한다. 이로부터 $Q$가 준콤팩트임이 쉽게 따라온다.
실제로 $Q$는 힐베르트 다항식 $P$를 갖는 $\mathcal F_K$의 몫층들의 족을
매개하는, $S$ 위에서 유한형인 준스킴 $S'$의 상이다. 따라서 $Q$는
$Q_\nu$들 가운데 하나와 같고, 그러므로 $S$ 위에서 준사영적이다. $Q$가
$S$ 위에서 사영적임을 증명하려면 이제 $S$ 위에서 고유임을 증명하면
된다. 이를 위해서는 [5], II, 7.3.8의 형태로 고유성의 값매김 판정법을
적용하면 충분하다. 즉 다음을 확인하면 충분하다.

\result{보조정리}{3.7} $S$를 이산 값매김환의 스펙트럼, $s$를 그 일반점,
$X$를 $S$ 위의 준스킴, $\mathcal F$를 $X$ 위의 준연접 가군층이라 하고,
$\mathcal G_s$를 $X_s$ 위에서
$\mathcal F_s=\mathcal F\otimes_{\mathcal O_S}k(s)$의 준연접 몫가군층이라 하자.
그러면 $\mathcal F$의 준연접 몫가군층 $\mathcal G$로서 $S$ 위에서 평탄하고
$X_s$로 제한한 것이 $\mathcal G_s$인 것은 유일하게 존재한다.
\src{265}
실제로 $\mathcal G_s=\mathcal F_s/\mathcal K_s$라면, $\mathcal K_s$를
유도하는 $\mathcal F$의 가장 큰 부분층 $\mathcal K$를 생각하고 [5], I,
9.4.2,
\[
\mathcal G=\mathcal F/\mathcal K
\]
로 두면 충분하다. 이 층이 요구한 조건을 만족함은 쉽게 확인된다.

이로써 정리 3.2와, 따라서 (3.1)이 완전히 증명되었다.

그 증명은 동시에 다음을 증명한다.

\result{명제}{3.8} (3.2)의 조건 아래에서
$Q=\operatorname{Quot}^P_{\mathcal F/X/S}$,
$X_Q=X\times_SQ$,
$\mathcal F_Q=\mathcal F\otimes_{\mathcal O_S}\mathcal O_Q$라 하고,
$(Q,\mathcal G)$가 함자
$\underline{\operatorname{Quot}}^P_{\mathcal F/X/S}$를 표현하도록 하는,
$Q$ 위에서 평탄하고 상대 힐베르트 다항식이 $P$인 $\mathcal F_Q$의
연접 몫 $\mathcal G$를 생각하자. 그러면 정수 $\nu$가 존재하여,
$n\geq\nu$이면 $(f_Q)_*(\mathcal G(n))$은 $Q$ 위의 계수 $P(n)$인
국소 자유 가군층이고 $S$에 대해 매우 풍부하다. 즉, $Q$에서 $S$ 위의
그라스만 스킴 $\operatorname{Grass}_{P(n)}(M)$으로 가는 몰입 사상을
정의한다. 특히 $n\geq\nu$이면 $Q$ 위의 층
\[
\bigwedge^{P(n)}(f_Q)_*(\mathcal G(n))
\]
은 가역이고 $S$에 대해 매우 풍부하다.

(3.2)에서와 같이 $\mathcal F$가 $S$ 위에서 평탄한 경우로 환원할 수
있으며, 그때에는 앞의 표기에서 $A_\nu=A$가 되는 정수 $\nu$를 택하면
충분하다.

(3.2)가 적용되는 가장 중요한 경우는 $\mathcal F=\mathcal O_X$인
경우이다. 이때
\[
\operatorname{Quot}_{\mathcal O_X/X/S}=\operatorname{Hilb}_{X/S},
\qquad
\operatorname{Quot}^P_{\mathcal O_X/X/S}=\operatorname{Hilb}^P_{X/S}
\]
라고 쓰므로, 다음과 같은 분해를 얻는다.
\[
\operatorname{Hilb}_{X/S}=\coprod_P\operatorname{Hilb}^P_{X/S}.
\]
정의에 따라 $\operatorname{Hilb}_{X/S}$는 함자
$\underline{\operatorname{Hilb}}_{X/S}$를 표현한다. 여기서
$\underline{\operatorname{Hilb}}_{X/S}(S')$는
$X'=X\times_SS'$의 닫힌 부분준스킴 가운데 $S'$ 위에서 평탄한 것들의
집합이다. 또한 $\operatorname{Hilb}^P_{X/S}$는 주어진 힐베르트
다항식 $P$를 갖는 닫힌 부분준스킴들에 대응하는 부분함자를 표현한다.
이 준스킴들은 각각 $S$ 위의 $X$의 \emph{힐베르트 준스킴},
\emph{지표 $P$인 힐베르트 준스킴}이라고도 한다. 이 용어는 이 이론에서
힐베르트 다항식이 하는 역할로 정당화된다. 이것들과 고전적인 차우
다양체들(다양체가 아니라 순환을 매개하기 위한 것)의 본성상의 차이는,
\src{266}
한 다양체의 순환류들로 이루어진 차우 환과 그 다양체 위의 층들의 류로
이루어진 환(리만--로흐 정리 [1]에서 도입되는 것과
같다\footnote{인쇄본에는 여성형 « telle »라고 되어
있으나, 그 선행사인 « l'anneau »는 남성명사이다.}) 사이의 차이와 같은
정도이다. 실제로 $X=\mathbf P^r_S$이고 $S$가 어떤 체의 스펙트럼일 때,
$X$ 위의 연접 가군층 $\mathcal F$의 힐베르트 다항식을 아는 것은
$\mathcal F$의 천 특성류들을 아는 것과도 동치이고, 또는 $X$ 위의 연접층들의
류환에서 $\mathcal F$의 류를 아는 것과도 동치임에 유의하자.

\result{주의}{3.9} 또한 $\operatorname{Quot}_{\mathcal F/X/S}$와
$\operatorname{Quot}^P_{\mathcal F/X/S}$의 구성을
$X=\mathbf P^r_S$, $\mathcal F=\mathcal O_X^N$이고
$\mathcal O_X(1)$이 통상적인 매우 풍부한 층인 경우로 환원했음에
유의하자. 더 정확히 말해, 일반적인
$\operatorname{Quot}_{\mathcal F/X/S}$는 앞의 것들의 닫힌
부분준스킴으로 실현된다. $\operatorname{Quot}^P_{\mathcal F/X/S}$의
형성은 기저변환 $S'\longrightarrow S$와 분명히 양립하므로, 나아가
$S=\operatorname{Spec}(\mathbf Z)$인 경우로 환원된다. 결국에는
$\mathbf Z$ 위의 다음 사영 스킴들을 연구하는 것으로 환원되는 셈이다.
\[
\mathfrak Q^P_{r,N}=
\operatorname{Quot}^{P}_{(\mathcal O_{\mathbf P^r_{\mathbf Z}})^N/
\mathbf P^r_{\mathbf Z}/\operatorname{Spec}(\mathbf Z)},
\]
그리고 특히 다음 절대 힐베르트 스킴들을 연구하는 것으로 환원된다.
\[
\operatorname{Hilb}^P_r=\mathfrak Q^P_{r,1}.
\]
이 스킴들의 연결 성분들(과연 연결인가?)을 결정하는 것부터
시작하여,\footnote{1962년 정오표는 Hartshorne이 대수적으로 닫힌 체 위의
$\operatorname{Hilb}^P_r$가 연결임을 증명하고
$\operatorname{Hilb}^P_r\neq\varnothing$인 순서쌍 $(r,P)$를 결정했다고
밝힌다.} 이들의 기약 성분들을 결정하는 것까지 더 자세히 연구한다면 매우
흥미로울 것이다(Serre [7]에 따르면 영이 아닌 소수 $p$ 위에 전부 놓이는
기약 성분이 있을 수 있다). $\operatorname{Spec}(\mathbf Z)$의
$s$ 위에서 $\operatorname{Hilb}^P_r$의 섬유가 갖는 기약 성분들이
소체의 ``정칙'' 확대들에 대응하는지, 즉 이들이 ``상대적으로 연결''인지
묻는 Weil의 문제를 상기하자. 이러한 문제들은 ``차우 다양체''보다
힐베르트 스킴에 대해서 더 다루기 쉬울 가능성이 있다.

\fsection{4}{변형}

\noindent a. (3.1)의 조건 아래에서 $U$를 $X$의 열린집합이라 하고,
$A=\underline{\operatorname{Quot}}_{\mathcal F/X/S}$의 부분함자 $A'$를
다음과 같이 정하자. $A'(S')$는 $\mathcal F'$의 몫가군층 $\mathcal G$
가운데 $S'$ 위에서 평탄하고 그 지지집합이 $U'$에 포함되는 것들의
집합이다. 곧바로 $A'$는 $A$를 표현하는 준스킴
$\operatorname{Quot}_{\mathcal F/X/S}$의 열린 부분집합으로 표현됨을
알 수 있다. 따라서 정리 (3.1)과 (3.2)는
\src{267}
$X$가 $S$ 위에서 사영이라는 가정을 준사영이라는 가정으로 바꾸어도
여전히 성립한다. 다만 결론에서도 ``사영''을 ``준사영''으로 바꾸고,
이제 $\underline{\operatorname{Quot}}_{\mathcal F/X/S}(S')$로
$\mathcal F'$의 연접 몫 $\mathcal G$ 가운데 $S'$ 위에서 평탄하고
그 지지집합이 $S'$ 위에서 고유인 것들의 집합을 뜻하도록 한다.

\medskip
\noindent b. 일반적으로는
$\mathcal F'=\mathcal F\otimes_{\mathcal O_S}\mathcal O_{S'}$의
$S'$-평탄 몫 $\mathcal G$들에, 기저변환 아래에서 안정적인 여러 가지
자연스러운 추가 조건을 부과할 수 있다. 이로써
$\underline{\operatorname{Quot}}_{\mathcal F/X/S}$의 여러 부분함자를
얻고, 이들을 표현하고자 하는 문제를 생각하게 된다. 통상적인 판정법으로
많은 경우 이 함자들 역시 $\operatorname{Quot}_{\mathcal F/X/S}$의
열린 부분집합들로 표현됨을 증명할 수 있다. 특히 다음 추가 성질 가운데
하나를 부과하는 경우가 그러하다.

\smallskip
\noindent $1^\circ$ 섬유 $X'_{s'}$ 위에 유도되는 가군층
$\mathcal G_{s'}$ ($s'\in S'$)에 결부된 소순환들의 차원이 주어진
정수 집합에 속한다.\footnote{인쇄본에는
$\mathcal G_s$라고 되어 있으나, 양화된 점은 $s'$이다.}

\smallskip
\noindent $2^\circ$ ($\mathcal F=\mathcal O_X$여서 $\mathcal G$가
$X'$의 닫힌 부분준스킴 $Y$에 대응하는 경우) $Y$는 $S'$ 위의 매끄러운
준스킴이고\footnote{기저변환 뒤에 나오는 이 조건의
여섯 곳에서 인쇄본은 잘못하여 $S$ 또는 $X$를 그대로 두었지만, 그 족은
$X'=X\times_SS'$ 위에서 $S'$에 대하여 놓인다.}, 각각 $S'$ 위에서
정규이거나(즉, 섬유 $Y_{s'}$가 ``$k(s')$ 위에서''
정규이다\footnote{인쇄본에는 $k(s)$라고 되어 있으나,
점 $s'$의 잉여체는 $k(s')$이다.}; 다시 말해 기초체를 어떻게 확대해도
정규이다), 또는 ($X$가 $S$ 위에서 평탄한 경우) $S'$에 대해 $X'$ 안의
국소 완전 교차이다.\footnote{인쇄본에는
``$k$-intersections''라고 되어 있다. 이 조건은 기저에 관한 것이므로 여기서는 기저변환된 $S'$를 써야 한다.}
마지막 조건은 섬유 $Y_{s'}$가 $X'_{s'}$ 안의 국소 완전 교차라는 뜻이다.

다른 조건들은 $X'_{s'}$ 위에 유도되는 가군층 $\mathcal G_{s'}$의
코호몰로지적 성질 등을 사용하게 될 것이다. 물론, 각각이
$\operatorname{Quot}_{\mathcal F/X/S}$의 열린집합 $U_i$로 표현되는
조건들의 논리곱은 그 열린집합들의 교집합으로 표현된다. 예를 들어,
$S$ 위의 각 $S'$에 $X'=X\times_SS'$의 닫힌 부분준스킴 $Y$ 가운데
$S'$의 주어진 계수 $r$인 에탈 피복([3], I)인 것들의 집합을 대응시키면,
$S'$에 관한 표현 가능한 반변 함자를 얻는다.

\medskip
\noindent c. [2], II, C, n\textsuperscript{o} 2에서 정의한 준스킴
$\operatorname{Hom}_S(X,Y)$, $\Pi_{X/S}(Z/X)$,
$\operatorname{Isom}_S(X,Y)$는 적절한 사영성 가정 아래에서 존재하며,
적절한 힐베르트 준스킴들의 열린 부분집합으로 실현된다. 등식
\[
\operatorname{Hom}_S(X,Y)=\Pi_{X/S}((X\times_SY)/X)
\]
이 성립하므로, $\operatorname{Hom}_S(X,Y)$의
경우는\footnote{인쇄본에는 여기에서
$\operatorname{Hom}_S(X,X)$라고 되어 있으나, 이는 바로 앞의 표시된
등식과 양립하지 않는다.} $\Pi_{X/S}(Z/X)$의 경우로 환원된다.
$S$ 위의 각 $S'$에 대해, $X'=X\times_SS'$ 위의
$Z'=Z\times_SS'$의 절단들의 집합은, $Z'$에서 $X'$로 가는 사상으로부터
유도되는 사상 $\Gamma\longrightarrow X'$가 동형사상인 $Z'$의
부분준스킴 $\Gamma$들의 집합과 일대일로 대응한다.\footnote{
인쇄본에는 $Z$라고 되어 있으나, 여기서 생각하는
그래프들은 $Z'=Z\times_SS'$의 절단들의 그래프이다.} 이 부분준스킴들은
$Z$가 $X$ 위에서 분리되어 있으면 반드시 닫혀 있다. 따라서
\src{268}
$X$가 $S$ 위에서 평탄하고 고유이며 $Z$가 $S$ 위에서 준사영이면
$\Pi_{X/S}(Z/X)$가 존재하고 $\operatorname{Hilb}_{Z/S}$의 열린
부분준스킴으로 실현된다. 그러므로 $X$가 $S$ 위에서 사영이고 평탄하며
$Y$가 $S$ 위에서 준사영이면 $\operatorname{Hom}_S(X,Y)$가 존재하고
$\operatorname{Hilb}_{(X\times_SY)/S}$의 열린 부분준스킴으로 실현된다.
$X$와 $Y$가 모두 $S$ 위에서 사영이면, 곧바로
$\operatorname{Isom}_S(X,Y)$도 존재하고
$\operatorname{Hom}_S(X,Y)$의 열린 부분집합으로 표현된다. 마찬가지로
$X$가 $S$ 위에서 평탄하고 사영이며 $Y$가 $S$ 위에서 준사영이면,
$S'$-사상 $X'\longrightarrow Y'$ 가운데 몰입 사상인 것들에 대응하는
$\operatorname{Hom}_S(X,Y)$의 부분함자를 나타내는 $S$-준스킴
$\operatorname{Imm}_S(X/Y)$도 $\operatorname{Hom}_S(X,Y)$의 열린
부분집합으로 표현된다.

$\mathcal L$을 $X$ 위의, $\mathcal M$을 $Y$ 위의 가역층이라 하고,
각각 $S$에 대해 매우 풍부하다고 하자.\footnote{인쇄본에는 괄호 안의 대응 대상을 동반한 단수 주어 앞에 복수형
« Soient »가 쓰였다.} 그러면
$\mathcal L\otimes_{\mathcal O_S}\mathcal M$은 $X\times_SY$ 위에서
$S$에 대해 매우 풍부한 층이다. 따라서 유리수 계수의 각 다항식 $P$에
대해 $\operatorname{Hilb}^P_{(X\times_SY)/S}$가 정의되며, 이는 $S$
위의 준사영 준스킴이다. 따라서 이 준스킴은
$\operatorname{Hom}_S(X,Y)$ 위에 $S$ 위에서 준사영인 동시에 열리고
닫힌 부분집합을 유도하며, 이를 $\operatorname{Hom}_S(X,Y)^P$라고
쓰겠다. 그러므로 $S$ 위에서 $\operatorname{Hom}_S(X,Y)^P$의 절단들은,
모든 정수 $n$에 대해
\[
\chi\!\left(\left(\mathcal L\otimes_{\mathcal O_X}g^*(\mathcal M)\right)^{
\otimes n}\right)=P(n)
\]
을 만족하는 $S$-사상 $g:X\longrightarrow Y$이다.

이 방법으로, ``편극된'' 사영다양체의 자기동형들이 대수군을 이룬다는
마쓰사카 정리의 일반화들을 얻는다. 여기에서는 이 군을 어떤 보편 문제의
해로 정의할 수 있으므로, 이 명제가 분명히 더 정확한 의미를 갖는다.
또한 대수적으로 닫힌 체 위에서 예전에 생각하던 자기동형군은 여기서
정의한 ``참된'' 자기동형군에서 멱영원들을 제거하여 얻는 것임에 유의하자.
기초체를 확대하고 나서 나타나는 멱영원들의 아이디얼이 반드시 ``$k$ 위에서
정의''되는 것은 아니므로, 이는 더 오래된 구성들이 비완전 기초체 위에서
이루어질 가능성이 거의 없는 까닭을 설명한다. 같은 지적은 옛 방식의
구성들 대부분에도 적용된다.

\fsection{5}{힐베르트 스킴의 미분적 연구}

이는 다음 결과에서 따른다.

\src{269}
\result{명제}{5.1} $S$를 준스킴, $S_0$를 제곱이 영인 준연접
아이디얼 $\mathcal J$로 정의된 부분준스킴, $X$를 $S$-준스킴,
$\mathcal F$를 $X$ 위의 준연접 가군이라 하자. 또한
$X_0=X\times_SS_0$,
$\mathcal F_0=\mathcal F\otimes_{\mathcal O_S}\mathcal O_{S_0}$라 하고,
$\mathcal G_0=\mathcal F_0/\mathcal K_0$를 $S_0$ 위에서 평탄인
$\mathcal F_0$의 준연접 몫가군이라 하자. $X$의 각 열린집합 $U$에
대하여 $\mathcal E(U)$를 다음 조건을 만족하는 $\mathcal F|U$의 준연접
몫가군 $\mathcal G$들의 집합이라 하자. 즉 $\mathcal G$는 $S$ 위에서
평탄이고
\[
\mathcal G\otimes_{\mathcal O_S}\mathcal O_{S_0}=\mathcal G_0
\]
이다. 그러면 $U$가 변할 때 $\mathcal E(U)$들은 $X$ 위의 한 층
$\mathcal E$의 단면들을 이룬다. 이제 군층
\[
\mathcal A=\operatorname{Hom}_{\mathcal O_{X_0}}
(\mathcal K_0,\mathcal G_0\otimes_{\mathcal O_{S_0}}\mathcal J)
\]
은 $\mathcal E$에 자연스럽게 작용하며, 이로써 $\mathcal E$는
``$\mathcal A$ 아래에서 형식적으로 주동차적인 층''이 된다. 즉 $X$의
모든 열린집합 $U$에 대하여 $\mathcal E(U)$는 공집합이거나
$\mathcal A(U)$ 아래의 주동차 집합이다.

이로부터 다음을 얻는다.

\result{따름정리}{5.2} $X$ 위에서 국소적으로 $\mathcal G_0$를
$\mathcal F$의 $S$-평탄 몫 $\mathcal G$로 연장할 수 있다고 가정하자.
즉 층 $\mathcal E$의 섬유들이 공집합이 아니라고 가정하자. 그러면 표준
장애류
\[
c(\mathcal G_0)\in H^1(X,\mathcal A)
\]
가 존재하며, 이 류가 영일 필요충분조건은 $\mathcal G_0$를
$\mathcal F$의 $S$-평탄 몫인 전역 가군 $\mathcal G$로 연장할 수 있는
것이다. 이 류가 영이면 가능한 모든 연장들의 집합 $\mathcal E(X)$는
\[
\mathcal A(X)=\operatorname{Hom}_{\mathcal O_{X_0}}
(\mathcal K_0,\mathcal G_0\otimes_{\mathcal O_{S_0}}\mathcal J)
\]
아래의 주동차 집합이다.
\footnote{이 마지막 공식에서 인쇄본은 바로 앞의 $\mathcal A$ 정의에
고정되어 있는 $\mathcal O_{X_0}$ 대신 $\mathcal O_X$를 쓴다.}
특히 $H^1(X,\mathcal A)=0$이면 전역 연장이 존재한다.

\result{따름정리}{5.3}
$Q=\operatorname{Quot}_{\mathcal F/X/S}$가 존재한다고 가정하자
(4 (a) 참조). 예를 들어 $S$가 국소 뇌터이고 $X$가 $S$ 위에서
준사영이며 $\mathcal F$가 연접이면 그러하다. $x\in Q$가 어떤
$s\in S$에 대하여 $k(s)$의 잉여 확대 $K=k(x)$에 대응한다고 하자.
\src{270}
따라서 $x$는 $K$-준스킴 $X_K$ 위의 가군
$\mathcal F_0=\mathcal F_K$의 연접 몫가군
$\mathcal G_0=\mathcal F_0/\mathcal K_0$에 의해 정의된다. $\mathcal A$를
다음과 같이 정의되는 $X_K$ 위의 연접층이라 하자.
\[
\mathcal A=\operatorname{Hom}_{\mathcal O_{X_0}}(\mathcal K_0,\mathcal G_0).
\]
그러면 점 $x$에서 섬유 $Q_s$의 자리스키 접공간, 곧
$\mathcal O_{Q_K,x}$의 극대 아이디얼을 $\mathfrak m$이라 할 때
$\mathfrak m/\mathfrak m^2$의 $K$-쌍대는 $H^0(X_K,\mathcal A)$와
표준적으로 동형이다.

자리스키 접공간을 주는 이 결과는 다음과 같이 일반화할 수 있다. 주어진
$S$-사상 $g:S'\longrightarrow Q$, 즉 $S'$ 위의
$Q'=Q\times_SS'$의 단면 $g'$에 대하여 가군
\[
\Omega=g^*(\Omega^1_{Q/S})=g'^*(\mathcal J/\mathcal J^2)
\]
을 생각하자. 여기서 $\mathcal J$는 $S'$ 위의 $Q'$의 단면 $g'$가
정의하는 $Q'$ 위의 아이디얼이다. 이때 $\Omega$는 $S'$ 위의 연접 가군
$\mathcal M$에 관하여 함자적인 다음 공식으로 특징지어진다.
\[
\operatorname{Hom}_{\mathcal O_{S'}}(\Omega,\mathcal M)
\simeq H^0(X',\mathcal A),
\]
여기서 $X'=X\times_SS'$이고 $\mathcal A$는 $X'$ 위의 다음 가군이다.
\[
\mathcal A=\operatorname{Hom}_{\mathcal O_{X'}}
(\mathcal K,\mathcal G\otimes_{\mathcal O_{S'}}\mathcal M).
\]
\footnote{인쇄본은 이 텐서곱에 $\mathcal O_S$를 쓴다. 그러나 여기서
$\mathcal G$와 $\mathcal M$은 $S'$ 위의 가군들이며, 제곱영 변형의
구성에는 $\mathcal O_{S'}$ 위의 텐서곱이 필요하다.}
여기서 $\mathcal G=\mathcal F'/\mathcal K$는 $g$에 대응하는
$\mathcal F'=\mathcal F\otimes_{\mathcal O_S}\mathcal O_{S'}$의
몫가군이다. 실제로 (5.1)에서 $S_0$를 $S'$으로 바꾸고 $S$를 준스킴
$D(\mathcal M)=(S',\mathcal O_{S'}+\mathcal M)$으로 바꾸면 된다. 이때
$\mathcal M$은 제곱이 영인 아이디얼로 본다.

(5.1)에서 $\mathcal F=\mathcal O_X$인 경우, $\mathcal G_0$를 주는 것은
아이디얼 $\mathcal J_0=\mathcal K_0$로 정의되고 $S_0$ 위에서 평탄인
$X_0$의 닫힌 부분준스킴 $Y_0$를 주는 것과 같다. 이때 (*)는
\[
\mathcal A=\operatorname{Hom}_{\mathcal O_{X_0}}
(\mathcal J_0/\mathcal J_0^2,
\mathcal O_{Y_0}\otimes_{\mathcal O_{S_0}}\mathcal J)
\]
를 준다. 여기서 $\mathcal J_0/\mathcal J_0^2$는 $Y_0$의 $X_0$ 안에서의
코노멀층으로 해석하며, 이를
\src{271}
$\check{\mathcal N}_{Y_0/X_0}$라고도 쓴다. 이 경우에는 $\mathcal A$를
$Y_0$ 위의 가군으로 보고 $H^0$와 $H^1$을 $Y_0$ 위에서 계산하는 것이
유용하다. 더구나 $X$가 $S$ 위에서 평탄이고 $Y_0$가 $X_0$ 안에서
국소 완전교차이면, (5.1)의 국소 연장이 존재한다. 또한
$\mathcal J_0/\mathcal J_0^2$는 $Y_0$ 위에서 국소 자유이고 다음과 같이
쓸 수 있다.
\[
\mathcal A=\mathcal N_{Y_0/X_0}\otimes_{\mathcal O_{S_0}}\mathcal J,
\]
여기서 오른쪽의 첫째 인자는 $Y_0$의 $X_0$ 안에서의 법층이다.
매끄러움의 기본 판정법([3], III, 3.1)을 사용하면 예를 들어 다음을
얻는다.

\result{따름정리}{5.4} (5.3)의 조건 아래에서
$\mathcal F=\mathcal O_X$이고 $X$가 $S$ 위에서 평탄이며,
$\mathcal G_0$에 대응하는 $X_0$의 닫힌 부분준스킴 $Y_0$가 국소
완전교차라고 가정하자. 그러면 점 $x$에서 $Q_s$의 자리스키 접공간은
$H^0(Y_0,\mathcal N_{Y_0/X_0})$와 표준적으로 동형이다. 만일
$H^1(Y_0,\mathcal N_{Y_0/X_0})=0$이면 힐베르트 준스킴
$\operatorname{Hilb}_{X/S}$은 점 $x$에서 $S$ 위로 매끄럽다.
($\mathcal N_{Y_0/X_0}$는 $Y_0$의 $X_0$ 안에서의 법층이다.)

\result{비고}{5.5} 이 명제는 특히 $Y_0$가 한 방정식으로 정의되는
$X_0$ 안의 완전교차, 즉 양의 ``카르티에 인자''인 경우에 적용된다.
그러면 $\mathcal N_{Y_0/X_0}$는 인자 $Y_0$가 정의하는 $X_0$ 위의
가역층 $\mathcal J_0^{-1}$에서 유도되는 $Y_0$ 위의 층과 동형이다. 특히
비특이 사영 다양체 $X_0$ 위의 양의 인자들의 족을 연구할 때 이러한
상황이 나타난다. $Q$의 점 $x$에서의 자리스키 접공간과
$H^0(Y_0,\mathcal N_{Y_0/X_0})$ 사이의 동형은, 원한다면 인자들에
대응하는 $Q$의 열린부분집합 $D$의 점 $x$에서의 자리스키 접공간과의
동형으로 보아도 되는데, 고전 대수기하학에서는 첫째 공간에서 둘째
공간으로 가는 ``특성 준동형''이라는 이름으로 알려져 있었다. 이
준동형은 인자들의 ``완전 연속족''의 매개변수 다양체 $T$에서 $x$가
매끄러운 점일 때에만 정의되었다. 우리의 관점에서 $T$는 유도된 축소
구조를 갖춘 스킴 $D$의 한 기약 성분이다. 이때 $x$에서 $T$의 접공간은
$x$에서 $D$의 접공간의 부분공간이므로, 고전적인 특성 준동형은 실제로
단사이다. 그러나 이는 추가 조건 아래에서만 전사이며, 예를 들어
$\mathcal O_{D,x}$가 정역이면 전사이다. 실제로 Zappa [8]는
\src{272}
복소수체 위의 비특이 사영 곡면 $X$에 대하여 $T$의 일반점에서조차
특성 준동형이 전사가 아닌 예를 구성하였다. 따라서 고려 중인 기약
성분의 일반점에서 $D$의 국소환조차 정역이 아니다. 이는 가장 고전적인
곡면론의 현상들을 이해하는 데 멱영원을 지닌 다양체들이 얼마나
필수적인지를 특히 인상적으로 보여 준다.\footnote{1962년의 정오표는
Mumford가 $\mathbf P^3_{\mathbf C}$의 힐베르트 스킴에서, 일반점들이 차수
14이고 종수가 24인 비특이 곡선들을 나타내며 그 일반점에서 비축소인 한
기약 성분을 막 구성했다고 덧붙인다. 그러한 곡선 가운데 하나를
블로업하면 그 형식적 모듈라이 스킴이 일반점에서 비축소인 3차원 정칙
사영 스킴도 얻어진다.}

\subsection*{5.6}

(5.4)에서 특히 약수 스킴들에 적용되는 매끄러움 판정법을 제시하였다.
고다이라는 [6]에서 다른 판정법, 즉 $\mathcal L=\mathcal J_0^{-1}$이
약수 $Y_0$가 정하는 $X_0$ 위의 가역층일 때
$H^1(X_0,\mathcal L)$이 0이라는 판정법을 제시하였다. 이 판정법은 $S$가
표수 0인 체의 스펙트럼일 때 성립하며, 기초체가 $\mathbf C$인 경우에는
[6]에서 초월적 방법으로 증명되었다. 여기서 일반적으로 $S$를 다시 임의로
두면, 고다이라의 조건은 표준 사상
\[
D\longrightarrow\operatorname{Pic}_{X/S}
\]
이, 즉 $X/S$의 약수 준스킴에서 피카르 준스킴으로 가는 사상이 문제의 점
$x$에서 매끄럽기 위한 충분조건임을 지적해 두자
($\operatorname{Pic}_{X/S}$의 존재가 확보되면 통상적인 매끄러움 판정법으로
쉽게 확인된다). 따라서 더 나아가 $\operatorname{Pic}_{X/S}$가 $x$의 상인
점에서 $S$ 위로 매끄럽다면(예를 들어 $\operatorname{Pic}_{X/S}$가 $S$ 위로
매끄럽다면), $D$는 $x$에서 $S$ 위로 매끄럽다. 한편 카르티에는 표수 0인
체 $k$ 위에서 국소 유한형인 모든 군 준스킴이 $k$ 위로 매끄럽다는 것을
증명하였다. 이 두 결과를 결합하면 고다이라의 결과를 다시 얻는다. 이
논의로부터, 표수가 $p>0$인 체 $K$ 위에서
$\operatorname{Pic}_{X/S}$가 $k$ 위로 매끄럽지 않을 때($X$가 이구사
곡면이면 그러하다), $H^1(X_0,\mathcal L)=0$이라는 조건은 오히려 $D$가
$x$에서 매끄럽지 않음을 뜻하며, $K$가 대수적으로 닫혀 있으면 $x$에서
축소되지 않기까지 함을 알 수 있다.

끝으로 섬유공간의 미분론적 연구에서 중요한 역할을 하는 다음 결과를
기록해 두자.

\result{명제}{5.7} $X$를 국소 뇌터 준스킴 $S$ 위에서 유한하고 평탄한
준스킴이라 하고, $Z$를 $X$ 위의 준스킴이라 하자.
$\Pi_{X/S}(Z/X)$가 존재한다고 가정한다($Z$가 $X$ 위에서 준사영이면
그러하다). 이때 $Z$가 $X$ 위로 매끄러우면 $\Pi_{X/S}(Z/X)$는 $S$ 위로
매끄럽다.

이는 정의와 통상적인 매끄러움 판정법([3], III, 3.1)에서 곧바로 따른다.
$X$가 $S$ 위에서 유한하고 평탄할 때에는
$\Pi_{X/S}(Z/X)$의 존재 문제\footnote{인쇄본에는 여기서 $Z/K$라고 되어
있다. $K$는 정의되어 있지 않고, (5.7)부터 고찰한 단면 함자는
$\Pi_{X/S}(Z/X)$이다.}를 힐베르트 스킴 이론을 사용하지 않고도 매우
초등적으로 다룰 수 있음에 유의하자.
\src{273}
예를 들어 $X$가 $S$ 위에서 라디시엘이면 $Z$에 아무런 제한을 두지 않아도
$\Pi_{X/S}(Z/X)$가 존재한다. 또 다른 예로, $T$를 $S$-준스킴이라 하고
$T_n$을 $T\times_ST$ 안에서 그 대각선의 ``$n$차 무한소 근방''이라 하자.
두 사영이 유도하는 사상
$p_1,p_2:T_n\longrightarrow T$를 갖추어 둔다. $p_1$을 통해 $T_n$을 $T$
위의 유한 준스킴으로 보고, 더 나아가 $T$ 위에서 평탄하다고 가정하자
($T$가 $S$ 위로 매끄러우면 그러하다). $T$ 위의 모든 준스킴 $X$에 대하여
\[
(X/T/S)^{(n)}=
\Pi_{T_n/T}\bigl(p_2^*(X/T)/T_n\bigr)
\]
로 놓는다.
\footnote{인쇄본에서는 $\Pi$ 기호 아래에 $T_n/S$를 둔다. 여기서는 유한
평탄 사상 $p_1:T_n\to T$에 명제 5.7을 적용하며, 그 결과는 명시적으로
$T$ 위의 준스킴이다. 따라서 분모는 $T_n/T$이다.}
이는 $T$ 위의 준스킴으로서 ($S$에 대한) \emph{$T$ 위의 $X$의 $n$차
단면 싹 다발}이라 한다. 이 준스킴은 $X$에 함자적으로 의존하며, $X$가
매끄러우면 $T$ 위로 매끄럽다.

\fsection{6}{노름과 대칭곱 사이의 관계}

$S$를 준스킴, $X$와 $Y$를 $S$-준스킴이라 하고,
\[
u:(X/S)^n\longrightarrow Y
\]
를 $X/S$의 $n$중 데카르트곱에서 $Y$로 가는 대칭인 $S$-사상이라 하자.
단순하게 하기 위해 $S$는 국소 뇌터이고 $X$와 $Y$는 $S$ 위에서
유한형이라고 가정하자. 이때 $X$ 위의 연접 가군층 $\mathcal F$로서 그
지지집합이 $S$ 위에서 유한하고, $S$ 위에서 평탄하며, $S$ 위의 상대 계수가
$n$인 것(즉 $f_*(\mathcal F)$가 $S$ 위의 계수 $n$인 국소 자유 가군층인
것)에 자연스럽게 $S$ 위의 $Y$의 한 단면을 대응시킬 수 있다.
\[
\pi^u_{X/S}(\mathcal F)\in\Gamma(Y/S).
\]
여기서는 엄밀한 정의를 제시하지 않고, 이렇게 얻는 형식체계가 통상적인
노름과 대각합의 형식체계를 자연스럽게 일반화한 것임만 지적하겠다. 임의의
분리 사상 $X/S$에 대해 노름 법칙의 표준 표적은 분할 거듭제곱 공간
$\Gamma^n_S(X)$이다. $X$가 $S$ 위에서 평탄하고 $n$차 대칭곱이 존재하면
(예를 들어 $(X/S)^n$에서 $\mathfrak S_n$의 궤도들이 아핀 열린집합들에
포함될 때), 표준 사상
$\operatorname{Symm}^n_S(X)\to\Gamma^n_S(X)$은 동형사상이고, 다음 표준
원소를 얻는다.
\[
\pi_{X/S}(\mathcal F)\in\Gamma(\operatorname{Symm}^n_S(X)/S).
\]
\src{274}
이 원소로부터 $\pi^u_{X/S}(\mathcal F)$들을 다시 얻을 수 있다. 또 하나의
중요한 경우는 $X$가 $S$ 위의 가환 모노이드이고 $X=Y$이며 $u$가 $X$의
합성 법칙에서 오는 경우이다. 이때에는 $X$ 위의 가군층 $\mathcal F$에
대응하는 $S$ 위의 $X$의 단면을 간단히 $\pi(\mathcal F)$로 쓴다.

이제 $X$ 위의 연접 가군층 $\mathcal F$가 주어져서
$\operatorname{Quot}_{\mathcal F/X/S}$가 존재한다고 하자. 또는 적어도
$S$ 위의 각 $S'$에 $\mathcal F'=\mathcal F\otimes_{\mathcal O_S}
\mathcal O_{S'}$의 연접 몫층 $\mathcal M$들 가운데 $S'$ 위에서 평탄하고
상대 계수가 $n$인 것들의 집합을 대응시키는 함자
$\underline{\operatorname{Quot}}^n_{\mathcal F/X/S}$가 $S$-준스킴
$\operatorname{Quot}^n_{\mathcal F/X/S}$로 표현된다고 하자. ($X$가 $S$
위에서 준사영이면 $\operatorname{Quot}^n_{\mathcal F/X/S}$는 실제로
존재하며, 제3절의 표기로는 $P$가 상수항 $n$뿐인 다항식일 때의
$\operatorname{Quot}^P_{\mathcal F/X/S}$와 다르지 않다.)
$\pi^u_{X/S}(\mathcal M)$의 구성이 밑변경과 양립하므로 다음 표준 사상을
얻는다.
\[
\pi^u_{X/S}:\operatorname{Quot}^n_{\mathcal F/X/S}\longrightarrow Y,
\]
그리고 일반적으로
\[ \pi_{X/S}:\operatorname{Quot}^n_{\mathcal F/X/S}
\longrightarrow\Gamma^n_S(X). \]
$X$가 $S$ 위에서 평탄하면 이 표적을 $\operatorname{Symm}^n_S(X)$와
동일시한다. 앞의 동일시와 무관하게, $X$가 $S$ 위에서 분리되어 있으면
$\mathcal F=\mathcal O_X$인 경우의 유한 평탄 족을 통한 구성은 다음 사상을
준다.
\[
\pi_{X/S}:\operatorname{Hilb}^n_{X/S}
\longrightarrow\operatorname{Symm}^n_S(X).
\]
이는 $n=0$ 또는 $n=1$이면 분명히 동형사상이다. 그러나
$n>1$이면\footnote{인쇄본은 방금 $n=1$에 대해 단언한 동형사상과
모순되게 $n\geq1$을 되풀이한다.}, $S$가 체 $k$의 스펙트럼이고 $X$가
$S$ 위로 매끄러운 경우에도 일반적으로 동형사상이 아니며 단사인 사상조차
아니다. 실제로 $X$의 0차원 부분스킴(예를 들어 $X$의 닫힌점 $x$에 대해
국소환 $\mathcal O_{X,x}$의 극대 아이디얼에 대한 준소 아이디얼
$\mathfrak i$에 대응하는 것)은 그것이 정하는 순환을 알아도 정해지지
않는다(이 경우에는 $\mathcal O_{X,x}$ 안에서 $\mathfrak i$의 $k$-여차원을
알아도 정해지지 않는다). $S$를 다시 임의로 두면 다음 사실만을 말할 수
있다.

\medskip
\noindent a. $X$가 $S$ 위로 매끄러우면, 노름 사상은 $X$에 포함된 계수
$n$의 에탈 피복들을 분류하는 $\operatorname{Hilb}^n_{X/S}$의 열린집합
(제4절 (b) 참조)과, 중복 성분이 없는 $n$-순환들에 대응하는
$\operatorname{Symm}^n_S(X)$의 열린집합 사이의 동형사상을 정한다.
\src{275}

\medskip
\noindent b. 더 나아가 $X$의 $S$ 위의 상대 차원이 1이면, 노름 사상은
$\operatorname{Hilb}^n_{X/S}$와 $\operatorname{Symm}^n_{X/S}$ 사이의
동형사상을 정한다.

마지막 사실은 비특이 대수곡선의 0차원 부분스킴은 그에 대응하는 순환을
알면 정해지기 때문이다. 같은 논의는 더 일반적으로 비특이 대수 스킴 위의
유효 카르티에 약수에도 적용된다(그리고 이 매우 특별한 경우에는 차우
다양체가 힐베르트 스킴과 같은 것을 줄 가능성도 배제되지 않는다).

\fsection{7}{보충과 문제}

J.-P.~세르가 지적했듯이, 나가타의 잘 알려진 예로부터 다음과 같은 것들을
찾을 수 있다. 즉 체 $k$의 스펙트럼인 스킴 $S$, $k$의 이차 확대체 $k'$의
스펙트럼인 $S$-스킴 $S'$, 그리고 차원이 3이고 고유하고 매끄럽지만
사영이지 않은 $S'$-스킴 $X$로서 $\Pi_{S'/S}(X/S')$가 존재하지 않는 것을
찾을 수 있다. 따라서 더구나 힐베르트 스킴
$\operatorname{Hilb}^2_{X/S}$도 존재하지 않는다($X$에 포함된 계수 2의
$S$의 에탈 피복들을 표현하는 $k$-스킴조차 존재하지 않으며, 더구나 $X$의
2차 대칭곱도 존재하지 않는다. 앞 절을 참조하라). 이는 대수기하학에서
비사영적 구성을 할 가능성에 심각한 제한을 부과한다. (그러나 형식기하학에서
그러한 제한이 나타나지 않는 것처럼([2], II 참조), 해석기하학에서도 그러한
제한이 나타나지 않으리라고 보는 것이 타당하다.) 한편 $X$가 체 $k$의
스펙트럼 $S$ 위의 고유 스킴이고 $Z$가 $X$ 위에서 준사영이면
$\Pi_{X/S}(Z/X)$가 존재하며, 이는 $S$ 위에서 준사영인 일련의 스킴들의
서로소 합인 스킴이다(사영인 경우 (3.1)과 같다). 실제로 $X$를 사영
$S$-스킴 $X'$로 지배하여 $X$ 자체가 사영인 경우로 환원한다. 유한 사상의
인수분해에 관한 [2], I, A, 제2절 (b)의 결과도 사용하는 이 증명의 자세한
내용은 여기서 제시하지 않겠다. 이 방법이 성공하는 까닭은 $S$가 체의
스펙트럼일 때 차우 보조정리에 등장하는 $X'$이 자동으로 $S$ 위에서
평탄하기 때문이다.

$S$에 아무 가정도 두지 않고 $X$가 $S$ 위에서 고유하고 평탄하며 $Z$가
$X$ 위에서 준사영이라고만 가정할 때에도 이 결과가 성립하는지는 알지
못한다. 응용에서 중요한 한 경우는 $Z$가 $X$의 닫힌 부분스킴인 경우이다.
이때 $\Pi_{X/S}(Z/X)$가 존재하면 이는 반드시 $S$의 닫힌 부분스킴이다.
그리고 $X$가 $S$ 위에서 사영이면 힐베르트 스킴 이론을 사용하지 않고도
이를 상당히 간단하게 직접 구성할 수 있다.
\src{276}
그 방법은 더 일반적으로 $Z$가 $X$ 위에서 아핀이면
$\Pi_{X/S}(Z/X)$가 존재하고 $S$ 위에서 아핀임을 보인다. 또한 $X$가 $S$
위에서 고유하고 평탄하면($S$ 위에서 사영일 필요는 없다), $X$ 위의 모든
국소 자명 벡터 다발 $Z$에 대하여 $\Pi_{X/S}(Z/X)$가 존재하고 $S$ 위의
벡터 다발임도 보인다. 이 결과들을 다시 다루어 통일하는 것이 바람직하다.

\section*{참고문헌}
\addcontentsline{toc}{section}{참고문헌}

\begin{description}
\item[{[1]}] BOREL (A.) et SERRE (J.-P.). -- Le théorème de Riemann--Roch,
\emph{Bull. Soc. math. France}, 제86권, 1958, 97--136쪽.

\item[{[2]}] GROTHENDIECK (Alexander). -- Technique de construction et
théorèmes d'existence en géométrie algébrique, I, II, III, \emph{Séminaire
Bourbaki}, 제13권, 1959/60, 제190호, 29쪽, 제195호, 22쪽; 그리고 제14권,
1960/61, 제212호, 20쪽.

\item[{[3]}] GROTHENDIECK (Alexander). -- Séminaire de Géométrie algébrique,
I, II, III, IV. -- Paris, Institut des hautes Études scientifiques, 1960/61
(등사판).

\item[{[4]}] GROTHENDIECK (Alexander). -- Technique de construction en
géométrie analytique, IV, V, \emph{Séminaire Cartan}, 제13권, 1960/61,
제11호 및 제12호.

\item[{[5]}] GROTHENDIECK (A.) et DIEUDONNÉ (J.). -- \emph{Éléments de
géométrie algébrique}, I : Le langage des schémas. -- Paris, Presses
universitaires de France, 1960 (Institut des hautes Études scientifiques,
Publications mathématiques, 4); II, III, IV (출간 예정).

\item[{[6]}] KODAIRA (K.). -- Characteristic linear systems of complete
continuous systems, \emph{Amer. J. of Math.}, 제78권, 1956, 716--744쪽.

\item[{[7]}] SERRE (Jean-Pierre). -- Exemples de variétés projectives en
caractéristique $p$ non relevables en caractéristique 0, \emph{Proc. Nat. Acad.
Sc. U.S.A.}, 제47권, 1961, 108--109쪽.

\item[{[8]}] ZAPPA (G.). -- Sull'esistenza, sopra le superficie algebriche, di
sistemi continui completi infiniti, la cui curva è a serie caratteristica
incompleta, \emph{Pont. Acad. Sc. Acta}, 제9권, 1945, 91--93쪽.
\end{description}

\section*{추기}
\addcontentsline{toc}{section}{추기}
\begin{center}
{\small [교정쇄를 교정하면서 추가함]}
\end{center}

이제 [2], III, 제8절의 추측들은 몫의 존재와 준사영성 양쪽에 관하여 표수
0인 체 위의 비특이 다양체에 대해서도 거짓이며, $G$가 닫힌 그래프를 갖고
작용하는 경우에도 거짓임이 드러났다.

\par\smallskip
\begingroup\footnotesize




\noindent\textit{5.7의 인쇄본은 $Z$를 $S$ 위의 준스킴이라고 부른다. 그러나 $\Pi_{X/S}(Z/X)$와 $X$ 위의 매끄러움 가정을 정의하려면 $Z\to X$가 주어져야 한다.}\par
\noindent\textit{2.2의 마지막 주장에서는 유계성과 다음 계수의 하계를 동시에 귀납해야 한다. 하나의 연접층의 몫이라는 사실만으로 그 몫들의 족이 유계인 것은 아니다. 제6절에서 일반적인 분리 사상 $X/S$에 대한 노름의 표적은 $\Gamma^n_S(X)$이며, $X/S$가 평탄할 때 대칭곱과 동일시된다. 힐베르트 스킴의 사상은 유한 평탄 족에서 별도로 구성된다.}\par
\endgroup

\par
\renewcommand{\sourceunit}{e221}
\clearpage
\chapter*{강연 221의 정오표}
\addcontentsline{toc}{chapter}{강연 221의 정오표}
\setcounter{footnote}{3}
\src{302}

\noindent\textbf{6. TDTE IV: 힐베르트 스킴.}

\begin{center}
[부르바키 세미나, 제13권, 1960/61, 제221호, 28쪽.]
\end{center}

\noindent\textbf{6쪽, 정리 2.1, 3행.} --- ``필요하다'' 대신
``필요충분하다''로 읽어야 한다.

\noindent\textbf{15쪽, 13행.} --- ``$P\neq Q$'' 대신
``$P<Q$''로 읽어야 한다.

\noindent\textbf{15행.} --- ``반대로 $P=Q$이면'' 대신
``그렇지 않은 경우에는 충분히 큰 $n$에 대하여 $P(n)\geq Q(n)$이다''로
읽어야 한다.

\noindent\textbf{18행, 공식 (*).} --- ``$<$''와 ``$>$'' 대신
``$\leq$''와 ``$\geq$''로 읽어야 한다.

\noindent\textbf{아래에서 5행.} --- ``$r$'' 대신 ``$Z$''로
읽어야 한다.

\noindent\textbf{18쪽, 주 3.9.} --- 대수적으로 닫힌 체 위의 힐베르트
스킴들의 연결 성분에 대한 연구는 하츠혼이 수행했음을 지적해 둔다. 그는
$\operatorname{Hilb}^{P}_{r}$가 연결임을 증명하고,
$\operatorname{Hilb}^{P}_{r}\neq\varnothing$이 되는 쌍 $(r,P)$를
결정한다\footnote{HARTSHORNE (R.). --- Connectedness of the Hilbert
scheme (Thesis, Harvard, 1962).}.

\noindent\textbf{23쪽, 아래에서 5행.} --- ``한 성분'' 대신
``한 성분의''로 읽어야 한다.

\noindent\textbf{23쪽과 24쪽, 주 5.5.} --- 자파의 예와 관련하여,
멈퍼드가 최근 $\mathbf P^3_{\mathbf C}$에 대한 힐베르트 스킴의 한 기약
성분을 구성하기까지 했음을 지적해 둔다. 이 성분의 일반적인 점들은 차수
14, 종수 24인 비특이 곡선들을 나타내며, 이 성분은 그 일반점에서
비축약이다. 이렇게 얻은 곡선들 가운데 하나를 중심으로
$\mathbf P^3_{\mathbf C}$를 블로업함으로써, 그는 또한 $\mathbf C$ 위의
3차원 정칙 사영 스킴을 얻는다. 그 형식적 모듈라이 스킴은
그 일반점에서 비축약이다. 동치인 표현으로, $\mathbf C$ 위의
복소해석기하학의 의미에서 그 국소 모듈라이 다양체(cf. 카르탕 세미나,
제13권, 1960/61)는 모든 점에서 비축약이다.

\par
\renewcommand{\sourceunit}{232}
\unit{강연 232 — 피카르 스킴: 존재 정리}
\sid{232}
\src{143}
\seminar{제14년차, 1961/62, 제232호\\1962년 2월}
\paperhead{대수기하학에서의 하강 기법과 존재 정리\\
V. 피카르 스킴: 존재 정리}%
  {알렉상드르 그로텐디크 지음}

\fsection{1}{상대 피카르 군과 함자}

모든 준스킴(더 일반적으로는 모든 환 달린 공간) $X$에 대하여, $X$ 위의
가역 가군(즉 $\mathcal O_X$와 국소적으로 동형인 가군)들의 동형류로
이루어진 군을 $X$의 (절대) \emph{피카르 군}이라 하고
$\operatorname{Pic}(X)$로 나타낸다. 따라서 표준 동형사상
\begin{equation}
\operatorname{Pic}(X)\xrightarrow{\sim}H^1(X,\mathcal O_X^*)
\tag{1.1}
\end{equation}
이 있다. 여기서 $\mathcal O_X^*$는 $\mathcal O_X$의 가역원층을 나타낸다.
(실제로 이는 표준 가역 가군 $\mathcal O_X$의 자기동형사상들로 이루어진
층과 동일시된다.) 분명히 $X\longmapsto\operatorname{Pic}(X)$는 $X$에 관한
반변 함자이며, 동형사상 (1.1)은 함자적이다.

$X$가 다른 준스킴 $S$ 위의 준스킴이라고 하자. $S$ 위의 준스킴들의 범주
$(\mathrm{Sch})/S$에서 $S'$이 움직일 때, 앞의 논의에 따라
\[
S'\longmapsto\operatorname{Pic}(X\times_S S')
\]
라는 반변 함자를 얻는다. 분류하려는 가역 가군들이 자명하지 않은
자기동형사상을 가지므로 이 함자는 “국소적 성격”을 갖지 않으며([5], IV,
5.4), 따라서 “표현 가능”할 가능성이 전혀 없다([4], II, A). 그러므로 이
함자를 “국소적으로 만들어야” 한다. 일반적으로 모든 상대 준스킴 $X/S$에
대하여 상대적인 군
\begin{equation}
\operatorname{Pic}'(X/S)=H^0\bigl(S,R^1f_*(\mathcal O_X^*)\bigr)
\tag{1.2}
\end{equation}
을 도입하자. 여기서 $f:X\longrightarrow S$는 구조사상이다([4], II, C 3과
비교하라). 인용한 곳에서는 이 군을 상대 피카르 군이라 부르지만, 곧
드러날 이유로 여기서는 $X/S$의 \emph{제한 상대 피카르 군}이라 부르는
편이 낫다. $S'$이 $(\mathrm{Sch})/S$에서 움직이면
\[
S'\longmapsto\operatorname{Pic}'(X\times_S S'/S')
\]
은 $S'$에 관한 반변 함자이다. 이 함자도
$\operatorname{Pic}'_{X/S}$로 나타내며, 본질적으로 다음 식으로 정의된다.
\begin{equation}
\operatorname{Pic}'_{X/S}(S')=
\operatorname{Pic}'(X\times_S S'/S').
\tag{1.3}
\end{equation}

\src{144}
필요한 조치를 취했으므로 이제 이 함자는 “국소적 성격”을 갖는다.
직관적으로 (1.3)의 우변은 매개변수 준스킴 $S'/S$로 지표화된, $X/S$의
(섬유들 위의) 가역층 동형류들의 “대수적 족”의 집합으로 해석된다. 함자
$\operatorname{Pic}'$가 표현 가능할 때, 이를 표현하는 $S$ 위의 준스킴을
$\operatorname{Pic}'_{X/S}$로 나타내고 $S$ 위의 $X$의
\emph{피카르 준스킴}이라 한다. 따라서 이 경우
\begin{equation}
\operatorname{Hom}_S(S',\operatorname{Pic}'_{X/S})
\simeq\operatorname{Pic}'_{X/S}(S')
=\operatorname{Pic}'(X\times_S S'/S')
\tag{1.4}
\end{equation}
을 갖는다.

그러나 $\operatorname{Pic}'_{X/S}$가 표현 가능하지 않은데도 상대 피카르
준스킴을 자연스럽게 정의할 수 있는 중요한 경우들이 있다. 예를 들어 체
$k$ 위에 $k$-유리점이 없는 “브라우어--세베리” 다양체가 그러하다. 이는
절대 피카르 군 $\operatorname{Pic}(X\times_S S'/S')$으로부터
$\operatorname{Pic}'$를 정의할 때 아직 충분히 국소화하지 않았기
때문이다. 정확히 말하면 $\operatorname{Pic}'$는 일반적으로 “충실 평탄
하강과 양립”하지 않는다. 이를 자세히 설명하자.

$\mathfrak M$을 \emph{충실 평탄이고 준콤팩트인} 준스킴 사상들의 모임이라
하자. 이 모임은 밑변경과 합성에 대해 안정하다. $P$를
$(\mathrm{Sch})/S$에서 집합의 범주로 가는 반변 함자라 하고, 모든
$S$-사상 $u:T'\longrightarrow T$, $u\in\mathfrak M$에 대하여 도식
\begin{equation}
P(T)\longrightarrow P(T')\rightrightarrows P(T'\times_TT')
\tag{1.5}
\end{equation}
을 생각하자. 이는 도식
\[
T\longleftarrow T'\mathrel{\substack{\xleftarrow{\,\mathrm{pr}_1\,}\\[-0.35ex]
\xleftarrow[\,\mathrm{pr}_2\,]{}}}T'\times_TT'
\]
에 $P$를 적용하여 얻는다. $P$가 표현 가능하면 하강 이론([4], I, B,
정리 2)에 따라 모든 $u\in\mathfrak M$에 대하여 도식 (1.5)가 완전하다.
이를 $P$가 $\mathfrak M$과 양립한다고 말한다. 지금의 경우에는 $P$가
“충실 평탄 하강과 양립”한다고도 하고, $(\mathrm{Sch})/S$ 위의 “준층”
$P$가 $\mathfrak M$이 주는 국소화 개념에 관한 “층”이라고도 한다. $P$가
임의의 함자일 때, 보통의 위상적 국소화에서 잘 알려진 표준 절차를 통해
“층” $\mathcal P$와 함자 준동형 $P\longrightarrow\mathcal P$를 대응시킬
수 있으며, 이 준동형은 자명한 의미에서 보편적이다. $\mathcal P$는 다음과
같이 구체적으로 계산할 수 있다. $\mathcal P(T)$를 정의하기 위하여,
$u:T'\longrightarrow T$가 $\mathfrak M$에 속하는 모든 $T$ 위의 $T'$에
대해 다음 집합을 생각한다.

\src{145}
$\overline H^0(T'/T,P)$를 $P(T')$의 원소 $\xi$ 가운데, 그 두 상
$\xi_1,\xi_2\in P(T'\times_TT')$가 다음 성질을 갖는 것들로 이루어진
부분집합이라 하자. 즉 $v\in\mathfrak M$인 사상
$v:T''\longrightarrow T'\times_TT'$가 존재하여, $\xi_1$과 $\xi_2$의
$P(T'')$에서의 상이 같아야 한다. [주의. 이와 같이 정의한 집합
$\overline H^0$는 [4], I, A, 4(a)에서 도입한 집합 $H^0(T'/T,P)$보다
크다.] $T$를 고정하고 $T'$을 움직이면(여전히 $u\in\mathfrak M$이라
하고, $T'$들의 모임에는 지배관계가 정하는 준순서를 준다)
$\overline H^0(T'/T,P)$들은 귀납계를 이루며, 다음과 같이 놓는다.
\begin{equation}
\mathcal P(T)=\varinjlim_{T'}\overline H^0(T'/T,P).
\tag{1.6}
\end{equation}
이 식의 $T$에 관한 함자 법칙은 분명하게 구체화된다.

\[
P(T)=\operatorname{Pic}(X\times_ST)
\]
일 때, (1.6)이 정의하는 $(\mathrm{Sch})/S$ 위의 반변 함자를 $S$ 위의
$X$의 \emph{상대 피카르 함자}라 하고 $\operatorname{Pic}_{X/S}$로
나타낸다. 또한 군 $\operatorname{Pic}_{X/S}(S)$를 $S$ 위의 $X$의
\emph{상대 피카르 군}이라 하고 $\operatorname{Pic}(X/S)$로 나타낸다.
그러면 분명한 전단사
\begin{equation}
\operatorname{Pic}_{X/S}(T)\simeq\operatorname{Pic}(X\times_ST/T)
\tag{1.7}
\end{equation}
를 얻는다.

따라서 $\operatorname{Pic}(X/S)$의 한 원소는 다음과 같은 자료로
정의된다. 충실 평탄이고 준콤팩트인 사상 $S'\longrightarrow S$와 군
$\operatorname{Pic}(X\times_SS')$의 원소 $\xi'$를 택하되, 충실 평탄이고
준콤팩트인 사상 $S''\longrightarrow S'\times_SS'$가 존재하여 $\xi'$의
두 역상이 $\operatorname{Pic}(X\times_SS'')$에서 같아야 한다. 방금
설명한 조건을 만족하는
\[
\xi'\in\operatorname{Pic}(X\times_SS'),\qquad
\xi_1\in\operatorname{Pic}(X\times_SS_1)
\]
이 $\operatorname{Pic}(X/S)$의 같은 원소를 정의할 필요충분조건은,
충실 평탄이고 준콤팩트인 사상
$S'_1\longrightarrow S'\times_SS_1$이 존재하여 문제의 두 원소의 상이
$\operatorname{Pic}(X\times_SS'_1)$에서 같은 것이다. 앞에서 도입한 함자
$P'=\operatorname{Pic}'_{X/S}$로 계속 작업하는 편이 흔히 편리하다.
표준 사상 $P\longrightarrow P'$이 동형사상
\begin{equation}
\mathcal P\xrightarrow{\sim}\mathcal P'
\tag{1.8}
\end{equation}
을 유도함을 곧바로 알 수 있다. 이는
$\operatorname{Pic}'_{X/S}=P'$를 이용하여
$\operatorname{Pic}_{X/S}$를 기술하는, 흔히 더 편리한 방법을 준다. 아래
2.3에 따라, 방금 설명한 $\operatorname{Pic}(X/S)$의 기술에서 $P$를
$P'$로 바꾸면 실제로

\src{146}
$S''=S'\times_SS'$ 및 $S'_1=S'\times_SS_1$로 잡을 수 있다. 적어도 인용한
곳에 명시된 조건 아래에서는 그렇다.

함자 $\operatorname{Pic}_{X/S}$가 표현 가능하면 $X/S$가 피카르 준스킴을
갖는다고 말한다. 이 함자를 표현하는 $S$ 위의 준스킴을 $S$ 위의 $X$의
\emph{피카르 준스킴}이라 하고, 여전히 $\operatorname{Pic}_{X/S}$로
나타낸다. 이를 위해서는 $P'=\operatorname{Pic}'_{X/S}$가 표현
가능하기만 하면 됨이 분명하다. 실제로 이 경우 $P'$는 이미 “층”이고,
(1.8)에 따라 $P'\longrightarrow\mathcal P'$은 표준 사상
\begin{equation}
\operatorname{Pic}'_{X/S}\longrightarrow\operatorname{Pic}_{X/S}
\tag{1.9}
\end{equation}
과 동일시되며, 이 사상은 동형사상이다. 따라서 우리의 용어는 (1.4)와
함께 앞에서 도입한 용어와 양립한다. 일반적으로
$\operatorname{Pic}_{X/S}$가 존재할 때, 이는 함자적 동형사상
\begin{equation}
\operatorname{Hom}_S(S',\operatorname{Pic}_{X/S})\xrightarrow{\sim}
\operatorname{Pic}(X\times_SS'/S')
\tag{1.10}
\end{equation}
으로 정의된다.

\fsection{2}{여러 상대 및 절대 피카르 군 사이의 관계}

\result{명제}{2.1} $f:X\longrightarrow S$가
$\mathcal O_S\xrightarrow{\sim}f_*(\mathcal O_X)$를 만족하는 사상이라고
하자. 그러면 완전열
\[
0\longrightarrow\operatorname{Pic}(S)\longrightarrow\operatorname{Pic}(X)
\longrightarrow\operatorname{Pic}'(X/S)
\]
이 있다. $X$가 $S$ 위의 단면을 가지면 마지막 사상은 전사이다. 즉
동형사상
\[
\operatorname{Pic}'(X/S)\xrightarrow{\sim}
\operatorname{Pic}(X)/\operatorname{Pic}(S)
\]
이 있다.

이 완전열은 $f$와 $\mathcal O_X^*$에 관한 르레 스펙트럴 열에 대응하는
저차 완전열로 볼 수 있다. 두 번째 주장도 마찬가지로 형식적이다.

\result{명제}{2.2} $f:X\longrightarrow S$가 준콤팩트 분리 사상이고
$\mathcal O_S\xrightarrow{\sim}f_*(\mathcal O_X)$를 만족한다고 하자. 또한
$S'\longrightarrow S$를 충실 평탄이고 준콤팩트인 사상이라 하자. 그러면
다음이 성립한다.

\noindent (i) $\operatorname{Pic}'(X/S)\longrightarrow
\operatorname{Pic}'(X\times_SS'/S')$은 단사이다.

\noindent (ii) $X$가 $S$ 위에서 국소적으로 단면을 갖는다고 하자. 즉 모든
$s\in S$가 열린 근방 $U$를 가져서 $X|U$가 $U$ 위의 단면을 갖는다고
하자. 그러면 도식
\[
\operatorname{Pic}'(X/S)\longrightarrow
\operatorname{Pic}'(X\times_SS'/S')\rightrightarrows
\operatorname{Pic}'(X\times_SS''/S'')
\]

\src{147}
은 완전하다. 여기서 $S''=S'\times_SS'$이다.

첫 번째 주장은 충실 평탄 하강의 초등적 성질과 다음의 일반적 관찰로부터
따라온다. $f:X\longrightarrow S$가
$\mathcal O_S\xrightarrow{\sim}f_*(\mathcal O_X)$를 만족하는 사상이라
하자. $S$ 위의 유한 계수 국소 자유 가군들의 범주에서 $X$ 위의 유한 계수
국소 자유 가군들의 범주로 가는 함자
\[
\mathcal F\longmapsto f^*(\mathcal F)
\]
는 충만하고 충실하다. 그 본질적 상은 $X$ 위의 가군 $\mathcal G$ 가운데
$f_*(\mathcal G)$가 국소 자유이고 표준 준동형
\[
f^*\bigl(f_*(\mathcal G)\bigr)\longrightarrow\mathcal G
\]
이 동형사상인 것들로 이루어진다. 두 번째 주장은 [4], I, B, 4에서 하강
이론으로 증명되었다.

2.2의 결과는 다음과 같이도 진술할 수 있다.

\result{따름정리}{2.3} 2.2의 조건 아래에서 표준 준동형 (1.9)
\[
\operatorname{Pic}'_{X/S}\longrightarrow\operatorname{Pic}_{X/S}
\]
은 단사이고, $X$가 $S$ 위에서 국소적으로 단면을 가지면 전단사이기까지
하다. (따라서 후자의 경우 상대 피카르 군 $\operatorname{Pic}(X/S)$과
제한 상대 피카르 군 $\operatorname{Pic}'(X/S)$은 동일하다.)

이를 2.1과 결합하면 다음을 얻는다.

\result{따름정리}{2.4} 2.2의 조건 아래에서 완전열
\[
0\longrightarrow\operatorname{Pic}(S)\longrightarrow\operatorname{Pic}(X)
\longrightarrow\operatorname{Pic}(X/S)
\]
이 있다. $X$가 $S$ 위의 단면을 가지면 마지막 준동형은 전사이다. 즉
동형사상
\[
\operatorname{Pic}(X/S)\xrightarrow{\sim}
\operatorname{Pic}(X)/\operatorname{Pic}(S)
\]
이 있다.

\result{비고}{2.5} $f:X\longrightarrow S$가
$\mathcal O_S\xrightarrow{\sim}f_*(\mathcal O_X)$를 만족하는 사상이고,
$g$가 $S$ 위의 $X$의 단면이라고 하자. $\mathcal L$을 $X$ 위의 가역
가군이라 하자. 동형사상
\[
\mathcal O_S\xrightarrow{\sim}g^*(\mathcal L)
\]
을 $\mathcal L$의 $g$-강체화라 하고, $g$-강체화가 주어진 $X$ 위의 가역
가군 $\mathcal L$을 \emph{$g$-강체화된 가역 가군}이라 하자. 이러한
구조의 모든 자기동형사상은 자명하며, $\operatorname{Pic}'(X/S)$은 $X$
위의 $g$-강체화된 가역 가군들의 동형류 집합과
동일시된다.\footnote{인쇄본에는 “$S$ 위”라고 되어 있다. 그러나 단면
$g:S\to X$에 의한 강체화는 여기서 $X$ 위의 가군 $\mathcal L$에 관한
것이다.} (바로 이 사실 덕분에 2.2(ii)를 증명할 때 하강 이론을 사용할 수
있었다.) 이는 적어도 $f$가 더 나아가 준콤팩트이고 분리인 경우
$\operatorname{Pic}(X/S)$에 대한 새로운 해석을 준다. 실제로 이 경우
2.3에 따라
\[
\operatorname{Pic}(X/S)\simeq\operatorname{Pic}'(X/S)
\]
이다.

\src{148}
\result{비고}{2.6} $f:X\longrightarrow S$를 2.2에서와 같은 사상이라
하고, $S'\longrightarrow S$를 다음 조건을 만족하는 충실 평탄이고
준콤팩트인 사상이라 하자. 즉 $S$-사상 $S'\longrightarrow X$가 존재하며,
동치로 $X'=X\times_SS'$가 $S'$ 위의 단면을 갖는다고 하자.
$S''=S'\times_SS'$, $X''=X\times_SS''$로 놓고 완전열
\[
\operatorname{Pic}(X/S)\longrightarrow\operatorname{Pic}(X'/S')
\rightrightarrows\operatorname{Pic}(X''/S'')
\]
을 생각하자. 2.4를 적용하면 완전열
\[
\operatorname{Pic}(X/S)\longrightarrow
\operatorname{Pic}(X')/\operatorname{Pic}(S')
\rightrightarrows
\operatorname{Pic}(X'')/\operatorname{Pic}(S'')
\]
을 얻는다. 특히 상대 피카르 군의 모든 원소는 이미
$\operatorname{Pic}(X')$의 어떤 원소에서 “온다”. 따라서 이는 앞
번호에서 설명한 상대 피카르 군의 기술을 크게 단순화하며, $S$ 위의 $X$의
피카르 함자에 대해서도 마찬가지다. 실제로 $S$ 위의 모든 $T$에 대하여
$X\times_ST/T$와 사상 $T'=S'\times_ST\longrightarrow T$에 앞의 논의를
적용할 수 있다. 예를 들어 $f$ 자체가 충실 평탄이면 $S'=X$로 잡을 수
있다. 더 나아가 $f$가 유한형(각각 매끄러운 등)이면, 상대 피카르 함자
\[
T\longmapsto\operatorname{Pic}(X\times_ST/T)
\]
를 기술할 때 유한형(각각 매끄러운 등)인 밑변경
$T'\longrightarrow T$만을 사용해도 된다. $f$가 사영이고 평탄하며 $S$가
국소 뇌터일 때에는, 위의 $S'\longrightarrow S$를 다음과 같이 잡을 수
있음이 증명된다. $S$를 덮는 열린집합 $S_i$들의 평탄 피복 $S'_i$를 택해
$S'$이 이들의 서로소 합이 되게 할 수 있다. $f$가 더 나아가 분리
가능하면 $S'_i$를 $S_i$ 위에서 에탈이 되게 잡을 수 있다.

\fsection{3}{주요 존재 정리: 진술}

알려진 모든 경우를 포괄하는 피카르 준스킴의 존재 명제는 추측의 형태로도
아직 없다. 말하자면 “실질적으로 필요한” 조건은
$f:X\longrightarrow S$가 고유이고(이는 본질적인 유한성 성질들을
보장한다) 평탄인 것이다. 그러나 $S$가 체 $k$ 위의 쌍대수
$k[t]/(t^2)$의 스펙트럼이고(가령 $k$를 복소수체 $\mathbf C$로 잡아도
된다) $X$가 1차원인 경우에도 이 조건들은 충분하지 않다. 이 강연을 쓰는
현재, 피카르 준스킴에 관한 가장 중요한 존재 정리들은 다음 정리에서
따라온다.

\result{정리}{3.1} $f:X\longrightarrow S$를 국소 뇌터 준스킴들의
사상이라 하자. 다음을 가정한다.

\noindent (i) $f$는 사영이다.

\src{149}
\noindent (ii) $f$는 평탄이다.

\noindent (iii) $f$의 기하적 섬유들은 정역 준스킴이다.

이 조건들 아래에서 $\operatorname{Pic}_{X/S}$가 존재한다.

뒤의 두 번호에서 개략적으로 제시할 증명은 동시에 다음도 보일 것이다.
$\xi$를 $S$에 관해 매우 풍부한 층 $\mathcal O_X(1)$에 대응하는
$\operatorname{Pic}_{X/S}$의 단면이라 하자. (즉 이 층은 사영 몰입 사상
$X\longrightarrow\mathbf P(\mathcal E)$에 의해 유도된다.) 그러면
$\operatorname{Pic}_{X/S}$의 열린 부분 $U$가 존재하여, $U$는 $S$ 위에서
준사영인 열린 부분들의 서로소 합이고 $\xi$에 의한 평행이동에 대해
안정하며, $\operatorname{Pic}_{X/S}$는 열린집합 $U-n\xi$들의 증가하는
합집합이다. (이 열린집합들은 모두 $U$와 동형이다.) 특히 3.1의 조건
아래에서 $\operatorname{Pic}_{X/S}$는 $S$ 위에서 분리되어 있다.

\result{비고}{3.2} 예들을 통해 다음을 알 수 있다. (예를 들어 $S$가 이산
값매김환의 스펙트럼이고 $X$가 $S$ 위에서 상대 차원 1인 경우를 생각할 수
있다.) 3.1에서 가정 (iii)을 생략하고, 모든 $s\in S$에 대하여 준동형
\[
k(s)\longrightarrow H^0(X_s,\mathcal O_{X_s})
\]
이 동형이라는 더 약한 가정으로 바꾸면 $\operatorname{Pic}_{X/S}$는
$S$ 위에서 분리되어 있지 않을 수도 있다. 이는 $f$의 기하적 섬유들이
축소이지만 정역인 기하적 일반섬유 하나가 특수화될 때 두 기약 성분으로
``갈라지는'' 경우에도, 기하적 섬유들이 기약이지만 정역인 기하적
일반섬유 하나가 ``다중 섬유''로 특수화되는 경우에도 일어난다. 첫 번째
경우는 예를 들어 원뿔곡선이 한 점에서 만나는 두 직선으로 퇴화할 때
나타난다. 두 번째 경우의 예는 D. 멈퍼드가 알려 주었는데, 타원곡선이
이중 타원곡선으로 퇴화하는 예이다. 이 예들은 임의의 표수에서 성립한다.

\result{비고}{3.3} 3.1의 조건 아래에서 $\operatorname{Pic}_{X/S}$가
$S$ 위에서 유한형이고 따라서 준사영인 열린 부분들의 서로소 합인지는
나는 알지 못한다.\footnote{1962년 정오표는 D. 멈퍼드가 여기서 전개한
피카르 스킴 이론을 이용하여 이 문제를 긍정적으로 해결했다고 전한다.}
힐베르트 스킴의 경우와 마찬가지로([4], IV), 힐베르트 다항식
$Q\in\mathbf Q[t]$를 생각하면 $\operatorname{Pic}_{X/S}$를 열린집합
$\operatorname{Pic}^{Q}_{X/S}$들의 서로소 합으로 분해할 수 있다. 이
열린집합들은 $S$ 위에서 준사영일 것으로 보이며, 적어도 다음 강연에서
$f$가 매끄러운 사상일 때 이를 보게 될 것이다. 가정 (i)을 ``$X$는
$S$ 위에서 국소적으로 사영이다''라는 가정으로 바꾸는 경우에는 주의해야
한다. ($\operatorname{Pic}_{X/S}$의 존재 문제는 분명히 $S$ 위에서
국소적이므로, 이 가정만으로도 3.1은 성립한다.) 반면 이 경우에는

\src{150}
$\operatorname{Pic}_{X/S}$가 $S$ 위에서 유한형이 아닌 연결 성분들을
포함하는 예를 쉽게 들 수 있다. 예를 들어 $X_0$를 대수적으로 닫힌 체
$k$ 위의 비특이 사영 대수다양체라 하고, $u$를 그 자동사상, $\xi$를
$X_0$의 네론--세베리 군의 원소라 하되 $u^n(\xi)$들이 서로 다르다고
하자. 이를테면 $X_0$로 타원곡선 $E$와 자기 자신의 곱을 취하고, $u$로는
$E\times E$의 자동사상
\[
(x,y)\longmapsto(x,y+x)
\]
을 취할 수 있다. $S$를 두 점 $a$와 $b$에서 만나는 두 비특이 기약곡선의
합집합이라 하자. $S$ 위에는 연결 $\mathbf Z$-토서 $P$가 있다.
$u$가 정하는 $\mathbf Z$의 $X_0$ 위의 작용들을 이용하면, 섬유가 $X_0$인
$S$ 위의 연관 다발을 얻는다. 이 다발은 $S-a$와 $S-b$ 위에서 자명하고,
지금 생각하는 특별한 경우에는 실제로 $S$ 위의 아벨 스킴이다. 또한
$\operatorname{Pic}_{X/S}$는 $P$와, $u$를 통해
$\operatorname{Pic}_{X_0/k}$ 위에 주어지는 $\mathbf Z$의 작용들에
연관된 다발이며, 다음과 동형인 연결 성분을 포함함을 쉽게 알 수 있다.
\[
P\times\operatorname{Pic}^{0}_{X_0/k}.
\]
(여기서 $\operatorname{Pic}^{0}$는 $\operatorname{Pic}$ 안에서 항등원의
연결 성분을 나타낸다.) 이 성분은 $S$ 위에서 유한형이 아니다. [$S$ 위에서
분리되지 않은 피카르 준스킴의 여러 경우에도 3.2에서 생각한 것과 비슷한
현상이 일어난다.]

\fsection{4}{상대 카르티에 인자와 사영 다발}

우리는 양의 인자만 사용할 것이며, 이 번호의 나머지 부분에서는 이
덧붙인 수식어를 생략한다.

$X$를 준스킴이라 하자. $X$ 위의 \emph{카르티에 인자}, 또는 간단히
$X$ 위의 인자란 가역 가군인 아이디얼 $\mathcal J$에 의해 정의되는
$X$의 닫힌 부분준스킴 $D$를 말한다. 즉 $\mathcal J$는
$\mathcal O_X$의 영인자가 아닌 한 절단에 의해 국소적으로 생성된다.
$D$에 가역 가군
\[
\mathcal L(D)=\mathcal J^{-1}
\]
을 결부한다. 그리고 표준 포함 사상
$\mathcal J\longrightarrow\mathcal O_X$는 표준 준동형
\[
s_D:\mathcal O_X\longrightarrow\mathcal J^{-1}=\mathcal L(D),
\qquad\text{즉}\quad s_D\in\Gamma(X,\mathcal L(D))
\]
을 준다. 더 나아가 인자 하나를 주는 것은 본질적으로 $X$ 위의 가역 가군
$\mathcal L$과 모든 점에서 영인자가 아닌 절단 $s$를 주는 것과 동치이다.
역으로 이러한 쌍 $(\mathcal L,s)$에는 $s$의 ``인자''를 결부하며, 이를
$\operatorname{div}(s)$로 나타낸다. $X$ 위의 주어진 가역 가군
$\mathcal L$을 정의하는 인자 $D$들의 집합은 몫집합

\src{151}
\[
\Gamma(X,\mathcal L)^*/\Gamma(X,\mathcal O_X^*)
\]
과 전단사로 대응한다. 여기서 $\Gamma(X,\mathcal L)^*$는
$\Gamma(X,\mathcal L)$의 절단들 가운데 모든 점에서 영인자가 아닌 것들로
이루어진 부분집합을 나타낸다.

이제 국소 유한형 사상 $f:X\longrightarrow S$가 주어졌다고 하고, 간단히
$S$가 국소 뇌터라고 가정하자. $\mathcal J$를 $X$ 위의 연접 아이디얼층,
$D$를 이것이 정의하는 $X$의 부분스킴이라 하고, $x\in X$ 및 $s=f(x)$라
하자. 다음 조건들은 서로 동치이다.

\noindent (i) $\mathcal J$는 $x$에서 가역이고(즉 $\mathcal J_x$는
$\mathcal O_{X,x}$의 정칙 원소 하나로 생성된다), $D$는 $x$에서
$S$ 위로 평탄이다.

\noindent (ii) $X$와 $D$는 $x$에서 $S$ 위로 평탄이고, $D_s$는 점
$x$에서 섬유 $X_s$ 위의 카르티에 인자이다.

\noindent (iii) $X$는 $x$에서 $S$ 위로 평탄이고, $\mathcal J_x$는
$X_s$ 위에 영인자가 아닌 싹을 유도하는 원소 $f_x$로 생성된다.

이때 $D$를 생각하는 점에서 $X/S$ 위의 \emph{상대 카르티에 인자}(또는
간단히 $X/S$ 위의 상대 인자)라고 한다. (i)에서 $D$가 $x$의 근방에 있는
점들에서도 상대 인자임을 알 수 있다. 따라서 $X$와 $D$가 $S$ 위에서
평탄이고 $D$가 $S$ 위에서 고유이면, $D_s$가 $X_s$ 안의 카르티에
인자인 $s\in S$들의 집합(즉 $D$가 $X_s$의 점들에서 상대 카르티에
인자인 $s$들의 집합)은 $S$의 열린 부분이다. 한편 앞의 정의는 상대
카르티에 인자라는 개념이 임의의 밑변환 $S'\longrightarrow S$ 아래에서
안정하도록 정해졌다. 이제 $X/S$ 위의 상대 인자들의 집합을
$\operatorname{Div}(X/S)$라 하고, $S$ 위에서 움직이는 $S'$에 관한 반변
함자를 다음과 같이 정의하자.
\[
\operatorname{Div}_{X/S}(S')=
\operatorname{Div}(X\times_SS'/S').
\]

$X$가 $S$ 위에서 평탄이고 고유라고 가정하자. 그러면 상대 카르티에 인자의
특성화 (ii)에 따라 $\operatorname{Div}_{X/S}$를 [4], IV에서 정의한 함자
$\operatorname{Hilb}_{X/S}$의 부분함자로 볼 수 있으며, 포함 준동형
\[
\operatorname{Div}_{X/S}\longrightarrow\operatorname{Hilb}_{X/S}
\]
은 앞의 비고들에 따라 ``열린 몰입 사상들에 의해 표현 가능''하다
([5], IV, 3.13과 비교하라). [4], IV의 주요 존재 정리를 이용하면 다음을
얻는다.

\src{152}
\result{명제}{4.1} $f:X\longrightarrow S$가 사영이고 평탄하다고 하자.
그러면 함자 $\operatorname{Div}_{X/S}$는 표현 가능하며, 더 정확히는
$\operatorname{Hilb}_{X/S}$의 열린 부분준스킴으로 표현된다.

$X/S$ 위에 주어진 매우 풍부한 층 $\mathcal O_X(1)$에 대하여,
$\operatorname{Hilb}_{X/S}$를 힐베르트 다항식
$Q\in\mathbf Q[t]$에 대응하는 열린 부분준스킴
$\operatorname{Hilb}^{Q}_{X/S}$들의 합으로 분해하는 표준 분해를 사용하면,
마찬가지로 분해
\[
\operatorname{Div}_{X/S}=
\coprod_{Q\in\mathbf Q[t]}\operatorname{Div}^{Q}_{X/S}
\]
를 얻는다. 이는 $S$ 위에서 준사영인 서로소 열린 부분준스킴들의 합이다.

한편 대응 $D\longmapsto\mathcal L(D)$를 사용하면 함자적 준동형
\begin{equation}
\operatorname{Div}_{X/S}\longrightarrow\operatorname{Pic}_{X/S}
\tag{+}
\end{equation}
을 얻는다. 이제 이를 연구하고자 한다. 상당히 일반적인 조건 아래에서 이
준동형은 상대적으로 표현 가능함이 드러날 것이다([5], IV, 3). 따라서
$\operatorname{Pic}_{X/S}(S')$의 한 원소 $\xi$에서 출발하자. 표기를
간단히 하기 위하여 $S'=S$라고 가정하고, 이에 대응하는
$\operatorname{Div}_{X/S}$의 부분함자가 표현 가능함을 보이자. 먼저
$\xi$가 $X$ 위의 가역 가군 $\mathcal L$로 정의되는 경우를 생각한다.
$X$가 $S$ 위에서 고유이고 평탄하며, $S$ 위의 $X$의 기하적 섬유들이
정수적이라고 가정하자. 그러면 ([2], III, 7절)
$\mathcal O_S\xrightarrow{\sim}f_*(\mathcal O_X)$이고, 이 관계는 임의의
밑변경 $S'\longrightarrow S$ 뒤에도 성립한다. $X/S$ 위의 상대 카르티에
인자 $D$ 가운데 $\mathcal L(D)$와 $\mathcal L$이
\[
\operatorname{Pic}(X/S)=\operatorname{Pic}_{X/S}(S)
\]
의 같은 원소를 정의하는 것들, 즉 2.4에 따라 $\mathcal L(D)$와
$\mathcal L$이 $S$ 위에서 국소적으로 동형인 것들은 몫층
$f_*(\mathcal L)^*/\mathcal O_S^*$의 단면들과 일대일로 대응한다. 이
대응은 밑변경과 양립한다. 한편 인용한 곳의 “퀸네트” 유형의 일반적
고찰([6]도 보라)에 따르면, $X/S$의 고유성과 $S$ 위에서의
$\mathcal L$의 평탄성으로부터 $S$ 위의 연접 가군 $\mathcal Q$와 층의
동형사상
\[
f_*(\mathcal L)\xrightarrow{\sim}
\operatorname{Hom}_{\mathcal O_S}(\mathcal Q,\mathcal O_S)
\]
이 존재한다. 여기서 $\mathcal Q$는 유일한 동형사상을 제외하고
정의되며, 그 형성은 밑변경과 양립한다.
\footnote{인쇄본에서는 $\operatorname{Hom}$의 아래첨자에
$\mathcal O_X$를 두었다. 그러나 $\mathcal Q$는 $S$ 위의 가군이므로
그 쌍대는 $\mathcal O_S$ 위에서 취한다.} 여기서
$f_*(\mathcal L)^*$는 $f_*(\mathcal L)$의 집합들의 부분층으로서, $U$
위의 단면들은 $f^{-1}(U)$ 위의 $\mathcal L$의 단면들 가운데
$f^{-1}(U)/U$ 위의 상대 카르티에 인자를 정의하는 것들이다. 즉 이들은

\src{153}
$X_s$ 위에 영인자가 아닌 단면을 유도한다($s\in U$). 섬유 $X_s$가
정수적이라는 가정을 사용하면, 이는 유도된 단면이 각 섬유 $X_s$ 위에서
항등적으로 영이 아니라는 뜻일 뿐이다. 또는 국소 준동형
$\mathcal Q\longrightarrow\mathcal O_S$의 말로는 이 준동형들이
전사라는 뜻이다(나카야마 보조정리). 따라서
$f_*(\mathcal L)^*/\mathcal O_S^*$의 단면들의 집합은 $\mathcal Q$의
가역 몫가군들의 집합과 일대일로 대응한다. 동치로, 연접 가군
$\mathcal Q$에 대응하는 사영 다발 $\mathbf P(\mathcal Q)$의 정의에
따라([5], V, 2), 이는 $S$ 위의 $\mathbf P(\mathcal Q)$의 단면들의
집합과 일대일로 대응한다. 이 기술은 역상 형성과 양립하므로 다음 정리를
얻는다.

\result{정리}{4.3} $f:X\longrightarrow S$를 기하적 섬유가 정수적인
고유 평탄 사상이라 하고, $S$는 국소 뇌터라고 하자. 또한 $\mathcal L$을
$X$ 위의 가역 가군이라 하자. $S$ 위의 모든 $S'$에 대하여 $T(S')$를
$X\times_SS'/S'$ 위의 상대 인자 $D$ 가운데 $\mathcal L(D)$가
$S'$ 위에서 국소적으로
$\mathcal L\otimes_{\mathcal O_S}\mathcal O_{S'}$와 동형인 것들의
집합이라 하자. 즉 $\mathcal L(D)$와
$\mathcal L\otimes_{\mathcal O_S}\mathcal O_{S'}$가
$\operatorname{Pic}(X\times_SS'/S')$의 같은 원소를 정의하도록 한다.
그러면 $S$ 위의 연접 가군 $\mathcal Q$가 존재하고 유일한 동형사상을
제외하고 결정되며, 함자 $T$는 사영 다발 $\mathbf P(\mathcal Q)$로
표현된다.

\result{따름정리}{4.4} $f$가 더 나아가 사영이라고 가정하면, 함자적
준동형
$\operatorname{Div}_{X/S}\longrightarrow\operatorname{Pic}_{X/S}$은
사영 사상들로 표현된다.

실제로 $X$가 $S$ 위의 단면을 가지면(각각 국소적으로 단면을 가지면),
4.3과 2.1에 따라 앞의 준동형은 사영 다발들로(각각 국소 사영 다발들로)
표현된다. $f$가 준사영인 경우에는 $S$ 위에서 국소적으로 $X$의 유한
평탄 준단면들을 사용하는 하강 방법으로 앞의 경우에 쉽게 환원된다.

\result{비고}{4.5} 4.3의 조건 아래에서 $S$ 위의 가군 $\mathcal Q$는
일반적으로 국소 자유가 아니다. 이는 $s\in S$가 변할 때
$\mathcal Q$의 축소 섬유의 차원, 즉 $H^0(X_s,\mathcal L_s)$의 차원이
뛸 수 있다는 사실로 나타난다. 국소 뇌터 준스킴 $S$ 위의 연접 가군
$\mathcal Q$가 주어졌다고 하자. 주어진 $s\in S$에 대하여
$\mathcal Q$가 $s$에서 자유일 필요충분조건은
$\mathbf P(\mathcal Q)$가 $s$ 위의 점들에서 $S$ 위로 평탄한 것이다.
이 경우에는 그 점들에서 $S$ 위로 매끄럽기까지 하다. 여기서는
$\mathcal Q$가 위와 같이 $\mathcal L$로 정의되므로, 이는 직상
$f_*(\mathcal L)$의 형성이 “$s$의 근방에서 밑변경과 가환한다”는 뜻과도
같다. 다시 말해

\src{154}
\[
f_*(\mathcal L)_s\longrightarrow H^0(X_s,\mathcal L_s)
\]
가 전사라는 뜻이다. 예를 들어 $H^1(X_s,\mathcal L_s)=0$이면 이렇게
된다. 문제의 준스킴들이 존재한다는 전제 아래에서, 이 판정법들은 특히
보편적 상황
$\operatorname{Div}_{X/S}\longrightarrow\operatorname{Pic}_{X/S}$에
적용된다. 그리하여 이 사상이 $\operatorname{Div}_{X/S}$의 주어진 한
점에서 매끄러울 필요충분조건, 각각 충분조건을 준다.

\fsection{5}{주요 존재 정리의 증명}

3.1의 조건 아래에서 $X/S$ 위의 매우 풍부한 가군
$\mathcal O_X(1)$을 택하고, $\xi$를 이에 대응하는
$\operatorname{Pic}(X/S)$의 원소라 하자. 줄여서
\[
P(S')=\operatorname{Pic}(X\times_SS'/S')
\]
로 놓고, 먼저 표기를 간단히 하기 위하여 $X/S$가 단면을 가진다고
가정하자. $P^+(S')$를 $X\times_SS'$ 위의 가역 가군 $\mathcal L$의
류들 가운데
\[
R^if'_*(\mathcal L(n))=0\quad\text{$i>0$이고 모든 $n\geq0$에 대하여},
\qquad
f'_*(\mathcal L(n))\neq0\quad\text{모든 $n\geq0$에 대하여}
\]
를 만족하는 것들로 이루어진 $P(S')$의 부분집합이라 하자. 이 조건들은
밑변경에 대해 안정하므로 $P$의 부분함자 $P^+$를 정의하며, 이는 분명히
$\xi$에 의한 평행이동에 대해 안정하다. 세르의 “정리 A와 B”([2], III,
2)와 일반적 결과들([5], IV, 5)을 사용하면 $P$가 표현 가능할
필요충분조건은 $P^+$가 표현 가능한 것임을 쉽게 알 수 있다. 이 경우
$P^+$는 $P$를 표현하는 준스킴 $\operatorname{Pic}_{X/S}$의 열린
부분준스킴 $U$로 표현되고, 이 준스킴은 열린 부분준스킴 $U-n\xi$들의
증가 합집합이다.

줄여서 $D=\operatorname{Div}_{X/S}$로 놓고, $D^+$를 표준 사상
$D\longrightarrow P$에 의한 $P^+$의 역상이라 하자. 따라서 사상
\begin{equation}
D^+\longrightarrow P^+
\tag{+}
\end{equation}
이 있다. 또한 $D^+$는 이미 준스킴 $D=\operatorname{Div}_{X/S}$의
열린 부분준스킴 $D^+$로 표현됨을 알고 있다. 이 준스킴의 존재는 4.1이
보장한다. 4.3으로부터 사상 (+)가 매끄럽고 사영이며 전사인 사상들로
표현됨이 쉽게 따른다. 더 정확히는, 모든 점에서 영이 아닌 국소 자유
가군에 대응하는 사영 다발들로 표현된다. 실제로
$X\times_SS'$ 위의 $\mathcal L$이 이 번호의 첫머리에서와 같으면,
$f'_*(\mathcal L)$은 영이 아닌 국소 자유 가군이고 그 형성은 밑변경과
가환한다. 4.3의 표기로 $\mathcal Q$는 $f'_*(\mathcal L)$의 쌍대와
동형이다. 이제 일반 판정법([5], IV, 4.7)을 사용하여 $P^+$의 표현
가능성을 이끌어 내겠다. 인용한 곳에서 $\mathfrak C$를 충실 평탄이고

\src{155}
준콤팩트인 준스킴 사상들의 모임으로 잡는다. 이들은 [4], I B에 따라
유효 에피사상이다. 인용한 곳의 조건 (a), 즉 (+)가
$\mathfrak C$에 속하는 사상들로 표현된다는 조건은 방금 본 바와 같이
성립한다. 마찬가지로 조건 (b), 즉 함자 $P^+$가 충실 평탄이고
준콤팩트인 하강과 양립한다는 조건도 즉시 성립한다. 이제 인용한 곳의
조건 (c)를 증명하는 일만 남는다. 즉 $\mathfrak C$-표현 가능 사상
(+)로부터 준스킴 $D^+$ 위에 유도된 동치관계 $R$이
$\mathfrak C$-유효임을, 다시 말해 $R$이 유효이고
$D^+\longrightarrow D^+/R$이 $\mathfrak C$에 속함을 증명해야 한다.
이를 위해 먼저, 여러 가상 힐베르트 다항식 $Q\in\mathbf Q[t]$에 대응하는
$D^+$의 열린 부분준스킴 $D^{+Q}$가 $R$에 대해 안정함을 주목한다.
이는 $R$의 섬유들이 연결이기 때문이다. 따라서 모든 $Q$에 대하여
$D^{+Q}$ 위에 유도된 동치관계 $R^Q$가 $\mathfrak C$-유효임을 보이는
문제로 환원된다. 그런데 이제 $D^{+Q}$는 준사영이고, 한편 동치관계
$R^Q$는 사영이고 평탄하다. 그러므로 [4], III, 6.1을 적용할 수 있는
조건 아래에 있으며, 이 결과로 원하는 결론이 따른다.

$X/S$가 반드시 단면을 가지지는 않는 일반적인 경우에는 하강 기법으로
앞의 경우에 쉽게 환원한다. 또는 $P^+$의 정의에 필요한 수정을 가하여
앞의 증명을 되풀이한다.

\result{비고}{5.1} 여기서 따른 방법은 피카르 다양체를 사영적으로 구성하기
위한 마쓰사카의 방법과 본질적으로 같다. 유효 몫으로 넘어갈 수 있음을 위해
인용한 [4], III의 결과는 힐베르트 스킴의 존재 정리에서도 쉽게 도출할 수
있다(예컨대 [6]을 보라). (고전적으로 이러한 몫은 차우 좌표를 이용하여
구성했다.) $\operatorname{Pic}_{X/S}$의 열린 부분
$\operatorname{Pic}^{+}_{X/S}$의 형성과, 고려하는 가역 가군들을 정하는
약수들의 힐베르트 다항식에 따라 이를 $S$ 위에서 준사영인 열린 부분들
$\operatorname{Pic}^{+Q}_{X/S}$로 분해하는 것은 밑변경과 양립함에
주의하자. 이 사실 덕분에 하강 기법을 적용할 수 있다.

\result{비고}{5.2} $f:X\longrightarrow S$가 고유이고 평탄이며 준동형
\[
k(s)\longrightarrow H^0(X_s,\mathcal O_{X_s})\qquad(s\in S)
\]
들이 동형사상일 때마다 $\operatorname{Pic}_{X/S}$가 존재할 가능성을
배제할 수 없다. (이 마지막 조건은 이때
$\mathcal O_S\xrightarrow{\sim}f_*(\mathcal O_X)$이고 이 관계가 임의의
밑변경 $S'\longrightarrow S$ 뒤에도 유지된다는 뜻이기도 하다.) 적어도
$f$가 더 나아가 사영인 해석공간의 경우에는, [5], IX, 3.1에서 설명한
다른 방법(내가 틀리지 않았다면 차우의 방법)으로 이 사실을
증명한다.\footnote{1962년 정오표는 이 비고의 첫머리에서 제시한 존재
추측이 거짓이지만, 멈퍼드의 방법으로 그보다 약간 약한 정리를 얻을 수
있다고 지적한다.} 그 방법에서는 약수 스킴의 한 열린 부분 위에서 고유이고
평탄인 동치관계로 몫을 취하는 대신, $X$를 $\mathbf P^r_S$ 안에
몰입시키는 몰입 사상들의 스킴에서 사영군으로 몫을 취한다.

\src{156}
이 방법은 사영군으로 몫을 취하는 것에 관한 멈퍼드의 결과 [8]을 이용하면
스킴의 경우에도 아마 적용할 수 있을 것이다. 현재로서는 $X$가 $S$ 위에서
``많은 국소 단면''을 갖는 경우에만 증명이 쓰여 있다. 예컨대 $X$가
잉여체가 대수적으로 닫힌 완비 국소환 위에서 분리 가능한 경우가
그렇다.\footnote{인쇄본에는 ``잉여체가 대수적으로 닫힌''이라는 말이
빠져 있다. 정오표는 멈퍼드의 반례 때문에 이 제한이 반드시 필요하다고
명시한다.} 원칙적으로 이 방법은 더 일반적으로 적용될 것이다. 실제로
피카르 준스킴이 분리되지 않는 경우에도 그 존재를 보이므로, 그런 경우에는
첫 번째 방법이 필연적으로 실패한다. (기술적인 어려움은 $f$의 기하적
섬유들이 더 이상 정역이 아닐 때 4.3에서 고려한 함자가 사영 다발
$\mathbf P(\mathcal Q)$ 자체로 표현되는 것이 아니라 그 안의 열린 부분으로
표현된다는 데서 생긴다. 이 때문에 평탄이지만 고유이지 않은 동치관계로
몫을 취하는 까다로운 문제가 나타난다.)

\result{비고}{5.3} 여기서 제시한 증명은 곡선 또는 곡선족의 야코비안을
미리 구성하는 일에도, 아벨 다양체나 아벨 스킴의 이론에도 의존하지 않는다.
따라서 A. 베유가 개략적으로 제시한 길을 따르는 랑의 책 [7]이나 슈발레의
강연 [1] 같은 전통적인 설명과 본질적으로 다르다. 대수적으로 닫힌 체
(가령 복소수체) 위의 비특이 곡선의 야코비안인 경우에도, 여기서 제시한
야코비안의 구성은 1번에서 정의로 삼은 매우 강한 성질들을 보장하는 것으로
알려진 유일한 구성이다. (이는 본질적으로 슈발레의 성질들이지만, 멱영원을
갖는 ``매개변수 다양체''까지 고려한 것이다.) 피카르 스킴의 구성은 아벨
다양체 이론 뒤에 와야 하는 것이 아니라 그에 앞서야 한다. 이는 일반적으로
피카르 스킴이 아벨 스킴도 아니고 아벨 스킴으로 환원되지도 않는다는
사실에서 선험적으로 분명하다. 대수적으로 닫힌 체 위의 특이 곡선에서 이미
이를 볼 수 있다. 그 경우에는 아벨 다양체가 아닌 로젠리히트의 ``일반화
야코비안''이 나타난다. 더 나아가 피카르 스킴의 일반 이론을 갖추고 나면
아벨 다양체 이론, 더 일반적으로는 아벨 스킴 이론이 훨씬 간단해진다.
특히 아벨 스킴의 쌍대성 이론, 그중에서도 카르티에형 결과들은 이로부터
거의 형식적인 귀결이 된다(예컨대 [10]을 보라).

\src{157}
\result{비고}{5.4} 특이 곡선으로 퇴화하는 곡선의 야코비안에 관한
이구사의 ``양립성 원리''는, 반드시 $S$ 위에서 매끄럽지는 않은 곡선을
섬유로 갖는 상대 스킴 $X/S$의 피카르 스킴에 대한 존재 정리로 이해해야만
제대로 이해할 수 있다. 따라서 특수화된 곡선이 정역일 때, 즉 고전적
용어로 기약이고 중복도가 1일 때 이는 주요 존재 정리 3.1의 특수한 경우다.
현재로서는 특수 곡선이 가약인 경우는 알려진 존재 정리들로 다룰 수 없다.
각 성분의 중복도가 1인 경우, 즉 특수 곡선이 잉여체 위에서 분리 가능한
경우에도 마찬가지다. 다만 잉여체가 대수적으로 닫힌 완비 이산 값매김환
위에서는 예외이며, 5.2와 비교하라. 종수 $g$인 곡선들의 모듈라이 스킴에
대한 베일리--사타케 ``콤팩트화''를 스킴 이론 안에서 기하적으로 구성할
때 이 존재 문제는 분명 제기될 것이다. (이 콤팩트화는 현재 이구사의
연구 덕분에 $g=1$인 경우에만 알려져 있다.)

\fsection{6}{상대 존재 정리}

우리는 여기서 어떤 피카르 스킴들의 존재가 다른 피카르 스킴들의 존재를
함의하는 몇 가지 유용한 경우를 개략적으로 살펴본다. 이로써 주요 정리
3.1에서 여러 다른 존재 정리를 도출할 수 있다.

\result{명제}{6.1} $f:X\longrightarrow S$를 평탄 사영 사상이라 하자. 스타인
인수분해 $f=f''f'$에서 사상 $f':X\longrightarrow S'$은 평탄이고 그 기하적
섬유들이 정역이며(따라서 3.1의 가정들을 만족한다), 유한 사상
$f'':S'\longrightarrow S$은 평탄하다고 하자. 그러면
$\operatorname{Pic}_{X/S}$가 존재하고, [4], II, C 2에서 도입한 표기를 쓰면
표준 동형사상
\[
\operatorname{Pic}_{X/S}\xrightarrow{\sim}
\Pi_{S'/S}\operatorname{Pic}_{X/S'}
\]
이 있다.

이를 보이기 위해 먼저 함자들의 동형사상
\[
\operatorname{Pic}_{X/S}\xrightarrow{\sim}
\Pi_{S'/S}\operatorname{Pic}_{X/S'}
\]
을 확립하고, $\operatorname{Pic}_{X/S'}$가 표현 가능함을 함의하는 3.1을
사용한다. 이어서 제3절에서 지적한 $\operatorname{Pic}_{X/S'}$의 구조, 즉
$S$ 위의 $\operatorname{Pic}_{X/S'}$의 한 섬유의 모든 유한 부분집합이 어떤
아핀 열린집합에 포함된다는 사실을 사용하여
$\Pi_{S'/S}\operatorname{Pic}_{X/S'}$의 존재를 얻는다.

\src{158}
예를 들어 $X$가 3.1의 조건들을 만족하는 $S$ 위의 스킴 $X_i$들의 서로소
합인 스킴이면, 명제 6.1은 자명한 명제
\[
\operatorname{Pic}_{X/S}\xrightarrow{\sim}
\prod_i\operatorname{Pic}_{X_i/S}
\]
로 환원된다.

\result{따름정리}{6.2} $f:X\longrightarrow S$를 기하적 섬유들이 국소 정역인
사영 평탄 사상이라 하자(예를 들어 사영이고 정규인 사상). 그러면
$\operatorname{Pic}_{X/S}$가 존재한다.

더욱이 이 경우 $S'$은 $S$의 에탈 피복이다. 이는 $f$가 분리 가능할 때, 즉
평탄이고 그 기하적 섬유들이 축소일 때마다 참이다. 또한
$\operatorname{Pic}_{X/S'}$의 유사한 구조 덕분에, 제3절에서
$\operatorname{Pic}_{X/S}$에 대해 진술한 구조정리가 여전히 성립함을 확인할
수 있다.

한편 하강 절차는 상대 존재정리를 준다. 그 적용 범위는 [4], I, A (c)에서
제기한 비평탄 하강 문제들의 해결에 달려 있는데, 여기서는 다음 특수한
경우만 명시하겠다.

\result{명제}{6.3} $f:X\longrightarrow S$를 고유 사상이라 하고, $X_1$과
$X_2$를 $S$ 위에서 평탄인 $X$의 두 부분준스킴이라 하자. 이들이 두 연접
아이디얼 $\mathcal I_1$과 $\mathcal I_2$로 정의되고
$\mathcal I_1\cap\mathcal I_2=(0)$이며
$\mathcal O_X/(\mathcal I_1+\mathcal I_2)$가 $S$ 위에서 평탄이라고 하자.
즉 $X$의 부분준스킴 가운데 $X_1$과 $X_2$의 상한은 $X$이고, 그 하한 $Z$는
$S$ 위에서 평탄이라고 하자.\footnote{마지막 글자는 인쇄본에서 ``$X$''로
되어 있으나, 바로 앞 조건은 $\mathcal O_X/(\mathcal I_1+\mathcal I_2)$가
$S$ 위에서 평탄하다는 것이므로 $S$ 위에서 평탄인 것은 $Z$이다.} 또한 모든
$s\in S$에 대하여 준동형사상
\[
k(s)\longrightarrow H^0(X_{i,s},\mathcal O_{X_{i,s}}),\qquad i=1,2
\]
가 전단사라고 가정하자. 그러면 함자들의 자연 준동형사상
\[
\operatorname{Pic}_{X/S}\longrightarrow
\operatorname{Pic}_{X_1/S}\times\operatorname{Pic}_{X_2/S}
\]
은 아핀 사상으로 표현 가능하다. 따라서 $\operatorname{Pic}_{X_1/S}$와
$\operatorname{Pic}_{X_2/S}$가 존재하면 $\operatorname{Pic}_{X/S}$도
존재하고, 표준 사상
\[
\operatorname{Pic}_{X/S}\longrightarrow
\operatorname{Pic}_{X_1/S}\times\operatorname{Pic}_{X_2/S}
\]
은 아핀이다.

충실 평탄 하강에 의해 $Z$가 $S$ 위의 단면을 갖는 경우로 환원할 수 있다.
이 단면은 $X$, $X_1$, $X_2$의 $S$ 위 단면들을 정하므로 2.5에서 설명한
대로 지금 고려하는 구조들에서 자기동형사상들을 제거할 수 있다. 이제 증명은
$X$ 위의 ``강체화된'' 가역 가군 $\mathcal L$을 주는 것이 삼중항

\src{159}
$(\mathcal L_1,\mathcal L_2,u)$를 주는 것과 동치임을 지적하는 데 있다.
여기서 $\mathcal L_i$는 $X_i$ 위의 ``강체화된'' 가역 가군이고, $u$는
강체화들과 양립하는 $\mathcal L_1|Z$에서 $\mathcal L_2|Z$로의
동형사상이다. $\mathcal L_1$, $\mathcal L_2$가 고정되었을 때 $u$를 주는
것이 $S$ 위에서 아핀인 적당한 $S$ 위 스킴의 한 단면으로 표현됨을 확인하기만
하면 되는데, 이는 쉽다. 6.3에서 곧바로 다음을 얻는다.

\result{따름정리}{6.4} $X$를 체 $k$ 위의 고유이고 분리 가능한 스킴이라
하고, $X_i$들을 $X$의 기약 성분들이라 하자. 모든
$\operatorname{Pic}_{X_i/k}$가 존재하면 $\operatorname{Pic}_{X/k}$도
존재하고, 표준 사상
\[
\operatorname{Pic}_{X/k}\longrightarrow\prod_i\operatorname{Pic}_{X_i/k}
\]
은 아핀이다.

이를 6.2와 결합하면, 예를 들어 $X$가 체 $k$ 위의 사영이고 분리 가능한
스킴일 때마다 $\operatorname{Pic}_{X/k}$가 존재함을 알 수 있다. $X$가 더는
$k$ 위에서 분리 가능하지 않을 때에도 Oort [11]의 논증을 사용하면 환원
결과를 얻는다. 이 방법은 임의의 밑스킴 위에서도 적용된다(예를 들어 다음
강연에서 3.3에 예고한 유한성 결과를 증명할 때 유용하다). 지나치게 기술적인
진술을 피하기 위해 기초체 위의 경우로 한정하겠다.

\result{명제}{6.5} $X$를 체 $k$ 위의 고유 스킴이라 하고, $X_0$를 $X$와
같은 바탕집합을 갖는 부분스킴이라 하자(따라서 $X$ 위의 한 멱영 아이디얼로
정의된다). 그러면 함자적 사상
\[
\operatorname{Pic}_{X/k}\longrightarrow\operatorname{Pic}_{X_0/k}
\]
은 아핀 사상으로 표현 가능하다. 특히 $\operatorname{Pic}_{X_0/k}$가
존재하면 $\operatorname{Pic}_{X/k}$도 존재하고, 사상
$\operatorname{Pic}_{X/k}\longrightarrow\operatorname{Pic}_{X_0/k}$는
아핀이다.

이를 6.4와 결합하면 곧바로 다음을 얻는다.

\result{따름정리}{6.6} $X$를 체 $k$ 위의 사영 준스킴이라 하자. 그러면
$\operatorname{Pic}_{X/k}$가 존재한다.

\result{비고}{6.6} 체 $k$ 위의 모든 고유 스킴 $X$에 대하여
$\operatorname{Pic}_{X/k}$가 존재한다는 것은 지극히 그럴듯하다.\footnote{
1962년 정오표는 Murre가 이 문제를 막 긍정적으로 해결했다고 알린다.} 앞의
결과들은 이 문제를 $k$가 대수적으로 닫혀 있고 $X$가 정역인 경우로
환원한다. 그러면 차우 보조정리에 의해 $k$ 위에서 사영인 정역 스킴 $X'$과
우세 사상 $g:X'\longrightarrow X$가 존재한다. 따라서 대응하는 함자적 사상

\src{160}
\[
\operatorname{Pic}_{X/k}\longrightarrow\operatorname{Pic}_{X'/k}
\]
이 표현 가능함을(그리고 아마도 아핀 사상으로 표현 가능함을) 보이면 충분할
것이다. 실제로 $\operatorname{Pic}_{X'/k}$가 표현 가능함은 이미 알고 있기
때문이다. 이는 현재 해결되지 않은 비평탄 하강 문제들을 제기한다. $X$가
대수적으로 닫힌 체 $k$ 위의 고유 정역 스킴일 때, 함자
$\operatorname{Pic}_{X/k}$를 축소 준스킴들에 제한하면 표현 가능한 함자를
얻는다는 점에 유의하자. $X$가 정규인 경우에는 Chevalley [1]가 이를
보였으며, 일반적인 경우에는 Seshadri [12]가 하강 방법으로 보였다. 그러나
우리 표기에서 이 저자들이 구성한 스킴은 $\operatorname{Pic}_{X/k}$가 아니라
그에 대응하는 축소 스킴
$(\operatorname{Pic}_{X/k})_{\mathrm{réd}}$이다.

\result{비고}{6.7} 더 일반적으로 $f:X'\longrightarrow X$를 $k$ 위의 고유
준스킴들의 전사 사상이라 하자. 그러면 비평탄 하강을 고찰하면
\[
\operatorname{Pic}_{X/k}\longrightarrow\operatorname{Pic}_{X'/k}
\]
가 아핀 사상일 것이라고 추측하게 된다.\footnote{1962년 정오표는 이 문제가
긍정적으로 해결되었다고 알리며 다음 강연의 주석 마지막 문단을 참조하도록
한다.} 이는 특히 항등원의 연결 성분들로 나눈 뒤 대응하는 네롱--세베리
군들의 준동형사상이 꼬임을 법으로 단사임을 함의할 것이다. $X$와 $X'$이
비특이이면 교차이론으로 이를 확인할 수 있다. 그 밖의 모든 경우에는 답이
알려져 있지 않은 듯하다.

\result{비고}{6.8} 생각과 달리 멱영원을 갖는 대수적 스킴들의 피카르
스킴을 고찰하는 일은 여러 문제에서 유용하고, 심지어 필수적이다. 예를 들어
$X$가 $k$ 위의 사영 스킴, 이를테면 매끄러운 스킴이고,\footnote{인쇄본에는
``정칙''이라고 되어 있으나 정오표는 ``$k$ 위에서 매끄러운''으로 바로잡도록
한다.} $Y$가 한 초평면 단면이면, $Y$의 모든 차수의 ``무한소 근방'' $X_n$과
피카르 스킴 $\operatorname{Pic}_{X_n/k}$를 고려해야 한다. $X$가 차원
$\geq4$(각각 $\geq3$)인 기약 스킴이면 표준 사상
\[
\operatorname{Pic}_{X/k}\longrightarrow\operatorname{Pic}_{X_n/k}
\qquad(n\text{이 충분히 클 때})
\]
은 동형사상이다(각각 네롱--세베리 꼬임 부분군들의 역상들 사이에
동형사상을 유도한다). 이 결과는 다음 강연에서 피카르 스킴을 정성적으로
연구할 때 유용할 것이다. 마찬가지로 멱영원을 갖는 몇몇 곡선들의 피카르
스킴을 고찰하고 형식기하학의 기본정리들 [3]을 사용하면, 동표수인 경우에
Mumford의 한 추측을 확인할 수 있다. 즉 차원 2인 모든 완비 정규 뇌터 국소환
$A$에 대하여, $A$의 약수류군을 잉여체 $k$ 위의 한 대수군 $G$의 $k$-유리점
집합으로 볼 수 있다.

\src{161}
더욱이 $A$ 안에서 계수체 하나를 택하면 $G$는 표준적으로 결정된다. $A$의
차원이 임의일 때에는 위의 $G$와 같은 역할을 하는 $k$ 위의 프로대수군이
존재할 법하다. $\operatorname{Spec}(A)$의 ``특이점을 해소''할 수 있는
경우, 이 군은 $k$ 위의 적당한 사영 스킴들(멱영원을 갖는)의 피카르 스킴들의
사영극한으로 구성된다.

\section*{참고문헌}
\addcontentsline{toc}{section}{참고문헌}

\begin{description}
\item[{[1]}] CHEVALLEY (Claude). -- Sur la théorie de Picard,
\emph{Amer. J. of Math.}, 제82권, 1960, 435--490쪽.

\item[{[2]}] GROTHENDIECK (A.) et DIEUDONNÉ (J.). -- Éléments de géométrie
algébrique, 제I장 이하. -- Paris, Presses universitaires de France,
1960--1961 (Institut des Hautes Études Scientifiques, Publications
mathématiques 4, 8, \ldots).

\item[{[3]}] GROTHENDIECK (Alexander). -- Géométrie formelle et géométrie
algébrique, \emph{Séminaire Bourbaki}, 제11권, 1958/59,
제182호, 28쪽.

\item[{[4]}] GROTHENDIECK (Alexander). -- Techniques de descente et théorèmes
d'existence en géométrie algébrique, I--IV, \emph{Séminaire Bourbaki}, 제12권,
1959/60, 제190호, 29쪽 및 제195호, 22쪽; 제13권, 1960/61, 제212호, 20쪽 및
제221호, 28쪽.

\item[{[5]}] GROTHENDIECK (Alexander). -- Techniques de construction en
géométrie analytique, I--X, \emph{Séminaire Cartan}, 제13권, 1960/61:
Déformation des structures analytiques complexes, 제7--17호.

\item[{[6]}] GROTHENDIECK (Alexander). -- Séminaire de géométrie algébrique,
Lichtenbaum 집필, Harvard University, 1961 (출간 예정).

\item[{[7]}] LANG (Serge). -- \emph{Abelian varieties}. -- New York,
Interscience Publishers, 1959 (Interscience Tracts in pure and applied
Mathematics, 7).

\item[{[8]}] MUMFORD (D.). -- An elementary theorem in geometric invariant
theory, \emph{Bull. Amer. math. Soc.}, 제67권, 1961, 483--486쪽.

\item[{[9]}] MUMFORD (D.). -- The topology of normal singularities of an
algebraic surface and a criterion for simplicity. -- Paris, Presses
universitaires de France, 1961 (Institut des Hautes Études Scientifiques,
Publications mathématiques, 9).

\item[{[10]}] MUMFORD (D.) and TATE (J.). -- Séminaire de géométrie
algébrique, Harvard University, 1962년 봄학기 (출간 예정).

\item[{[11]}] OORT (F.). -- Sur le schéma de Picard,
\emph{Bull. Soc. math. France}, 제90권, 1962 (출간 예정).

\item[{[12]}] SESHADRI (C. S.). -- \emph{Annali di Matematica}에 게재 예정인
박사학위논문.
\end{description}

\par\smallskip
\begingroup\footnotesize
\noindent\textit{인쇄본의 3.3 참조가 가리키는 명제는 3.2이다.}\par
\endgroup

\par
\renewcommand{\sourceunit}{e232}
\clearpage
\chapter*{강연 232의 정오표}
\addcontentsline{toc}{chapter}{강연 232의 정오표}
\setcounter{footnote}{9}
\src{303}

\noindent\textbf{7. TDTE V: 피카르 스킴. 존재 정리.}

\begin{center}
[부르바키 세미나, 제14권, 1961/62, 제232호, 19쪽.]
\end{center}

\noindent\textbf{7쪽, 비고 3.3.} --- 여기서 제기한 문제는 멈퍼드에
의해 긍정적으로 해결되었다(다음 강연에 대한 논평을 보라).

\noindent\textbf{13쪽, 비고 5.2.} --- 다음 강연의 첫머리에서
지적하듯이, 여기서 제기한 존재 추측은 거짓이다. 그러나 멈퍼드는 자신의
방법으로 이보다 약간 약한 정리를 증명하는 데 성공했다.

\noindent\textbf{14쪽, 7행.} --- ``완비'' 뒤에 ``잉여체가 대수적으로
닫힌''이라는 조건을 덧붙여야 한다. 곧 ``완비 국소환''을 ``잉여체가
대수적으로 닫힌 완비 국소환''으로 읽어야 한다. 앞서 언급한 멈퍼드의
반례는 실제로 이 제한이 필수적임을 보여 준다.

\noindent\textbf{17쪽, 비고 6.6.} --- 다음 강연의 첫머리에서
지적하듯이, 여기서 제기한 문제에 Murre가 방금 긍정적인 해답을 주었다.

\noindent\textbf{18쪽, 비고 6.7.} --- 여기서 제기한 문제는
긍정적으로 해결된다(다음 강연에 대한 논평의 마지막 문단을 보라).

\noindent\textbf{18쪽, 비고 6.8, 4행.} --- ``정칙'' 대신
``$k$ 위에서 매끄러운''으로 읽어야 한다.


\par
\renewcommand{\sourceunit}{236}
\unit{강연 236 — 피카르 스킴: 일반 성질}
\sid{236}
\src{221}
\seminar{제14년차, 1961/62, 제236호\\1962년 5월}
\paperhead{대수기하학에서의 하강 기법과 존재 정리\\
VI. 피카르 스킴: 일반 성질}%
  {알렉상드르 그로텐디크 지음}

\fsection{0}{앞 강연의 보충 ([3], V)\textsuperscript{(1)}}

[3], V에서 제기한 피카르 준스킴의 존재 문제와 관련하여 약간의 진전이
있었다.

\medskip
\noindent a. (멈퍼드). $f:X\longrightarrow S$가 사영이고 분리 가능한
사상(= 평탄하고 섬유들이 분리 가능)일 때 피카르 준스킴
$\operatorname{Pic}_{X/S}$가 존재한다는 명제는 일반적으로 참이 아니다.
이는 $f$의 섬유들이 1차원이고 $S$가 완비 이산 값매김환의 스펙트럼인
경우에도 마찬가지이다. 예를 들려면
$S=\operatorname{Spec}\mathbf R[[t]]$로 놓고, $X$를 다음 방정식으로
정의되는 $\mathbf P^2_S$의 부분스킴으로 잡으면 된다(여기서 $x,y,z$는
동차변수이다).
\[
x^2+y^2=tz^2.
\]
따라서 이 식은 기하적으로는 서로 만나는 두 직선으로 퇴화하는 이차곡선을
나타내지만, 체 $\mathbf R$ 위의 특수 섬유는 여전히 기약이다(그 섬유는
$\mathbf R$ 위의 방정식 $x^2+y^2=0$으로 주어진다). 에탈 확대
$S'\longrightarrow S$, $S'=\operatorname{Spec}\mathbf C[[t]]$를 한 뒤에는
$X'/S'$의 피카르 준스킴이 존재함을 쉽게 알 수 있으며, 이를
$\widetilde S'$의 복사본들의 서로소 합으로 명시적으로 기술할 수 있다.
여기서 $\widetilde S'$은 $S'$의 원점을 무한히 여러 번 둘로 갈라서 얻는다.
$S'\longrightarrow S$에 대한 $\operatorname{Pic}_{X'/S'}$ 위의 하강 자료는
(여기서는 $S'$의 $S$ 위 갈루아 군 $\mathbf Z/2\mathbf Z$의 작용으로
주어진다) 유효하지 않음을 쉽게 확인할 수 있다. 이 군이 둘로 갈라진 몇몇
점을 서로 교환하므로 어떤 궤도들은 아핀 열린집합에 들어 있지 않기
때문이다. 그러나 멈퍼드는 $f:X\longrightarrow S$가 사영이고 분리 가능한
사상이며, 모든 $s\in S$에 대하여 $X_s$의 기약성분들이 $k(s)$에 관하여
기하적으로 기약이면 $\operatorname{Pic}_{X/S}$가 존재함을 보일 수 있다.
그 증명은 자신의 몫 정리를 정밀화한 결과에 의존한다([7] 참조). 한편
섬유 $X_s$의 기약성분들에 아무 가정도 두지 않더라도, 아래에서 도입할
스킴 $\operatorname{Pic}^{\tau}_{X/S}$가 존재할 가능성은 여전히 남아
있음에 유의하라.

\src{222}
\noindent b. (뮈르). $X$가 한 체 위의 고유 스킴이면
$\operatorname{Pic}_{X/k}$가 존재한다. 그 증명은 슈발레의 증명 [1]을
일부 되풀이하며, 피카르 함자의 군 구조를 본질적으로 이용한다.

피카르 스킴 이론에 관한 몇 가지 보충, 특히 아벨 스킴과 관련된 내용은
[7]을 참조하는 것이 유용할 것이다. 끝으로 이 강연에는 주목할 만한
공백이 하나 있다. 사영 스킴의 피카르 스킴과 그 초평면 단면들의 피카르
스킴을 비교하게 해 주는 ``동치 판정법''이 없다는 점이다. 그러한 판정법을
전개하는 데 필요한 핵심 정리들은 [5]에 있으며, 여기에 피카르 스킴의
존재 정리들을 함께 사용해야 한다.

\fsection{1}{가환 군 준스킴의 위상적 성질}

먼저 $k$를 체, $G$를 $k$ 위의 군 준스킴이라 하자. 항등원 $e$는
$k$-유리이므로 반드시 닫혀 있고, 따라서 $G\times_kG$의 대각선은 곧바로
닫혀 있다. 그러므로 $G$는 분리되어 있다. 즉 체 위의 모든 군 준스킴은
분리되어 있다. 항등원 $e$의 연결성분을 $G^\circ$로 나타내자. $e$가
$k$-유리이므로 $G^\circ$는 실제로 기하적으로 연결되어 있다. 다시 말해
$G^\circ$를 만드는 과정은 기초체의 확대와 양립한다. 또한 $G^\circ$는
집합으로서 곱셈에 안정하다. $G$가 국소 뇌터이면 $G^\circ$가 열려
있으므로, $G^\circ$를 $G$의 열린 부분군으로 볼 수 있다. 이하에서는
$G$가 $k$ 위에서 국소 유한형이라고 가정한다. 그러면 $G^\circ$는
기하적으로 기약이고 $k$ 위에서 유한형이다. 실제로 $k$가 대수적으로
닫혀 있다고 가정해도 되고, 더 나아가 $G$가 축약이라고 가정해도 된다.
(이 경우 $G_{\mathrm{réd}}\times_kG_{\mathrm{réd}}$도 여전히 축약이므로
$G_{\mathrm{réd}}$는 $G$의 부분군이다.) 따라서 $G$는 비어 있지 않은
열린집합 위에서 $k$에 대하여 매끄럽고, 그 열린집합을 평행이동하면 모든
곳에서 매끄럽다. 이제 $G$는 국소적으로 기약이므로 그 기약성분들은
연결성분들과 일치하며, 따라서 $G^\circ$는 기약이다. $U$를 $G^\circ$에서
$e$의 아핀 근방이라 하자. $G^\circ$가 기약이라는 사실을 이용하면
$U\cdot U=G^\circ$임을 곧바로 알 수 있다. 이는 $G^\circ$가
준콤팩트이고, 따라서 $k$ 위에서 유한형임을 보인다.

간단히 하기 위해 $G$가 가환이라고 가정하자. 모든 정수 $n>0$에 대하여
$G$의 $n$제곱 준동형 $\varphi_n$에 의한 $G^\circ$의 역상을
$G^{(n)}$이라 하자. 그러면 $G^{(n)}$은 $G$의 열린 부분군이다. 다음과
같이 놓는다.
\[
G^\tau=\bigcup_nG^{(n)}.
\]

\src{223}
\[
G^\sigma=\bigcup_{(n,p)=1}G^{(n)},
\qquad
G^\rho=\bigcup_hG^{(p^h)},
\]
여기서 $p$는 체 $k$의 표수 지수이다. 이로써 $G$의 열린 부분군들을
얻으며, 이들은
\[
G^\sigma\cap G^\rho=G^\circ,
\qquad
G^\sigma G^\rho=G^\tau
\]
를 만족한다.

\result{비고}{} 몫 군 스킴 $G/G^\circ=N$을 구성할 수 있다([3], IV 참조).
그러면 $G^\tau$, $G^\sigma$, $G^\rho$를 각각 $N$의 유한 위수 원소들로
이루어진 부분군, 그 $p$-일차 성분의 자연스러운 여부분(즉 소수
$q\neq p$에 대한 $q$-일차 성분들의 합), 그리고 그 $p$-일차 성분의
$G$ 안에서의 역상으로 정의할 수 있다. 이와 관련하여 $N$은 $k$ 위의
이산 분리 가능 군 스킴임에 유의하라. 따라서 $k$의 대수적 폐포
$\overline k$를 하나 고르면(이에 따라 갈루아 군 $\pi$가 생긴다), $N$은
$\pi$가 자기동형사상으로 작용하는 보통의 이산군과 동일시된다. 이
관점에서 유한 위수 원소들의 부분군을 구성하고 이를 $q$-일차 성분들로
분해하는 것을 명백하게 해석할 수 있다. $G$가 $k$ 위의 고유 스킴 $X$의
피카르 스킴일 때, $N$을 $k$ 위의 $X$의 (축약)
\emph{네론--세베리 스킴}이라 부를 수 있다. $G^\circ_{\mathrm{réd}}$가
$G^\circ$의 군 부분스킴일 때에도 몫
$N'=G/G^\circ_{\mathrm{réd}}$를 도입하는 것이 적절하다. 이는 특히
$k$가 완전체이거나 $G^\circ$가 $k$ 위에서 고유일 때(예를 들어 $X$가
기하적으로 정규일 때) 성립한다. 이 몫은 특수화의 관점, 즉 $X$가
대수적 스킴들의 족 안에서 변할 때 $N$보다 더 잘 거동하는 경향이 있다.

이제 $S$를 국소 뇌터 준스킴, $G$를 $S$ 위에서 국소 유한형인 $S$ 위의
군 준스킴이라 하자. $G$가 $S$ 위에서 유한형이거나 분리되어 있다고
가정하지 않는다. 이때
\[
G^\circ=\bigcup_{s\in S}(G_s)^\circ
\]
로 놓고, $G$가 가환일 때에는
\[
G^\tau=\bigcup_{s\in S}(G_s)^\tau,
\qquad
G^\sigma=\bigcup_{s\in S}(G_s)^\sigma,
\qquad
G^\rho=\bigcup_{s\in S}(G_s)^\rho
\]
로 놓는다.

\src{224}
이들은 $G$의 곱셈에 안정한 부분집합들이지만, 그렇다고 이들을 $G$의 군
부분준스킴으로 정의할 수 있는 것은 물론 아니다. 특히 바탕집합이
$G^\circ$인 $G$의 군 부분준스킴은 일반적으로 존재하지 않는 듯하다.
물론 이 집합들 가운데 하나가 열려 있다면, 유도된 구조를 부여했을 때
$G$의 열린 군 부분준스킴이 된다. 우리는 $G^\tau$의 경우에는 언제나
그러함을 보일 것이다. 따라서 표현 가능 함자의 관점, 특히 ``특수화''의
관점에서는 수치적 동치가 대수적 동치보다 더 만족스럽게 거동한다. 이제
막 정의한 집합들의 주된 일반 성질을 서술하자.

\result{정리}{1.1} $G^\circ$, $G^\tau$, $G^\sigma$, $G^\rho$는 국소
구성가능이다. 더 나아가 다음이 성립한다.

\noindent (i) $G^\circ$는 $S$ 위에서 준콤팩트이다. $G_s^\circ$들이
고유이고 $G$가 $S$ 위에서 분리되어 있으면 $G^\circ$는 $S$ 위에서
고유이며, 따라서 $G$에서 닫혀 있다.

\noindent (ii) $G^\tau$는 열려 있다. $G^\circ$가 닫혀 있으면
$G^\tau$도 닫혀 있다.

\noindent (iii) $G^\circ$가 닫혀 있으면, 동표수인 경우, 즉 $S$의 모든
잉여체가 같은 표수를 갖는 경우에는 $G^\sigma$도 닫혀 있다.
$G^\circ$가 닫혀 있고 $G\longrightarrow S$가 $G^\sigma$의 점들에서
보편적으로 열려 있으면(아래의 따름정리 1.5 참조),
$G^\sigma\longrightarrow S$는 보편적으로 열린다.

\noindent (iv) $G^\circ$가 닫혀 있으면 $G^\rho$도 닫혀 있다.
동표수라고 가정하고, $(n,p)=1$인 모든 정수 $n>0$에 대하여 $G$의
$n$제곱 준동형이 열린 사상이라고 가정하자. 그러면 $G^\rho$는 열려 있다.

증명의 개요를 제시하자. $G^\circ$가 국소 구성가능이고 $S$ 위에서
준콤팩트라는 사실은 다음 보조정리에 포함되어 있다.

\result{보조정리}{1.2} $S\neq\varnothing$이라고 가정하자. 그러면
$S$ 안에 비어 있지 않은 열린집합 $U$가 있고, $U$ 위에 연결인 섬유들을
갖는 유한형 군 스킴 $H$가 있으며, 상이 $G^\circ|U$인 열린 몰입 사상인
군 스킴 준동형 $H\longrightarrow G|U$가 있다.

보조정리를 증명하기 위해 $S$가 기약이라고 가정해도 된다. 그 일반점을
$\eta$라 하고 밑변경
$S'=\operatorname{Spec}(\mathcal O_{S,\eta})\longrightarrow S$를 하자.
그러면 국소 아르틴 환 $\mathcal O_{S,\eta}=A$ 위의 군 스킴 $G'$을 얻고,
앞서 말한 바와 같이 그 안에는 $A$ 위에서 유한형인 열린 군 부분스킴
$G'{}^\circ$가 있다. 따라서 이는 $\eta$의 어떤 열린 근방 $U$ 위의
유한형 군 스킴 $H$에서 유도되며, 표준 몰입 사상
$G'{}^\circ\longrightarrow G'$은 몰입 사상

\src{225}
$H\longrightarrow G|U$에서 유도된다. $U$를 충분히 작게 잡으면 이
몰입 사상은 열려 있고 군 스킴 준동형이다. 또한 $U$를 충분히 작게 잡으면
$H$의 섬유들은 연결되어 있고 모두 같은 차원을 가지며, 그 차원은 $G$의
섬유들의 차원과 같다. 따라서 모든 $s\in U$에 대하여 $G_s$ 안에서
$H_s$의 상은 정확히 $G_s^\circ$이다(필요하면 $U$를 다시 충분히 작게
잡는다). 이로써 1.2가 증명된다.

(i)의 두 번째 명제는 다음 보조정리에 포함되어 있다(이를 $G$ 안에서
$G^\circ$의 준콤팩트 열린 근방에 적용한다).

\result{보조정리}{1.3} $S$가 국소 뇌터이고, $X$가 $S$ 위에서 유한형이고
분리된 준스킴이며, $g$가 $S$ 위의 $X$의 단면이라고 하자. $X^\circ$를
각 $X_s$에서 $g(s)$를 포함하는 연결성분들의 합집합이라 하자. 어떤
$s\in S$에 대하여 $X_s^\circ$가 $k(s)$ 위에서 고유이면, $s$의 열린
근방 $U$가 있어서 $X^\circ|U$는 $U$ 위에서 고유이고, 특히 $X|U$에서
닫혀 있다.

밑공간의 충실 평탄 하강에 의해 $S$가 완비 국소환의 스펙트럼이고 $s$가
그 닫힌점인 경우로 환원된다. [2], III, 5.5.1을 적용하면 $X$가 서로소인
두 열린집합 $X'$과 $X''$의 합으로 분해됨을 알 수 있다. 여기서 첫째 것은
$S$ 위에서 고유이고 $X'_s=X_s^\circ$를 만족한다. 그러므로 $X=X'$, 즉
$X$가 $S$ 위에서 고유인 경우로 환원된다. 이 경우에는 부분집합의 고유성에
대한 값매김 판정법([2] 제II장에서 빠져 있다)을 이용하는 표준적인 증명으로
결론을 얻는다.

$G^\tau$가 열려 있음을 증명하자. $G^\tau$를 만드는 과정($G^\circ$,
$G^\sigma$, $G^\rho$의 경우도 마찬가지이다)이 밑변경과 양립함을 감안하면,
이는 모든 $S$ 위의 $G$의 단면 $g$에 대하여 $g^{-1}(G^\tau)$가 열려
있음을 보이는 것과 같다. 이는 다음 두 명제를 뜻한다.

\medskip
\noindent a. $y\in S$이고 $g(y)\in G_y^\tau$라 하자. 그러면
$y$에 충분히 가까운 모든 $y'\in\overline{\{y\}}$에 대하여
$g(y')\in G_{y'}^\tau$이다.

\medskip
\noindent b. $y'\in S$이고 $g(y')\in G_{y'}^\tau$라 하자. 그러면
$y'$의 모든 일반화 $y$에 대하여 $g(y)\in G_y^\tau$이다.

(a)를 위해 $g^n(y)\in G_y^\circ$가 되는 정수 $n>0$이 존재함에 유의하자.
$G^\circ$가 구성가능이므로 $(g^n)^{-1}(G^\circ)$도 구성가능하다. 따라서
$y$에 충분히 가까운 $y'\in\overline{\{y\}}$에 대하여
$g^n(y')\in G_{y'}^\circ$이다. (b)를 위해서는 $g(y')^n=g^n(y')$들이
$G_{y'}$의 한 준콤팩트 열린집합 안에 머묾에 유의하자. 실제로 이들은
$G_{y'}^\circ$를 법으로 한 유한 개의 잉여류 안에 들어 있다. 그러므로
$g^n(y')$들을 포함하는 $G$의 준콤팩트 열린집합 $U$가 있고, 이들의
일반화 $g^n(y)$들도 $U$에 들어 있다. 따라서 $g(y)$의 거듭제곱들은
$G_y$의 한 준콤팩트 열린집합 안에 머물며, 이로부터
$g(y)\in G_y^\tau$임을 쉽게 얻는다.

\src{226}
$G^\circ$가 닫혀 있다고 가정하고 $G^\tau$도 닫혀 있음을 증명하자.
$G^\tau$가 열려 있고 따라서 국소 구성가능임은 이미 알고 있으므로,
특수화에 안정함만 보이면 된다. 이는 $G^\tau$가 닫힌집합들의 합집합,
곧 $n$제곱 준동형 $\varphi_n$에 의한 닫힌집합 $G^\circ$의 역상들의
합집합이라는 사실에서 따른다.

$G^\circ$가 닫혀 있으면 같은 논증으로 $G^\sigma$와 $G^\rho$도 닫혀
있음을 보일 수 있다(첫째 경우에는 동표수라는 조건을 둔다). 그러려면
$G^\sigma$와 $G^\rho$가 국소 구성가능임을 먼저 보이면 된다.
$x\in G^\tau$에 대하여, $n$제곱 준동형 $\varphi_n$이 $x$를
$G^\circ$로 보내는 가장 작은 정수 $n>0$을 $\nu(x)$라 하자. 그러면
$G^\sigma$(각각 $G^\rho$)는 $\nu(x)$가 $p$와 서로소인(각각 $p$의
거듭제곱인) $x\in G$\footnote{인쇄본에는 여기에 ``$X$''라고 되어 있으나,
이 대목에서 $X$는 정의되어 있지 않다. 지금 살피는 네 장소는 모두
$G$의 부분집합이다.}로 이루어진다. 이제 구성가능성에 관한 명제는 다음의
더 정밀한 명제에서 따른다.

\result{보조정리}{1.4} $G^\tau$ 위의 함수 $\nu$는 국소 구성가능이다.

실제로 이는 모든 정수 $n>0$에 대하여 $\nu(x)=n$인 $x\in G^\tau$들의
집합이 국소 구성가능이라는 뜻이다. 그런데 이 집합은
$\varphi_n^{-1}(G^\circ)$와, $d$가 $n$의 진약수들을 움직일 때의
$\varphi_d^{-1}(G^\circ)$들의 합집합의 차집합이다. $G^\circ$가 국소
구성가능이므로 $\varphi_d^{-1}(G^\circ)$들도 국소 구성가능이고, 따라서
앞의 차집합도 국소 구성가능이다.

$G^\circ$가 닫혀 있고 $G\longrightarrow S$가 $G^\sigma$의 점들에서
보편적으로 열려 있다고 가정하자. $G^\sigma\longrightarrow S$가 열린
사상임을, 즉 $G^\sigma$ 안에서 $x\in G^\sigma$의 한 근방을
$y=f(x)$의 한 근방으로 보냄을 증명하자. $G^\sigma$가 국소
구성가능이므로, $y$의 모든 일반화 $y'$에 대하여 $y'$ 위에 놓이는
$G^\sigma$ 안의 $x$의 일반화 $x'$이 존재함을 보이면 된다. 밑변경을
하면 $S$가 이산 값매김환의 스펙트럼이고 $y,y'$이 각각 닫힌점과
일반점인 경우로 환원된다. $G\longrightarrow S$가 $x$에서 열려 있으므로
$y'$ 위에 놓이는 $G$ 안의 $x$의 일반화 $x_1$이 존재한다. 다시 밑변경을
할 수 있으므로, $x=g(y)$를 만족하는 $S$ 위의 $G$의 단면 $g$가 있다고
더 가정해도 된다. $k(y')$의 표수가 0이면 $y'$ 위에 놓이는 $G$ 안의
$x$의 아무 일반화 $x'$이나 택하면 된다. $G^\tau$가 열려 있으므로
$x'\in G^\tau$이고, $G_{y'}^\sigma=G_{y'}^\tau$이므로
$x'\in G^\sigma$이다. $k(y')$의 표수가 $p>0$이면
\[
\nu(g(y'))=p^hm\qquad\text{이며 }(m,p)=1
\]
로 놓는다.

\src{227}
$ap^h+bm=1$을 만족하는 두 정수 $a,b$를 잡고
\[
g_1=g^{ap^h},\qquad g_2=g^{bm},\qquad\text{따라서}\qquad g=g_1g_2
\]
로 놓자. 그러면 정의에 따라 $g_1(y')\in G^\sigma$이고
$g_2(y')\in G^\rho$이다. $G^\circ$가 닫혀 있으므로
$g_1(y)\in G^\sigma$이고 $g_2(y)\in G^\rho$이다. 그런데
\[
g(y)=g_1(y)g_2(y)\in G^\sigma
\]
이므로 $g_2(y)\in G^\sigma$이기도 하다. 따라서
\[
g_2(y)\in G_y^\sigma\cap G_y^\rho=G_y^\circ.
\]
가정과 $S$가 이산 값매김환의 스펙트럼이라는 사실에서
$G-(G_y-G_y^\circ)$는 $G$의 열린집합이고, 그 위에서
$G\longrightarrow S$가 열린 사상을 유도함을 얻는다. 따라서
$G_y^\circ$의 모든 점을 지나는 ``준단면''이 존재한다. 그러므로 필요하면
밑공간 $S$를 다시 확대하여, $g'_2(y)=g_2(y)$를 만족하는 $S$ 위의
$G^\circ$의 단면 $g'_2$가 있다고 가정해도 된다. $g'=g_1g'_2$로 놓자.
그러면 정의에 따라 $g'(y)=g(y)=x$이고
\[
g'(y')=g_1(y')g'_2(y')\in G^\sigma
\]
이다. 따라서 $g'(y')=x'$\footnote{인쇄본에는 여기에 ``$g(y')=x'$''라고
되어 있으나, 앞 문장에서 구성하고 지금 일반 섬유를 계산하는 단면은
$g'$이다.}은 $y'$ 위에 놓이는 $G^\sigma$ 안의 $x$의 일반화이다. 이로써
(iii)이 증명된다.

끝으로 (iv)의 마지막 명제를 증명하자. $x\in G^\rho$이면 $x$의 모든
일반화 $x'$이 $G^\rho$에 속함을 보이면 된다. 적당한 $h$에 대하여
$\varphi_{p^h}$로 상을 취하면\footnote{인쇄본에는 $\varphi_{n^h}$라고
되어 있다. $x\in G^\rho$이면 p.~223의 정의에 따라 표수 지수 $p$의 한
거듭제곱이 $x$를 $G^\circ$로 보낸다.} $x\in G^\circ$이라고 더 가정해도
된다. 이제 표수와 서로소인 모든 정수 $n$에 대하여 $x$는
$\varphi_n$의 상에 속한다(표수가 $n$과 서로소인 체 위의 연결 유한형
군에서 $n$제곱 사상은 전사이기 때문이다). $\varphi_n$이 열려 있으므로
$x'$도 $\varphi_n$의 상에 속한다. 더 정확히 말해, $G_y^\circ$를 포함하는
$G^\tau$의 준콤팩트 열린집합 $U$를 잡자. 그러면 $p$와 서로소인 모든
$n$에 대하여 $x'\in\varphi_n(U)$이다. $z\in U$에 대한 $\nu(z)$의
$p$와 서로소인 인수들의 공배수를 $n$으로 택하면 $x'\in G^\rho$를
얻는다.

정리 (1.1), (iii)는 다음과 같이 보충된다.

\result{따름정리}{1.5} $n>1$을 $n$제곱 준동형
$\varphi_n:G\longrightarrow G$가 보편적으로 열린(예를 들어 에탈인)
정수라 하고
\[
G^{(n)}=\varphi_n^{-1}(G^\circ)
\]
로 놓자. 모든 기하적 섬유 $G_{\bar s}$에서 부분집합

\src{228}
$\bigcup_{h\geq0}|{}_{n^h}G_{\bar s}|$가
$\bigcup_{h\geq0}G_{\bar s}^{(n^h)}$에서 조밀하다고 가정하자. 그러면
$G\longrightarrow S$는 $G^{(n^h)}$들의 점들에서 보편적으로 열려 있다.
특히 $G_s^\circ$들이 가법 성분을 포함하지 않고, 모든 정수 $n>1$에
대하여 $\varphi_n$이 잉여표수가 $n$과 서로소인 점들에서 보편적으로
열려 있다고 가정하자\footnote{인쇄본에는 ``$p$와 서로소''라고 되어 있다.
$n$으로 양화된 이 명제에서 에탈성 판정법과 핵 ${}_nG$를 이용하는 논증은
잉여표수가 $n$과 서로소일 것을 요구한다.}. 그러면
$G\longrightarrow S$는 $G^\sigma$의 점들에서 보편적으로 열려 있다.
따라서 (1.1), (iii)에 의해, $G^\circ$가 닫혀 있으면
$G^\sigma\longrightarrow S$는 보편적으로 열린다. 이 경우들에는
$s\longmapsto\dim G_s$가 $S$ 위의 국소 상수 함수이다.

실제로 가정에 따라 $\varphi_n$의 핵 ${}_nG$는 $S$ 위에서 보편적으로
열려 있다. 따라서 $G\longrightarrow S$는 ${}_nG$의 점들에서 보편적으로
열려 있고, $n$을 $n^h$으로 바꾸면 ${}_{n^h}G$의 점들에서도 보편적으로
열려 있다. 조밀성 가정은 섬유마다 위수가 $n$의 거듭제곱인 점들이
$G^{(n^h)}$들의 합집합에서 조밀하다는 것을 정확히 뜻한다. 특별한 명제의
경우, 가법 성분이 없다는 사실은 잉여체에서 가역인 정수 $n$들에 대하여
이 조밀성을 함의한다. 이로부터 $G\longrightarrow S$가 이 합집합의 모든
점에서, 특히 항등원 단면을 따라 보편적으로 열려 있음을 쉽게 얻으며,
따라서 $s\longmapsto\dim G_s$가 국소 상수 함수임을 쉽게 얻는다.

\result{비고}{1.6} 사상 $X\longrightarrow S$가 모든 열린집합을
열린집합으로 보내고, 임의의 밑변경 $S'\longrightarrow S$ 뒤에도 이 훌륭한
성질을 유지할 때 이 사상을 보편적으로 열려 있다고 한다는 것을 상기하자. 이는
또한 ($S$가 국소 뇌터이고 $X\longrightarrow S$가 국소 유한형일 때)
$X$의 모든 기약 성분이 $S$를 지배하고, 이 성질이 임의의 밑변경
$S'\longrightarrow S$ 뒤에도 보존됨을 뜻한다. 더구나 두 진술 모두에서
(위의 유한성 가정 아래) $S'$이 이산 값매김환의 스펙트럼인 밑변경
$S'\longrightarrow S$만으로 검사해도 충분하다(원한다면 이 환을 완비이고
잉여체가 대수적으로 닫힌 것으로 취할 수 있다~...). 이 정의는 $X$의 부분집합
$Z$에도 자명하게 확장된다(예를 들면 $G$의 부분집합 $G^\sigma$). 심지어
기하적 섬유가 연결인 매끄러운 사영 사상 $X\longrightarrow S$에서 시작해도
(예컨대 $S=\operatorname{Spec}\mathbf C[j]$ 위의 타원곡선 모듈러 족의
섬유제곱), $\operatorname{Pic}_{X/S}$가 $S$ 위에서 보편적으로 열려 있지
않는 것은 지극히 정상적인 현상이다. 즉
$\operatorname{Pic}_{X/S}$의 어떤 기약 성분들이 $S$의 한 점 위에만 완전히
놓일 수 있다. 이는 $X/S$의 섬유들의 네론--세베리 군의 계수(階數)가
위쪽으로 도약할 수 있다는 사실(``복소곱셈'' 현상)과 관련된다. 반면 (1.5)는
좋은 경우 $\operatorname{Pic}_{X/S}^\sigma$가(더 나아가 흔히
$\operatorname{Pic}_{X/S}^\tau$까지도 그런 듯하다) $S$ 위에서
보편적으로 열려 있음을 보장한다.

\src{229}
마지막으로 예외적으로 $G^\circ$가 열린집합이 되어 주는 유용한 경우가
다음과 같다.

\result{따름정리}{1.7} $G\longrightarrow S$가 $G^\circ$의 점들에서
보편적으로 열려 있고(cf. (1.5)), 섬유 $G_s$가 분리 가능하여 $k(s)$ 위에서
매끄럽다고 가정하자(카르티에의 한 결과에 따라 뒤의 조건은 잉여 표수가 0이면
자동으로 충족된다). 그러면 $G^\circ$는 $G$에서 열린집합이다. 더 나아가
$S$가 축소이면 $G^\circ$는 $S$ 위에서 매끄럽다(특히 평탄하다).

첫 번째 진술은 다음과 같이 더 정밀하게 할 수 있다. 주어진 $s\in S$에 대하여
$G\longrightarrow S$가 $G_s^\circ$의 점들에서 보편적으로 열려 있고
$G_s^\circ$가 $k(s)$ 위에서 분리 가능하면, $G^\circ$는
$G_s^\circ$의 근방이다(더구나 $s$에서 한 가정은 이웃한 점들에서도 계속
성립한다). 이 진술은 어떤 군 구조에도 의존하지 않으며 [2], IV, 제7절에서
증명될 것이다. (1.7)의 마지막 진술도 어떤 군 구조나 연결 성분에 관한
문제와 무관하며, [2], IV, 제5절에 주어진 평탄성 판정법의 특수한 경우이다.
그 판정법은 더 일반적으로 다음을 함의한다.

\result{따름정리}{1.8} $U$를 섬유 $G_s$의 열린 부분집합이라 하고,
$G\longrightarrow S$가 $U$의 점들에서 보편적으로 열려 있다고 하자(cf.
(1.5)). $G_s$가 $k(s)$ 위에서 분리 가능하고 $S$가 $s$에서 축소이면,
$G$는 $U$의 점들에서 $S$ 위에 평탄하다(따라서 이 경우 그 점들에서
$S$ 위에 매끄럽다).

\result{비고들}{1.9} $G$가 $S$ 위에서 분리되어 있을 때 $G^\circ$나
$G^\tau$가 언제나 $G$에서 닫혀 있는지는 알지 못한다. 그럴 가능성은 작아
보인다. 어쨌든 분리 가정을 버리면 자명한 반례들이 있다. 예컨대 아핀 직선의
원점을 가산 무한 번 중복시키면, 모든 섬유가 항등군 하나로 축약되되 한 섬유만
$\mathbf Z$인 아핀 직선 위의 군 스킴 $G$를 얻는다. 이 군 스킴은 두 직선이
한 점에서 만나도록 퇴화하는 원뿔곡선족에 대응하는 $S$-스킴 $X$의 피카르
준스킴 안에서, 항등 단면의 폐포인 열린 부분군 스킴이다.

더구나 이산 값매김환 $V$의 스펙트럼 $S$ 위의 유한 평탄 군 스킴 $G$에서
출발하여 $G^\circ$가 항등 단면 하나로 축약되고 따라서 닫혀 있는 경우에도,
몇몇 가정을 버리면 여러 결론이 거짓이 된다. $V$의 동일 표수가 $p>0$이라고
하고, 함자들의 준동형
\[
f\longmapsto f^p-tf
\]
로 정의되는 준동형 $\mathbf G_a\longrightarrow\mathbf G_a$의 핵을 $G$라
하자. 여기서 $t$는 $V$의 균등화원이다.

\src{230}
(정의에 따라 $S$ 위의 ``가법군'' $\mathbf G_a$는 함자
$S'\longmapsto\Gamma(S',\mathcal O_{S'})$를 표현함을 상기하라.) 그러면
$S$-스킴으로서
\[
G=\operatorname{Spec}V[x]/(x^p-tx)
\]
는 한 점에서 만나는 $p$개의 ``직선''의 합집합이며, 가법형 무한소군으로
퇴화하는 군 $\mathbf Z/p\mathbf Z$를 나타낸다. 이 예에서는
$G^\circ=G^\sigma$이고(이는 $p$개 직선 가운데 하나이다), 따라서
$G$에서 열려 있지 않다. 이제 잉여 표수가 $p>0$인 혼합 표수의 경우에는
$\mathbf G_m$의 $p$제곱 사상의 핵인 군 스킴 $H=\symbf{\mu}_p$에서
출발하자. 따라서
\[
H=\operatorname{Spec}V[x]/(x^p-1) ;
\]
이것 역시 한 점에서 만나는 $p$개의 직선의 합집합이며, 표수 0에서의 군
$\mathbf Z/p\mathbf Z$가 표수 $p$에서 곱셈형 무한소군으로 퇴화하는 모습을
나타낸다. 통상적인 유한군 $\mathbf Z/p\mathbf Z$가 정의하는 ``상수 군
스킴''을 $H'$이라 하자. 이는 $\mathbf Z/p\mathbf Z$개인 $S$의 서로소 합,
달리 말해 $\operatorname{Spec}V^{\mathbf Z/p\mathbf Z}$이다. 그러면
$G=H\times_S H'$은 표수 0에서 군 $(\mathbf Z/p\mathbf Z)^2$을 나타내고,
표수 $p$에서는 무한소군과 $\mathbf Z/p\mathbf Z$의 곱으로 퇴화한다.
여기서 $G^\rho$는 항등 단면과 특수 섬유의 합집합이므로 열려 있지 않은
반면, $G^\sigma$는 항등 단면과 일반 섬유의 합집합이므로 닫혀 있지 않다.
이는 (1.1), (iii), (iv)의 동일 표수 경우에 한 주장과 대조된다. 물론 이는
표수 $p>0$과 관련된 현상이다. 실제로 앞의 결과들은 다음을 준다.

\result{따름정리}{1.10} $S$의 모든 잉여 표수가 0이라고 가정하자. 따라서
$G^\tau=G^\sigma$이고 $G^\circ=G^\rho$이다. 그러면
$G^\tau=G^\sigma$는 열려 있다. $G$가 $S$ 위에서 분리되어 있고
$G_s^\circ$들이 고유이면 이는 더 나아가 열리고 닫혀 있다. 같은 가정 아래
$G^\circ=G^\rho$는 $S$ 위에서 고유이므로 $G$에서 닫혀 있다. 마지막으로
모든 정수 $n>1$에 대하여 $G$의 $n$제곱 준동형이 보편적으로 열려 있다고
가정하면 $G^\circ$는 열려 있다. 더 나아가 $G_s^\circ$들이 가법 성분을
갖지 않으면(예컨대 위와 같이 $G_s^\circ$들이 고유이면),
$G^\tau\longrightarrow S$는 보편적으로 열려 있으며, $S$가 축소이면 더
나아가 매끄럽다.

마지막으로 다음의 쉬운 결과를 언급하자.

\result{명제}{1.11} $S$의 열린 부분집합 $U$가 존재하여, $G$가 $S$ 위에서
매끄러운(각각 평탄한) $G$의 점들의 집합 $G'$은 $G|U$에서 유도된 열린
부분군 스킴의 바탕 집합이다. 더 나아가 $U$ 위의 $G$의 모든 단면은 $U$
위의 $G'$의 단면이다.

\src{231}
\result{따름정리}{1.12} $G$가 항등 단면의 점들에서 $S$ 위에
매끄럽다면(각각 평탄하다면), $S$ 위의 $G$의 모든 단면의 점들과
$G^\circ$의 모든 점에서도 그러하다. 더 나아가 모든 정수 $n>0$에 대하여
$n$제곱 준동형 $\varphi_n:G\longrightarrow G$가 표수가 $n$과 서로소인
점들에서 에탈이면, $G$는 $G^\sigma$의 모든 점에서 $S$ 위에
매끄럽다(각각 평탄하다).

\fsection{2}{피카르 스킴의 국소적 성질에 대한 응용}

\result{정리}{2.1}

\noindent(i) $f:X\longrightarrow S$를 고유하고 매끄러운 사상이라 하고,
$\operatorname{Pic}_{X/S}$가 존재한다고 가정하자(예를 들어 $f$가 사영
사상인 경우). 그러면 $\operatorname{Pic}_{X/S}$는 $S$ 위에서 분리되어
있고, $S$ 위에서 유한형인 $\operatorname{Pic}_{X/S}$의 모든 닫힌 부분집합
$Z$는 $S$ 위에서 고유이다.

\noindent(ii) $X$를 체 $k$ 위에서 고유이고 기하적으로 정규인 준스킴이라
하자. 그러면 $\operatorname{Pic}_{X/k}^\circ$\footnote{인쇄본에는
$\operatorname{Pic}_{X/S}^\circ$라고 되어 있으나, 체 $k$ 위의 이 명제에는
$S$가 나타나지 않는다.}는 $k$ 위에서 고유이다.

\begin{proof}
(i) 값매김 판정법([2], II, 7절)에 따라 다음을 증명하면 충분하다. $S$를
완비 이산 값매김환의 스펙트럼이라 하고, $U$를 $S$의 일반점 하나로 이루어진
열린 부분집합이라 하자. 그러면 $S$ 위의 $\operatorname{Pic}_{X/S}$의 모든
유리 단면, 즉 $U$ 위의 모든 단면은 $S$ 위의 단면으로 유일하게 연장된다.
$\operatorname{Pic}_{X/S}$의 정의를 고려하면 이는 다음 명제와 동치이다.
$V=f^{-1}(U)$ 위의 모든 가역 가군 $\mathcal L$에 대하여, $\mathcal L$을
연장하는 $X$ 위의 가역 가군이 존재하고 동형사상을 제외하고 유일하다. 그런데
$V$ 위와 $X$ 위의 가역 가군을 각각 ``카르티에'' 인자류로 기술하면 이 사실을
쉽게 얻는다. 실제로 $X$의 국소환들은 정칙이다($X$가 정칙인 $S$ 위에서
매끄럽기 때문이다). 따라서 아우슬란더의 결과에 의해 이 국소환들은
유일인수분해정역이고, 이에 따라 $X$ 위의 모든 인자\footnote{인쇄본에는
여기서 ``$S$ 위''라고 되어 있으나, 사용되는 유일인수분해성은 $X$의
국소환들의 성질이다.}는 카르티에 인자이다. 실제로 $V$ 위의 모든 인자는 그
``폐포''를 취함으로써 $X$ 위의 인자로 연장할 수 있다.
\end{proof}

\result{비고}{2.2} $f$가 평탄하고 그 섬유 $X_s$들이 국소 완전교차이며
여차원 $\leq2$에서 $k(s)$ 위로 매끄럽다고만 가정해도 위의 증명은 유효하다.
이는 [5]에서 증명된 다음 결과에 따른다. 여차원 $\leq3$에서 정칙인 뇌터
국소 완전교차환은 유일인수분해정역이다(``사무엘의 추측''). 다음 사실에
유의하자.

\src{232}
``여차원 $\leq2$''를 ``여차원 $\leq1$'', 즉 ``정규''라는 가정으로
바꾸면 이 결과는 거짓이 된다. 비특이 이차 초곡면들의 족이 이차 원뿔로
퇴화하는 예에서 이를 확인할 수 있다.

(ii) 차우 보조정리를 사용하면 $X$가 사영인 경우로 환원된다. 따라서
$X\subset\mathbf P_k^n$이라 할 수 있고, $X$가 연결이라고 가정해도 된다.
$\dim X=1$이면 $X$는 $k$ 위에서 매끄러우므로 (i)을 적용한다.
$\dim X\geq2$이면 알려진 ``동치 판정법''의 한 변형을 사용할 수 있다. 이에
따르면 $\mathbf P_k^n$의 적절한 선형 부분공간들과 $X$의 교차로 얻어지는,
$X$ 위의 매끄러운 곡선 $Y_i$들이 유한 개 존재하여
\[
\operatorname{Pic}_{X/k}^\tau\longrightarrow
\prod_i\operatorname{Pic}_{Y_i/k}^\tau
\]
가 단사사상이 되고, 따라서 특히 항등원 연결 성분들 사이에서도 단사사상을
유도한다. 우변의 항등원 연결 성분은 앞의 결과에 의해 $k$ 위에서 고유이다.
또한 지금의 사상은 군 스킴의 준동형이므로 반드시 닫힌 몰입 사상이다. 따라서
$\operatorname{Pic}_{X/k}^\circ$도 $k$ 위에서 고유이다. 대수적으로 닫힌 체
위의 가환 대수군의 구조(슈발레--보렐에 의한 것)를 사용하면 까다로운 동치
판정법에 의존하지 않을 수 있다. 원점을 뺀 아핀 직선\footnote{인쇄본에는
``아핀 직선에서''라고 되어 있다. 이에 연결된 303쪽의 정오표는 ``원점을 뺀
아핀 직선''으로 고치도록 한다.}에서 $\operatorname{Pic}_{X/k}^\tau$로 가는
모든 사상이 상수임을 증명하는 문제로 환원된다. 이는
$X[t,t^{-1}]$ 위의 모든 가역 가군\footnote{인쇄본에는 $X[t]$라고 되어
있다. 이에 연결된 303쪽의 정오표는 $X[t,t^{-1}]$로 고치도록 한다.}이 $X$
위의 한 가역 가군에서 온다는 것과 동치이다. 이 사실은 잘 알려져 있고 매우
초보적이며, $X$가 $k$ 위에서 고유라는 사실조차 사용하지 않는다($X$가
정규라는 가정으로부터 곧바로 $X$가 정칙인 경우로 환원할 수 있다).

\result{따름정리}{2.3} $f:X\longrightarrow S$를 고유한 정규 사상(즉
평탄하고 기하적으로 정규인 섬유들을 갖는 사상)이라 하고,
$\operatorname{Pic}_{X/S}$가 존재한다고 가정하자. 그러면
$\operatorname{Pic}_{X/S}^\circ$는 $S$ 위에서 고유이고, 따라서
$\operatorname{Pic}_{X/S}$ 안에서 닫혀 있다. 더 나아가
$\operatorname{Pic}_{X/S}^\tau$와 $\operatorname{Pic}_{X/S}^\rho$도
닫혀 있고, ``동표수''인 경우에는 $\operatorname{Pic}_{X/S}^\sigma$도
닫혀 있다.

(1.1)과 (2.1), (ii)를 적용하면 된다.

\result{따름정리}{2.4} $f:X\longrightarrow S$를 고유하고 매끄러운
사상이라 하고, $\operatorname{Pic}_{X/S}$가 존재하며 $S$ 위에서 유한형인
스킴 $P^{(i)}$들의 서로소 합이라고 가정하자(cf. [3], V, 4.1). 그러면 각
$P^{(i)}$는 $S$ 위에서 고유이다.

이는 (2.1), (i)에 따른다. (2.2)에서 지적했듯이 $f$의 섬유들에 덜 제한적인
가정을 두어 이 결과를 일반화할 수 있지만, $f$가 정규라고만 가정하면 거짓이
된다. 이 경우에는 $\operatorname{Pic}_{X/S}^\tau$가 $S$ 위에서 유한형이라고
가정하더라도 여전히 $S$ 위에서 고유인지 나는 알지 못한다.

\src{233}
\result{정리}{2.5} $f:X\longrightarrow S$를 고유 평탄 사상이라 하고,
$\operatorname{Pic}_{X/S}$가 존재한다고 하자. 모든 정수 $n$에 대하여 이
군 준스킴의 $n$제곱 준동형을 $\varphi_n$이라 하자. 그러면 $\varphi_n$은
잉여 표수가 $n$과 서로소인 모든 점
$x\in\operatorname{Pic}_{X/S}$\footnote{인쇄본에는 ``$x\in X$''라고
되어 있으나, $\varphi_n$은 피카르 준스킴의 거듭제곱 자기준동형이다.}에서
에탈이다.

에탈 사상의 무한소적 특징짓기에 따르면 이 주장은 다음 명제와 동치이다.

\result{보조정리}{2.6} $S$를 국소 아르틴 환 $A$의 스펙트럼이라 하고, 그
극대 아이디얼 $\mathfrak m$이 $\mathfrak m^{\nu+1}=0$을 만족한다고 하자.
다음과 같이 놓자.
\[
A_{\nu-1}=A/\mathfrak m^\nu,
\qquad X_{\nu-1}=X\otimes_A A_{\nu-1}.
\]
$\mathcal L$을 $X$ 위의 가역 가군이라 하고, $\mathcal L'_{\nu-1}$을
$X_{\nu-1}$ 위의 가역 가군으로서 그 $n$차 텐서 거듭제곱이
\[
\mathcal L_{\nu-1}=\mathcal L\otimes_A A_{\nu-1}
\]
과 동형인 것이라 하자. 그러면 $n$이 $k=k(A)$의 잉여 표수와 서로소일 때,
$n$차 텐서 거듭제곱이 $\mathcal L$과 동형인 $X$ 위의 가역 가군
$\mathcal L'$이 존재한다.

보조정리를 증명하자. 다음과 같이 놓는다.
\[
V=\mathfrak m^\nu=\mathfrak m^\nu/\mathfrak m^{\nu+1}.
\]
이는 $k=k(A)$ 위의 벡터공간이다. 우선 $\mathcal L'_{\nu-1}$을 $X$ 위의
어떤 가역 가군 $\mathcal L'$으로 연장해 보자. 그러한 연장의 장애물은
\[
H^2(X_0,\mathcal O_{X_0})\otimes_k V
\]
에 속한다. 그런데 $\mathcal L'_{\nu-1}$에 둔 가정과
$\mathcal L_{\nu-1}$이 연장될 수 있다는 사실에 따라, 이 장애물에 $n$을
곱하면 영이 된다. $n$이 표수와 서로소이므로 장애물 자체가 영이다. 한편
선택한 연장의 불확정성은
\[
H^1(X_0,\mathcal O_{X_0})\otimes_k V
\]
에 속하고, $\mathcal L'^{\otimes n}$과 $\mathcal L$ 사이의 차이 $\xi$도
같은 가군에 속한다. 이 차이가 영이 되도록 $\mathcal L'$을 고치려면, 그
가군 안에서 $n\eta=\xi$를 만족하는 $\eta$를 찾으면 된다. 이는 역시 $n$이
표수와 서로소이므로 가능하다.

\result{따름정리}{} (2.5)의 조건 아래에서, 섬유 $X_s$들의 피카르 스킴이
가법 성분을 갖지 않는다고 더 가정하자(예를 들어 $X_s$들이 기하적으로
정규인 경우, cf. (2.1), (ii)). 그러면
$\operatorname{Pic}_{X/S}\longrightarrow S$는
$\operatorname{Pic}_{X/S}^\circ$의 점들에서 보편적으로 열려 있다.
$\operatorname{Pic}_{X/S}^\circ$가 닫혀 있으면(예를 들어 $X_s$들이
기하적으로 정규인 경우, cf. (2.3)),
$\operatorname{Pic}_{X/S}^\sigma$도 $S$ 위에서 보편적으로 열려 있다.
마지막으로 동표수인 경우
$\operatorname{Pic}_{X/S}^\rho\longrightarrow S$는 보편적으로 열려 있다.

(1.5)와 (1.1)을 적용하면 된다.

\src{234}
\result{따름정리}{2.7} $f:X\longrightarrow S$\footnote{인쇄본에는
``$f:X\longrightarrow Y$''라고 되어 있으나, 뒤따르는 함수와 피카르
준스킴은 모두 $S$를 밑으로 한다.}를 고유 평탄 사상이라 하고,
$\operatorname{Pic}_{X/S}$가 존재한다고 하자. 그러면 $S$ 위의 함수
\[
s\longmapsto\dim\operatorname{Pic}_{X_s/k(s)}
\]
는 상반연속이다(즉 위쪽으로 도약할 수는 있지만 아래쪽으로 도약할 수는 없다).
또한 $\operatorname{Pic}_{X_s/k(s)}$들이 가법 성분을 갖지 않으면 이 함수는
연속이기까지 하다(즉 국소 상수이다).

첫 번째 주장은 국소 뇌터 밑 위에서 국소 유한형인 모든 군 준스킴에 대하여
자명하거나 거의 자명하다. 항등 단면을 따라 살펴보는 것으로 충분하기
때문이다. 두 번째 주장은 (2.5)에 따른다.

\result{비고}{2.8} $s,s'\in S$이고 $s$가 $s'$의 특수화라고 하자. 그러면
(2.7)은 $X_{s'}$의 피카르 스킴과 그 ``특수화'' $X_s$의 피카르 스킴의
차원 사이의 부등식(각각 등식)과 동치이다. 더구나 세르는 피카르 스킴이
구성되기 전부터, 매끄러운 사상 $f:X\longrightarrow S$의 경우 $X_s$들의
피카르 스킴 차원의 불변성이 기본군의 특수화 이론([4], X)과 피카르 다양체
위의 유한 위수 점들과 아벨화한 기본군 사이의 고전적인 쿠머 관계([4], XI)의
형식적 귀결임을 지적했다. 일반적으로
$\operatorname{Pic}_{X_s/k(s)}^\circ$의 아벨 부분, 곱셈적 부분, 가법적
부분의 차원을 각각 $\alpha,\mu,\lambda$라 하고,
$\alpha',\mu',\lambda'$도 같은 방식으로 정의하자. 알려진 관계들은 다음
부등식으로 표현된다.
\[
\tag{*}\alpha+\mu+\lambda\geq\alpha'+\mu'+\lambda'
\]
(이는 $\operatorname{Pic}_{X/S}$가 존재하기만 하면 성립하고, 따라서 아마
항상 성립할 것이다). $\lambda=\lambda'=0$이면 이 부등식은 같은 존재 가정
아래에서 성립하는 등식
\[
\alpha+\mu=\alpha'+\mu'
\]
으로 바뀐다. 한편 다음 부등식도 있다.
\[
\tag{**}2\alpha+\mu\leq2\alpha'+\mu'
\]

\src{235}
$X_s$들이 분리가능이면 세르의 논증에 따라 이 부등식이 성립한다
($\operatorname{Pic}_{X/S}^\tau$의 존재조차 가정할 필요가 없다). 또는
${}_n\operatorname{Pic}_{X/S}$, 즉 피카르 함자 안에서 $\varphi_n$의 핵들이
$n$이 가역인 $S$의 열린 부분집합 위에서 분리되어 있는 경우에도 성립한다.
실제로 (2.5)에 따라 이 핵들은 그 열린 부분집합 위에서 에탈이다. (*)가 항상
등식이거나, 적어도 $X_s$들이 분리가능일 때에는 등식이라고 추측하고 싶다.
또한 다음 부등식들이
\[
\tag{***}\alpha\leq\alpha',\qquad\lambda\geq\lambda'
\]
$S$가 국소 뇌터일 때, $S$ 위에서 국소 유한형이고 섬유의 차원이 일정한 군
준스킴이 주어질 때마다 성립할 것으로 추측하고 싶다(이 방향의 긍정적인
결과는 (1.3)을 보라).

\result{비고}{2.9} 알려진 모든 경우에
$\operatorname{Pic}_{X/S}^\tau$는 $S$ 위에서 보편적으로 열려 있다. 그러나
$f:X\longrightarrow S$가 매끄럽더라도 이를 근거로 지나친 기대를 품어서는
안 될 듯하다. 어쨌든 멈퍼드는 $\operatorname{Pic}_{X/S}^\tau$가 $S$
위에서 평탄하지 않은 예를 구성했다. 이 예에서 $S$는 축소가 아니며, 실제로
아르틴 환의 스펙트럼이고, 이구사 곡면을 무한소적으로 변형하여 얻어진다.
멈퍼드가 고찰한 점은 $\operatorname{Pic}_{X/S}^\rho$에 속한다. 한편
$f:X\longrightarrow S$가 매끄러울 때 $\operatorname{Pic}_{X/S}^\tau$가
$\operatorname{Pic}_{X/S}^\sigma$의 점들에서 $S$ 위에 평탄할 가능성은
남아 있다. 그러나 강연자는 $\operatorname{Pic}_{X/S}^\circ$의 점들로
제한하고 $S$가 이산 값매김환의 스펙트럼이라고 가정하더라도 이것이 항상
성립할지는 의심한다. 이 문제는 유한 자기동형군의 작용 아래 아벨 스킴의
고정점을 연구하는 문제와도 관련되는데, 이 상황에 대해서는 예들이 부족한
듯하다.

$f:X\longrightarrow S$가 매끄럽고 사영이라고 제한하더라도, 이 절에서 말한
$\operatorname{Pic}_{X/S}$의 국소적 정칙성 결과들과 (2.8)에서 언급한
추측들은 $f$의 섬유의 성질에 더 특별한 가정을 두지 않고 이 문제에 관해 말할
수 있는 것을 거의 다 포괄하는 듯하다. 그러나
$\operatorname{Pic}_{X/S}$의 기하적 섬유들이 축소이고 가법 성분을 갖지
않으면, (1.8)과 (2.5)에 따라 $S$가 축소일 때
$\operatorname{Pic}_{X/S}$는 $\operatorname{Pic}_{X/S}^\sigma$의 점들에서
$S$ 위에 매끄럽다는 사실을 상기하자. $f:X\longrightarrow S$가 정규이면
$S$에 아무 가정도 두지 않고 같은 결과가 성립하며, 이는 (3.5)에서 보게 될
것이다. 이와 관련하여 다음을 언급하자.

\result{명제}{2.10}

\noindent(i) $\operatorname{Pic}_{X/S}$가 항등 단면의 점들에서 $S$ 위에
매끄럽다면(각각 평탄하다면), $\operatorname{Pic}_{X/S}^\sigma$의 모든
점과 $S$ 위의 $\operatorname{Pic}_{X/S}$의 모든 단면의 점들에서도 그러하다.

\src{236}
\noindent(ii) $s\in S$가
\[
H^2(X_s,\mathcal O_{X_s})=0
\]
을 만족한다고 하자. 그러면 $s$의 열린 근방 $U$가 존재하여
$\operatorname{Pic}_{X/S}|U$는 $U$ 위에서 매끄럽다.

\noindent(iii) $X$를 체 $k$ 위의 고유 스킴이라 하자. 그러면
\[
\dim\operatorname{Pic}_{X/k}\leq\dim H^1(X,\mathcal O_X)
\]
이고, 등호가 성립할 필요충분조건은 $\operatorname{Pic}_{X/k}$가 $k$ 위에서
매끄러운 것이다. $k$의 표수가 0이면 이는 항상 성립한다.

(i)은 (2.5)와 (1.12)에 따르고, (ii)는 매끄러운 사상의 무한소 판정법과 잘
알려진 장애물 계산에 따른다. 실제로 ``반연속성 정리''에 따라 $s$에서의
가정은 이웃한 점들에서도 계속 성립한다. 마지막으로 (iii)은
$H^1(X,\mathcal O_X)$가 $\operatorname{Pic}_{X/k}$의 항등원에서의 자리스키
접공간과 동형이라는 사실에 따른다. 마지막 주장은 표수 0의 ``형식군''은
$k$ 위에서 형식적으로 매끄럽다는 카르티에 정리의 특수한 경우이다.

\fsection{3}{\texorpdfstring{$\operatorname{Pic}_{X/S}$}{Pic(X/S)}의 표준 아벨 부분스킴과 알바네세
스킴}

\result{명제}{3.1} $k$를 체라 하고, $G$를 $k$ 위의 유한형 군 스킴으로서
가환이고 ``가법 성분이 없는'' 것이라 하자. 그러면
$G_{\mathrm{réd}}^\circ$는 $k$ 위에서 분리가능하며, 따라서 $k$ 위의
매끄러운 부분군 스킴이다.

$k$가 완전체일 때에는, 특히 $\bar k$가 $k$의 대수적 폐포일 때
$G_{\bar k}$에 대해서는 이 주장이 자명하므로,
$(G_{\bar k}^\circ)_{\mathrm{réd}}$가 $G$의 한 부분스킴에서 유래함을 보이면
충분하다. 그런데 $G_{\bar k}$에 대한 가정(가법 성분이 없음)\footnote{인쇄본에는
``Or l'hypothèse ...''라고 되어 있으나, 이에 연결된 정오표 303쪽은
``Or, de l'hypothèse ...''로 고치도록 규정한다.}에서 쉽게 다음이 따른다.
어떤 정수 $m$이 존재하여 $(G_{\bar k}^\circ)_{\mathrm{réd}}$는
$G_{\bar k}$에서의 $m$제곱 준동형의 ``스킴론적'' 상이다. 이 준동형은
$G$에 대한 같은 준동형에서 유래하므로, 후자의 스킴론적 상이 요구한
부분스킴이다.

\result{따름정리}{3.2} $X$를 $k$ 위에서 기하적으로 정규이고 고유인 스킴이라 하고
$\operatorname{Pic}_{X/k}$가 존재한다고 하자. 그러면
$\operatorname{Pic}_{X/k}$의 아벨 부분스킴 $A$로서 바탕 집합이
$\operatorname{Pic}_{X/k}^\circ$인 것이 존재한다.

실제로 (2.1), (ii)에 따라 $\operatorname{Pic}_{X/k}^\circ$는 $k$ 위에서
고유이므로 (3.1)의 조건을 만족한다.

\src{237}
따라서 앞의 결과는 어떤 경우에는 고전적인 ``피카르 다양체''(현대 이론에서는
$(\operatorname{Pic}_{X/k})_{\mathrm{réd}}^\circ$로 표기할 대상)가
체 $k$가 완전체라고 가정하지 않아도 ``$k$ 위에서 정의된다''는 것을 보인다.

이제 $f:X\longrightarrow S$를 고유하고 평탄한 상대 스킴이라 하고, 간단히
하기 위해
\[
\mathcal O_S\xrightarrow{\sim}f_*(\mathcal O_X)
\]
라고 하자. 또한 $\operatorname{Pic}_{X/S}$가 존재하고
$\operatorname{Pic}_{X/S}^\circ$가 $S$ 위에서 고유라고 하자. (3.3), (ii)를
위해서는 $\operatorname{Pic}_{X/S}^\circ$를 포함하고 $S$ 위에서
준사영인 $\operatorname{Pic}_{X/S}$의 한 열린부분이 존재한다고 더
가정한다. 앞에서 보았듯이 $f$가 사영이고 그 기하 섬유들이 분리가능하며
기약이면 이 조건이 성립한다. $S$ 위의 아벨 스킴이란 $S$ 위의 군
스킴으로서 $S$ 위에서 고유하고 매끄러우며 기하 섬유들이 연결인 것을
말함을 상기하자. 우리는 $\operatorname{Pic}_{X/S}$의 부분군 스킴 $A$로서
아벨 스킴이고 바탕 집합이 $\operatorname{Pic}_{X/S}^\circ$인 것이
존재하는지를 살펴보고자 한다. $S$가 한 체의 스펙트럼이면 그러한 것이
항상 존재함을 방금 보았다. 여기서 생각하는 일반적인 경우에는 다음을
말할 수 있다.

\result{정리}{3.3} 앞의 조건 아래에서 다음이 성립한다.

\noindent(i) $\operatorname{Pic}_{X/S}$의 아벨 부분스킴으로서 바탕 집합이
$\operatorname{Pic}_{X/S}^\circ$인 것이 존재하면, 그것은 유일하다. 따라서
그 형성은 밑변경과 양립한다.

\noindent(ii) 그러한 아벨 부분스킴이 존재하기 위해서는, $S'$가 아르틴
국소 스킴인 모든 밑변경 $S'\longrightarrow S$ 뒤에 그러한 것이 존재하면
충분하다. $S$가 한 국소환의 스펙트럼이면
$A_n=A/\mathfrak m^{n+1}$로 놓았을 때
$S'=\operatorname{Spec}(A_n)$인 경우들만 시험해도 충분하다. $S$가
축소이면, $S'$가 이산 값매김환의 스펙트럼인 밑변경
$S'\longrightarrow S$들만 시험해도 마찬가지로 충분하다(원한다면 그 환은
완비이고 그 잉여체는 대수적으로 닫혀 있다고 해도 된다).

\noindent(iii) $A$가 존재한다고 가정하고,
$B=\operatorname{Alb}^0(X/S)$를 그 쌍대 아벨 스킴, 즉
$B=\operatorname{Pic}_{A/S}^\circ$ [7]이라 하자. 그러면 $B$에 대한
주동차 공간 $P=\operatorname{Alb}^1(X/S)$와 $S$-사상
$X\longrightarrow P$를 표준적으로 구성할 수 있으며, 이 사상은 $X$에서
준아벨 스킴들(즉 아벨 스킴에 대한 주동차 공간들)로 가는 $S$-사상들에
대하여 보편적이다. $\operatorname{Alb}^0(X/S)$,
$\operatorname{Alb}^1(X/S)$ 및
$X\longrightarrow\operatorname{Alb}^1(X/S)$의 형성은 밑변경과 가환한다.

증명을 개략적으로 설명하자.

(i)은 가환 군 스킴들의 아벨 부분스킴에 관한 일반적인 강직성 성질이다.
그러한 두 부분스킴이 한 점 $s\in S$에서 집합론적으로 일치하면, $s$가
속한 연결 성분

\src{238}
전체 위에서 서로 일치한다 [7]. (이 결과는 Chow의 고전적인 한 정리를
일반화한다.)

(ii) 힐베르트 스킴을 사용하면, $S$ 위의 각 $S'$에
$(\operatorname{Pic}_{X/S})\times_S S'$의 표준 아벨 부분스킴들의 집합
(원소가 하나 또는 0개인 집합)을 대응시키는 함자가 $S$ 위의 유한형 스킴
$T$로 표현됨을 알 수 있다. (i)에 따라 $T\longrightarrow S$는
단사사상이고, (3.2)에 따라 전사이다. $\operatorname{Pic}_{X/S}$의 표준
아벨 부분스킴이 존재한다는 것은 $T$가 $S$ 위의 절단을 갖는다는 것, 다시
말해 $T\longrightarrow S$가 동형사상이라는 뜻이다. 이는
$T\longrightarrow S$가 에탈이라는 것과도 동치이고, $S$가 축소인
경우에는 $T\longrightarrow S$가 고유하다는 것과도 동치이다. 이로부터
(ii)가 곧 따른다.

(iii) $\operatorname{Pic}_{X/S}$의 정의만을 사용하면, $S$ 위의 임의의
아벨 스킴 $C$에 대하여 $X$에서 $C$에 대한 한 주동차 공간으로 가는
$S$-사상을 주는 것은 군 준동형
$C'\longrightarrow\operatorname{Pic}_{X/S}$를 주는 것과 동치임을 알 수
있다. 여기서 $C'$는 $C$의 쌍대 아벨 스킴이다. 그런데
$\operatorname{Pic}_{X/S}$의 표준 아벨 부분스킴 $A$가 존재하면, 이러한
준동형들은 반드시 $A$를 통해 인수분해된다(유한 위수의 점들을 사용해서도
이를 알 수 있다). 이로부터 (iii)가 곧 따른다.

\result{비고}{3.4} $\operatorname{Pic}_{X/S}^{\circ\circ}$로
$\operatorname{Pic}_{X/S}$의 표준 아벨 부분스킴을, 그것이 존재할 때,
나타내겠다. 불행히도 이것은 항상 존재하지는 않으며, 이구사 곡면을
무한소적으로 변형시키면(1차 모듈라이 변형을 통해) 이를 알 수 있다.
그럼에도 $\operatorname{Pic}_{X/S}^{\circ\circ}$는 적어도 $S$가 축소일
때에는 존재할 가능성이 있다. 또는 (ii)에 따르면 같은 말이지만, $S$가
한 이산 값매김환의 스펙트럼일 때에는 존재할 가능성이 있다. $X_0,X_1$을
$X$의 특수섬유와 일반섬유라 하고,
\[
A_1=\operatorname{Pic}_{X_1/K}^\circ
\]
라 하자. 여기서 $K$는 값매김환 $V$의 분수체이다. Koizumi의 결과에 따라
일반섬유가 $A_1$인, $S$ 위의 본질적으로 유일한 아벨 스킴 $A$가 존재한다.
이제 $X$가 $S$ 위에서 매끄럽다고 가정하면, (2.1), (i)에서와 마찬가지로
$A_1$의 항등사상이 다음 사상으로 연장됨을 쉽게 알 수 있다.
\[
A\longrightarrow\operatorname{Pic}_{X/S}.
\]
따라서 준동형
\[
\tag{*}A_0\longrightarrow\operatorname{Pic}_{X_0/k}^{\circ\circ}
\]
을 얻는다. 이것은 핵이 유한 $p$-군인 전사 준동형임을 쉽게 보일 수 있다.
여기서 $p$는 잉여체 $k$의 표수 지수이다.

\src{239}
(여기서도 유한 위수의 점들을 사용한다.) 이제 $X/S$에 대한 다음 조건들은
서로 동치이다.

\medskip
\noindent a. 앞의 준동형 (*)는 동형사상이다(Shimura의 표현으로는
``피카르 다양체''의 형성이 ``특수화와 양립한다'').

\noindent b. $\operatorname{Pic}_{X/S}^{\circ\circ}$가 존재한다(그때 그것은
$A$와 다름없다).

\noindent c. (기록을 위해) 모든
$\operatorname{Pic}_{X_n/S}^{\circ\circ}$가 존재한다.

(*)의 핵에 관해 한 비고에 따르면, 잉여 표수가 0이면 조건 (a)가 성립한다.
그러나 이 결과는 (3.5)에서 상당히 일반화될 것이다.

물론 $\operatorname{Pic}_{X/S}$가 $\operatorname{Pic}_{X/S}^\circ$의
점들에서 $S$ 위로 매끄러우면, 후자는 $\operatorname{Pic}_{X/S}$에서
열려 있다(cf. (1.7)). 따라서 유도된 구조를 주면 그것은
$\operatorname{Pic}_{X/S}$의 아벨 부분스킴이고, 그러므로 이 경우에
존재하는 $\operatorname{Pic}_{X/S}^{\circ\circ}$와 같다. 그러나 훨씬
더 좋은 결과가 성립한다.
\result{정리}{3.5} (3.3)의 조건 아래, $s\in S$이고
$\operatorname{Pic}_{X_s/k(s)}$가 $k(s)$ 위에서 매끄럽다고 하자(또는,
동치인 조건으로,
\[
\dim\operatorname{Pic}_{X_s/k(s)}=\dim H^1(X_s,\mathcal O_{X_s})).
\]
그러면 $s$의 열린 근방 $U$가 존재하여
$\operatorname{Pic}_{X/S}$는 $\operatorname{Pic}_{X/S}^\circ|U$의
점들에서 $S$ 위로 매끄럽다. 따라서 $\operatorname{Pic}_{X/S}^\circ|U$는
$\operatorname{Pic}_{X/S}|U$의 열린 아벨 부분스킴이다. 특히
$\operatorname{Pic}_{X|U/U}^{\circ\circ}$가 존재한다.

증명의 원리를 설명하자. 앞의 논의에 따라 $S$가 아르틴 국소환 $A$의
스펙트럼인 경우로 환원할 수 있고, $A$의 무한소 차수에 관하여 귀납한다.
따라서 $\operatorname{Pic}_{X_n/A_n}^\circ$가 $A_n$ 위에서 매끄럽다고
가정하고, $\operatorname{Pic}_{X_{n+1}/A_{n+1}}^\circ$가 $A_{n+1}$ 위에서
매끄러움을 보이는 것으로 충분하다. 이를 위해서는
\[
P_n=\operatorname{Pic}_{X_n/A_n}^\circ
\]
을 연장하는 $A_{n+1}$ 위의 아벨 스킴 $P_{n+1}$과 동시에,
$X_n\times_{A_n}P_n$ 위의 가역층 $\mathcal L_n$을 연장하는
$X_{n+1}\times_{A_{n+1}}P_{n+1}$ 위의 가역층 $\mathcal L_{n+1}$을
구성하면 충분함을 확인할 수 있다. 여기서 $\mathcal L_n$은 보편 문제의
해로서 피카르 스킴 $\operatorname{Pic}_{X_n/A_n}$을 정의할 때 나타나는
가역층이다. (주의 --- $X$가 $S$ 위의 단면을 갖는다고 가정할 수 있다~\ldots)
이 구성을 위해 다음 핵심 결과를 써야 한다. 아르틴 국소환의 몫환 위에
정의된 모든 아벨 스킴은 그 아르틴 국소환 위로 연장될 수 있다. 다시 말해,
대수적으로 닫힌 체 위의 아벨 스킴에 대한 절대적인 ``형식적 모듈라이
스킴''([3], II)은 $k$ 위의

\src{240}
비트 벡터 환 위에서 매끄럽다. 이 결과 자체는
올림 장애의 일반적인 형식적 성질과 군 연산들을 사용하여 얻는다. 이 결과를
인정하고 먼저 $P_n$을 임의의 방식으로 $P_{n+1}$로 연장한다. 그러면
$\mathcal L_n$을 올리는 데 대한 장애가 나타나며, 이 장애는
\[
H^2(X_0\times P_0,\mathcal O_{X_0\times P_0})\otimes_k V
\qquad(\text{여기서 }V=\mathfrak m^{n+1}/\mathfrak m^{n+2})
\]
에 속하고, 실제로는 그 부분공간
\[
H^1(X_0,\mathcal O_{X_0})\otimes
H^1(P_0,\mathcal O_{P_0})\otimes_k V
\]
에 속한다($\mathcal L_n$을 두 인자 $X_n$과 $P_n$에 제한하면 자명하기
때문이다). 그런데 이 마지막 공간은 바로
\[
H^1(P_0,\mathcal T_{P_0/k})\otimes V
\]
이다. 여기서 $\mathcal T_{P_0/k}$는 $P_0/k$의 접다발이다. 따라서 이는
$P_n$을 $P_{n+1}$로 올릴 때의 비유일성을 나타내는 공간이기도 하다([3],
II). 그러므로 이 올림을 (마땅히 그래야 하듯 유일한 한 가지 방식으로)
수정하여 $\mathcal L_n$을 올리는 데 대한 장애를 없앨 수 있다.

\result{따름정리}{3.6} (3.5)의 조건 아래,
$R^1f_*(\mathcal O_X)$는 $s$의 근방에서 $S$ 위의 국소 자유 가군이고,
그 구성은 밑변환과 가환한다.

실제로 이 가군은 $\operatorname{Pic}_{X/S}$의 단위원 단면을 따른
접가군이다.

\result{비고}{3.7} (3.5)에서와 같은 논증을 사용하면 다음을 보일 수 있다.
$S'$이 연접인 멱영 아이디얼층으로 정의되는 $S$의 부분스킴이고,
$\operatorname{Pic}_{X'/S'}^\circ$가 존재하며 그 섬유들이 매끄럽다고만
가정하면, $\operatorname{Pic}_{X/S}^\circ$는 반드시 존재하며 $S$ 위의
아벨 스킴이다. 이에 따라 사영성 가정이 없는 일부 경우에도 피카르 스킴을
구성할 수 있다. 예를 들어 아르틴 환 위의 아벨 스킴의 쌍대 아벨 스킴은
항상 존재한다. 한편 $X$가 $S$ 위의 아벨 스킴인 경우에 (3.6)을 적용하고,
$H^1(X_s,\mathcal O_{X_s})$ 위의 외대수로서
$H^*(X_s,\mathcal O_{X_s})$가 갖는 알려진 구조(로젠리히트--세르)를
사용하면, 모든 $i$에 대하여 $R^if_*(\mathcal O_X)$가 국소 자유이고,
더 정확히는 $R^1f_*(\mathcal O_X)$의 $i$차 외승과 동형임을 얻는다.

\result{비고}{3.8} $f:X\longrightarrow S$가 매끄러운 사영 사상이고,
$S$가 축약이며 모든 잉여 표수가 영인 경우, (3.6)의 결과는 호지 이론의
귀결로서 초월적 방법에 의해 알려져 있었다. 실제로 모든
$R^pf_*(\Omega^q_{X/S})$가 국소 자유이다. 반면 혼합 표수인 경우에는

\src{241}
세르 다양체를 이용하여 ``$R^1f_*(\mathcal O_X)$가 국소 자유이다''라는
명제의 반례를 얻는다([8], p.~50). 동표수인 경우에는 반례가 하나도 알려져
있지 않은 듯하다.

\result{따름정리}{3.8} (3.5)의 조건 아래,
$\operatorname{Pic}_{X/S}|U$는 $\operatorname{Pic}_{X/S}^\sigma|U$의
모든 점에서 $U$ 위로 매끄럽다.

(2.10)의 (i)를 적용하면 된다. 특히 (2.10)의 (iii)을 고려하면 다음을 얻는다.

\result{따름정리}{3.9} $S$의 모든 잉여 표수가 영이라고 가정하자. 그러면
$\operatorname{Pic}_{X/S}^\tau$는 $S$ 위에서 매끄럽다.

\result{비고}{3.10} 예를 들어 $\operatorname{Pic}_{X/S}^\tau$가 $S$
위에서 고유이면, $f$의 기하적 섬유들의 네론--세베리 군의 꼬임 부분군은
$S$의 각 연결 성분 위에서 일정하다는 결론을 얻는다($f$가 매끄럽고
사영이면 이는 초월적 방법으로도 명백하다). (2.5)를 직접 사용하면 더
일반적으로 다음을 보일 수 있다. $\operatorname{Pic}_{X/S}^\tau$가 $S$
위에서 고유이고(예를 들어 $f:X\longrightarrow S$가 매끄럽고 사영이면),
$q$가 $S$의 모든 잉여 표수와 다른 소수이면, $X/S$의 기하적 섬유들의
네론--세베리 군의 꼬임 부분군의 $q$-일차 성분은 $S$의 각 연결 성분
위에서 일정하다. 그러나 표수 $p>0$에서는 꼬임 부분군의 $p$-일차 성분에
대하여 더 이상 같은 결과가 성립하지 않는다. 그렇지만
\[
\operatorname{Pic}_{X_s/k(s)}^\tau/
\operatorname{Pic}_{X_s/k(s)}^{\circ\circ}=T_{X_s/k(s)}
\]
의 $k(s)$ 위의 계수가 국소 상수일 가능성은 남아 있다. $S$가 축약이면,
이것은 $\operatorname{Pic}_{X/S}^{\circ\circ}$가 존재하고
$\operatorname{Pic}_{X/S}^\tau$가 $S$ 위에서 평탄하다는 것과 동치임을
보일 수 있으며, $S$가 이산 값매김환의 스펙트럼인 경우에만 이 문제를
시험하면 충분하다. 내가 살펴본 몇 가지 예에서는 이를 확인했다. 그러나
$S$가 아르틴인 경우의 대응하는 명제는 거짓이므로((2.9)와 (3.4)의 비고
참조), 아마 지나친 기대를 품어서는 안 될 것이다.

\fsection{4}{피카르 스킴의 유한성 정리}

$f:X\longrightarrow S$를 $\operatorname{Pic}_{X/S}$가 존재하는 사영 평탄
사상이라 하자. ``힐베르트 다항식'' $Q\in\mathbf Q[t]$들을 생각하면
$\operatorname{Pic}_{X/S}$를 열린 부분스킴
$\operatorname{Pic}_{X/S}^Q$들의 서로소 합으로 분해할 수 있다. 예를 들어
$\operatorname{Pic}_{X/S}$가 $S$ 위에서 분리임을 보장하는 추가 가정을
두지 않으면, 일반적으로 이 열린 부분스킴들은 $S$ 위에서 유한형이지 않다.
$X$가 ``두 직선으로 퇴화하는 이차곡선''일 때 반례를 얻는다.

\src{242}
그렇지만 $f$가 분리가능이고 기하학적으로 기약인 섬유들을 가지면 이들이
유한형일 가능성이 있다. 이 문제는 $\operatorname{Pic}_{X/S}^\tau$가 $S$
위에서 유한형인지의 문제와 관련되며, 후자는 $X/S$의 섬유에 아무런 가정도
두지 않아도 참일지 모른다. $f:X\longrightarrow S$가 매끄러우면
$\operatorname{Pic}_{X/S}^\tau$가 $\operatorname{Pic}_{X/S}^Q$ 가운데
하나에 포함됨에 유의하자. 이는 비특이 사영 다양체 위에서 인자에 관한
``꼬임'' 동치가 수치적 동치보다 더 세밀하기 때문이다(실제로 Matsusaka에
따르면 두 동치는 ``동일하다''). 따라서 $\operatorname{Pic}_{X/S}^Q$들이
$S$ 위에서 유한형이면 $\operatorname{Pic}_{X/S}^\tau$도 그러하다. 방금
제기한 유한성 문제들은 $\operatorname{Pic}_{X/S}$의 존재조차 가정하지
않아도 명백한 의미를 갖는다. 실제로 이 문제들은 어떤 가역층들의 족이
[3]의 IV의 의미에서 ``유계''라고 말하는 것으로 표현된다. 모든 대수적으로
닫힌 체 $k$에 대하여, 차원이 $n$이고 차수가 $d$인 $\mathbf P_k^r$의 정역
부분스킴들과, 힐베르트 다항식이 $Q$인 그러한 준스킴 위의 가역층들을
생각하자(여기서 $r,n,d,Q$는 주어져 있다). 이 가역층들을
$\mathbf P_k^r$ 위의 연접 가군으로 볼 때, 그 족이 유계임, 즉 $\mathbf Z$
위의 유한형 스킴으로 매개될 수 있음을 보여야 한다.

Matsusaka의 방법 [6]을 사용하면, 상당히 기술적인 증명(``동치 판정법''을
알맞은 형태로 사용한다)을 통해 $\mathbf P_k^r$의 비특이 부분다양체들로
한정할 때에는 이 문제에 긍정적으로 답할 수 있다. 더 정확히는 다음 결과를
얻는다.
\result{정리}{4.1} $f:X\longrightarrow S$\footnote{인쇄본에는
``다음을 두자: $X\longrightarrow S$''라고 되어 있으나, 이하에서 사용하는
사상은 $f$이다.}를 기하학적으로 연결된 섬유를 갖는 매끄러운 사영 사상이라
하고, $S$는 뇌터라고 하자. $\mathcal O_X(1)$을 $S$에 관해 $X$ 위에서 매우
풍부한 가군이라 하자. $E$를 $\operatorname{Pic}_{X/S}$의 부분집합이라 하며,
이는 $X/S$의 기하학적 섬유 $\overline X_{s_i}$ 위의 가역 가군들의 족
$(\mathcal L_i)$에 대응한다고 하자. $D_i$를 $\overline X_{s_i}$ 위에서
$\mathcal L_i$를 정의하는 인자(반드시 유효일 필요는 없다)라 하고,
\[
a_n^{(i)}t^n+\cdots+a_0^{(i)}
\]
를 $\mathcal L_i$의 힐베르트 다항식이라 하며, $\xi=\xi_i$를
$\mathcal O_X(1)$을 정의하는 인자, 즉 초평면 절단이라 하자. 다음 조건들은
서로 동치이다.

\noindent a. $E$\footnote{인쇄본에는 여기서 ``$Q$''라고 되어 있다. 명제에서
도입되었으며 그 준콤팩트성이 이 족의 유계성을 나타내는 부분집합은
$E$이다.}는 준콤팩트이다. 즉 $S$ 위에서 유한형인 열린집합에 포함된다.
다시 말해 $(\mathcal L_i)$의 족은 유계이다.

\noindent b. $E$는 유한 개의 집합 $\operatorname{Pic}_{X/S}^Q$,
$Q\in\mathbf Q[t]$의 합집합에 포함된다. 즉 $\mathcal L_i$들의 힐베르트
다항식들로 이루어진 집합은 유한하다.

\noindent b'. ($X/S$의 모든 섬유의 차원이 같은 $n$인 경우.)
$\mathcal L_i$들의 힐베르트 다항식의 두 계수 $a_{n-1}^{(i)}$와
$a_{n-2}^{(i)}$는 유한한 집합 안에서만 변한다.

\src{243}

\result{따름정리}{4.2} $f:X\longrightarrow S$를 기하학적으로 연결된 섬유를
갖는 매끄러운 사영 사상이라 하자. 그러면 스킴
$\operatorname{Pic}_{X/S}^Q$와 $\operatorname{Pic}_{X/S}^\tau$
($Q\in\mathbf Q[t]$)는 $S$ 위에서 사영이다.

이 스킴들은 (4.1)에 따라 $S$ 위에서 유한형이므로 (2.4)를 적용할 수 있다.

\section*{참고문헌}
\addcontentsline{toc}{section}{참고문헌}

\begin{description}
\item[{[1]}] CHEVALLEY (Claude). -- Sur la théorie de Picard,
\emph{Amer. J. of Math.}, 제82권, 1960, 435--490쪽.

\item[{[2]}] GROTHENDIECK (A.) 및 DIEUDONNÉ (J.). -- \emph{Éléments de
géométrie algébrique}, 제I장 이하. -- Paris, Presses universitaires de France,
1960--1961 (Institut des Hautes Études Scientifiques, 4, 8, 11, \ldots).

\item[{[3]}] GROTHENDIECK (Alexander). -- Technique de descente et théorèmes
d'existence en géométrie algébrique, I--V, \emph{Séminaire Bourbaki}, 제12권,
1959/60, 제190호, 29쪽 및 제195호, 22쪽; 제13권, 1960/61, 제212호, 20쪽 및
제221호, 28쪽; 제14권, 1961/62, 제232호, 19쪽.

\item[{[4]}] GROTHENDIECK (Alexander). -- \emph{Séminaire de géométrie
algébrique}, 1960/61. -- Paris, Institut des Hautes Études Scientifiques, 1961.

\item[{[5]}] GROTHENDIECK (Alexander). -- \emph{Séminaire de géométrie
algébrique}, 1961/62, 청강자 집단이 정리함. -- Paris, Institut des Hautes
Études Scientifiques (출간 예정).

\item[{[6]}] MATSUSAKA (T.). -- The criteria for algebraic equivalence and the
torsion group, \emph{Amer. J. Math.}, 제79권, 1957, 53--66쪽.

\item[{[7]}] MUMFORD (D.) 및 TATE (J.). -- \emph{Séminaire de géométrie
algébrique}, Harvard University, 1962년 봄 학기 (출간 예정).

\item[{[8]}] SERRE (Jean-Pierre). -- Sur la topologie des variétés algébriques
en caractéristique $p$, \emph{Symposium internacional de topología algebraica}
[1956, Mexico], 24--53쪽. -- Mexico, Universidad nacional autónoma, 1958.
\end{description}

\par\smallskip
\begingroup\footnotesize





\noindent\textit{4.1의 인쇄본 조건 $b^{\prime\prime}$은 $D_i$와 $D_i^2$의 차수에 상계만을 요구하며, 나머지 유계성 조건들과 일반적으로 동치가 아니다. 예를 들어 $\mathbf P^1\times\mathbf P^1$에서 $H=(1,1)$, $D_m=(m,-m)$이면 $D_m\cdot H=0$이고 $D_m^2=-2m^2$이지만, $\mathcal O(m,-m)$의 힐베르트 다항식 $(t+1)^2-m^2$은 서로 다르다.}\par
\noindent\textit{1.5의 증명에는 각 기하적 섬유에서 $n$의 거듭제곱으로 소멸하는 꼬임점들의 조밀성이 필요하다. 가법 성분이 없다는 가정은 $n$이 잉여체에서 가역일 때 이 조밀성을 보장한다.}\par
\noindent\textit{2.8에서 ${}_n\operatorname{Pic}_{X/S}$에 대한 에탈성 주장은 $n$이 가역인 열린 부분집합에 적용된다.}\par
\noindent\textit{3.2의 인쇄본은 기하적 정규성을 명시하지 않고 $X$가 $k$ 위에서 정규이고 고유라고 가정한다. 3.10의 인쇄본에서 두 고유성 가정은 $\operatorname{Pic}_{X/S}^{\tau}$가 아니라 피카르 스킴 전체에 관한 것이다. 2.10(i)와 3.8의 인쇄된 기호는 $\sigma$이며, $\circ$와 $\tau$는 필사상의 오기이다.}\par
\endgroup

\par
\renewcommand{\sourceunit}{e236}
\clearpage
\chapter*{강연 236의 정오표 및 보충}
\addcontentsline{toc}{chapter}{강연 236의 정오표 및 보충}
\src{303}

\noindent\textbf{8. TDTE VI: 피카르 스킴. 일반 성질.}

\begin{center}
[부르바키 세미나, 제14권, 1961/62, 제236호, 23쪽.]
\end{center}

\noindent\textbf{12쪽, 16행.} --- ``아핀 직선'' 대신
``원점을 뺀 아핀 직선''으로 읽어야 한다.

\noindent\textbf{18행.} --- ``$X[t]$'' 대신
``$X[t,t^{-1}]$''로 읽어야 한다.

\noindent\textbf{16쪽, 명제 3.1, 6행.} --- ``그런데 가정은 …'' 대신
``그런데 가정으로부터 …''로 읽어야 한다.

\noindent\textbf{21쪽, § 4: 피카르 스킴에 대한 유한성 정리.} ---
이 절에서 제기한 종류의 유한성 문제들은 이 강연을 작성한 뒤 거의 완전히
해결되었다. 이제 이 방향에서 알려진 주요 사실들을 제시하자. 명제를
간단히 하기 위하여, 명제들에 나타나는 모든 피카르 준스킴이 존재한다고
암묵적으로 가정한다. 다만 이 명제들을 명백하게 수정하면 어떠한 존재
가정도 둘 필요가 없다. 이하에서 $S$는 뇌터 스킴을, $X,Y$는 $S$ 위에서
고유한 스킴들을 나타낸다.

\medskip
\noindent(i) $f:X\longrightarrow Y$를 전사 사상이라 하자. 그러면
\[
f^*: \operatorname{Pic}_{Y/S}\longrightarrow\operatorname{Pic}_{X/S}
\]
는 유한형 사상이다.

\src{304}
\noindent(ii) $Y$가 $S$ 위에서 사영이고 $S$에 관해 풍부한 가역 가군을
갖는다고 하자. $X$를 이 가군의 임의의 단면의 영점 스킴이라 하고,
$f:X\longrightarrow Y$를 표준 몰입 사상이라 하자. 끝으로 $Y/S$의
섬유들의 기약 성분들이 차원 $\geq3$이라고 가정하자. 그러면
\[
f^*: \operatorname{Pic}_{Y/S}\longrightarrow\operatorname{Pic}_{X/S}
\]
는 유한형이다.

\noindent(iii) $X$가 $S$ 위에서 사영이고 그 모든 기하적 섬유가 차원
$n$인 정역 스킴이라고 하자. $\mathcal O_X(1)$을 $X$ 위의 풍부한 가역
가군으로서 힐베르트 다항식들을 정의할 수 있게 하는 것으로 잡자. $M$을
$\operatorname{Pic}_{X/S}$의 한 부분집합이라 하자. 그러면 $M$이
준콤팩트일 필요충분조건은 $M$의 원소들의 힐베르트 다항식
\[
a_0x^n+a_1x^{n-1}+a_2x^{n-2}+\cdots
\]
에서 계수 $a_1$과 $a_2$가 유계로 남는 것이다. 여기서는 그 단항식의
차수가 음수가 되는 모든 계수를 생략한다. 따라서 $n=1$이면 $a_1$만
관여하고, $n=0$이면 둘 다 관여하지 않는다.

\noindent(iv) 모든 정수 $n\neq0$에 대하여, 군 준스킴
$\operatorname{Pic}_{X/S}$의 $n$제곱 준동형은 유한형 사상이다.

(i)과 (ii)는 또한 (각각의 명제에서 고려한 가정 아래)
$\operatorname{Pic}_{Y/S}$의 부분집합 $M$이 준콤팩트일 필요충분조건이
$\operatorname{Pic}_{X/S}$ 안에서 그 상이 준콤팩트인 것임을 뜻한다.
따라서 $Y$ 위의 가역 가군 $\mathcal L$이 $0$과 $\tau$-동치일
필요충분조건은 $X$로의 그 역상이 그러한 것이다. 다시 말해
$\operatorname{Pic}_{Y/S}^\tau$는
$\operatorname{Pic}_{X/S}^\tau$의 역상이다. 특히 앞의 준스킴이 $S$
위에서 유한형임을 보이려면 뒤의 준스킴에 대하여 이를 증명하면 충분하다.
실제로 $\operatorname{Pic}_{Y/S}\longrightarrow
\operatorname{Pic}_{X/S}$는 유한형이다. 이제 (i), 차우 보조정리와
(iii)을 사용하면 다음을 얻는다.

\medskip
\noindent(v) $\operatorname{Pic}_{X/S}^\tau$는 $S$ 위에서 유한형이다.

일반적으로 유한 사상에 대한 (i)과 (ii)를 함께 사용하면 대부분의 유한성
문제를 $X/S$가 차원 $\leq2$인 기하적으로 정역이고 정규인 섬유들을 갖는
경우로 환원할 수 있다. 흔히 유한일 필요가 없는 전사 사상에 (i)을
적용하고 대수 곡면의 특이점 해소(임의의 표수에서 아브향카르에 의한
결과)를 사용하면, $X/S$가 매끄러워서 기하적 섬유들이 차원 2인
비특이 스킴인 경우로까지 환원된다. 예를 들어 (v)와 비특이 사영 곡면의
네론--세베리 군의 계수(階數)에 상계를 주는 피카르--이구사 부등식을
함께 사용하면, 네론의 유한성 정리의 다음 일반화를 증명할 수 있다.

\medskip
\noindent(vi) $X/S$가 $S$ 위에서 고유하고 그 밖에는 임의적이라고 하자.
그러면 $X/S$의 기하적 섬유 $X_i/k_i$의 네론--세베리 군
\[
\operatorname{Pic}_{X_i/k_i}/\operatorname{Pic}_{X_i/k_i}^\circ
\]

\src{305}
은 유한 생성이며, 그 계수(階數)와 꼬임 부분군의 위수에는 공통 상계가
있다.

비특이 곡면의 경우로 환원하는 같은 방법과 이 경우에 알려진 정리들, 즉
네론의 유한성 정리와 네론--세베리 군 위의 교차 형식이 비퇴화라는 사실은
다음을 함의한다.

\medskip
\noindent(vii) $X$를 대수적으로 닫힌 체 위의 고유 스킴이라 하자.
그러면 $X$ 안에는 다음 성질을 갖는 유한 개의 닫힌 정역 곡선
$C_i$ ($1\leq i\leq r$)가 존재한다. $\operatorname{Pic}_{X/k}$의
모든 부분집합 $M$에 대하여, $M$이 준콤팩트일 필요충분조건은 정수들
\[
\deg_{C_i'}\mathcal L\qquad(\mathcal L\in M)
\]
이 유계로 남는 것이다. 여기서 $C_i'$는 $C_i$의 정규화를 나타낸다.

여기서 $r$은 네론--세베리 군의 계수(階數)로 잡을 수 있다. 이 군이 유한
생성임을 알고 나면, (vii)은 $X$ 안의 곡선 $C$가 네론--세베리 군 위에
정의하는 선형 형식들이 동시에 영이 되는 원소는 네론--세베리 군의 꼬임
원소뿐이라는 사실로 환원된다. $X$가 사영이고 비특이인 경우에는 이 결과와
(v)가 마쓰사카에 의한 것이었다. (vii)을 사용하면 $X$ 위에서 $0$과
$\tau$-동치인 가역 가군들의 다음 특성화를 쉽게 얻는다.

\medskip
\noindent(viii) $X/k$를 한 체 위의 고유 스킴이라 하고, $\mathcal L$을
$X$ 위의 가역 가군이라 하자. 다음 조건들은 서로 동치이다.

\noindent a. $\mathcal L$은 $0$과 $\tau$-동치이다.

\noindent b. 모든 연접 가군 $F$에 대하여
\[
\chi(F\otimes\mathcal L)=\chi(F)
\]
이다. 여기서 $\chi$는 오일러--푸앵카레 특성을 나타낸다.

\noindent b'. (b)에서와 같되 $F=\mathcal O_Y$로 한다. 여기서 $Y$는
$X$ 안의 차원 1인 닫힌 정역 부분스킴이다.

\noindent c. 위와 같은 모든 $Y$에 대하여, $Y'$을 그 정규화 곡선이라
하면
\[
\deg\mathcal L_{Y'}=0
\]
이다.

$X/k$가 사영이고 풍부한 가역 가군 $\mathcal O_X(1)$을 갖는다면, 앞의
조건들은 다음 조건들과도 서로 동치이다.

\noindent d. 모든 정수 $m$에 대하여 $\mathcal L^{\otimes m}(1)$은
풍부하다.

\noindent e. ($X$가 정역인 경우.) 모든 정수의 쌍 $m,n$에 대하여
\[
\chi(\mathcal L^{\otimes m}(n))=\chi(\mathcal O_X(n))
\]

\src{306}
이다. 즉 $F=\mathcal L^{\otimes m}(n)$으로 놓으면 (b)가 성립한다.

마지막 조건이 충분함을 보려면, 이 조건이 $\mathcal L^{\otimes m}$의
힐베르트 다항식들이 모두 같다는 뜻임에 유의한다. 따라서 멈퍼드의 판정법
(iii)에 따라 $\mathcal L^{\otimes m}$들은
$\operatorname{Pic}_{X/k}$의 한 준콤팩트 부분집합 안에 머물며, 이는
(a)이다. 조건 (b), (b'), (c)는 보통 비특이 사영 다양체 위에서 정의하는
수치적 동치 개념을 임의의 고유 스킴 위로 옮긴 변형으로 보아야 한다.
비특이 사영 다양체의 경우 (a)와 (c)의 동치는 물론 잘 알려져 있었다
(마쓰사카).

판정법 (e)는 또한 다음 결과를 함의한다.

\medskip
\noindent(ix) $f:X\longrightarrow S$를 기하적 섬유들이 정역인 사영
평탄 사상이라 하자. 그러면 $\operatorname{Pic}_{X/S}^\tau$는
$\operatorname{Pic}_{X/S}$ 안에서 열리고 닫혀 있다.

핵심 결과 (i), (ii), (iii)의 증명에 관한 몇 가지 논평만 하기로 한다.
결과 (iv)는 다른 결과들과 다소 별개이며, 유한하고 전사이며 라디시엘인
사상, 더 정확히는 프로베니우스 사상에 대하여 (i)만을 사용해서 증명된다.
(i)에는 비평탄 하강의 발상들을 본질적으로 사용한다(특히 TDTE I, 9쪽을
보라). 유한성 결과만을 목표로 하므로 하강 자료의 유효성 판정법이 없는
것은 아무런 문제가 되지 않는다. 한편 멈퍼드는 최근 (iii)의 약간 약한
형태, 즉 힐베르트 다항식의 모든 계수를 사용하는 판정법을 증명했다. 그의
논증은 극히 간단하며, 나카이의 풍부성 판정법의 증명에서 착안한 것이다
(나카이는 비특이 곡면에 대하여 이를 진술했고, 멈퍼드는 임의의 사영
사상으로 일반화했다). 다만 내가 보기에는 이 논증이 $X/S$의 섬유에
약간의 추가 제한, 곧 세르의 조건 $(S_2)$를 둘 때에만 타당하다. 기하적
섬유들이 정규이면 이 조건은 성립한다. 이제 이 제한된 판정법을 (ii)의
증명에 사용한다. 실제로 판정법 (i)에 의해 $Y/S$가 평탄하고 기하적으로
정역이며 정규인 섬유들을 갖는 경우로 환원할 수 있고, 멈퍼드의 판정법을
적용하면 $X/S$가 같은 조건들을 만족하는 경우로 쉽게 환원할 수 있다.
그러면 차원에 대한 가정으로부터 $Y$와 $X$의 기하적 섬유들은 그 닫힌점들
에서 깊이 $\geq2$임을 알 수 있다. 따라서 SGA 1962(강연 XII)에 제시된
형태의 ``동치 판정법''을 적용하여 (ii)의 증명을 끝낼 수 있다. (i)과
(ii)를 얻은 뒤에는 멈퍼드의 판정법에서 섬유들의 정규성에 관한 모든
가정을 제거하고, (iii)에 제시된 더 강한 형태로도 이를 증명하기 어렵지
않다.

\src{307}
끝으로 (i)과 (ii)의 증명은 $S$가 체 $k$의 스펙트럼인 경우 사상
\[
\operatorname{Pic}_{Y/k}\longrightarrow\operatorname{Pic}_{X/k}
\]
이 유한형일 뿐만 아니라 아핀임을 보여 준다.

\par\smallskip
\begingroup\footnotesize
\noindent\textit{조건 (iii)의 계수는 단항식의 차수가 음수가 되지 않는 경우에만 고려한다. 따라서 $n=1$이면 $a_1$만, $n=0$이면 어느 계수도 고려하지 않는다.}\par
\endgroup

\par
\renewcommand{\sourceunit}{com}
\unit{1962년 논평}
\sid{com}
\src{297}
\seminar{제14년차, 1961/62, 보충\\1962년 5월}
\paperhead{대수기하학의 기초.\\논평}%
  {알렉상드르 그로텐디크 지음}

\section*{일러두기}
\addcontentsline{toc}{section}{일러두기}

우리는 1957년부터 1962년까지 부르바키 세미나에서 대수기하학의 기초에
관한 여덟 강연을 하였다. 첫 번째를 제외한 이 강연들은 스킴의 틀 안에서
서술되어 있다. 여기서 진술한 모든 결과는 장 디외도네와 알렉상드르
그로텐디크의 『Éléments de géométrie algébrique』에 실릴 것이다. 그러나
이 강연들 가운데 어느 것의 내용도 현재 이미 출판되었거나 준비 중인 어느
장에도, 다른 어떤 책이나 논문에도 다루어져 있지 않으며, 아마 앞으로도
몇 해 동안은 그러할 것이다. 이것이 우리가 이 강연들을 한데 모은 주된
이유다. 이 강연들은 정식 집필본을 기다리는 동안 이용자들에게 스킴 이론의
몇몇 개념과 핵심 결과를 제공할 것이다. 더욱이 이 글들을 읽으면 체계적
논고에서 불가피하게 번잡해지는 세부에 방해받지 않고 그 결과와 개념에
빠르게 익숙해질 수 있으며, 그러한 논고를 공부하는 데 반드시 필요한
동기도 얻을 수 있다.

독자의 편의를 위해 강연별로 묶은 몇 가지 논평과 정정 사항을 여기에
모았다. 특히 이들은 각 글이 작성된 뒤 이루어진 진전을 곳곳에서 알리고
몇몇 추가 참고문헌을 제시한다.

이 강연들에 나오는 여러 결과는 IHES 대수기하학 세미나와 1961/62년에
하버드에서 연이어 열린 두 세미나에서 자세히 다루어졌다. 첫 번째 세미나는
내가, 두 번째 세미나는 멈퍼드--테이트가 진행했으며, 이 두 세미나의 등사판
강의록(리히텐바움 씨 작성)은 현재 준비 중이다.

\section*{사용한 약어}
\addcontentsline{toc}{section}{사용한 약어}

\begin{description}
\item[SGA :] 대수기하학 세미나, Institut des Hautes Études Scientifiques,
1960--1962.

\src{298}
\item[SHG :] 그로텐디크 세미나, 하버드 대학교, 1961/62.

\item[SHMT :] 멈퍼드--테이트 세미나, 하버드 대학교, 1961/62.

\item[EGA :] GROTHENDIECK (A.)와 DIEUDONNÉ (J.). --- Éléments de géométrie
algébrique. --- Paris, Presses universitaires de France, 1960, 1961, …
(Institut des Hautes Études Scientifiques, Publications mathématiques, 4, 8,
11, …).

\item[TDTE :] GROTHENDIECK (Alexander). --- Technique de descente et théorèmes
d'existence en géométrie algébrique, I--VI, \emph{Séminaire Bourbaki}, 제12권,
1959/60, 제190호와 제195호; 제13권, 1960/61, 제212호와 제221호; 제14권,
1961/62, 제232호와 제236호.
\end{description}


\par
\end{document}
